% Generated by build_chapter.py from the GY1975 files. Do not edit by hand; edit the files and rebuild.
\chapter{The academic year 1975--76}\label{ch:1975}
The 1975--76 material comprises 34 catalogue records and 56 scanned pages,
\GSpage{P000603}--\GSpage{P000658}. The occurrence index below contains
256 distinct items. Record 378 is a year heading and has no invented
mathematical content. Catalogue scan envelopes are retained for talks;
the narrower page references on items identify the actual passages.
This distinction matters because several scans contain successive
speakers, and the catalogue does not consistently identify interventions.
In particular, record 383 is catalogued under Novikov, but its manuscript
heading names Drinfeld. Shnirelman's substantial spectral discussion is
embedded within record 391, catalogued under Krichever.

\begin{center}
\begin{tabular}{lrrrrrrr}
\toprule
Class&A&B&C&D&E&F&Total\\
\midrule
Items&45&16&5&0&174&16&256\\
\bottomrule
\end{tabular}
\end{center}

The principal mathematical line is the passage from stationary integrable
equations to commuting ordinary differential operators, spectral curves,
and line bundles. The autumn talks explain finite-gap and rank-one
constructions; the February--March discussions articulate the role of
filtered modules, the punctured-curve ring, acyclic line bundles, and
translation on the Jacobian. Drinfeld's 1977 note is an unusually precise
publication match, while Krichever's subsequent treatment separates
higher rank from merely choosing two operators of noncoprime orders
\citep{drinfeld1977commutative,krichever1978commutative}.
Earlier finite-gap occurrences remain essential context, but a recurrence
of a subject is not automatically the recurrence of an identical problem.

Polyakov's sequence moves from perturbative quantum fields and particle
models to gauge theory, static charges, canonical reduction, and a
lattice integral. The corrections repeatedly concern the distinction
between a group and its Lie algebra, gauge-invariant quantities and
gauge-dependent potentials, formal functional integrals and finite
measures, and strong-coupling lattice results and the four-dimensional
continuum confinement problem. The lattice construction was already
published before these talks \citep{wilson1974confinement}; that historical
fact is separated from the unresolved continuum claim. The spring
Hamiltonian calculation does not derive confinement simply by averaging
an instantaneous Coulomb kernel.

The other lines are not marginal appendices. Manin's arithmetic discussion
contains exact formulas whose numerical normalizations can be tested;
Kirillov's talks connect local Lie brackets, deformation, and
supergeometry; Roiter's arithmetic root criterion has an exact later
match in Proposition 4 of his 1978 paper \citep{roiter1978roots}.
Novikov's radial inverse problem is identifiable as Brodskii's scalar
model, rather than a theorem recovering every parameter of an arbitrary
elastic solid. The concluding algebraic discussion belongs to the
period of the solution of Serre's projective-module problem.

Two bounded developments are supplied. For the Fuchs--Perchik question,
the relative Hamiltonian Euler product is normalized, its finite cutoff
is proved, and two exact algorithms agree through weight 300. There are
135 nonzero positive-weight Euler coefficients in that range, and the
notebook's lower bound of 57 classes is recovered. These finite facts do
not prove infinite nonvanishing. The relevant modern computation retains
the exact infinitude question \citep[\S2]{mohammadzadehfuchs2014cohomology}.
For Shnirelman's Abel spectral-phase expression, a Fourier convergence
criterion, a complete circle calculation, and a qualified sphere
counterexample identify the role of a uniform bound. Neither development
is presented as a solution of its unbounded general question or as a
priority claim for the elementary arguments.

There are 20 indexed problems and 91 correction or
qualification entries. Class F marks an identification limit, not an
undeveloped research programme; no class D is assigned merely because a
search failed. Literal errors, omitted hypotheses, and publication matches
are assessed independently. The report and source ledgers explain the
remaining uncertainty, distinguish manuscript submission from public
circulation, and record which references were read in full and which
were checked only bibliographically. All present-status statements are
bounded by the documented searches as of 10 November 2026.

\GStalk{363}{1975-09-15}{A. A. Kirillov}{P000603-P000605}{Deforming the Poisson bracket}
The introductory announcement of Polyakov's physics lectures is followed by Kirillov's report of Vey's deformation theorem. The page distinguishes local finite-order brackets from formal bidifferential deformations, and ends with the suggestive relation \(pq-qp=t\). The displayed planar formula suppresses necessary mixed derivatives; its omission is mathematical shorthand or error in the notebook, not an error introduced by translation.

\begin{GSitem}{1975-363.1}{A}{Local Lie brackets and the locality hypothesis}
\GSfield{Kind; source}{claim; \GSpage{P000603}}
\GSfield{The notes}{On a symplectic manifold \((M,\omega)\), there are ``essentially no other Lie algebra structures'' besides the Poisson bracket, with Kirillov and others named.}
\GSfield{Statement}{A meaningful precise version concerns local Lie brackets on sections of a real or complex line bundle over a smooth manifold. Such brackets are differential operators of order at most one in each entry and are described locally by Jacobi brackets, including Poisson and contact brackets. This does not classify arbitrary Lie algebra structures on the vector space \(C^\infty(M)\).}
\GSfield{Reconstruction}{Locality, the line-bundle framework, and the qualification by Jacobi brackets are supplied to delimit the claim; the notes do not state a classification theorem with these hypotheses.}
\GSfield{Fate}{\citet[abstract]{kirillov1976local} is a close subsequent publication by the speaker (received 25 November 1975; published 1976). Only its publisher abstract and metadata were verified here, so an exact match to the notebook's shorthand is not asserted.}
\GSfield{Present status}{As of 10 November 2026: Established in the stated scope.}
\GSfield{Assessment}{literal statement: not determinable; publication match: partial; verification: publisher metadata checked; development: none; confidence: medium}
\end{GSitem}
\begin{GSitem}{1975-363.2}{E}{Formal deformation and the reported third-cohomology condition}
\GSfield{Kind; source}{theorem; \GSpage{P000603}}
\GSfield{The notes}{The theorem attributed to J. Vey asserts an inequivalent formal deformation \([u,v]_t=\sum_{k\geq0}t^k B_k(u,v)\), where \(B_0\) is Poisson and \(B_k\) has order \(2k+1\) in each argument. Global definition is said to follow from \(H^3(M,\mathbb R)=0\).}
\GSfield{Statement}{For a smooth symplectic manifold, seek skew bidifferential operators \(B_k\) such that the displayed series satisfies Jacobi coefficient by coefficient, with equivalence specified by invertible formal differential changes of functions. The cohomological condition in the notes is a sufficient condition reported for Vey's construction; it is not a necessary condition for existence of deformation quantization.}
\GSfield{Reconstruction}{The coefficientwise interpretation and equivalence by formal differential operators make explicit what ``formal'' and ``inequivalent'' must mean. The notes's narrower phrase ``change of coordinates'' alone does not specify the customary equivalence relation.}
\GSfield{Fate}{\citet[pp.~421--454]{vey1975deformation} is the named 1975 publication; its exact obstruction argument was not obtained here. The report treats it as an existing theorem, but the publication month was not verified. The later Fedosov construction gives star products on arbitrary smooth symplectic manifolds; its mechanism is checked in \citet[Sections~1--2]{emmrichweinstein1993differential}.}
\GSfield{Present status}{As of 10 November 2026: Formal symplectic deformation quantization is established without the displayed vanishing assumption; the exact Vey formulation in this shorthand remains only partially verified.}
\GSfield{Assessment}{literal statement: correct with missing hypotheses; publication match: partial; verification: indirect; development: none; confidence: medium}
\end{GSitem}
\begin{GSitem}{1975-363.3}{E}{The planar Moyal bracket requires mixed derivatives}
\GSfield{Kind; source}{construction; \GSpage{P000603}}
\GSfield{The notes}{The coefficient \(B_k\) is drawn as the difference of the two pure \((2k+1)\)-st derivative products, divided by \((2k+1)!\), followed by dots called lower terms; on the plane or torus the lower terms are said to vanish.}
\GSfield{Statement}{For a constant Poisson tensor \(P=\partial_p\otimes\partial_q-\partial_q\otimes\partial_p\), one consistent formal normalization is \[ [f,g]_t=\sum_{k\geq0}\frac{t^k}{(2k+1)!}\,\mu P^{2k+1}(f\otimes g). \] Here \(\mu\) multiplies the tensor factors. Already \(P^3\) contains the mixed terms \(-3f_{ppq}g_{qqp}+3f_{pqq}g_{qpp}\); these have the same differential order as the two displayed pure terms. The formula descends to a flat symplectic torus using periodic functions.}
\GSfield{Reconstruction}{This is an explicit flat-model reconstruction. Writing \(t=h^2\), it is \((f\star_h g-g\star_h f)/(2h)\) for \(\star_h=\mu\exp(hP)\). Curvature corrections can vanish in a flat model; the binomial mixed derivatives cannot.}
\GSfield{Fate}{The flat Weyl/Moyal model is described in \citet[Section~2]{emmrichweinstein1993differential}. No assertion that the notebook's two pure terms alone define a Lie bracket is retained.}
\GSfield{Present status}{As of 10 November 2026: Established in the stated scope.}
\GSfield{Assessment}{literal statement: false as written; publication match: partial; verification: read the result; development: annotated; confidence: high}
\end{GSitem}
\begin{GSitem}{1975-363.4}{F}{The final canonical-commutator suggestion}
\GSfield{Kind; source}{remark; \GSpage{P000603}}
\GSfield{The notes}{The page ends with a question mark followed by \(pq-qp=t\).}
\GSfield{Statement}{This can be read as the formal Weyl-algebra relation for two generators over \(\mathbb C[[t]]\). It suggests an associative quantization whose commutator recovers a Poisson bracket after division by the parameter. The line does not specify whether existence, representations, global gluing, or comparison with the preceding deformation is being asked.}
\GSfield{Reconstruction}{The Weyl-algebra interpretation is a possible identification, not a recovered definite research question. The use of the same letter \(t\) as in the previous expansion need not preserve the normalization.}
\GSfield{Fate}{The relation is compatible with the local Weyl construction in \citet[Section~2]{emmrichweinstein1993differential}, but that compatibility does not identify a unique problem or its fate.}
\GSfield{Present status}{As of 10 November 2026: No definite question can be recovered from this line.}
\GSfield{Assessment}{literal statement: not determinable; publication match: thematic only; verification: unresolved lead; development: none; confidence: low}
\end{GSitem}

\GStalk{364}{1975-09-15}{A. M. Polyakov}{P000603-P000611}{Spinors, gauge invariance, and perturbative quantum electrodynamics}
This extended introductory lecture moves from the particle table to Weyl and Dirac fields, local phase covariance, electromagnetic connections, propagators, Feynman parameters, and charge screening. Gelfand's Gaussian-integral analogy interrupts the discussion. The notebook combines exact finite-dimensional identities with formal functional integrals and schematic diagrams; the companion preserves that distinction.

\begin{GSitem}{1975-364.1}{E}{The 1975 particle table and interaction-strength heuristic}
\GSfield{Kind; source}{remark; \GSpage{P000604}}
\GSfield{The notes}{The table lists photons; electrons, muons and two massless neutrinos; protons, neutrons and pions. Approximate masses are \(0.5,105,940,140\), and gravitational, weak, electromagnetic and strong strengths are ranked by \(10^{-38},10^{-5},10^{-2},1\).}
\GSfield{Statement}{The masses are approximate values in \(\mathrm{MeV}/c^2\). The table is a historical teaching selection, not an exhaustive particle classification. A dimensionless interaction strength must specify the process and energy scale; \(E_{\rm int}/m\) is not a universal constant for each interaction. Statements ordering the number of violated conservation laws by strength have no general theorem behind them. Massless neutrinos are an approximation of the then minimal model.}
\GSfield{Reconstruction}{Units, scale dependence, and the distinction between a model assumption and a measured mass are supplied. No numerical modernization of this historical table is attempted.}
\GSfield{Fate}{The later oscillation evidence in \citet[abstract]{fukuda1998oscillation} rules out treating all neutrinos as exactly massless. The strength ranking is retained as an introductory heuristic, not a theorem awaiting development.}
\GSfield{Present status}{As of 10 November 2026: The particle table is historical; the neutrino masslessness and universal strength assertions are not valid literal modern statements.}
\GSfield{Assessment}{literal statement: false as written; publication match: partial; verification: publisher metadata checked; development: none; confidence: high}
\end{GSitem}
\begin{GSitem}{1975-364.2}{E}{Spin one-half, Weyl spinors, and Lorentz boosts}
\GSfield{Kind; source}{construction; \GSpage{P000604-P000605}}
\GSfield{The notes}{Rotations act by \(\exp(i\boldsymbol\sigma\cdot\boldsymbol\Theta/2)\); the two boost laws are \(e_L\mapsto e^{\boldsymbol\sigma\cdot\boldsymbol\psi/2}e_L\) and \(e_R\mapsto e^{-\boldsymbol\sigma\cdot\boldsymbol\psi/2}e_R\), with \(v=\tanh|\boldsymbol\psi|\).}
\GSfield{Statement}{The Pauli matrices generate the two-dimensional representation of \(SU(2)\). The two inequivalent Weyl representations of \(SL_2(\mathbb C)=\operatorname{Spin}^+(1,3)\) restrict to this rotation representation and have opposite boost generators. The signs depend on active/passive and left/right conventions; rapidity has \(|\mathbf v|=c\tanh|\boldsymbol\psi|\). They are representations of the spin cover, not single-valued spinor representations of \(SO^+(1,3)\).}
\GSfield{Reconstruction}{The covering group and units are supplied; the notebook boost convention is retained.}
\GSfield{Fate}{\citet[Section~4.1]{tong2006qft} supplies the standard spinor construction. This is exposition of pre-existing representation theory.}
\GSfield{Present status}{As of 10 November 2026: Established in the stated scope.}
\GSfield{Assessment}{literal statement: correct with missing hypotheses; publication match: exact; verification: read the result; development: none; confidence: high}
\end{GSitem}
\begin{GSitem}{1975-364.3}{E}{Gelfand’s covering-group explanation}
\GSfield{Kind; source}{remark; \GSpage{P000605}; I. M. Gelfand}
\GSfield{The notes}{Gelfand draws \(SO(3)\longrightarrow SU(2)\) and explains the spinor rotation action and its extension to Lorentz transformations.}
\GSfield{Statement}{The canonical surjective covering map is \(SU(2)\to SO(3)\), with kernel \(\{\pm I\}\). A rotation lifts to two elements, so an arrow in the opposite direction cannot be a globally defined inverse homomorphism. Likewise \(SL_2(\mathbb C)\to SO^+(1,3)\) is a double cover.}
\GSfield{Reconstruction}{The arrow is interpreted as a covering-group explanation because of its surrounding text, not silently reversed in the diplomatic transcription.}
\GSfield{Fate}{The covering description is standard; see \citet[Section~4.1]{tong2006qft}.}
\GSfield{Present status}{As of 10 November 2026: Established in the stated scope.}
\GSfield{Assessment}{literal statement: false as written; publication match: exact; verification: read the result; development: annotated; confidence: high}
\end{GSitem}
\begin{GSitem}{1975-364.4}{E}{Parity, the Weyl equations, and the Dirac mass}
\GSfield{Kind; source}{construction; \GSpage{P000605-P000606}}
\GSfield{The notes}{Parity interchanges \(e_L,e_R\); the massless equations imply \(E^2=\mathbf p^2\). A bilinear \(m(e_L^\dagger e_R+e_R^\dagger e_L)\) is added and rewritten as \(\bar e(i\gamma^\mu\partial_\mu-m)e\).}
\GSfield{Statement}{With \(\eta=\operatorname{diag}(1,-1,-1,-1)\) and \(\{\gamma^\mu,\gamma^\nu\}=2\eta^{\mu\nu}\), the Dirac equation implies \((\Box+m^2)e=0\). Plane waves have \(E=\pm\sqrt{\mathbf p^2+m^2}\); the displayed positive root is the positive-energy branch. A parity-invariant charged Dirac field combines the two Weyl modules. Up to normalization, phase choices, and the stated parity and charge symmetries, its real scalar mass couples opposite chiralities.}
\GSfield{Reconstruction}{The factors of \(i\), relative signs, both energy branches, and the symmetries needed for the claimed uniqueness of the mass term are supplied. Unqualified uniqueness would ignore pseudoscalar and, for neutral fields with different symmetries, Majorana masses.}
\GSfield{Fate}{\citet[Sections~4.2--4.4]{tong2006qft} treats the free Dirac construction.}
\GSfield{Present status}{As of 10 November 2026: Established in the stated scope.}
\GSfield{Assessment}{literal statement: correct with missing hypotheses; publication match: partial; verification: read the result; development: none; confidence: high}
\end{GSitem}
\begin{GSitem}{1975-364.5}{E}{Local phase covariance introduces a connection}
\GSfield{Kind; source}{construction; \GSpage{P000606-P000607}}
\GSfield{The notes}{The global phase \(e\mapsto e^{i\beta}e\) is made position dependent; the connection transforms as \(A_\mu\mapsto A_\mu-\partial_\mu\beta\), and \(\nabla_\mu=\partial_\mu+iA_\mu\). The first displayed coupling omits a gamma contraction.}
\GSfield{Statement}{For \(D_\mu=\partial_\mu+iA_\mu\), \(D^{\prime}_\mu(e^{i\beta}e)=e^{i\beta}D_\mu e\). Consequently \(\bar e(i\gamma^\mu D_\mu-m)e\) is invariant under the local \(U(1)\) transformation. Its interaction is \(-\bar e\gamma^\mu A_\mu e\); a separate charge can instead be absorbed into \(A\).}
\GSfield{Reconstruction}{The contraction and placement of parentheses are repaired mathematically; the malformed displayed coupling is present in the scan.}
\GSfield{Fate}{\citet[Section~6.3.1]{tong2006qft} gives this standard gauge-covariant construction.}
\GSfield{Present status}{As of 10 November 2026: Established in the stated scope.}
\GSfield{Assessment}{literal statement: false as written; publication match: exact; verification: read the result; development: annotated; confidence: high}
\end{GSitem}
\begin{GSitem}{1975-364.6}{E}{Curvature and the Maxwell kinetic term}
\GSfield{Kind; source}{construction; \GSpage{P000607}}
\GSfield{The notes}{To give equations for the connection, the notes add \((4e^2)^{-1}F_{\mu\nu}^2\), with \(F_{\mu\nu}=\partial_\mu A_\nu-\partial_\nu A_\mu\).}
\GSfield{Statement}{With the preceding Lorentz signature and \(D=\partial+iA\), a standard real Minkowski Lagrangian is \[ \mathcal L=\bar e(i\gamma^\mu D_\mu-m)e-\frac{1}{4e^2}F_{\mu\nu}F^{\mu\nu}. \] Variation in \(A\) gives Maxwell's equations with the Dirac current. A positive quadratic sign instead belongs to the usual Euclidean action after Wick rotation.}
\GSfield{Reconstruction}{Raising the tensor indices, selecting Lorentz signature consistently, and distinguishing the Euclidean sign are necessary additions; no Wick rotation appears at this point in the notes.}
\GSfield{Fate}{\citet[Section~6.1]{tong2006qft} verifies the conventional Maxwell and QED action.}
\GSfield{Present status}{As of 10 November 2026: Established in the stated scope.}
\GSfield{Assessment}{literal statement: false as written; publication match: partial; verification: read the result; development: none; confidence: high}
\end{GSitem}
\begin{GSitem}{1975-364.7}{E}{The scattering functional integral is formal}
\GSfield{Kind; source}{method; \GSpage{P000607}}
\GSfield{The notes}{A transition amplitude is written \(\int_{\Phi_{\rm in}}^{\Phi_{\rm out}}D\varphi\,e^{iS[\varphi]}\); \(D\varphi\) is called a natural Lebesgue measure. Free theory is Gaussian and interaction gives perturbation theory.}
\GSfield{Statement}{A finite-mode or lattice regulator produces finite-dimensional oscillatory bosonic integrals and Berezin integrals for fermions, with boundary and gauge-fixing choices. Their normalized perturbative expansion defines formal Green functions and, after suitable limiting and scattering prescriptions, perturbative amplitudes. There is no nonzero locally finite translation-invariant Lebesgue measure on an infinite-dimensional topological vector space of fields.}
\GSfield{Reconstruction}{The regulated/formal interpretation is supplied. It is neither a construction of an unregularized four-dimensional QED measure nor a theorem that pointwise field limits as \(t\to\pm\infty\) exist.}
\GSfield{Fate}{\citet[Sections~3.1--3.3, 6.1]{tong2006qft} checks the perturbative scattering framework; it does not verify the notebook's literal measure claim.}
\GSfield{Present status}{As of 10 November 2026: The perturbative interpretation is established; the written infinite-dimensional Lebesgue integral is not a literal measure-theoretic definition.}
\GSfield{Assessment}{literal statement: false as written; publication match: partial; verification: read the result; development: none; confidence: high}
\end{GSitem}
\begin{GSitem}{1975-364.8}{E}{Gelfand’s Gaussian-moment analogy and divergent cubic integral}
\GSfield{Kind; source}{method; \GSpage{P000607}; I. M. Gelfand}
\GSfield{The notes}{Gelfand compares the theory with \(\int e^{i(x^2+\varepsilon x^3)}dx\) and \(\int e^{-\sum a_{ij}x_ix_j+\varepsilon\sum x_i^3}dx\). The last expansion instead displays \(e^{ix^2(1+\varepsilon x^3)}\).}
\GSfield{Statement}{For a positive definite real quadratic form \(Q\), every coefficient \(\int P(x)^k e^{-Q(x)}dx/k!\) exists for a polynomial \(P\); these coefficients define a formal perturbation series. If \(P=x^3\) and real \(\varepsilon\ne0\), the displayed real-line integral diverges on one tail. Oscillatory analogues require a contour or distributional prescription. The final written phase is quintic, not cubic.}
\GSfield{Reconstruction}{The rigorous content retained is the coefficientwise Gaussian-moment calculation, not convergence of the formal series or the cubic exponential integral.}
\GSfield{Fate}{The Wick-contraction method is standard; \citet[Section~3.3]{tong2006qft} gives the corresponding perturbative mechanism. The convergence objection follows directly from the odd leading term in the notebook's real exponent.}
\GSfield{Present status}{As of 10 November 2026: Established in the stated scope.}
\GSfield{Assessment}{literal statement: false as written; publication match: partial; verification: read the result; development: annotated; confidence: high}
\end{GSitem}
\begin{GSitem}{1975-364.9}{E}{A logarithmically divergent model loop integral}
\GSfield{Kind; source}{claim; \GSpage{P000607}}
\GSfield{The notes}{The sample integral is \(\int_{\mathbb R^4}d^4x/((|x|^2+m_1)(|x|^2+m_2))\), without a regulator.}
\GSfield{Statement}{For \(m_1,m_2>0\), the integrand is locally integrable, but its radial integral at infinity behaves like \(\int^\infty r^{-1}dr\); hence the unregularized integral diverges logarithmically. A cutoff makes it a legitimate starting point for parameterization and subtraction.}
\GSfield{Reconstruction}{Positivity and the ultraviolet cutoff are supplied. The radial comparison is an elementary check of this specific notebook example, not a claim about all loop integrals.}
\GSfield{Fate}{The general perturbative setting is reviewed in \citet[Section~3.2]{tong2006qft}; the explicit divergence is determined directly from the written integrand.}
\GSfield{Present status}{As of 10 November 2026: Established in the stated scope.}
\GSfield{Assessment}{literal statement: correct with missing hypotheses; publication match: thematic only; verification: read the result; development: annotated; confidence: high}
\end{GSitem}
\begin{GSitem}{1975-364.10}{E}{Feynman parameters and the missing beta normalization}
\GSfield{Kind; source}{method; \GSpage{P000607-P000608}}
\GSfield{The notes}{The identity \(1/(ab)=\int_0^1(sa+(1-s)b)^{-2}ds\) is extended to arbitrary exponents and to a simplex without gamma factors.}
\GSfield{Statement}{For \(a_i>0\), \(\operatorname{Re}\lambda_i>0\), and \(\Lambda=\sum_i\lambda_i\), \[ \prod_i a_i^{-\lambda_i}=\frac{\Gamma(\Lambda)}{\prod_i\Gamma(\lambda_i)}\int_{s_i\geq0}\frac{\delta(1-\sum_i s_i)\prod_i s_i^{\lambda_i-1}\,ds_i}{(\sum_i s_i a_i)^\Lambda}. \] The two-denominator formula is its special case. Continuation to Lorentzian denominators additionally requires consistent branches and the pole prescription.}
\GSfield{Reconstruction}{The gamma ratio and domain are restored. At \(a=b=1\), \(\lambda=\mu=2\), the notebook's right-hand side is \(1/6\), while the left-hand side is \(1\); the missing factor is \(6\).}
\GSfield{Fate}{\citet[Section~5.12, equations~5.12.1 and~5.12.4]{nist2026beta} verifies the beta-integral normalization. The simplex identity follows by introducing positive Schwinger parameters and separating their sum.}
\GSfield{Present status}{As of 10 November 2026: Established in the stated scope.}
\GSfield{Assessment}{literal statement: false as written; publication match: exact; verification: read the result; development: annotated; confidence: high}
\end{GSitem}
\begin{GSitem}{1975-364.11}{E}{Dirac Green function and the Feynman pole prescription}
\GSfield{Kind; source}{construction; \GSpage{P000608}}
\GSfield{The notes}{Although the prose says ``according to Maxwell'', the equation is \((i\gamma^\mu\partial_\mu-m)K=\delta\), and the Fourier inverse is \((\gamma^\mu p_\mu+m)/(p^2-m^2+i0)\).}
\GSfield{Statement}{For the convention \(K_F(x)=\int d^4p\,(2\pi)^{-4}e^{-ipx}(\gamma^\mu p_\mu+m)/(p^2-m^2+i0)\), \((i\gamma^\mu\partial_\mu-m)K_F=\delta\) distributionally. The usual time-ordered two-point function is \(iK_F\) in this convention. The two displaced poles encode positive-frequency propagation forward and negative-frequency propagation backward in time.}
\GSfield{Reconstruction}{The Fourier sign, normalization, matrix identity on the delta function, and distinction between Green function and two-point function are supplied. ``Maxwell'' here must not be used to identify the displayed electron kernel as a photon propagator.}
\GSfield{Fate}{\citet[Section~5.5]{tong2006qft} gives the free Dirac propagator.}
\GSfield{Present status}{As of 10 November 2026: Established in the stated scope.}
\GSfield{Assessment}{literal statement: correct with missing hypotheses; publication match: partial; verification: read the result; development: none; confidence: high}
\end{GSitem}
\begin{GSitem}{1975-364.12}{E}{Single-photon exchange and the diagram expansion}
\GSfield{Kind; source}{method; \GSpage{P000608}}
\GSfield{The notes}{The diagram joins two electron lines by a photon. Each emission or absorption carries \(e\), but the displayed integral is a product of four electron kernels and one photon kernel without explicit vertices or \(e^2\). More diagrams contain loops and pair creation.}
\GSfield{Statement}{After a fixed gauge and propagator convention, the order-\(e^2\) electron-exchange kernel contains \[ e^2\!\int\! [S_F(x_3-z_1)\gamma^\mu S_F(z_1-x_1)]\otimes[S_F(x_4-z_2)\gamma^\nu S_F(z_2-x_2)]D_{\mu\nu}(z_1-z_2)\,d^4z_1d^4z_2, \] up to the conventional factors of \(i\), external-state contractions, and exchange antisymmetrization for identical fermions. Higher terms are organized by Wick contractions with loop signs and symmetry factors.}
\GSfield{Reconstruction}{The displayed formula is reconstructed as a schematic kernel, not asserted to be a fully normalized scattering amplitude. Unclear connections in the hand-drawn third and fourth diagrams are not assigned a unique graph.}
\GSfield{Fate}{\citet[Section~6.5]{tong2006qft} supplies the QED vertex and diagram rules.}
\GSfield{Present status}{As of 10 November 2026: Established in the stated scope.}
\GSfield{Assessment}{literal statement: correct with missing hypotheses; publication match: partial; verification: read the result; development: none; confidence: high}
\end{GSitem}
\begin{GSitem}{1975-364.13}{E}{Antiparticles do not make the Feynman inverse the only Green function}
\GSfield{Kind; source}{remark; \GSpage{P000608}}
\GSfield{The notes}{The final lines link negative-energy levels, backward time propagation, antiparticles, and an alleged unique consistent way to bypass the poles; the terminal Russian words remain doubtful.}
\GSfield{Statement}{The time-ordered vacuum propagator has the Feynman \(i0\) prescription. Retarded and advanced Green functions are also well-defined distributional inverses with different boundary conditions. Antiparticles arise in the quantized complex Dirac field; uniqueness is meaningful only after the state, ordering, and boundary conditions have been specified.}
\GSfield{Reconstruction}{This separates the unambiguous pole discussion from the unreadable terminal assertion. No claim that the seminar denied antiparticles is extracted from the doubtful English phrase.}
\GSfield{Fate}{\citet[Sections~5.2--5.5]{tong2006qft} gives the particle/antiparticle interpretation and the time-ordered propagator.}
\GSfield{Present status}{As of 10 November 2026: Established in the stated scope.}
\GSfield{Assessment}{literal statement: not determinable; publication match: partial; verification: read the result; development: none; confidence: medium}
\end{GSitem}
\begin{GSitem}{1975-364.14}{E}{Charge screening and the opposite nonabelian sign}
\GSfield{Kind; source}{claim; \GSpage{P000609}}
\GSfield{The notes}{The leading exchange is of order \(e^2\), the polarization correction of order \(e^4\); the notes write \(e_{\rm phys}^2<e^2\) and the reverse inequality for the strong interaction.}
\GSfield{Statement}{For a cutoff description of QED, vacuum polarization screens electric charge at long distances. In a mass-independent perturbative convention with one unit-charge Dirac fermion, \(\beta(e)=e^3/(12\pi^2)+O(e^5)\). For \(SU(N_c)\) with \(N_f\) fundamental Dirac flavours, \(\beta(g)=-g^3(11N_c-2N_f)/(48\pi^2)+O(g^5)\), so the coupling decreases toward high energy when \(N_f<11N_c/2\).}
\GSfield{Reconstruction}{The comparison is reconstructed as the sign of scale dependence, rather than a scheme-independent inequality between undefined bare and physical numbers. The negative beta function by itself does not prove confinement.}
\GSfield{Fate}{\citet[p.~1344, equations~(8)--(9)]{grosswilczek1973ultraviolet} gives the nonabelian and fermion contributions in June 1973, before this lecture; the comparison with the Abelian sign is made on the same page.}
\GSfield{Present status}{As of 10 November 2026: Established in the stated scope.}
\GSfield{Assessment}{literal statement: correct with missing hypotheses; publication match: partial; verification: indirect; development: none; confidence: high}
\end{GSitem}

\GStalk{365}{1975-09-22}{I. M. Gelfand}{P000609-P000611}{Essentially normal operators and extension classes}
Gelfand reports the Brown--Douglas--Fillmore classification and its connection with odd \(K\)-homology, under a heading mentioning Atiyah. Essential normality and the equivalence relation defining \(\operatorname{Ext}\) are omitted in the shorthand. The intervening line ``Moser---Korteweg--de Vries'' contains no mathematical assertion and receives no separate item.

\begin{GSitem}{1975-365.1}{E}{Classification modulo compact operators}
\GSfield{Kind; source}{question; \GSpage{P000610}}
\GSfield{The notes}{For the Calkin algebra \(\mathfrak B=B(H)/K(H)\), when are \(B_1,B_2\) unitarily equivalent? The answer given is equality of essential spectra and of the indices on every component of their complement.}
\GSfield{Statement}{Let \(T_1,T_2\) be bounded essentially normal operators on a separable infinite-dimensional complex Hilbert space. They are unitarily equivalent modulo compacts, meaning \(UT_1U^*-T_2\) is compact for a unitary \(U\) on \(H\), precisely when \(\sigma_{\rm ess}(T_1)=\sigma_{\rm ess}(T_2)=X\) and \(\operatorname{ind}(T_1-\lambda)=\operatorname{ind}(T_2-\lambda)\) for every \(\lambda\notin X\). Essential normality means \([T_i,T_i^*]\) is compact. The index is constant on each complementary component.}
\GSfield{Reconstruction}{Essential normality and the meaning of lifted unitary equivalence are indispensable additions. The theorem does not classify arbitrary nonnormal elements of the Calkin algebra by these data.}
\GSfield{Fate}{The cited 1973 chapter, \citet[pp.~58--128]{browndouglasfillmore1973unitary}, predates the seminar; its publisher metadata was opened, not its full proof. The claim is identified as the BDF theorem with the omitted hypothesis restored, rather than as a new Atiyah theorem.}
\GSfield{Present status}{As of 10 November 2026: Established in the stated scope.}
\GSfield{Assessment}{literal statement: correct with missing hypotheses; publication match: partial; verification: publisher metadata checked; development: none; confidence: high; problem \hyperref[prob:GP-1975-01]{GP-1975-01}}
\end{GSitem}
\begin{GSitem}{1975-365.2}{E}{Extensions as odd analytic K-homology}
\GSfield{Kind; source}{theorem; \GSpage{P000610}}
\GSfield{The notes}{The notes write \(\operatorname{Ext}(X)\cong K_1(X)\) and explain \(\operatorname{Ext}(X)\) as embeddings of \(C(X)\) in the Calkin algebra.}
\GSfield{Statement}{For compact metrizable \(X\), the stable unitary-equivalence classes of extensions \(0\to K(H)\to E\to C(X)\to0\), with split extensions treated as zero, form the BDF group \(\operatorname{Ext}(C(X))\). In modern notation it is naturally the odd analytic \(K\)-homology group \(KK^1(C(X),\mathbb C)\). This is a group of extension classes, not the unquotiented set of embeddings and not generally the group \(K_1(C(X))\).}
\GSfield{Reconstruction}{Metrizability, stabilization, equivalence, and the distinction between homology of the space and \(K\)-theory of its function algebra are supplied. The notation \(K_1(X)\) is retained only when explicitly understood homologically.}
\GSfield{Fate}{The 1973 precursor is \citet{browndouglasfillmore1973unitary}; the full paper is \citet[pp.~265--324]{browndouglasfillmore1977extensions}. Both publisher records were checked. The modern \(KK\)-notation is a reconstruction, not claimed to occur in the seminar.}
\GSfield{Present status}{As of 10 November 2026: Established in the stated scope.}
\GSfield{Assessment}{literal statement: correct with missing hypotheses; publication match: partial; verification: publisher metadata checked; development: none; confidence: high}
\end{GSitem}

\GStalk{366}{1975-09-22}{B. Kostant}{P000609-P000611}{Euler products, root systems, and Kostant’s character formula}
The Euler product and its familiar powers introduce the Macdonald identities and Kostant's character-theoretic form. Manin anticipates his next discussion of several expressions for \(\tau(n)\), and Gelfand asks only the suggestive question \(\psi(x)^z\). The spectral-zeta fragment is too compressed to identify a unique formula.

\begin{GSitem}{1975-366.1}{E}{Euler’s product, partitions, and the pentagonal expansion}
\GSfield{Kind; source}{theorem; \GSpage{P000610}}
\GSfield{The notes}{For \(\psi(q)=\prod_{n\geq1}(1-q^n)\), the reciprocal generates partitions. The exponent in the displayed pentagonal expansion is \((n^2-3n)/2\).}
\GSfield{Statement}{For \(|q|<1\), or as formal power series, \[ \psi(q)^{-1}=\sum_{n\geq0}p(n)q^n,\quad p(0)=1,\qquad \psi(q)=\sum_{k\in\mathbb Z}(-1)^kq^{k(3k-1)/2}. \] Each coefficient is finite in the formal interpretation.}
\GSfield{Reconstruction}{The correct pentagonal exponent and the bilateral index set are supplied. The notebook exponent would even give negative powers at small positive indices.}
\GSfield{Fate}{These classical formulas are recorded in \citet[Section~27.14(ii), equations~27.14.3--5]{nist2026dlmf}; no 1975 novelty is suggested.}
\GSfield{Present status}{As of 10 November 2026: Established in the stated scope.}
\GSfield{Assessment}{literal statement: false as written; publication match: exact; verification: read the result; development: annotated; confidence: high}
\end{GSitem}
\begin{GSitem}{1975-366.2}{E}{Jacobi’s cubic Euler-product identity}
\GSfield{Kind; source}{theorem; \GSpage{P000610}}
\GSfield{The notes}{The cube is \(\psi(q)^3=\sum_{n\geq0}(-1)^n(2n+1)q^{n(n+1)/2}\), with Jacobi named.}
\GSfield{Statement}{The displayed identity holds for \(|q|<1\) and coefficientwise in \(\mathbb Z[[q]]\). It is the rank-one member of the eta-product/root-system family, after the appropriate leading power of \(q\) is separated.}
\GSfield{Reconstruction}{Only the convergence/formal interpretation and the distinction between \(\eta\) and \(\psi\) are supplied.}
\GSfield{Fate}{\citet[Section~2, pp.~636--638]{dyson1972missed} explicitly identifies Jacobi's power \(3\) in the established sequence. This revisits the eta-power discussion in 1971-269.3.}
\GSfield{Present status}{As of 10 November 2026: Established in the stated scope.}
\GSfield{Assessment}{literal statement: correct; publication match: exact; verification: indirect; development: none; confidence: high}
\end{GSitem}
\begin{GSitem}{1975-366.3}{E}{The discriminant and the exponent 24}
\GSfield{Kind; source}{definition; \GSpage{P000610}}
\GSfield{The notes}{The notation is \(\eta(q)=q^{1/24}\psi(q)\), with \(\eta(q)^{24}=\sum\tau(n)q^n\); a similar formula for power \(10\) is mentioned.}
\GSfield{Statement}{For \(z\) in the upper half-plane and \(q=e^{2\pi iz}\), \(\eta(z)=q^{1/24}\prod_{n\geq1}(1-q^n)\) and \(\Delta(z)=\eta(z)^{24}=\sum_{n\geq1}\tau(n)q^n\). Thus \(\tau(1)=1\). Powers \(10\) and \(24\) have explicit lattice-sum identities, but the notes state no particular power-\(10\) summation.}
\GSfield{Reconstruction}{The variable \(z\), branch convention \(q^{1/24}=e^{2\pi iz/24}\), and lower index \(1\) are explicit here. The phrase about power \(10\) is catalogued without inventing its omitted formula.}
\GSfield{Fate}{\citet[equations~(1)--(3), pp.~636--637]{dyson1972missed} gives the discriminant and explains Winquist's power \(10\). Earlier occurrences are 1971-269.1--3.}
\GSfield{Present status}{As of 10 November 2026: Established in the stated scope.}
\GSfield{Assessment}{literal statement: correct; publication match: exact; verification: read the result; development: none; confidence: high}
\end{GSitem}
\begin{GSitem}{1975-366.4}{A}{Macdonald identities in Kostant’s character form}
\GSfield{Kind; source}{theorem; \GSpage{P000610}}
\GSfield{The notes}{The power is \(\dim K\), and the sum has terms \(\chi_\lambda(a)\dim V_\lambda\) with \(\chi_\lambda(a)\in\{0,\pm1\}\); the exponent is left unspecified.}
\GSfield{Statement}{For a compact connected simply connected simple Lie group \(K\), let \(\rho\) be half the sum of positive roots and use the Killing-form normalization to identify weights with the Cartan subalgebra. For Kostant's principal element \(a=\exp(2\pi i\,2\rho)\), \[ \prod_{n\geq1}(1-q^n)^{\dim K}=\sum_{\lambda\in P_+}\chi_\lambda(a)\dim V_\lambda\,q^{(\lambda+\rho,\lambda+\rho)-(\rho,\rho)},\qquad \chi_\lambda(a)\in\{0,1,-1\}. \] The inner product here is the one in the cited theorem; arbitrary root-length normalization changes the exponents.}
\GSfield{Reconstruction}{Simplicity, simple connectedness, the principal element, and the Casimir exponent replace the omitted hypotheses and dots. ``Dimensions of compact groups'' alone is too broad a description of the precise theorem.}
\GSfield{Fate}{Macdonald's underlying identities appeared in 1972 \citep{macdonald1972affine}. \citet[Theorems~0.1 and~2.1]{kostant2003euler} restates the character formula and explicitly identifies its original 1976 Theorem~3.1. This supports class A for Kostant's particular formulation, not a new priority claim for the older product identities.}
\GSfield{Present status}{As of 10 November 2026: Established in the stated scope.}
\GSfield{Assessment}{literal statement: correct with missing hypotheses; publication match: partial; verification: read the result; development: none; confidence: high}
\end{GSitem}
\begin{GSitem}{1975-366.5}{F}{The spectral-zeta fragment}
\GSfield{Kind; source}{remark; \GSpage{P000610}}
\GSfield{The notes}{After the character identity, the page records ``Kostant: \(\operatorname{Sp}(-\Delta+a)^{-s}\), etc.''}
\GSfield{Statement}{A spectral zeta function of a positive elliptic operator can be defined for sufficiently large \(\operatorname{Re}s\), and a group element can enter an equivariant trace. However, the notes do not specify the operator domain, the meaning of \(a\), whether it is added or acts by translation, or any asserted identity.}
\GSfield{Reconstruction}{An equivariant heat/zeta-trace connection is a possible interpretation only. The group element \(a\) from the preceding line cannot simply be treated as a scalar spectral shift without additional information.}
\GSfield{Fate}{The heat-kernel connection in \citet[Introduction]{kostant2003euler} provides thematic context, but does not verify this literal expression.}
\GSfield{Present status}{As of 10 November 2026: The intended spectral formula is not identifiable with sufficient precision.}
\GSfield{Assessment}{literal statement: not determinable; publication match: thematic only; verification: unresolved lead; development: none; confidence: low}
\end{GSitem}
\begin{GSitem}{1975-366.6}{F}{Gelfand’s complex-exponent suggestion}
\GSfield{Kind; source}{question; \GSpage{P000610}; I. M. Gelfand}
\GSfield{The notes}{Gelfand writes \(\psi(q)^z\), followed by a question mark, with \(z\in\mathbb C\).}
\GSfield{Statement}{There is a canonical analytic definition on \(|q|<1\):
     \[\psi(q)^z=\exp\!\left(z\sum_{n\geq1}\log(1-q^n)\right),\]
     using the branch equal to \(1\) at \(q=0\). Its formal coefficients are polynomials in \(z\). This observation does not determine whether Gelfand wanted analytic continuation, coefficient formulas, representation-theoretic interpolation, or another property.}
\GSfield{Reconstruction}{The elementary definition is supplied only to show that existence of a complex power is not itself a uniquely recoverable open problem. It is not equated with the earlier definite question about exponent \(26\), GP-1971-41.}
\GSfield{Fate}{\citet[Introduction]{kostant2003euler} studies coefficients of general powers, but the notebook specifies no assertion that can be matched to one of its theorems.}
\GSfield{Present status}{As of 10 November 2026: The intended question remains underdetermined; no new problem-index entry is manufactured.}
\GSfield{Assessment}{literal statement: not determinable; publication match: thematic only; verification: unresolved lead; development: none; confidence: low}
\end{GSitem}

\GStalk{367}{1975-09-22}{A. M. Polyakov}{P000609-P000614}{From lepton symmetries to electroweak gauge fields}
After reading recommendations, Polyakov extends the QED construction to chiral lepton doublets and local \(SU(2)\) symmetry. Charged currents, muon decay, neutral-field mixing, and the then four-flavour, three-colour quark count appear. The Higgs sector and hypercharge assignment are not supplied in these pages; a complete electroweak reconstruction must mark those additions.

\begin{GSitem}{1975-367.1}{E}{The enlarged symmetry of massless lepton kinetic terms}
\GSfield{Kind; source}{claim; \GSpage{P000611}}
\GSfield{The notes}{The kinetic terms for \(e_L,\nu_L,e_R\) have \(SU_L(2)\times U_R(1)\times U_L(1)\) symmetry; another lepton family is indicated. A preceding line still includes the electron mass term.}
\GSfield{Statement}{For a massless left Weyl doublet \(L=(\nu_L,e_L)^T\) and a right Weyl singlet \(e_R\), the free kinetic Lagrangian is invariant under \(U(2)_L\times U(1)_R\). At the Lie-algebra level this is \(\mathfrak{su}(2)_L\oplus\mathfrak u(1)_L\oplus\mathfrak u(1)_R\); the product-group notation suppresses a finite central quotient. A nonzero bare Dirac mass selecting the electron component breaks the independent left-doublet symmetry.}
\GSfield{Reconstruction}{The massless restriction and global central quotient are supplied. The doublet order is changed explicitly to the usual convention for later charge assignments.}
\GSfield{Fate}{\citet[p.~1264, equations~(1)--(4)]{weinberg1967leptons} starts from these kinetic symmetries and selects the gauged combination. It was published 20 November 1967.}
\GSfield{Present status}{As of 10 November 2026: Established in the stated scope.}
\GSfield{Assessment}{literal statement: correct with missing hypotheses; publication match: partial; verification: read the result; development: none; confidence: high}
\end{GSitem}
\begin{GSitem}{1975-367.2}{E}{Local SU(2) covariance and the three gauge fields}
\GSfield{Kind; source}{construction; \GSpage{P000611}}
\GSfield{The notes}{The doublet transforms by \(U(x)=\exp(i\boldsymbol\tau\cdot\boldsymbol\alpha(x)/2)\), while the infinitesimal line has the opposite sign. A three-component connection is introduced, yielding three massless particles at this stage.}
\GSfield{Statement}{Set \(D_\mu=\partial_\mu+igW_\mu\), \(W_\mu=W_\mu^a\tau_a/2\). Requiring \(D'_\mu(UL)=UD_\mu L\) gives \(W'_\mu=UW_\mu U^{-1}+(i/g)(\partial_\mu U)U^{-1}\). The curvature is \(F_{\mu\nu}=\partial_\mu W_\nu-\partial_\nu W_\mu+ig[W_\mu,W_\nu]\). The gauge-invariant Yang--Mills kinetic term describes three massless gauge fields before symmetry breaking.}
\GSfield{Reconstruction}{The coupling, complete connection transformation, and curvature are supplied. The infinitesimal sign is made consistent with the finite exponential; no mass-generation mechanism is inferred from this line.}
\GSfield{Fate}{The nonabelian construction is present in the electroweak Lagrangian of \citet[equation~(4)]{weinberg1967leptons}. The notebook is expository.}
\GSfield{Present status}{As of 10 November 2026: Established in the stated scope.}
\GSfield{Assessment}{literal statement: false as written; publication match: partial; verification: read the result; development: none; confidence: high}
\end{GSitem}
\begin{GSitem}{1975-367.3}{E}{Charged weak currents and muon decay}
\GSfield{Kind; source}{construction; \GSpage{P000612}}
\GSfield{The notes}{The off-diagonal fields \(A^1\pm iA^2\) change electron-type leptons into their neutrinos; a diagram represents muon decay.}
\GSfield{Statement}{With canonically normalized \(W^\pm=(W^1\mp iW^2)/\sqrt2\), the charged-current term is proportional to \(g\bar\nu_{eL}\gamma^\mu e_LW^+_\mu/\sqrt2\), plus its adjoint and the muon-family analogue. Exchange of \(W^-\) gives \(\mu^-\to e^-+\bar\nu_e+\nu_\mu\). At energies much below \(m_W\), the propagator reduces to the local four-fermion interaction with \(G_F/\sqrt2=g^2/(8m_W^2)\).}
\GSfield{Reconstruction}{The \(\sqrt2\) normalization, conjugate neutrino, chirality, and low-energy hypothesis are supplied. The low-energy coefficient is a reconstruction, not an extra formula recovered from the sketch.}
\GSfield{Fate}{\citet[equations~(8), (13), and~(16)]{weinberg1967leptons} gives the charged fields and their weak-coupling normalization, predating the lecture.}
\GSfield{Present status}{As of 10 November 2026: Established in the stated scope.}
\GSfield{Assessment}{literal statement: correct with missing hypotheses; publication match: partial; verification: read the result; development: none; confidence: high}
\end{GSitem}
\begin{GSitem}{1975-367.4}{E}{Neutral mixing requires a hypercharge assignment}
\GSfield{Kind; source}{claim; \GSpage{P000612}}
\GSfield{The notes}{A right-handed \(U(1)\) field \(B\) and \(A_\gamma=A^3\cos\theta+B\sin\theta\) are written, without a complete left-handed \(B\)-coupling or symmetry-breaking sector.}
\GSfield{Statement}{In the standard completion, \(SU(2)_L\times U(1)_Y\) has \(Y_L=-1/2\), \(Y_{e_R}=-1\), \(Q=T_3+Y\), and a scalar doublet of \(Y=1/2\) with a neutral nonzero vacuum value. Then \(A_\gamma=W^3\sin\theta_W+B\cos\theta_W\), \(\tan\theta_W=g'/g\), and \(e=g\sin\theta_W=g'\cos\theta_W\). The orthogonal neutral field is massive. A \(U(1)\) acting only on \(e_R\) cannot by itself reproduce these charges.}
\GSfield{Reconstruction}{Hypercharges, the scalar doublet, and its vacuum are explicitly added. The reversed sine and cosine in the notebook can be an angle convention; it is not alone an error.}
\GSfield{Fate}{\citet[equations~(4), (7), and~(10)--(15)]{weinberg1967leptons} verifies this completion. It is only a partial match to the incomplete notebook account.}
\GSfield{Present status}{As of 10 November 2026: Established in the stated scope.}
\GSfield{Assessment}{literal statement: correct with missing hypotheses; publication match: partial; verification: read the result; development: none; confidence: high}
\end{GSitem}
\begin{GSitem}{1975-367.5}{E}{Four flavours and three colours in the quark count}
\GSfield{Kind; source}{remark; \GSpage{P000612}}
\GSfield{The notes}{The last line counts twelve quarks as four flavours times three colours.}
\GSfield{Statement}{In the four-flavour model under discussion, the internal one-particle labels are \(\{u,d,s,c\}\times\{1,2,3\}\), giving twelve flavour--colour labels. This does not mean twelve flavours, and it does not count spin or antiquark labels.}
\GSfield{Reconstruction}{The interpretation of the multiplication is made explicit; the notebook is not recast as a claim that four flavours exhaust all future observations.}
\GSfield{Fate}{This is bookkeeping within the lecture's historical model, with no independent theorem or definite research question to trace.}
\GSfield{Present status}{As of 10 November 2026: Established in the stated scope.}
\GSfield{Assessment}{literal statement: correct; publication match: none found; verification: indirect; development: none; confidence: high}
\end{GSitem}

\GStalk{368}{1975-09-29}{I. M. Gelfand}{P000612-P000614}{The Lie algebra of differential polynomials}
Gelfand introduces the bracket on differential polynomials underlying evolutionary equations and the KdV hierarchy. Povzner identifies its action by derivations. On the same page Manin's opening definition of \(\tau\) precedes the Gelfand passage and his numbered formulas follow it; the catalogue's two Manin records are therefore kept distinct.

\begin{GSitem}{1975-368.1}{A}{Differential polynomials and total differentiation}
\GSfield{Kind; source}{definition; \GSpage{P000613}}
\GSfield{The notes}{The algebra consists of polynomials in \(u,u',u'',\ldots\), and differentiation is \(D=\sum u^{(k+1)}\partial/\partial u^{(k)}\).}
\GSfield{Statement}{Let \(\mathcal A=\mathbb R[u_0,u_1,\ldots]\), each element depending on finitely many variables, and \(D=\sum_{j\geq0}u_{j+1}\partial_{u_j}\). Evaluation at a smooth function, \(u_j=\partial_x^j u(x)\), intertwines \(D\) with \(\partial_x\). Each formal sum acts finitely on every element of \(\mathcal A\).}
\GSfield{Reconstruction}{Finite dependence and the evaluation interpretation are explicit additions. The variables are algebraically independent before evaluation.}
\GSfield{Fate}{The framework is developed in \citet[Sections~8--10]{gelfanddikii1976fractional}, received 22 June 1976 and published in 1976. The source's differential-polynomial construction was read.}
\GSfield{Present status}{As of 10 November 2026: Established in the stated scope.}
\GSfield{Assessment}{literal statement: correct; publication match: exact; verification: read the result; development: none; confidence: high}
\end{GSitem}
\begin{GSitem}{1975-368.2}{A}{The evolutionary Lie bracket}
\GSfield{Kind; source}{construction; \GSpage{P000613}}
\GSfield{The notes}{For two differential polynomials the notes write
     \[[P,Q]=\sum_{k\geq0}\left(D^kP\,\frac{\partial Q}{\partial u_k}-D^kQ\,\frac{\partial P}{\partial u_k}\right).\]}
\GSfield{Statement}{For \(P\in\mathcal A\), define the evolutionary derivation \(\mathcal D_P=\sum_{k\geq0}D^k(P)\partial_{u_k}\). Then \([P,Q]=\mathcal D_P(Q)-\mathcal D_Q(P)\) is a Lie bracket, and \([\mathcal D_P,\mathcal D_Q]=\mathcal D_{[P,Q]}\). The injection \(P\mapsto\mathcal D_P\) follows from \(\mathcal D_P(u_0)=P\); Jacobi follows from the derivation commutator.}
\GSfield{Reconstruction}{\ours\ The short derivation argument makes the Jacobi claim precise and fixes the sign convention, without introducing a Poisson bracket on the original jet variables.}
\GSfield{Fate}{\citet[Section~10, equation~(24)]{gelfanddikii1976fractional} gives the same evolutionary bracket in the subsequent joint paper.}
\GSfield{Present status}{As of 10 November 2026: Established in the stated scope.}
\GSfield{Assessment}{literal statement: correct; publication match: exact; verification: read the result; development: none; confidence: high}
\end{GSitem}
\begin{GSitem}{1975-368.3}{A}{Povzner’s identification with compatible evolution equations}
\GSfield{Kind; source}{remark; \GSpage{P000613}; A. Ya. Povzner}
\GSfield{The notes}{Povzner associates \(a\) with \(u_t=a\), explains its derivation on all derivatives of \(u\), and identifies the bracket of these derivations.}
\GSfield{Statement}{For autonomous differential-polynomial flows \(u_t=P[u]\) and \(u_s=Q[u]\), compatibility of the formal prolonged vector fields is \([P,Q]=0\). Indeed \(\partial_tQ[u]-\partial_sP[u]=\mathcal D_PQ-\mathcal D_QP\). This is a formal jet-space statement; it alone does not establish existence, uniqueness, or commuting solution maps for a chosen PDE function space.}
\GSfield{Reconstruction}{The distinction between algebraic compatibility and analytic well-posedness is supplied. The observation belongs to Povzner's intervention, not to a separately catalogued talk.}
\GSfield{Fate}{The derivation interpretation is incorporated in \citet[Section~10]{gelfanddikii1976fractional}.}
\GSfield{Present status}{As of 10 November 2026: Established in the stated scope.}
\GSfield{Assessment}{literal statement: correct; publication match: exact; verification: read the result; development: none; confidence: high}
\end{GSitem}

\GStalk{369}{1975-09-29}{Yu. I. Manin}{P000612-P000614}{Opening definition of Ramanujan’s tau function}
This catalogue record is represented by the short Manin heading and product definition at the top of P000613. Gelfand's differential-polynomial discussion follows immediately; Manin resumes below it in record 370. The mathematical content is therefore one definition, despite the archive's broad overlapping page range.

\begin{GSitem}{1975-369.1}{E}{Tau as the coefficients of the discriminant}
\GSfield{Kind; source}{definition; \GSpage{P000613}}
\GSfield{The notes}{Manin opens with \(x\prod_{n\geq1}(1-x^n)^{24}=\sum_{n\geq1}\tau(n)x^n\).}
\GSfield{Statement}{The identity defines integers \(\tau(n)\), \(n\geq1\), as the coefficients of \(\Delta(q)=q\prod_{m\geq1}(1-q^m)^{24}\), with \(\tau(1)=1\). It is valid formally and analytically for \(|q|<1\).}
\GSfield{Reconstruction}{Only the formal/analytic interpretation is supplied. This occurrence repeats 1975-366.3 and 1971-269.1, rather than introducing a second tau function.}
\GSfield{Fate}{\citet[equation~(1), p.~636]{dyson1972missed} gives the same definition before the seminar.}
\GSfield{Present status}{As of 10 November 2026: Established in the stated scope.}
\GSfield{Assessment}{literal statement: correct; publication match: exact; verification: read the result; development: none; confidence: high}
\end{GSitem}

\GStalk{370}{1975-09-29}{Yu. I. Manin}{P000612-P000617}{Eight descriptions of Ramanujan’s tau function}
Manin resumes with a list announced as six expressions but numbered through eight: the product, Dyson's root-lattice sum, divisor convolutions, lattice theta series, a period formula, partition asymptotics, Hecke multiplicativity, and a Galois representation. The subsequent explanation invokes eta transformations, the ring of modular forms, Rademacher contours, and Kostant's identity. Several formulas are faithfully transcribed mathematical errors; their corrected forms and the exact scope of independent finite checks are separated from the notebook.

\begin{GSitem}{1975-370.1}{F}{The unexplained nonalgorithmic qualification}
\GSfield{Kind; source}{remark; \GSpage{P000613}}
\GSfield{The notes}{The heading calls the expressions ``not obtained by a finite algorithm'', while the list itself begins with a finite coefficient-extraction procedure.}
\GSfield{Statement}{For any given positive integer \(n\), \(\tau(n)\) is computable by truncating \(q\prod_{m=1}^{n-1}(1-q^m)^{24}\) modulo \(q^{n+1}\). Therefore a literal assertion that no finite algorithm computes \(\tau(n)\) is false. The phrase may instead concern deriving one description from another or the conceptual diversity of the formulas; no complexity model or quantified claim is stated.}
\GSfield{Reconstruction}{The finite coefficient algorithm is an elementary consequence of the supplied definition. The possible alternative reading is explicitly conjectural and is not indexed as a research problem.}
\GSfield{Fate}{No exact noncomputability or complexity assertion was identified. The surrounding explicit coefficient formulas in \citet[Section~1.3]{manin1973periods} reinforce the need to retain this textual uncertainty.}
\GSfield{Present status}{As of 10 November 2026: The qualification is not a recoverable mathematical conjecture.}
\GSfield{Assessment}{literal statement: not determinable; publication match: none found; verification: unresolved lead; development: none; confidence: low}
\end{GSitem}
\begin{GSitem}{1975-370.2}{E}{The combinatorial expansion needs multiplicities}
\GSfield{Kind; source}{construction; \GSpage{P000613}}
\GSfield{The notes}{The first expression sums \((-1)^k\) over ordered-by-size partitions of \(n-1\), with at most 24 consecutive equality signs. No colours or binomial coefficients are shown.}
\GSfield{Statement}{The literal product expansion is \[ \tau(n)=\sum_{\substack{0\leq r_j\leq24\ (j\geq1)\\\sum j r_j=n-1}}(-1)^{\sum r_j}\prod_j\binom{24}{r_j}. \] Only \(j\leq n-1\) contribute. Equivalently, each selected part size comes with a subset of 24 labelled copies; it can occur at most 24 times, hence at most 23 equality signs within a run.}
\GSfield{Reconstruction}{The binomial multiplicities are mandatory. At \(n=2\), an unweighted partition count gives \(-1\), while the product gives \(-24\). This counterexample is independent of the ambiguity between copies and equality signs.}
\GSfield{Fate}{The reconstruction follows directly by the binomial theorem from the discriminant product in \citet[equation~(1)]{dyson1972missed}. The accompanying exact-arithmetic code independently computes the coefficients through \(n=50\).}
\GSfield{Present status}{As of 10 November 2026: Established in the stated scope.}
\GSfield{Assessment}{literal statement: false as written; publication match: partial; verification: read the result; development: annotated; confidence: high}
\end{GSitem}
\begin{GSitem}{1975-370.3}{E}{Dyson’s five-variable formula}
\GSfield{Kind; source}{theorem; \GSpage{P000613}}
\GSfield{The notes}{The second expression is \(\tau(n)=(1!2!3!4!)^{-1}\sum\prod_{i<j}(u_i-u_j)\), subject to \(\sum u_i=0\), \(\sum u_i^2=10n\), and \(u_i\equiv i\pmod5\).}
\GSfield{Statement}{For every positive integer \(n\), sum over all labelled integer 5-tuples satisfying those three conditions. Then the displayed Vandermonde sum equals \(\tau(n)\); its finiteness follows from the sum-of-squares condition. The congruences are attached to the labels, so no extra quotient by permutations is taken.}
\GSfield{Reconstruction}{The labelled summation convention and positivity of \(n\) are made explicit; no further mathematical repair is required.}
\GSfield{Fate}{\citet[equation~(2), pp.~636--637]{dyson1972missed} prints exactly this formula. His Gibbs lecture was delivered 17 January 1972 and the paper appeared in September 1972; this repeats 1971-269.2.}
\GSfield{Present status}{As of 10 November 2026: Established in the stated scope.}
\GSfield{Assessment}{literal statement: correct; publication match: exact; verification: read the result; development: none; confidence: high}
\end{GSitem}
\begin{GSitem}{1975-370.4}{E}{Ramanujan’s divisor-convolution identity}
\GSfield{Kind; source}{theorem; \GSpage{P000613}}
\GSfield{The notes}{The third formula is \(\tau(n)=\frac{65}{756}\sigma_{11}(n)+\frac{691}{756}\sigma_5(n)-\frac{691}{3}\sum_{m=1}^{n-1}\sigma_5(m)\sigma_5(n-m)\).}
\GSfield{Statement}{For \(\sigma_k(n)=\sum_{d\mid n}d^k\) and \(n\geq1\), this identity is correct. It follows by comparing coefficients in \(E_6^2=E_{12}-(762048/691)\Delta\), using \(E_6=1-504\sum\sigma_5(n)q^n\) and \(E_{12}=1+(65520/691)\sum\sigma_{11}(n)q^n\).}
\GSfield{Reconstruction}{The modular-form identity explains the otherwise opaque constants. The coefficient comparison is a reconstruction, not a second formula written in the notes.}
\GSfield{Fate}{The modular-form basis theorem is checked in \citet[Proposition~4.16]{milne2017modular}; the weight-12 normalization is checked independently in \citet[Section~3]{brouwerleech}. Direct Euler-product coefficients and a separate divisor recurrence agree with the notebook convolution for \(1\leq n\leq50\).}
\GSfield{Present status}{As of 10 November 2026: Established in the stated scope.}
\GSfield{Assessment}{literal statement: correct; publication match: exact; verification: read the result; development: annotated; confidence: high}
\end{GSitem}
\begin{GSitem}{1975-370.5}{E}{The congruence modulo 691}
\GSfield{Kind; source}{theorem; \GSpage{P000613}}
\GSfield{The notes}{The convolution formula is followed by \(\tau(n)\equiv\sigma_{11}(n)\pmod{691}\), attributed to Ramanujan.}
\GSfield{Statement}{For every positive integer \(n\), \(\tau(n)-\sigma_{11}(n)\) is divisible by \(691\). Indeed multiply the preceding identity by \(756\), reduce modulo \(691\), and use \(756\equiv65\pmod{691}\) with \(65\) invertible modulo \(691\).}
\GSfield{Reconstruction}{The elementary modular reduction is supplied; it does not require the later Galois-representation construction.}
\GSfield{Fate}{\citet[Section~1.3, p.~374]{manin1973periods} explicitly recovers this already classical congruence from his coefficient theorem. It also occurs in 1973-305.9.}
\GSfield{Present status}{As of 10 November 2026: Established in the stated scope.}
\GSfield{Assessment}{literal statement: correct; publication match: exact; verification: read the result; development: none; confidence: high}
\end{GSitem}
\begin{GSitem}{1975-370.6}{E}{The assertion that there are only two rank-24 forms}
\GSfield{Kind; source}{claim; \GSpage{P000614}}
\GSfield{The notes}{Integral quadratic forms of determinant one in 24 variables are said to number ``2 in all''; neither evenness nor positive definiteness is written.}
\GSfield{Statement}{Even after imposing positive definiteness and even unimodularity on the associated lattice, there are more than two classes. For example \(E_8^{\oplus3}\), \(D_{24}^+\), and the Leech lattice are distinguished by their numbers of norm-two vectors: \(720\), \(1104\), and \(0\). Here \(D_{24}=\{x\in\mathbb Z^{24}:\sum x_i\text{ is even}\}\) and \(D_{24}^+=D_{24}\cup((1/2,\ldots,1/2)+D_{24})\); the index-two extension is even and unimodular, and its new coset has minimal squared norm \(6\).}
\GSfield{Reconstruction}{This explicit three-lattice counterexample suffices; a full classification is not needed for the correction. The nearby number two may refer to the dimension of the weight-12 modular-form space, but that is an interpretation, not a legible correction of the scan.}
\GSfield{Fate}{\citet[Theorems~1.1--1.3 and Section~3]{brouwerleech} checks the Leech and theta-series facts and distinguishes the two rank-16 lattices from the two-dimensional weight-12 space. The \(D_{24}^+\) count above is the elementary root count \(2\cdot24\cdot23\).}
\GSfield{Present status}{As of 10 November 2026: Established in the stated scope.}
\GSfield{Assessment}{literal statement: false as written; publication match: partial; verification: read the result; development: annotated; confidence: high}
\end{GSitem}
\begin{GSitem}{1975-370.7}{E}{Recovering tau from a rank-24 lattice theta series}
\GSfield{Kind; source}{theorem; \GSpage{P000614}}
\GSfield{The notes}{The fourth formula expresses \(\tau(n)\) as a linear combination of the number of representations \(Q_{24}=n\) and \(\sigma_{11}(n)\), with constants obtained from small coefficients.}
\GSfield{Statement}{Let \(L\) be an even positive definite unimodular lattice of rank 24, set \(Q(v)=(v,v)/2\), and \(r_L(n)=\#\{v\in L:Q(v)=n\}\). Then \[ r_L(n)=\frac{65520}{691}\sigma_{11}(n)+c_L\tau(n),\qquad c_L=r_L(1)-\frac{65520}{691}. \] Thus \(\tau(n)\) is recovered by division by \(c_L\ne0\); nonvanishing follows because \(r_L(1)\) is an integer. This uses \(\Theta_L\in M_{12}(SL_2(\mathbb Z))=\mathbb C E_{12}\oplus\mathbb C\Delta\).}
\GSfield{Reconstruction}{The lattice hypotheses and half-norm convention are necessary. Determinant one alone is insufficient for this weight and full modular group. The constant coefficient and the coefficient of \(q\) already determine the two constants.}
\GSfield{Fate}{\citet[Theorem~1.2 and Section~3]{brouwerleech} verifies the theta-series argument and the explicit Leech instance. This is established modular-form theory, not a new classification of forms.}
\GSfield{Present status}{As of 10 November 2026: Established in the stated scope.}
\GSfield{Assessment}{literal statement: correct with missing hypotheses; publication match: partial; verification: read the result; development: none; confidence: high}
\end{GSitem}
\begin{GSitem}{1975-370.8}{E}{Manin’s finite coefficient formula and its boundary weights}
\GSfield{Kind; source}{theorem; \GSpage{P000614}}
\GSfield{The notes}{The fifth expression sums over \(n=\Delta\Delta'+\delta\delta'\), but its polynomial is printed as \((\Delta^8\delta^2-\Delta^2\delta^8)+(691/6)(\Delta^6\delta^4-\Delta^4\delta^6)\), subtracted from \(\sigma_{11}(n)\).}
\GSfield{Statement}{The corrected factorized form is \[ \tau(n)=\sigma_{11}(n)-\frac{691}{18}\sum{}'\Delta^2\delta^2(\Delta^2-\delta^2)^3. \] Sum over integers \(\Delta>\delta>0\) satisfying either \(\Delta'>\delta'>0\), \(n=\Delta\Delta'+\delta\delta'\), or \(\delta'=0\), \(\Delta\mid n\), \(\Delta'=n/\Delta\), and \(0<\delta/\Delta\leq1/2\). In the latter case the boundary \(\delta/\Delta=1/2\) has weight \(1/2\). Expanding the polynomial gives coefficient \(691/18\) on the first difference and \(-691/6\) on the second.}
\GSfield{Reconstruction}{Both coefficients and the boundary weight are restored. At \(n=2\), the corrected sum gives \(-24\); even granting the missing half weight, the notebook polynomial gives \(-841\). The factorized formula was independently checked for \(n\leq50\) against two exact coefficient algorithms; this finite check is not offered as a proof for all \(n\).}
\GSfield{Fate}{\citet[Section~1.3, pp.~373--374; proof in Section~5]{manin1973periods} is the direct earlier publication. The summation region and half-weight rule were read; the poor formula OCR was cross-checked by the factorization and exact computation. Russian publication was 1973, whereas the opened English translation carries copyright 1975; neither is inferred merely from receipt dates.}
\GSfield{Present status}{As of 10 November 2026: Established in the stated scope.}
\GSfield{Assessment}{literal statement: false as written; publication match: partial; verification: indirect; development: annotated; confidence: high}
\end{GSitem}
\begin{GSitem}{1975-370.9}{E}{Recursive computation using partition numbers}
\GSfield{Kind; source}{method; \GSpage{P000614}}
\GSfield{The notes}{The sixth expression uses \(P(q)=\sum p(n)q^n=\prod(1-q^n)^{-1}\), but writes the inner sum from \(n=1\) in \(\Delta=qP^{-24}\). It proposes recovering \(\tau\) recursively from \(p\).}
\GSfield{Statement}{The invertible formal series is \(P(q)=\sum_{n\geq0}p(n)q^n\), with \(p(0)=1\). Thus \(\Delta(q)=qP(q)^{-24}\) and ordinary series inversion gives a finite recursion for each coefficient. An independent recurrence from logarithmic differentiation is \((n-1)\tau(n)=-24\sum_{m=1}^{n-1}\sigma_1(m)\tau(n-m)\), \(n\geq2\).}
\GSfield{Reconstruction}{The lower index is repaired. The divisor recurrence is explicitly an additional elementary verification device, not asserted to be the recursion Manin wrote.}
\GSfield{Fate}{\citet[equations~27.14.3 and~27.14.7]{nist2026dlmf} verifies the partition normalization. The product and derivative recurrences independently agree through \(n=50\) in the supplied computation.}
\GSfield{Present status}{As of 10 November 2026: Established in the stated scope.}
\GSfield{Assessment}{literal statement: false as written; publication match: partial; verification: read the result; development: annotated; confidence: high}
\end{GSitem}
\begin{GSitem}{1975-370.10}{E}{The convergent Rademacher series and its phases}
\GSfield{Kind; source}{theorem; \GSpage{P000614}}
\GSfield{The notes}{A Hardy--Ramanujan--Rademacher series is written with a phase depending on \(1/k\) rather than \(h/k\), an asserted \(\omega_{h,k}^{24}=1\), and a derivative of a normalized hyperbolic sine. The exponential-only predecessor is contrasted with it.}
\GSfield{Statement}{For \(n\geq1\), put \(t=n-1/24\), \(C=\pi\sqrt{2/3}\), and \[ A_k(n)=\sum_{\substack{h\bmod k\\(h,k)=1}}e^{\pi i s(h,k)-2\pi inh/k},\quad s(h,k)=\sum_{r=1}^{k-1}\frac r k\left(\frac{hr}{k}-\left\lfloor\frac{hr}{k}\right\rfloor-\frac12\right). \] Then \[ p(n)=\frac1{\pi\sqrt2}\sum_{k\geq1}\sqrt{k}\,A_k(n)\left.\frac{d}{dt}\frac{\sinh(C\sqrt t/k)}{\sqrt t}\right|_{t=n-1/24}. \] The series converges. The factor \(e^{\pi i s(h,k)}\) does not in general have order dividing 24.}
\GSfield{Reconstruction}{This complete conventional normalization replaces the malformed formula. For instance \(s(1,5)=1/5\), so \((e^{\pi i/5})^{24}\ne1\). The Hardy--Ramanujan truncation is an asymptotic method, not the same convergent series with \(\sinh\) simply replaced by \(e^x\).}
\GSfield{Fate}{\citet[Section~27.14(iii), equations~27.14.8--11]{nist2026dlmf} gives the convergent series, phases, and historical distinction between the 1918 asymptotic and Rademacher's 1938 treatment.}
\GSfield{Present status}{As of 10 November 2026: Established in the stated scope.}
\GSfield{Assessment}{literal statement: false as written; publication match: partial; verification: read the result; development: annotated; confidence: high}
\end{GSitem}
\begin{GSitem}{1975-370.11}{E}{Hecke multiplicativity of tau}
\GSfield{Kind; source}{theorem; \GSpage{P000614}}
\GSfield{The notes}{The seventh entry underlines \(\tau(mn)=\tau(m)\tau(n)\) when \((m,n)=1\), noting that this is not apparent from the earlier formulas.}
\GSfield{Statement}{The normalized cusp form \(\Delta\), spanning the one-dimensional space \(S_{12}(SL_2(\mathbb Z))\), is a simultaneous Hecke eigenform. With the conventional normalized Hecke operators, \(T_n\Delta=\tau(n)\Delta\), and \(T_mT_n=T_{mn}\) for coprime \(m,n\), giving the displayed multiplicativity.}
\GSfield{Reconstruction}{The Hecke interpretation and normalization are supplied as the explanation. No assertion of multiplicativity for noncoprime arguments is made.}
\GSfield{Fate}{\citet[Proposition~5.18 and Corollary~5.20]{milne2017modular} verifies the standard argument. Earlier relevant occurrences are 1973-304.8--9.}
\GSfield{Present status}{As of 10 November 2026: Established in the stated scope.}
\GSfield{Assessment}{literal statement: correct; publication match: exact; verification: read the result; development: none; confidence: high}
\end{GSitem}
\begin{GSitem}{1975-370.12}{E}{Prime-power coefficients are complete symmetric sums}
\GSfield{Kind; source}{claim; \GSpage{P000614}}
\GSfield{The notes}{After \(\tau(p)=\theta_p+\bar\theta_p\) and \(\theta_p\bar\theta_p=p^{11}\), the notes underline \(\tau(p^a)=\theta_p^a+\bar\theta_p^a\).}
\GSfield{Statement}{If \(\alpha_p,\beta_p\) are the roots of \(T^2-\tau(p)T+p^{11}\), then \[ \sum_{a\geq0}\tau(p^a)T^a=\frac1{(1-\alpha_pT)(1-\beta_pT)},\qquad \tau(p^a)=\sum_{j=0}^a\alpha_p^{a-j}\beta_p^j. \] Equivalently \(\tau(p^{a+1})=\tau(p)\tau(p^a)-p^{11}\tau(p^{a-1})\). A sum of just the two powers describes a trace of a Frobenius power, not this Hecke coefficient.}
\GSfield{Reconstruction}{The distinction is decisive already at \(p=2,a=2\): the notebook expression is \(-3520\), while \(\tau(4)=-1472\). The complete sum, or its quotient \((\alpha^{a+1}-\beta^{a+1})/(\alpha-\beta)\) when distinct, repairs the formula.}
\GSfield{Fate}{\citet[Corollary~5.20]{milne2017modular} verifies the recurrence. The small counterexample is checked independently by direct product multiplication and the divisor recurrence.}
\GSfield{Present status}{As of 10 November 2026: Established in the stated scope.}
\GSfield{Assessment}{literal statement: false as written; publication match: partial; verification: read the result; development: annotated; confidence: high}
\end{GSitem}
\begin{GSitem}{1975-370.13}{E}{The two-dimensional l-adic representation}
\GSfield{Kind; source}{theorem; \GSpage{P000614}}
\GSfield{The notes}{The eighth entry says that \(\tau\) is the trace of a representation and names Serre, Deligne, Grothendieck, Weil, Kuga, and Sato. The target is drawn as \(\mathbb Q_\ell^2\).}
\GSfield{Statement}{For every prime \(\ell\), the eigenform \(\Delta\) has a continuous two-dimensional \(\ell\)-adic representation \(\rho_\ell:G_{\mathbb Q}\to GL_2(\mathbb Q_\ell)\), unramified outside \(\ell\), with \(\det(1-\rho_\ell(\operatorname{Frob}_p)T)=1-\tau(p)T+p^{11}T^2\) for \(p\ne\ell\), choosing Frobenius and dual conventions to give determinant \(p^{11}\). In particular \(\tau(p)\) is a trace at primes; the preceding item explains why \(\tau(p^a)\) is not usually the trace of the \(a\)-th power.}
\GSfield{Reconstruction}{The group-valued target, continuity, ramification restriction, and convention are supplied. A vector space is the representation space, not the target group of the homomorphism.}
\GSfield{Fate}{\citet[Sections~3--4, especially Theorem~4.9]{deligne1971modulaires} constructs the cohomological representation. The Bourbaki talk was given in February 1969; the collected publication is 1971, so this is established before the seminar.}
\GSfield{Present status}{As of 10 November 2026: Established in the stated scope.}
\GSfield{Assessment}{literal statement: false as written; publication match: partial; verification: read the result; development: none; confidence: high}
\end{GSitem}
\begin{GSitem}{1975-370.14}{E}{Deligne’s bound does not make the roots purely imaginary}
\GSfield{Kind; source}{theorem; \GSpage{P000614}}
\GSfield{The notes}{The text calls \(\theta_p\) ``purely imaginary (nonreal)'', and attributes the result to Deligne's proof of the Weil conjectures.}
\GSfield{Statement}{The roots of \(T^2-\tau(p)T+p^{11}\) have complex absolute value \(p^{11/2}\), and \(|\tau(p)|\leq2p^{11/2}\). Since \(\tau(p)\) is an integer and \(2p^{11/2}\) is irrational, equality cannot occur; the roots are a nonreal conjugate pair. They need not be purely imaginary: at \(p=2\), their real part is \(\tau(2)/2=-12\).}
\GSfield{Reconstruction}{The distinction between \ru{чисто мнимые} and \ru{невещественные} is mathematical, not an English mistranslation. The strictness and small-prime counterexample are elementary consequences of the established bound.}
\GSfield{Fate}{\citet[Theorem~5.6]{deligne1971modulaires} gives the implication Weil \(\Rightarrow\) Ramanujan; \citet[Theorem~1.6]{deligne1974weil} proves the needed weight statement. The latter was publicly published in 1974, not merely received then; it precedes this talk.}
\GSfield{Present status}{As of 10 November 2026: Established in the stated scope.}
\GSfield{Assessment}{literal statement: false as written; publication match: partial; verification: read the result; development: annotated; confidence: high}
\end{GSitem}
\begin{GSitem}{1975-370.15}{E}{Eta transformations and the invariant sixth differential}
\GSfield{Kind; source}{theorem; \GSpage{P000615}}
\GSfield{The notes}{Eta is rewritten in the upper-half-plane variable, with an extraneous \(n\) in the leading exponential. Translation and inversion laws are given, and \(\eta^{24}(z)(dz)^6\) is called invariant.}
\GSfield{Statement}{For \(\operatorname{Im}z>0\), \(\eta(z)=e^{2\pi iz/24}\prod_{n\geq1}(1-e^{2\pi inz})\), \(\eta(z+1)=e^{\pi i/12}\eta(z)\), and \(\eta(-1/z)=\sqrt{-iz}\,\eta(z)\), with the holomorphic square root positive at \(z=i\). Therefore \(\Delta(\gamma z)=(cz+d)^{12}\Delta(z)\) for \(\gamma\in SL_2(\mathbb Z)\), and \(\Delta(z)(dz)^6\) is invariant under pullback.}
\GSfield{Reconstruction}{The leading exponent, square-root branch, and group are specified. The multiplier of eta has finite order in the transformation law; this does not justify the different \(\omega_{h,k}^{24}=1\) assertion in the partition formula.}
\GSfield{Fate}{\citet[equations~27.14.12--14]{nist2026dlmf} verifies the eta definition and multiplier formula. This is classical exposition.}
\GSfield{Present status}{As of 10 November 2026: Established in the stated scope.}
\GSfield{Assessment}{literal statement: false as written; publication match: partial; verification: read the result; development: none; confidence: high}
\end{GSitem}
\begin{GSitem}{1975-370.16}{E}{The graded ring generated by E4 and E6}
\GSfield{Kind; source}{theorem; \GSpage{P000615}}
\GSfield{The notes}{The notes assert that the ring of modular forms is \(\mathbb C[E_4,E_6]\) and identify the nonconstant coefficients with \(\sigma_3(n)\) and \(\sigma_5(n)\), without their normalization factors.}
\GSfield{Statement}{For holomorphic scalar modular forms on \(SL_2(\mathbb Z)\), holomorphic also at the cusp, \[ M_* = \bigoplus_{k\geq0}M_k(SL_2(\mathbb Z))=\mathbb C[E_4,E_6], \] where \(E_4=1+240\sum_{n\geq1}\sigma_3(n)q^n\), \(E_6=1-504\sum_{n\geq1}\sigma_5(n)q^n\), and \(\Delta=(E_4^3-E_6^2)/1728\). Homogeneous polynomial identities therefore yield coefficient convolutions such as 1975-370.4.}
\GSfield{Reconstruction}{The group, cusp condition, grading, constant coefficients, and scaling factors are essential. The theorem is not a statement about all levels, meromorphic forms, or arbitrary multiplier systems.}
\GSfield{Fate}{\citet[Proposition~4.16 and p.~78]{milne2017modular} proves the structure and supplies the two normalized generators, using half the conventional weight in its subscripts. Earlier occurrences are 1973-304.6--7.}
\GSfield{Present status}{As of 10 November 2026: Established in the stated scope.}
\GSfield{Assessment}{literal statement: correct with missing hypotheses; publication match: partial; verification: read the result; development: none; confidence: high}
\end{GSitem}
\begin{GSitem}{1975-370.17}{E}{Rademacher contours and cusp singularities}
\GSfield{Kind; source}{method; \GSpage{P000615}}
\GSfield{The notes}{The explanation of formula six draws circles tangent at rational points and successive integration contours; it calls the boundary singularities poles of infinite order.}
\GSfield{Statement}{Coefficient extraction for \(P(q)=\prod(1-q^n)^{-1}\), followed by \(q=e^{2\pi iz}\), deforms a contour through arcs associated with rational cusps (Ford circles in the standard construction). Modular transformations give local expansions contributing the sums \(A_k(n)\). The analytic boundary singularities at roots of unity are not meromorphic poles of a finite order; ``pole of infinite order'' is informal terminology.}
\GSfield{Reconstruction}{The identification of the circles as Ford circles and the modular coefficient-extraction mechanism are reconstructed from the labelled Rademacher diagram. No error estimate for a particular drawn truncation is inferred.}
\GSfield{Fate}{The resulting convergent coefficient formula is checked in \citet[equations~27.14.9--11]{nist2026dlmf}; this verifies the endpoint, while the exact hand-drawn contour parametrization remains schematic.}
\GSfield{Present status}{As of 10 November 2026: Established in the stated scope.}
\GSfield{Assessment}{literal statement: correct with missing hypotheses; publication match: partial; verification: indirect; development: none; confidence: high}
\end{GSitem}
\begin{GSitem}{1975-370.18}{E}{Kostant’s identity revisited with a missing power}
\GSfield{Kind; source}{theorem; \GSpage{P000615}}
\GSfield{The notes}{The final formula again sums \(\chi_\lambda(a)\dim V_\lambda q^{c(\lambda)}\), now with \(c(\lambda)=(\lambda+\rho,\lambda+\rho)-(\rho,\rho)\), but the left side is only \(\psi(q)\).}
\GSfield{Statement}{With the compact simple simply connected group and Killing normalization of 1975-366.4, the left side is \(\psi(q)^{\dim K}\), and the displayed Casimir exponent is correct in that normalization. The root-system character values are again \(0,\pm1\).}
\GSfield{Reconstruction}{The omitted dimension power is restored by direct comparison with P000610 and with the verified theorem; the new occurrence is kept because the exponent is supplied here for the first time in these notes.}
\GSfield{Fate}{\citet[Theorem~0.1]{kostant2003euler} verifies the formula. It had already been reported at this seminar on 22 September; classification E describes this repeated exposition, while 1975-366.4 records the subsequent 1976 publication.}
\GSfield{Present status}{As of 10 November 2026: Established in the stated scope.}
\GSfield{Assessment}{literal statement: false as written; publication match: partial; verification: read the result; development: none; confidence: high}
\end{GSitem}

\GStalk{371}{1975-10-06}{J. Leray}{P000615-P000618}{Quadratic Fourier operators and the metaplectic group}
Leray explains the algebraic background of Maslov's method, beginning with the Schwartz space, its distribution dual, and linear position--momentum operators. Quadratic generating functions produce unitary integral operators. The notes then pass from their arbitrary scalar phases to the metaplectic double cover and its Maslov signs; this passage needs a phase convention beyond unitarity alone. The subject continues the metaplectic discussion in 1973-312.1--2 and 1973-315.1--4.

\begin{GSitem}{1975-371.1}{E}{Heisenberg operators and their symplectic normalizer}
\GSfield{Kind; source}{construction; \GSpage{P000616}}
\GSfield{The notes}{The operators \(a(x,\nu^{-1}\partial_x)\), linear in position and momentum, are normalized by automorphisms of \(S'(\mathbb R^\ell)\). The quotient by nonzero scalar operators is identified with \(\operatorname{Sp}(\ell)\).}
\GSfield{Statement}{Let \(Q_jf=x_jf\), \(P_jf=-i\hbar\partial_jf\), with \(\hbar>0\), on \(\mathcal S(\mathbb R^\ell)\subset L^2(\mathbb R^\ell)\subset\mathcal S'(\mathbb R^\ell)\). Their commutators encode the real symplectic form on \(\mathbb R^{2\ell}\). An invertible continuous complex-linear implementer preserving the common Schwartz domain and the real linear span of the \(Q_j,P_j\) induces a real symplectic transformation. In the irreducible Schrödinger representation, implementers of a fixed transformation differ by a scalar; allowing all nonzero scalar multiples gives kernel \(\mathbb C^\times\), while unitary implementers have kernel \(U(1)\).}
\GSfield{Reconstruction}{The real structure, continuity, irreducible Schrödinger representation, and common invariant Schwartz domain are supplied. These conditions delimit the phrase \emph{automorphisms of distributions}; arbitrary algebraic automorphisms of \(S'\) are not being classified.}
\GSfield{Fate}{The metaplectic implementers and their action on Schwartz functions and distributions are verified in \citet[Section~1]{degosson1990maslov}. That article explicitly reviews the earlier Weil--Leray framework; this is an expository occurrence, not a new 1975 normalizer theorem.}
\GSfield{Present status}{As of 10 November 2026: Established in the stated scope.}
\GSfield{Assessment}{literal statement: correct with missing hypotheses; publication match: partial; verification: read the result; development: none; confidence: high}
\end{GSitem}
\begin{GSitem}{1975-371.2}{E}{Quadratic generating functions and unitary Fourier kernels}
\GSfield{Kind; source}{construction; \GSpage{P000617}}
\GSfield{The notes}{For a real quadratic \(A(x,x')\), the notes write a kernel proportional to \(\Delta(A)e^{\nu A(x,x')}\), with \(\Delta(A)=\sqrt{\det A''}\), and the canonical relation \(p=A_x,\ p'=-A_{x'}\).}
\GSfield{Statement}{Assume \(A(x,y)\) is a real homogeneous quadratic form and its mixed Hessian \(B=\partial_x\partial_y A\) is invertible. The equations \(p=\partial_xA,\ p'=-\partial_yA\) define a linear symplectic map. For any phase \(\zeta\in U(1)\),
     \[ (T_Af)(x)=\zeta(2\pi\hbar)^{-\ell/2}|\det B|^{1/2}\int_{\mathbb R^\ell}e^{iA(x,y)/\hbar}f(y)\,dy \]
     defines a unitary operator on \(L^2\), initially on Schwartz functions and then by extension. It factors into quadratic phase multipliers, an invertible linear substitution with its Jacobian factor, and a Fourier transform.}
\GSfield{Reconstruction}{Here \(\nu=i/\hbar\); the determinant means the mixed Hessian, and its absolute value supplies the modulus. The complex square-root phase in the notebook must be chosen consistently. Symplectic maps for which this mixed-coordinate generating function fails require other charts or products of such operators.}
\GSfield{Fate}{The explicit kernels and elementary factorization appear in \citet[Section~1, equations~(1.1)--(1.5)]{degosson1990maslov}. The formula verifies the local chart construction, with its missing determinant and phase qualifications.}
\GSfield{Present status}{As of 10 November 2026: Established in the stated scope.}
\GSfield{Assessment}{literal statement: correct with missing hypotheses; publication match: exact; verification: read the result; development: none; confidence: high}
\end{GSitem}
\begin{GSitem}{1975-371.3}{E}{Unitarity alone does not select the metaplectic double cover}
\GSfield{Kind; source}{claim; \GSpage{P000617}}
\GSfield{The notes}{The line \ru{Унитарность} \(\Rightarrow\) \ru{метапл. группа} is followed by \(|c|=1\), then by the instruction to take \(c=\pm1\).}
\GSfield{Statement}{All unitary canonical implementers form a central \(U(1)\)-extension of \(\operatorname{Sp}(2\ell,\mathbb R)\). The metaplectic group is a distinguished subgroup defined by coherent phase choices and fits into
     \[1\longrightarrow\{\pm I\}\longrightarrow\operatorname{Mp}(2\ell,\mathbb R)\longrightarrow\operatorname{Sp}(2\ell,\mathbb R)\longrightarrow1.\]
     Multiplying a unitary implementer by an arbitrary phase preserves unitarity and its underlying symplectic map. Thus the displayed implication requires the additional normalization.}
\GSfield{Reconstruction}{The modern notation \(\operatorname{Sp}(2\ell,\mathbb R)\) denotes the group written \(\operatorname{Sp}(\ell)\) in the notes. The double cover is the corrected conclusion after phase normalization, not a consequence of \(|c|=1\) alone.}
\GSfield{Fate}{\citet[Section~1, Theorem~1]{degosson1990maslov} states the double-cover construction, using discrete phase indices in the generators. Earlier seminar occurrences are 1973-312.1--2.}
\GSfield{Present status}{As of 10 November 2026: Established in the stated scope.}
\GSfield{Assessment}{literal statement: false as written; publication match: partial; verification: read the result; development: annotated; confidence: high}
\end{GSitem}
\begin{GSitem}{1975-371.4}{E}{Composition phases and the Maslov sign}
\GSfield{Kind; source}{claim; \GSpage{P000617}}
\GSfield{The notes}{The notebook writes \(S_AS_{A'}=S_{A''}\) and, when \(s_As_{A'}s_{A''}=e\), the doubly underlined identity \(S_AS_{A'}S_{A''}=\pm E\), with its sign linked to the Maslov index.}
\GSfield{Statement}{For lifts \(\widehat s_1,\widehat s_2,\widehat s_3\in\operatorname{Mp}(2\ell,\mathbb R)\) with \(s_1s_2s_3=I\), their product lies in the kernel and is \(\pm I\). In a nondegenerate quadratic-kernel chart, Gaussian integration contributes the signature phase \(e^{i\pi\operatorname{sgn}(H)/4}\), where \(H\) is the Hessian in the integrated intermediate variables. The discrete phase indices must incorporate this contribution. For arbitrary unitary lifts, the corresponding product can instead be any scalar in \(U(1)\).}
\GSfield{Reconstruction}{The triple-product assertion is separated from a globally defined multiplication law for generating functions: compositions may leave a given chart, and an unqualified assignment \(A\mapsto S_A\) does not provide a global section of the cover.}
\GSfield{Fate}{\citet[Section~1, Theorem~2]{degosson1990maslov} gives the quadratic composition rule with its signature correction. The notes correctly point toward the Maslov obstruction but suppress the phase bookkeeping.}
\GSfield{Present status}{As of 10 November 2026: Established in the stated scope.}
\GSfield{Assessment}{literal statement: correct with missing hypotheses; publication match: partial; verification: read the result; development: none; confidence: high}
\end{GSitem}

\GStalk{372}{1975-10-13}{E. S. Fradkin}{P000616-P000622}{Integrable particle systems and constrained quantization}
After Fradkin's reading list, Perelomov gives a substantial intervening report on Toda and Calogero--Moser systems. Its table distinguishes finding conserved quantities, proving their involution, reconstructing the dynamics, and performing quadratures; elliptic involution and exceptional root systems are explicitly questioned. Gelfand asks about a larger group-theoretic construction. Fradkin then returns to Bose and Fermi canonical variables, constrained Hamiltonian systems, and a formal functional integral. The individual attributions below follow the page headings rather than assigning Perelomov's report to Fradkin.

\begin{GSitem}{1975-372.1}{E}{Liouville integrability and quadratures}
\GSfield{Kind; source}{theorem; \GSpage{P000618}; A. M. Perelomov; I. M. Gelfand}
\GSfield{The notes}{Perelomov defines complete integrability by \(n\) first integrals in involution for \(H=\frac12\sum p_k^2+U(q)\); Gelfand recalls Liouville's theorem that the system can then be solved by quadratures.}
\GSfield{Statement}{On a smooth \(2n\)-dimensional symplectic manifold, suppose \(F_1=H,F_2,\ldots,F_n\) Poisson commute and \(dF_1,\ldots,dF_n\) are independent on the regular region under consideration. The local integration procedure reduces the motion to quadratures. A compact connected regular common level is an \(n\)-torus, with action--angle coordinates in a neighborhood and linear motion on each torus. Without compactness this global torus conclusion does not follow.}
\GSfield{Reconstruction}{Functional independence, regularity and the separate compactness hypothesis are supplied; merely listing \(n\) possibly dependent conserved functions is insufficient.}
\GSfield{Fate}{The classical theorem is stated with completeness and regularity conditions in \citet[Section~4.4, Theorem~4.11]{cannas2007introduction}. This is background for the four distinct tasks subsequently listed by Gelfand.}
\GSfield{Present status}{As of 10 November 2026: Established in the stated scope.}
\GSfield{Assessment}{literal statement: correct with missing hypotheses; publication match: exact; verification: read the result; development: none; confidence: high}
\end{GSitem}
\begin{GSitem}{1975-372.2}{E}{The pendulum quadrature is not a second independent integral}
\GSfield{Kind; source}{remark; \GSpage{P000618}; A. M. Perelomov}
\GSfield{The notes}{For \(H=\frac12p^2+\sin q\), the next line reads \ru{2-й инт. --- эллиптич. ф-я}, rendered as a second integral and an elliptic function.}
\GSfield{Statement}{On an energy level \(H=E\), \(\dot q=p\) gives
     \[t-t_0=\pm\int_{q_0}^{q}\frac{ds}{\sqrt{2(E-\sin s)}}.\]
     Away from turning points this is a local quadrature; its inversion is expressed through elliptic functions for generic energy. The critical energies \(E=\pm1\) require separate limiting or equilibrium solutions. There cannot be a second functionally independent first integral on a regular two-dimensional phase region.}
\GSfield{Reconstruction}{The abbreviation is interpreted as the second integration step, not an additional independent conserved quantity. The quadrature follows directly from the stated Hamiltonian.}
\GSfield{Fate}{This is the classical one-degree-of-freedom example; \citet[Section~4.5]{cannas2007introduction} treats the equivalent shifted pendulum Hamiltonian and its critical levels.}
\GSfield{Present status}{As of 10 November 2026: Established in the stated scope.}
\GSfield{Assessment}{literal statement: not determinable; publication match: partial; verification: read the result; development: none; confidence: high}
\end{GSitem}
\begin{GSitem}{1975-372.3}{E}{The periodic Toda Hamiltonian}
\GSfield{Kind; source}{construction; \GSpage{P000618-P000619}; A. M. Perelomov}
\GSfield{The notes}{The first Toda potential is written as a difference of exponentials, with a summation index \(k\) but exponents using \(n\). P000619 instead gives \(U=g^2\sum_{k=1}^n e^{-(q_{k+1}-q_k)}\); P000618 states \(q_{k+n}=q_k\).}
\GSfield{Statement}{For \(n\geq3\), take the phase space \(T^*\mathbb R^n\), canonical brackets, and the periodic convention \(q_{n+1}=q_1\). The corrected Hamiltonian is
     \[H=\tfrac12\sum_{k=1}^np_k^2+g^2\sum_{k=1}^ne^{q_k-q_{k+1}}.\]
     The difference of neighboring exponentials belongs in the force equation \(\dot p_k=g^2(e^{q_{k-1}-q_k}-e^{q_k-q_{k+1}})\). Summing such differences as the potential would telescope to zero after correcting the index.}
\GSfield{Reconstruction}{The potential is restored from the explicit repetition on P000619. The open chain has a different boundary convention and is not silently substituted for the periodic system.}
\GSfield{Fate}{\citet[abstract]{henon1974integrals}, published 15 February 1974, explicitly reports \(n\) integrals for periodic Toda and covers fixed ends separately. Only that publisher statement and metadata were checked here.}
\GSfield{Present status}{As of 10 November 2026: Established in the stated scope.}
\GSfield{Assessment}{literal statement: false as written; publication match: partial; verification: publisher metadata checked; development: annotated; confidence: high}
\end{GSitem}
\begin{GSitem}{1975-372.4}{B}{Involution of the Toda integrals}
\GSfield{Kind; source}{question; \GSpage{P000618}; A. M. Perelomov}
\GSfield{The notes}{The chronology names Hénon, Flaschka, Manakov, and Kac--van Moerbeke, adding that integrals existed but involution had not been proved.}
\GSfield{Statement}{For the periodic Toda Hamiltonian of 1975-372.3, the issue is to exhibit \(n\) generically independent conserved functions \(I_1,\ldots,I_n\) with \(\{I_j,I_k\}=0\), not merely to produce an isospectral matrix evolution. Independence is understood on an open dense regular set; the center-of-mass motion remains a free degree of freedom.}
\GSfield{Reconstruction}{The precise family and boundary condition are fixed by the adjacent potential. The notebook's historical assertion about which proof was available in October 1975 has not been independently established.}
\GSfield{Fate}{A subsequent paper is specifically devoted to this claim: \citet[pp.~1668--1669]{kawasaki1976involution}, received 8 January and published 1 May 1976. Its publisher record was checked, but the proof was not accessed, so neither an earliest-proof priority nor an exact match to every named 1975 family is asserted.}
\GSfield{Present status}{As of 10 November 2026: Toda involution is established; the precise chronology claimed by this note remains only partially verified.}
\GSfield{Assessment}{literal statement: not determinable; publication match: partial; verification: publisher metadata checked; development: none; confidence: medium; problem \hyperref[prob:GP-1975-05]{GP-1975-05}}
\end{GSitem}
\begin{GSitem}{1975-372.5}{E}{Five pair-potential families and their phase spaces}
\GSfield{Kind; source}{construction; \GSpage{P000618}; A. M. Perelomov}
\GSfield{The notes}{The table lists \(q^{-2}\), \(a^2\sinh^{-2}(aq)\), \(a^2\sin^{-2}(aq)\), \(a^2\wp(aq)\), and \(q^{-2}+a^2q^2\), first for \(U=g^2\sum_{i<j}V(q_i-q_j)\). The hyperbolic and trigonometric cases are called limits of the elliptic case.}
\GSfield{Statement}{For \(n\) particles on the line, use canonical coordinates \((q_i,p_i)\) on the cotangent bundle of a collision-free configuration chamber; exclude all poles of \(V\). The trigonometric model may instead be interpreted on a circle, with collisions modulo its period excluded. In the elliptic model choose periods and a real slice on which the potential is real. The fifth pair potential leaves the center of mass free:
     \(\sum_{i<j}(q_i-q_j)^2=n\sum_i(q_i-Q_{\rm cm})^2\).
     Thus its relative motion has harmonic confinement, while a one-body trap \(\frac12\omega^2\sum q_i^2\) also confines the center of mass. Elliptic degenerations give the first three potentials up to scale and additive constants.}
\GSfield{Reconstruction}{The domains, reality conditions and separation of center-of-mass motion are supplied. No integrability claim is inferred for arbitrary singular continuations or boundary identifications.}
\GSfield{Fate}{The five types and their allowed chambers are explicitly listed in \citet[Section~2, equations~(2.1)--(2.6)]{olshanetskyperelomov1978radial}; their classical antecedents are summarized in \citet[Section~2]{sasaki2006explicit}. The named \(A\)-type models were already being reported as known here.}
\GSfield{Present status}{As of 10 November 2026: Established in the stated scope.}
\GSfield{Assessment}{literal statement: correct with missing hypotheses; publication match: partial; verification: read the result; development: none; confidence: high}
\end{GSitem}
\begin{GSitem}{1975-372.6}{A}{Classical and nonreduced root-system potentials}
\GSfield{Kind; source}{construction; \GSpage{P000618-P000619}; A. M. Perelomov}
\GSfield{The notes}{The second family adds \(V(q_i+q_j)\), \(V(q_i)\), and \(V(2q_i)\). P000619 rewrites this as \(\sum_{\alpha>0}g_\alpha^2V((\alpha,q))\), with equal couplings for equal root lengths, and identifies types \(A,B,C,BC,D\).}
\GSfield{Statement}{Let \(R\) be a finite crystallographic root system, or the standard nonreduced \(BC_r\), in a Euclidean space \(\mathfrak a\). On \(T^*\mathfrak a_{\rm reg}\), away from the potential's root hyperplanes or affine pole walls, set
     \[H=\tfrac12(p,p)+\sum_{\alpha\in R_+}c_\alpha V((\alpha,q)),\]
     with couplings constant on Weyl orbits. The roots \(e_i-e_j\) yield \(A_{n-1}\) for \(n\) particles after separating the center of mass; \(e_i\pm e_j,e_i,2e_i\) yield the \(BC\) family. In the displayed \(B\) formula the plus and minus pair terms must have compatible Weyl-invariant coefficients; an arbitrary independent \(\varepsilon\) is not covered by this statement.}
\GSfield{Reconstruction}{The shift from \(A_n\) to \(A_{n-1}\) records the usual rank convention. Weyl-orbit coupling and regular-domain hypotheses are made explicit. This item is the construction, distinct from the exceptional-case integrability question.}
\GSfield{Fate}{The speaker's joint work gives this root-system framework in \citet[Section~2]{olshanetskyperelomov1978radial}; that 1978 paper itself cites the classical 1976 construction. The match is to the stated family, not to every arbitrary coefficient suggested by the shorthand.}
\GSfield{Present status}{As of 10 November 2026: Established in the stated scope.}
\GSfield{Assessment}{literal statement: correct with missing hypotheses; publication match: partial; verification: read the result; development: none; confidence: high}
\end{GSitem}
\begin{GSitem}{1975-372.7}{B}{Commuting complete integrals in the elliptic case}
\GSfield{Kind; source}{question; \GSpage{P000618}; A. M. Perelomov}
\GSfield{The notes}{In row 4, conserved quantities for the elliptic potential are credited to Calogero and Olshanetsky--Perelomov, while the involution column contains two question marks, corresponding to families A and B.}
\GSfield{Statement}{For the unreduced \(n\)-particle elliptic Calogero--Moser system on \(T^*Q_{\rm reg}\), where \(Q_{\rm reg}\subset\mathbb R^n\) (or \(\mathbb C^n\)) excludes every pair difference at a potential pole, do there exist \(n\) generically independent meromorphic integrals in involution, including the free center-of-mass momentum? After removing the center-of-mass canonical pair, the relative \(A_{n-1}\) problem instead needs \(n-1\) integrals. The analogous root-system question requires orbit-constant couplings. The assertion concerns generic regular points, not a uniform action--angle chart through collisions or singular fibers.}
\GSfield{Reconstruction}{The two table columns are retained, but the completely verified answer below is restricted to reduced crystallographic root systems, including \(A,B,C,D\). No blanket identification with every nonreduced \(BC\) or arbitrary-\(\varepsilon\) variant is made.}
\GSfield{Fate}{\citet[pp.~3--5 and Table~I]{khastgirsasaki2001liouville} proves involution and independence for the root-system Calogero--Moser models, including the elliptic potentials in its scope. The preprint circulated 30 May 2000 and appeared in print in 2001; it supplies a verified later answer, not a claim to the first solution.}
\GSfield{Present status}{As of 10 November 2026: Affirmative for the reduced crystallographic elliptic models specified here; the shorthand B-family is broader than the exact formulation checked.}
\GSfield{Assessment}{literal statement: correct with missing hypotheses; publication match: partial; verification: read the result; development: none; confidence: high; problem \hyperref[prob:GP-1975-02]{GP-1975-02}}
\end{GSitem}
\begin{GSitem}{1975-372.8}{E}{Four separate tasks in solving an integrable model}
\GSfield{Kind; source}{method; \GSpage{P000618}; I. M. Gelfand}
\GSfield{The notes}{Gelfand lists: find the integrals, prove involution, solve the inverse problem, and express the solution by quadratures.}
\GSfield{Statement}{Conservation, \(\{I_j,H\}=0\), does not imply mutual involution or functional independence. Liouville integration requires a complete independent commuting family on a regular region. An inverse problem additionally requires specification of spectral data and its equivalences; quadrature or explicit inversion is a further constructive step. Consequently the table's question marks concern different mathematical assertions.}
\GSfield{Reconstruction}{The fourth phrase is interpreted as integrating the dynamics, rather than evaluating an already known constant of motion. No single unsolved theorem is extracted from this methodological list.}
\GSfield{Fate}{The distinction is consistent with the hypotheses of \citet[Theorem~4.11]{cannas2007introduction}; concrete inverse data are addressed separately in 1975-372.11.}
\GSfield{Present status}{As of 10 November 2026: Established in the stated scope.}
\GSfield{Assessment}{literal statement: correct; publication match: partial; verification: read the result; development: none; confidence: high}
\end{GSitem}
\begin{GSitem}{1975-372.9}{E}{Lax matrices and trace invariants}
\GSfield{Kind; source}{construction; \GSpage{P000619}; A. M. Perelomov}
\GSfield{The notes}{For types 1--4, \(I_k=\operatorname{Sp}(L^k)\), with \(L=P+iX\), \(P_{ij}=p_i\delta_{ij}\), and \(X_{ij}=(1-\delta_{ij})x(q_i-q_j)\). The table gives \(x=q^{-1},a\coth(aq),a\cot(aq),a\,\operatorname{cn}(aq)/\operatorname{sn}(aq)\).}
\GSfield{Statement}{An exact rational reconstruction is
     \[L_{ij}=p_i\delta_{ij}+\frac{ig(1-\delta_{ij})}{q_i-q_j},\qquad
     M_{ij}=-\frac{ig}{(q_i-q_j)^2}\ (i\ne j),\quad M_{ii}=ig\sum_{j\ne i}(q_i-q_j)^{-2}.\]
     For \(H=\frac12\sum p_i^2+g^2\sum_{i<j}(q_i-q_j)^{-2}\), Hamilton's equations give \(\dot L=[L,M]\), hence every \(\operatorname{tr}L^k\) is conserved and \(H=\frac12\operatorname{tr}L^2\). This argument alone does not prove involution. The nonrational kernels require compatible diagonal terms, coupling and elliptic-period conventions; the table alone does not specify these.}
\GSfield{Reconstruction}{The missing coupling and companion matrix are supplied only for the rational case. No assertion that the displayed scalar \(n\times n\) matrix directly represents every \(B,C,D\) model is retained.}
\GSfield{Fate}{The rational matrix is verified in \citet[Section~5, equations~(5.2)--(5.10)]{sasaki2006explicit}; the notebook's broader table is therefore a partial match.}
\GSfield{Present status}{As of 10 November 2026: Established in the stated scope.}
\GSfield{Assessment}{literal statement: correct with missing hypotheses; publication match: partial; verification: read the result; development: none; confidence: high}
\end{GSitem}
\begin{GSitem}{1975-372.10}{E}{The quadratic matrix invariant for harmonic confinement}
\GSfield{Kind; source}{construction; \GSpage{P000619}; A. M. Perelomov}
\GSfield{The notes}{For type 5 the notes use \(N=L^2+\omega^2Q^2\), \(Q=\operatorname{diag}(q_i)\), and \(I_k=\operatorname{Sp}(N^k)\).}
\GSfield{Statement}{For the rational Calogero Hamiltonian with a one-body harmonic trap, the rational matrices satisfy
     \(\dot L=[L,M]-\omega^2Q\), \(\dot Q=[Q,M]+L\).
     Therefore \(N=L^2+\omega^2Q^2\) satisfies \(\dot N=[N,M]\), and \(\operatorname{tr}N^k\) is conserved. For the pairwise quadratic potential on P000618 this applies to the relative variables after separating the free center of mass and matching the frequency; with the table's coefficients, \(\omega^2=2ng^2a^2\). Conservation of these traces is the assertion verified here, not a separate proof that this particular list is independent and in involution.}
\GSfield{Reconstruction}{The phase space is a collision-free chamber, with either the stated one-body trap or the explicitly reduced center-of-mass system. The matrix identity follows by the product rule from the two displayed evolution equations.}
\GSfield{Fate}{The required equations are read in \citet[Section~3, equations~(3.1)--(3.3)]{sasaki2006explicit}. Their elementary consequence validates the notebook formula once its phase-space convention is supplied.}
\GSfield{Present status}{As of 10 November 2026: Established in the stated scope.}
\GSfield{Assessment}{literal statement: correct with missing hypotheses; publication match: partial; verification: read the result; development: none; confidence: high}
\end{GSitem}
\begin{GSitem}{1975-372.11}{A}{Inverse spectral data require more than the traces}
\GSfield{Kind; source}{question; \GSpage{P000618-P000619}; A. M. Perelomov}
\GSfield{The notes}{Every inverse-problem column for types 1--5 carries a question mark. The text says to find matrix entries from traces of powers and says the inverse problem is solved only for Toda, with Stieltjes named.}
\GSfield{Statement}{The traces determine the characteristic polynomial of a finite matrix, not its entries or the full phase point. Even in the rational family, \(q_i\mapsto q_i+c\) leaves \(L\) unchanged. A well-posed reconstruction must add position or eigenvector data and specify permutation, conjugation and center-of-mass conventions. For the rational model, full initial matrices \(Q_0,L_0\) determine the particle positions as the eigenvalues of \(Q_0+tL_0\); for a harmonic trap replace this by \(Q_0\cos\omega t+\omega^{-1}L_0\sin\omega t\), on regular intervals.}
\GSfield{Reconstruction}{The literal trace-only reconstruction is underdetermined. The explicit initial-value formulas are a bounded corrected formulation, not a solution of recovering every matrix from its spectrum. Stieltjes data for a finite Jacobi matrix include spectral weights; the notebook gives no boundary specification making that name an exact periodic-Toda inverse theorem.}
\GSfield{Fate}{\citet[Sections~2--4, especially (2.17)--(2.18) and (3.7)--(3.8)]{sasaki2006explicit} verifies the augmented-data solution. Its historical discussion attributes the \(A\)-type matrix method to Olshanetsky--Perelomov's 1976 papers. Thus the precise initial-value problem was developed, while the trace-only wording remains false.}
\GSfield{Present status}{As of 10 November 2026: Explicit reconstruction is available with adequate initial or spectral data; traces alone cannot determine a phase point.}
\GSfield{Assessment}{literal statement: false as written; publication match: partial; verification: indirect; development: annotated; confidence: high; problem \hyperref[prob:GP-1975-04]{GP-1975-04}}
\end{GSitem}
\begin{GSitem}{1975-372.12}{E}{Higher commuting particle flows}
\GSfield{Kind; source}{claim; \GSpage{P000619}; A. M. Perelomov}
\GSfield{The notes}{Manakov's work is described as a finite-dimensional analogue of the higher KdV hierarchy, with the first two cases treated.}
\GSfield{Statement}{Once a mutually commuting complete collection of spectral Hamiltonians has been established, their Hamiltonian vector fields commute. For the rational \(A\)-type model, \(H_k=\operatorname{tr}(L^{k+1})/(k+1)\) gives higher flows with \(\partial_{t_k}L=[L,M_k]\). Their existence and commutativity require the Poisson statement; isospectrality alone is insufficient.}
\GSfield{Reconstruction}{The report does not identify the dissertation or the exact two cases. Only the rational \(A\)-type hierarchy is reconstructed explicitly; the historical priority claim remains unverified.}
\GSfield{Fate}{The concrete hierarchy and integration formulas are read in \citet[Section~5]{sasaki2006explicit}. This is a partial mathematical identification of the notebook's analogy, without replacing its uncertain Manakov reference by an invented bibliographic match.}
\GSfield{Present status}{As of 10 November 2026: Established in the stated scope.}
\GSfield{Assessment}{literal statement: correct with missing hypotheses; publication match: partial; verification: read the result; development: none; confidence: medium}
\end{GSitem}
\begin{GSitem}{1975-372.13}{A}{Quantum Hamiltonians as conjugated radial Laplacians}
\GSfield{Kind; source}{claim; \GSpage{P000619}; A. M. Perelomov}
\GSfield{The notes}{After \(p=-i\partial_q\), the quantum Hamiltonian is identified with the radial part of a Laplace operator on a symmetric space, particularly for the trigonometric system.}
\GSfield{Statement}{For a Riemannian symmetric space, restriction of the Laplace--Beltrami operator to invariant functions gives a radial differential operator with first-derivative terms. Conjugation by the square root of the radial density removes them and produces a Schrödinger operator of Calogero--Sutherland type, up to an additive constant and scale. The couplings are determined by the restricted-root multiplicities. The compact, noncompact and flat cases lead respectively to trigonometric, hyperbolic and rational interactions. Arbitrary continuous couplings need not come from an actual symmetric space.}
\GSfield{Reconstruction}{Conjugation, the energy shift, root multiplicities and radial domains are missing from the equality in the notes. The statement is a relation of differential expressions; self-adjoint domains and boundary conditions are additional analytic data.}
\GSfield{Fate}{\citet[Section~3, Lemma~1, Proposition~1 and Theorems~1--2]{olshanetskyperelomov1978radial} supplies the precise group-coupling formulation. Received 16 June 1977 and published in 1978, it is a subsequent publication by the intervening speaker.}
\GSfield{Present status}{As of 10 November 2026: Established in the stated scope.}
\GSfield{Assessment}{literal statement: correct with missing hypotheses; publication match: partial; verification: read the result; development: none; confidence: high}
\end{GSitem}
\begin{GSitem}{1975-372.14}{F}{Gelfand’s larger group-integrability programme}
\GSfield{Kind; source}{question; \GSpage{P000619}; I. M. Gelfand}
\GSfield{The notes}{Gelfand asks whether everything can be reduced to group integrability in a larger space.}
\GSfield{Statement}{The line suggests realizing the particle system by reduction of a simpler group or matrix motion. It does not specify the larger phase space, group, moment-map value, equivalence relation, allowed potentials, or meaning of solving the group motion. It therefore does not determine a universal theorem or a yes/no problem with fixed hypotheses.}
\GSfield{Reconstruction}{Hamiltonian reduction is an interpretation of the suggestion, not wording recovered from the note.}
\GSfield{Fate}{A concrete rational \(A\)-type realization is described in \citet[Section~7, equations~(7.1)--(7.4)]{sasaki2006explicit}. This confirms the theme but does not settle the unspecified word \emph{everything}.}
\GSfield{Present status}{As of 10 November 2026: The general programme is not a determinate problem in this record.}
\GSfield{Assessment}{literal statement: not determinable; publication match: thematic only; verification: indirect; development: none; confidence: medium}
\end{GSitem}
\begin{GSitem}{1975-372.15}{B}{Complete integrability for exceptional root systems}
\GSfield{Kind; source}{question; \GSpage{P000619}; A. M. Perelomov}
\GSfield{The notes}{The root-system construction ends: it is still unclear whether the systems for exceptional groups are completely integrable.}
\GSfield{Statement}{For \(R=G_2,F_4,E_6,E_7,E_8\) of rank \(r\), take the canonical phase space away from all poles of
     \[H=\tfrac12(p,p)+\sum_{\alpha\in R_+}c_\alpha V((\alpha,q)),\]
     with positive real couplings constant on Weyl orbits, and \(V\) rational, hyperbolic, trigonometric or untwisted elliptic, with a compatible real slice. The question is existence of \(r\) generically independent integrals in involution. Harmonic confinement is a separate variant and is not included in the verified elliptic theorem below.}
\GSfield{Reconstruction}{The statement fixes the root systems, couplings and regular domain, and distinguishes Liouville integrability from global nonsingular coordinates or an explicit solution formula.}
\GSfield{Fate}{\citet[pp.~3--5 and Table~I]{khastgirsasaki2001liouville} verifies Liouville integrability in this scope, with first public preprint 30 May 2000. Earlier group-coupling cases are covered in \citet[Theorem~1 and Corollary~1]{olshanetskyperelomov1978radial}. These are affirmative developments, not evidence of a remaining exceptional-root obstruction.}
\GSfield{Present status}{As of 10 November 2026: Affirmative in the stated exceptional-root scope; no open remainder is inferred from the 1975 uncertainty.}
\GSfield{Assessment}{literal statement: correct with missing hypotheses; publication match: exact; verification: read the result; development: none; confidence: high; problem \hyperref[prob:GP-1975-03]{GP-1975-03}}
\end{GSitem}
\begin{GSitem}{1975-372.16}{E}{The quantum sine-Gordon and massive Thirring correspondence}
\GSfield{Kind; source}{claim; \GSpage{P000619}; A. M. Perelomov}
\GSfield{The notes}{A short line names Coleman and Schroer for a connection between the Thirring model and sine-Gordon.}
\GSfield{Statement}{The identifiable claim concerns bosonization in \(1+1\) spacetime dimensions: a suitably renormalized sine-Gordon scalar theory corresponds to the zero-charge sector of the massive Thirring fermion theory, with coupling parameters related. At \(\beta^2=4\pi\), Coleman's formulation gives the free massive fermion case. This is a quantum correspondence with a sector and renormalization prescription, not equality of the two classical equations.}
\GSfield{Reconstruction}{Only Coleman's identifiable result is delimited here. The single surname Schroer does not identify a second paper or an independent theorem, and no exact coupling formula absent from the opened abstract is supplied.}
\GSfield{Fate}{\citet[abstract]{coleman1975equivalence} states this correspondence and its free-fermion specialization; the paper was publicly published 15 April 1975. The full proof was not accessed.}
\GSfield{Present status}{As of 10 November 2026: Established in the stated scope.}
\GSfield{Assessment}{literal statement: correct with missing hypotheses; publication match: partial; verification: publisher metadata checked; development: none; confidence: high}
\end{GSitem}
\begin{GSitem}{1975-372.17}{E}{Classical Grassmann coordinates and the graded Poisson bracket}
\GSfield{Kind; source}{construction; \GSpage{P000619-P000620}}
\GSfield{The notes}{Fradkin allows Bose and Fermi canonical variables, writes \(pq+qp=1\) after Grassmann variables, and then introduces a Poisson bracket using right and left derivatives and the parities \(n_A,n_B\).}
\GSfield{Statement}{Classical odd coordinates in a supercommutative algebra satisfy \(\theta_a\theta_b=-\theta_b\theta_a\), so their algebra product does not obey \(pq+qp=1\). That relation belongs to quantized creation and annihilation operators, or to an odd multiplication operator paired with differentiation. An even graded Poisson bracket instead satisfies
     \[\{A,B\}=-(-1)^{|A||B|}\{B,A\},\qquad
       \{A,BC\}=\{A,B\}C+(-1)^{|A||B|}B\{A,C\}.\]
     Canonical even pairs obey \(\{q_i,p_j\}=\delta_{ij}\), while an odd canonical pair has a symmetric bracket \(\{\theta_a,\pi_b\}=\{\pi_b,\theta_a\}=\delta_{ab}\), after choosing left/right derivative conventions.}
\GSfield{Reconstruction}{The algebra product, the classical graded bracket and the quantum anticommutator are separated. This removes the contradiction while preserving the intended mixed canonical formalism.}
\GSfield{Fate}{The basic distinction is algebraic and predates the talk; Berezin is explicitly named in the record. The later general constrained formalism in \citet[abstract, publication lead]{fradkinfradkina1978quantization} matches the theme, but its inaccessible full text is not used to certify the notebook's derivative signs.}
\GSfield{Present status}{As of 10 November 2026: Established in the stated scope.}
\GSfield{Assessment}{literal statement: false as written; publication match: partial; verification: indirect; development: annotated; confidence: high}
\end{GSitem}
\begin{GSitem}{1975-372.18}{E}{Gauge constraints and the two photon polarizations}
\GSfield{Kind; source}{claim; \GSpage{P000619}}
\GSfield{The notes}{Electromagnetism and Yang--Mills theory are examples of constrained systems; light has two states rather than three because gauge invariance removes an extra state.}
\GSfield{Statement}{For free Maxwell theory in \(3+1\)-dimensional Minkowski space, a nonzero null momentum has a transverse polarization space modulo gauge shifts \(\epsilon_\mu\mapsto\epsilon_\mu+\lambda k_\mu\). This quotient is two-dimensional, corresponding to two physical helicities. In Hamiltonian language the primary and Gauss constraints, together with gauge equivalence, remove the unphysical canonical variables. The assertion does not count polarizations of a massive vector field, which has a different constraint structure.}
\GSfield{Reconstruction}{The spacetime dimension, massless free field and gauge quotient are supplied to make the numerical assertion definite.}
\GSfield{Fate}{The constraint and polarization analysis is checked in \citet[Sections~6.1--6.2]{tong2006qft}; this is established Maxwell theory, used as an example rather than a new 1975 result.}
\GSfield{Present status}{As of 10 November 2026: Established in the stated scope.}
\GSfield{Assessment}{literal statement: correct with missing hypotheses; publication match: exact; verification: read the result; development: none; confidence: high}
\end{GSitem}
\begin{GSitem}{1975-372.19}{E}{The singular Legendre map and the constraint algorithm}
\GSfield{Kind; source}{method; \GSpage{P000620}}
\GSfield{The notes}{The momentum is defined by a right derivative of \(L\); failure to solve for \(\dot q\), expressed by a vanishing Hessian determinant, leads to additional constraints and the Dirac procedure.}
\GSfield{Statement}{For an ordinary finite-dimensional Lagrangian \(L(q,\dot q)\), an invertible velocity Hessian gives a local Legendre transformation by the inverse function theorem. A constant-rank singular Legendre map has a locally constrained image. Preservation of primary constraints under Hamiltonian evolution can impose secondary constraints, fix multipliers, or reveal gauge freedom. In a Bose--Fermi system the relevant invertibility is that of a graded linear map; an ordinary determinant is inadequate for a general mixed Hessian.}
\GSfield{Reconstruction}{Constant-rank and local regularity assumptions are supplied. Singular rank-changing Lagrangians are not covered by a statement that all constraints automatically form a finite-dimensional Lie algebra with constant structure coefficients.}
\GSfield{Fate}{The regular Legendre step is recalled in \citet[Appendix~C.4]{cannas2007introduction}. The mixed constrained extension is the subject of the later Fradkin--Fradkina publication lead, whose exact formula has not been read \citep{fradkinfradkina1978quantization}.}
\GSfield{Present status}{As of 10 November 2026: Established in the stated scope.}
\GSfield{Assessment}{literal statement: correct with missing hypotheses; publication match: partial; verification: indirect; development: none; confidence: high}
\end{GSitem}
\begin{GSitem}{1975-372.20}{E}{First- and second-class constraints require a rank condition}
\GSfield{Kind; source}{definition; \GSpage{P000620}}
\GSfield{The notes}{Second-class constraints are described by brackets \(\{\varphi_k,\psi_j\}\ne0\) after imposing the other constraints, with central constraints discarded.}
\GSfield{Statement}{On a regular constraint surface, a first-class function has bracket zero on that surface with every constraint. A selected family \(\chi_a\) is second class when its bracket matrix \(C_{ab}=\{\chi_a,\chi_b\}\) is invertible on the surface (in the graded sense when needed). A nonzero matrix entry alone is insufficient. For purely even second-class constraints,
     \[\{F,G\}_D=\{F,G\}-\{F,\chi_a\}(C^{-1})^{ab}\{\chi_b,G\}\]
     is the Dirac bracket, under the usual regularity assumptions. First-class constraints may close with phase-space-dependent structure functions rather than constants.}
\GSfield{Reconstruction}{The invertibility condition replaces the notebook's entrywise nonvanishing shorthand. The displayed bracket is explicitly restricted to the ungraded case so that no missing super-sign convention is hidden.}
\GSfield{Fate}{This is the Dirac distinction named in the note. The relation with general Bose--Fermi quantization is only indirectly identified through \citet[abstract, publication lead]{fradkinfradkina1978quantization}; the original Dirac book edition is not specified here.}
\GSfield{Present status}{As of 10 November 2026: Established in the stated scope.}
\GSfield{Assessment}{literal statement: false as written; publication match: partial; verification: indirect; development: annotated; confidence: high}
\end{GSitem}
\begin{GSitem}{1975-372.21}{E}{Sources and the formal generating functional}
\GSfield{Kind; source}{construction; \GSpage{P000620}}
\GSfield{The notes}{The displayed \(Z\) integrates fields and canonical momenta against constraint delta functions and an exponential of the Hamiltonian action plus sources. It is called the generating function of all Green functions and an analogue of a Fourier transform.}
\GSfield{Statement}{Formally, with fixed boundary or vacuum prescription and consistent factors of \(i\), one writes \(Z[J]=C^{-1}\int D\mu\,e^{iS+i\langle J,\Phi\rangle}\). Functional derivatives with respect to even and odd sources generate the corresponding ordered correlation functions. The measure \(D\mu\), constraint reduction or gauge fixing, regularization, and source parity must be specified. This symbolic integral is not a countably additive Lebesgue measure on an infinite-dimensional field space, and the line alone does not define all quantum theories.}
\GSfield{Reconstruction}{Only the formal generating-functional mechanism is reconstructed. The notebook's crossed-out measure and uncertain grouping are not silently converted into a rigorous path-integral theorem.}
\GSfield{Fate}{Formal source differentiation gives the generating mechanism; the Wick expansion read in \citet[Section~3.3]{tong2006qft} checks the corresponding perturbative correlation-function structure, not existence of the displayed measure. A more specific subsequent constrained construction is suggested by \citet[abstract, publication lead]{fradkinfradkina1978quantization}, but an exact formula-by-formula identification has not been verified.}
\GSfield{Present status}{As of 10 November 2026: Established in the stated scope.}
\GSfield{Assessment}{literal statement: not determinable; publication match: partial; verification: indirect; development: none; confidence: medium}
\end{GSitem}
\begin{GSitem}{1975-372.22}{F}{The constraint determinant and the ambiguous term A}
\GSfield{Kind; source}{claim; \GSpage{P000620}}
\GSfield{The notes}{The term \(A\) contains trace logarithms of even and odd constraint brackets and a Schur-complement expression. It appears after the exponent's integral, while another exponential and \(\delta(0)/2\) occur inside the definition of \(A\).}
\GSfield{Statement}{A finite-dimensional model for regular second-class reduction inserts the constraint delta function and the square root of the determinant of the bracket matrix; with even and odd constraints the corresponding object is a square root of a Berezinian, with conventions fixed. For a block matrix \(C=\left(\begin{smallmatrix}A&B\\ C_1&D\end{smallmatrix}\right)\) with invertible odd block \(D\),
     \[\operatorname{Ber}C=\frac{\det(A-BD^{-1}C_1)}{\det D}.\]
     This explains the visible Schur complement and difference of trace logarithms. It does not determine the notebook's misplaced exponential, numerical coefficient, sign, or distributional factor. First-class systems also require gauge fixing or its equivalent.}
\GSfield{Reconstruction}{The Berezinian is an explicit algebraic reconstruction of the visible pattern; the complete \(A\) formula remains unresolved. No universal determinant formula for arbitrary quantum field theories is claimed.}
\GSfield{Fate}{The abstract of the 1978 Fradkin--Fradkina lead concerns precisely mixed first- and second-class constraints, but its full text was inaccessible \citep{fradkinfradkina1978quantization}. The exact displayed expression therefore remains a partial, indirect identification rather than a verified publication match.}
\GSfield{Present status}{As of 10 November 2026: The regular determinant/Berezinian mechanism is standard; the exact notebook term A has not been recovered.}
\GSfield{Assessment}{literal statement: not determinable; publication match: thematic only; verification: indirect; development: none; confidence: medium}
\end{GSitem}
\begin{GSitem}{1975-372.23}{F}{The asserted cancellation of delta at zero}
\GSfield{Kind; source}{claim; \GSpage{P000620}}
\GSfield{The notes}{The final line says that \(\delta(0)\) later cancels in the \(S\)-matrix.}
\GSfield{Statement}{No theory, regularization, observable definition, order in perturbation theory, or cancellation mechanism is specified. Consequently the line does not state a determinate cancellation theorem. In particular, a formal trace logarithm at coincident points cannot be assigned a numerical value merely by writing \(\delta(0)\).}
\GSfield{Reconstruction}{The uncertainty concerns the missing model and prescriptions, beyond a possible transcription ambiguity in the preceding \(A\).}
\GSfield{Fate}{The related general constrained-quantization publication lead does not verify this specific cancellation \citep{fradkinfradkina1978quantization}. No exact source match was obtained for the isolated claim.}
\GSfield{Present status}{As of 10 November 2026: No definite universal cancellation statement can be extracted.}
\GSfield{Assessment}{literal statement: not determinable; publication match: none found; verification: unresolved lead; development: none; confidence: low}
\end{GSitem}

\GStalk{373}{1975-10-20}{I. M. Gelfand}{P000620-P000622}{Fractional differential operators and stationary KdV}
Gelfand introduces Krichever's report by recalling the differential parts of fractional Schrödinger powers and the stationary one-gap equation. His final suggestion of an unworked-out noncommutative theory remains a programme rather than a precise conjecture. This brief intervention continues the higher-KdV discussion of 1973 and the finite-gap work of 1974.

\begin{GSitem}{1975-373.1}{E}{Differential parts of fractional powers and the KdV hierarchy}
\GSfield{Kind; source}{construction; \GSpage{P000621}}
\GSfield{The notes}{For \(L=-D^2+u\), \(A=D^3+\cdots\) is described as \(L^{3/2}\) with nondifferential terms cut off. The same operation on \(L^{5/2},L^{7/2},\ldots\) leaves a commutator that is multiplication by a function.}
\GSfield{Statement}{Use the formal pseudodifferential algebra and choose \(R=(-L)^{1/2}=D+O(D^{-1})\). Then \(B_{2j+1}=(R^{2j+1})_+\), the part with nonnegative powers of \(D\), satisfies \([B_{2j+1},L]\) of differential order at most zero. For \(j=1\),
     \[B_3=D^3-\tfrac32uD-\tfrac34u',\qquad [L,B_3]=\tfrac32uu'-\tfrac14u'''.\]
     Thus the corresponding Lax equation gives KdV after a constant time rescaling. The formal higher flows commute; analytic existence for particular function spaces is a separate issue.}
\GSfield{Reconstruction}{The branch and sign are essential: \(L=-D^2+u\) has no formal square root with leading term \(D\). The construction is algebraic, not the spectral fractional power of an unbounded self-adjoint operator.}
\GSfield{Fate}{The hierarchy was already discussed in 1973-326.1 and 1973-327.1, so the present statement is expository. A systematic later fractional-operator treatment is read in \citet[Sections~8--12]{gelfanddikii1976fractional}, received 22 June 1976 and published that year.}
\GSfield{Present status}{As of 10 November 2026: Established in the stated scope.}
\GSfield{Assessment}{literal statement: correct with missing hypotheses; publication match: partial; verification: read the result; development: none; confidence: high; problem \hyperref[prob:GP-1973-53]{GP-1973-53}}
\end{GSitem}
\begin{GSitem}{1975-373.2}{E}{The stationary third-order equation and one-gap potentials}
\GSfield{Kind; source}{claim; \GSpage{P000621}}
\GSfield{The notes}{A one-zone potential is identified by \(K(u)=0\), so that the second- and third-order operators commute; this is called a hypothesis in the notes.}
\GSfield{Statement}{For a real smooth periodic potential, the existence of a normalized third-order commuting operator gives a stationary KdV equation
     \[u'''-6uu'+cu'=0\]
     for a constant \(c\). For example \(A=-4D^3+6uD+3u'-cD\) obeys \([A,-D^2+u]=6uu'-u'''-cu'\). Such periodic stationary solutions have at most one finite open spectral gap; the nonconstant nonsingular elliptic case has one, and constant or degenerate cases must be distinguished. A general one-gap potential need not obey the unshifted equation with \(c=0\).}
\GSfield{Reconstruction}{Real smooth periodicity, the lower-flow constant, and the meaning of one finite open gap are supplied. The algebraic commutator calculation by itself does not establish a band-spectrum statement for arbitrary complex or nonperiodic coefficients.}
\GSfield{Fate}{\citet[Section~2, Theorem~2.1 and the \(n=1\) example on p.~63]{novikov1974periodic} verifies the stationary-to-finite-gap implication before this talk. The converse characterization belongs to the earlier problem GP-1974-10, represented here by its one-gap occurrence; see also 1974-340.7 and 1974-343.2.}
\GSfield{Present status}{As of 10 November 2026: Established in the stated scope.}
\GSfield{Assessment}{literal statement: correct with missing hypotheses; publication match: partial; verification: read the result; development: none; confidence: high; problem \hyperref[prob:GP-1974-10]{GP-1974-10}}
\end{GSitem}
\begin{GSitem}{1975-373.3}{F}{An unspecified noncommutative theory}
\GSfield{Kind; source}{conjecture; \GSpage{P000621}}
\GSfield{The notes}{Gelfand calls the construction a piece of an \ru{несделанной некоммут. теории}; Krichever replies that he in fact started from this.}
\GSfield{Statement}{No algebra, objects, equivalences or proposed theorem are specified. The phrase might refer to noncommuting coefficient rings, the ambient differential-operator algebra, or a broader structural explanation of integrability; these are different mathematical programmes.}
\GSfield{Reconstruction}{None of these alternatives is selected as the uniquely intended problem. The following theory of commutative scalar differential-operator rings cannot be treated as the solution of every possible noncommutative extension.}
\GSfield{Fate}{The adjacent report is identified precisely in 1975-374.1--10. The historical thematic connection is also visible in \citet[introduction and Section~2]{krichever1977methods}, but this does not recover a definite conjecture.}
\GSfield{Present status}{As of 10 November 2026: No determinate conjecture can be extracted from this phrase.}
\GSfield{Assessment}{literal statement: not determinable; publication match: thematic only; verification: indirect; development: none; confidence: low}
\end{GSitem}

\GStalk{374}{1975-10-20}{I. M. Krichever}{P000620-P000625}{Algebraic curves and commuting ordinary differential operators}
Krichever passes from coefficient equations for commuting scalar operators to their algebraic spectral curve, then reverses the construction using a normalized Baker--Akhiezer function. The report includes the polynomial-relation theorem attributed to Shabat, pole-order gaps, a genus criterion, a centralizer intervention, and a theta formula. Several distinctions matter: algebraic dependence was known earlier; rank-one spectral data are required for the scalar eigenfunction construction; a pole divisor is a bound rather than an exact divisor at every \(x\); and the claimed algebraic independence and final theta normalization contain errors. The notebook's date is retained even though Krichever's 1977 retrospective mentions related seminar talks in November--December 1975.

\begin{GSitem}{1975-374.1}{A}{Coefficient recursion for a commuting operator}
\GSfield{Kind; source}{method; \GSpage{P000621}}
\GSfield{The notes}{Expanding \([L_n,L_m]=0\) gives coefficient equations. The first \(m-1\) equations express the \(v_i\) as differential polynomials in the \(u_i\), with integration constants; the first relation is \(v_{m-2}=(m/n)u_{n-2}+h_2\).}
\GSfield{Statement}{Work locally with scalar differential operators over a characteristic-zero differential field, with constant nonzero leading coefficients. After scaling, let
     \(L_n=D^n+\sum_{j=0}^{n-2}u_jD^j\), while \(L_m=D^m+cD^{m-1}+\sum_{j=0}^{m-2}v_jD^j\), with \(n,m\geq2\) and \(c\) constant. Equating the coefficient of \(D^{n+m-3}\) gives
     \(nv'_{m-2}-mu'_{n-2}=0\), hence the displayed relation. Recursive comparison determines the lower \(v_j\) as differential polynomials with constants; the remaining \(n-1\) equations constrain the \(u_j\). For \(n=2\) and odd \(m\), these are stationary higher-KdV equations.}
\GSfield{Reconstruction}{The monic normalization and vanishing \(D^{n-1}\) coefficient are supplied. One common scalar gauge can remove that coefficient of \(L_n\); it need not simultaneously remove the constant \(D^{m-1}\) coefficient of an arbitrary commuting partner.}
\GSfield{Fate}{This elimination is part of the systematic formal-operator framework in \citet[Sections~8--12]{gelfanddikii1976fractional}; Krichever's normalization is explicit in \citet[Section~2.1]{krichever1977methods}. The subsequent algebraic-geometric solution is more than this coefficient count.}
\GSfield{Present status}{As of 10 November 2026: Established in the stated scope.}
\GSfield{Assessment}{literal statement: correct with missing hypotheses; publication match: partial; verification: read the result; development: none; confidence: high}
\end{GSitem}
\begin{GSitem}{1975-374.2}{E}{Polynomial dependence of commuting scalar differential operators}
\GSfield{Kind; source}{theorem; \GSpage{P000621-P000622}}
\GSfield{The notes}{The theorem attributed to A. B. Shabat says \([L_n,L_m]=0\Rightarrow P(L_n,L_m)=0\). Its proof uses \(L_m\) on the solution space of \(L_ny=Ey\), but calls a polynomial of degree \(n+1\) the characteristic polynomial. Bernstein mentions an earlier polynomial-coefficient result without a title.}
\GSfield{Statement}{For commuting scalar ordinary differential operators of positive orders \(n,m\), with complex analytic germs as coefficients and nonvanishing leading coefficients, there is a nonzero \(P(E,w)\in\mathbb C[E,w]\) such that \(P(L_n,L_m)=0\). The solution space of \(L_ny=Ey\) has dimension \(n\). In its initial-jet basis, the induced matrix \(M(E)\) of \(L_m\) has entries polynomial in \(E\), by repeated substitution of the differential equation. Thus \(P(E,w)=\det(wI-M(E))\) has degree \(n\) in \(w\), and Cayley--Hamilton gives the operator identity. Coprimality of \(n,m\) is unnecessary.}
\GSfield{Reconstruction}{\ours\ The condensed initial-jet argument reconstructs the proof read in Krichever's Theorem~2.1; the notebook's exponential asymptotic sketch is not needed. The degree \(n+1\) is corrected to \(n\). The untitled 1967 reference mentioned by Bernstein was not independently identified.}
\GSfield{Fate}{\citet[Theorem~1]{amitsur1958commutative} already yields algebraic dependence. The notebook proof is reproduced as \citet[Theorem~2.1]{krichever1977methods}; the Shabat attribution does not establish general priority.}
\GSfield{Present status}{As of 10 November 2026: Established in the stated scope.}
\GSfield{Assessment}{literal statement: false as written; publication match: partial; verification: read the result; development: annotated; confidence: high}
\end{GSitem}
\begin{GSitem}{1975-374.3}{A}{The spectral curve and a normalized common eigenfunction}
\GSfield{Kind; source}{construction; \GSpage{P000622}}
\GSfield{The notes}{The eigenvalues satisfy \(P(E,w)=0\), and a common eigenfunction normalized by \(\psi(x_0,P)=1\) is viewed as a function on the resulting algebraic curve.}
\GSfield{Statement}{Take an irreducible component of the joint spectral curve and, when necessary, its normalization. If the generic simultaneous eigenspace is one-dimensional, the eigenline can locally be normalized by evaluation at \(x_0\), away from points where that evaluation vanishes. This gives a common eigenfunction meromorphic in the spectral point off the distinguished point at infinity, with a fixed divisor bounding its finite poles. A commuting pair of coprime positive orders is a sufficient rank-one situation. Higher generic rank requires vector-bundle data and does not yield one scalar normalized eigenfunction carrying the full information.}
\GSfield{Reconstruction}{Generic rank, normalization of a possibly singular affine curve, and the excluded normalization zeros are supplied. The scalar construction is not asserted for every commuting pair merely because a polynomial relation exists.}
\GSfield{Fate}{The rank-one and coprime-order argument is read in \citet[Section~2.2]{krichever1977methods}. It was developed in the speaker's 1976--1977 publications following these seminar reports.}
\GSfield{Present status}{As of 10 November 2026: Established in the stated scope.}
\GSfield{Assessment}{literal statement: correct with missing hypotheses; publication match: partial; verification: read the result; development: none; confidence: high}
\end{GSitem}
\begin{GSitem}{1975-374.4}{A}{Existence and uniqueness of the one-point Baker--Akhiezer function}
\GSfield{Kind; source}{theorem; \GSpage{P000622}}
\GSfield{The notes}{On a nonsingular algebraic curve \(X\) with marked \(P_0\), the notes prescribe an \(x\)-independent nonspecial divisor \(D\) of degree \(g\) and the expansion \(e^{k(x-x_0)}(1+\sum\xi_i(x)k^{-i})\). The local coordinate \(1/k\) is mistakenly said to be near \(x_0\).}
\GSfield{Statement}{Let \(X\) be a smooth compact complex curve of genus \(g\), \(P_\infty\in X\), \(z=k^{-1}\) a local coordinate there, and \(D\) a nonspecial effective divisor of degree \(g\), disjoint from \(P_\infty\). For \(x\) near \(x_0\), except possible exceptional parameter values, there is a unique function \(\psi(x,P)\), meromorphic on \(X\setminus\{P_\infty\}\) with poles bounded by \(D\), such that
     \[e^{-k(x-x_0)}\psi(x,P)=1+\sum_{j\geq1}\xi_j(x)k^{-j}\]
     is holomorphic at \(P_\infty\). At \(x=x_0\), nonspeciality and normalization give \(\psi=1\). For more general \(x\), the dependence may be meromorphic and the construction is used away from its parameter poles.}
\GSfield{Reconstruction}{Compactness, disjointness and exceptional \(x\)-values are made explicit. The divisor bounds the poles: it cannot be the exact polar divisor for every \(x\), since \(\psi(x_0,\cdot)=1\).}
\GSfield{Fate}{The lemma is read in \citet[pp.~291--292]{krichever1976zakharov}. Receipt on 8 October 1975 and presentation on 24 October do not independently prove public availability before the notebook's 20 October report.}
\GSfield{Present status}{As of 10 November 2026: Established in the stated scope.}
\GSfield{Assessment}{literal statement: correct with missing hypotheses; publication match: partial; verification: read the result; development: none; confidence: high}
\end{GSitem}
\begin{GSitem}{1975-374.5}{A}{The function ring acts by commuting differential operators}
\GSfield{Kind; source}{theorem; \GSpage{P000623}}
\GSfield{The notes}{Every \(E\in\Lambda(X,P_0)\), meromorphic only at \(P_0\), is assigned a unique \(L_E\) with \(L_E\psi=E\psi\). Matching principal parts determines its coefficients, and \(E\mapsto L_E\) is a ring homomorphism.}
\GSfield{Statement}{For the smooth Baker--Akhiezer data of 1975-374.4, let \(A(X,P_\infty)\) be the ring of meromorphic functions with poles only at \(P_\infty\). For pole order \(n\geq2\), if
     \(E=c_nk^n+c_{n-1}k^{n-1}+c_{n-2}k^{n-2}+\cdots\), there is a unique scalar differential operator
     \(L_E=c_nD^n+c_{n-1}D^{n-1}+\sum_{j=0}^{n-2}u_j(x)D^j\)
     satisfying \(L_E\psi=E\psi\), with coefficients meromorphic in the allowed \(x\)-domain. For example,
     \(u_{n-2}=c_{n-2}-nc_n\xi'_1\).
     Constants act by \(L_E=c_0\); for pole order one the formula is \(L_E=c_1D+c_0\).
     The map is injective and obeys \(L_{E+F}=L_E+L_F\), \(L_{EF}=L_EL_F\); hence its image is commutative and order equals pole order for nonconstant \(E\).}
\GSfield{Reconstruction}{After matching principal parts, the remainder has zero normalized leading term and the same allowed poles; Baker--Akhiezer uniqueness makes it vanish. The index bounds printed for \(u_i=Q(c_j,\xi_s^{(r)})\) cannot hold literally: already \(u_0\) for \(n=2\) depends on \(\xi'_1\).}
\GSfield{Fate}{The construction is read in \citet[pp.~291--292]{krichever1976zakharov}; the coefficient example and ring statement are explicit in \citet[Sections~1--2]{krichever1977methods}.}
\GSfield{Present status}{As of 10 November 2026: Established in the stated scope.}
\GSfield{Assessment}{literal statement: false as written; publication match: partial; verification: read the result; development: annotated; confidence: high}
\end{GSitem}
\begin{GSitem}{1975-374.6}{E}{Pole-order gaps and the genus sum}
\GSfield{Kind; source}{claim; \GSpage{P000623}}
\GSfield{The notes}{Let \(n\) be the least positive pole order and \(s+\mu_sn\) an available order in residue class \(s\); then \(\sum_{s=1}^{n-1}\mu_s=g\).}
\GSfield{Statement}{For a smooth genus-\(g\) curve and marked point, let \(H\subset\mathbb Z_{\geq0}\) be its Weierstrass semigroup of allowable pole orders. If \(n\) is its least positive element and \(w_s=s+\mu_sn\) is the \emph{least} element of \(H\) congruent to \(s\pmod n\), then the gaps in this residue class are \(s,s+n,\ldots,s+(\mu_s-1)n\). The gap theorem gives exactly \(g\) gaps, so \(\sum_s\mu_s=g\). Merely choosing an arbitrary available order in each residue class would not give this equality.}
\GSfield{Reconstruction}{\ours\ The residue-class count is supplied explicitly. Minimality in each residue class is essential and is restored. The same semigroup records orders of the operators from 1975-374.5.}
\GSfield{Fate}{This is the classical Riemann--Roch gap count applied to the function ring used in \citet[Section~2.2]{krichever1977methods}; the elementary residue-class count above explains the exact notebook formula.}
\GSfield{Present status}{As of 10 November 2026: Established in the stated scope.}
\GSfield{Assessment}{literal statement: correct with missing hypotheses; publication match: partial; verification: read the result; development: none; confidence: high}
\end{GSitem}
\begin{GSitem}{1975-374.7}{E}{The claimed algebraic independence is impossible}
\GSfield{Kind; source}{claim; \GSpage{P000623}}
\GSfield{The notes}{At a point in general position the notes assert \(g\) algebraically independent commuting operators.}
\GSfield{Statement}{The function field of a curve has transcendence degree one. Therefore any two elements of the associated scalar differential-operator ring are algebraically dependent, and for \(g>1\) the asserted \(g\)-fold algebraic independence is impossible. At an ordinary point the gaps are \(1,\ldots,g\); one can instead choose algebra generators with pole orders \(g+1,\ldots,2g+1\), subject to algebraic relations. An alternative module description uses one pole-order-\(g+1\) function and a basis over its polynomial ring.}
\GSfield{Reconstruction}{Algebra generators, module generators and algebraically independent elements are different notions. The notebook does not identify which of the first two was intended, so its numeral is not silently emended.}
\GSfield{Fate}{The contradiction already follows from \citet[Theorem~1]{amitsur1958commutative} and the polynomial dependence of 1975-374.2. The curve-ring construction supplies the corrected interpretation.}
\GSfield{Present status}{As of 10 November 2026: Established in the stated scope.}
\GSfield{Assessment}{literal statement: false as written; publication match: partial; verification: read the result; development: annotated; confidence: high}
\end{GSitem}
\begin{GSitem}{1975-374.8}{A}{A genus criterion requires more than a list of operator orders}
\GSfield{Kind; source}{claim; \GSpage{P000623-P000624}}
\GSfield{The notes}{An operator is called algebraic of genus \(g\) if it arises from the curve construction. An equivalence is then asserted with the existence of commuting operators of orders \(s+\mu_sn\), where \(\sum\mu_s=g\); the proof of sufficiency is only a short sketch.}
\GSfield{Statement}{In the smooth rank-one construction, the least pole orders in each residue class provide the stated commuting partners and recover the genus by the gap count. In the reverse direction, a commutative ring of scalar operators gives a spectral curve and eigenfunction data. To identify it with the preceding smooth-curve construction one must control rank, the full ring rather than an arbitrary subalgebra, the compactification and possible singularities, and the divisor and normalization at infinity. A list of available orders is not itself this spectral data, and arbitrary nonminimal \(\mu_s\) do not determine genus.}
\GSfield{Reconstruction}{This is a partial reconstruction of the claimed criterion. No exact if-and-only-if theorem with only the notebook's numerical hypotheses is asserted; in particular arithmetic and geometric genus cannot be silently interchanged in a singular case.}
\GSfield{Fate}{The relevant classification appears in \citet[Section~2.2, Theorem~2.2]{krichever1977methods}. It verifies the spectral-data framework, not the exact abbreviated numerical criterion here.}
\GSfield{Present status}{As of 10 November 2026: The rank-one spectral classification is established; the notebook’s bare genus criterion remains only partially reconstructed.}
\GSfield{Assessment}{literal statement: not determinable; publication match: partial; verification: indirect; development: none; confidence: medium}
\end{GSitem}
\begin{GSitem}{1975-374.9}{E}{Commutativity of the centralizer}
\GSfield{Kind; source}{theorem; \GSpage{P000624}; I. M. Krichever; participant recorded as Limanov}
\GSfield{The notes}{The report remarks that operators commuting with a pair commute with each other. A participant recorded as Limanov adds the polynomial-coefficient version: two operators commuting with a third commute mutually.}
\GSfield{Statement}{For scalar ordinary differential operators over a characteristic-zero differential field with algebraically closed constant field, the centralizer of a positive-order operator is commutative. Thus two operators commuting with such an operator commute with one another; a pair is unnecessary. The qualifier matters: if the third operator is merely a scalar constant, every operator commutes with it, and the conclusion is false. For example \(D\) and multiplication by \(x\) both commute with \(I\), but \([D,x]=1\).}
\GSfield{Reconstruction}{The hypothesis of a positive-order operator is supplied, covering the adjacent \(L_n\) setting. The tentatively read surname is retained without identifying a particular person or a new polynomial-coefficient paper.}
\GSfield{Fate}{\citet[Theorem~1]{amitsur1958commutative} proves this stronger centralizer statement in 1958. The seminar intervention is expository in that setting.}
\GSfield{Present status}{As of 10 November 2026: Established in the stated scope.}
\GSfield{Assessment}{literal statement: correct with missing hypotheses; publication match: exact; verification: read the result; development: none; confidence: high}
\end{GSitem}
\begin{GSitem}{1975-374.10}{A}{The normalized theta formula and its missing factors}
\GSfield{Kind; source}{construction; \GSpage{P000624}}
\GSfield{The notes}{The final formula has an exponential \(\exp\int_{P_0}^{P}\omega\) without \(x\), a theta quotient shifted by \(Ux\), and a theta series whose factors of \(2\pi i\) are inconsistent with the normalized periods.}
\GSfield{Statement}{Choose holomorphic differentials with \(a\)-periods \(\delta_{ij}\), period matrix \(B\), and Abel map \(\mathcal A(P_\infty)=0\). Set
     \[\theta(z\mid B)=\sum_{m\in\mathbb Z^g}\exp(\pi i\,m^TBm+2\pi i\,m^Tz),\qquad \operatorname{Im}B>0.\]
     Let \(d\Omega\) have zero \(a\)-periods, only a double pole at \(P_\infty\), and principal part \(dk=-dz/z^2\); normalize its regularized integral by \(\Omega(P)=k+O(k^{-1})\). Put \(U_j=(2\pi i)^{-1}\int_{b_j}d\Omega\), and choose \(Z\) so that \(\theta(\mathcal A(P)+Z)\) has zero divisor \(D\). With \(t=x-x_0\), on the set where the denominator is nonzero,
     \[\psi(x,P)=e^{t\Omega(P)}
       \frac{\theta(\mathcal A(P)+tU+Z)\,\theta(Z)}
       {\theta(\mathcal A(P)+Z)\,\theta(tU+Z)}.\]
     This has the required normalized exponential expansion and \(\psi(x_0,P)=1\).}
\GSfield{Reconstruction}{The parameter factor in the exponential, the normalization factor depending only on \(x\), the regularization at the pole, the sign of the second-kind principal part, and the consistent theta convention are all supplied explicitly. They cannot be recovered by treating an integral starting at its double pole as an ordinary convergent integral.}
\GSfield{Fate}{The theta construction and normalization at the marked point are read in \citet[Section~1, pp.~189--190]{krichever1977methods}; the notebook has omitted essential factors.}
\GSfield{Present status}{As of 10 November 2026: Established in the stated scope.}
\GSfield{Assessment}{literal statement: false as written; publication match: partial; verification: read the result; development: annotated; confidence: high}
\end{GSitem}

\GStalk{375}{1975-10-27}{B. A. Dubrovin}{P000623-P000628}{Matrix finite-gap operators and Jacobians}
The dated heading on P000625 begins Dubrovin's talk; P000623--P000624 continue Krichever's preceding talk. Dubrovin constructs matrix hierarchies, reduces their stationary equations to a polynomial matrix Lax equation, and identifies generic isospectral sets with a diagonal-group bundle over a Jacobian. The actual mathematical span is P000625--P000626. A sign reversal in the commutator equation is a notebook error, not a change of convention that can consistently be retained.

\begin{GSitem}{1975-375.1}{A}{Regular diagonal spectral pencil and diagonal gauge}
\GSfield{Kind; source}{construction; \GSpage{P000625}}
\GSfield{The notes}{The pencil is \(D+U-EA\), with \(A=\operatorname{diag}(a_1,\ldots,a_n)\), distinct \(a_i\), and \(U_i^i=0\), which the notes say can always be achieved.}
\GSfield{Statement}{On an interval, for a smooth complex matrix \(U\) and constant diagonal \(A\) with pairwise distinct entries, a diagonal invertible gauge \(g\) satisfying \(g_i^{\prime}=-U_{ii}g_i\) transforms \(g^{-1}(D+U-EA)g\) to a pencil with zero diagonal potential. A periodic gauge exists only when the diagonal holonomies permit it; the local assertion does not automatically preserve the periodic boundary condition.}
\GSfield{Reconstruction}{The gauge equation and its global obstruction are explicit reconstructions; neither qualification is written.}
\GSfield{Fate}{The regular-diagonal matrix setup appears in \citet[Ch.~III, \S2]{dubrovinmatveevnovikov1976nonlinear}; the paper was received 2 June 1975 and published in 1976. Receipt is not evidence of public circulation before this talk.}
\GSfield{Present status}{As of 10 November 2026, the result is established in the stated scope.}
\GSfield{Assessment}{literal statement: correct with missing hypotheses; publication match: exact; verification: read the result; development: annotated; confidence: high}
\end{GSitem}
\begin{GSitem}{1975-375.2}{A}{Formal diagonal resolvents and commuting matrix flows}
\GSfield{Kind; source}{construction; \GSpage{P000625}}
\GSfield{The notes}{For diagonal constant \(B\), the notes seek \(\lambda_B=B+\sum_{j\ge1}\lambda_{j,B}E^{-j}\) satisfying \(\lambda_B^{\prime}=[\lambda_B,U-EA]\), and set \(U_t=[A,\lambda_{N+1,B}]\).}
\GSfield{Statement}{Fix the formal diagonalization normalized to the identity at \(E=\infty\), and fix its diagonal integration constants. The equation determines the off-diagonal coefficients recursively because \(\operatorname{ad}A\) is invertible there. With the usual normalized diagonal resolvent, the resulting coefficients are differential polynomials in \(U\), and truncations give the matrix hierarchy. The displayed differential equation alone, without normalization, does not select a unique resolvent.}
\GSfield{Reconstruction}{Normalization and the distinction between off-diagonal recursion and diagonal constants have been supplied; the notebook asserts polynomial dependence without giving that construction.}
\GSfield{Fate}{Compare the normalized construction in \citet[Ch.~III, \S2]{dubrovinmatveevnovikov1976nonlinear} and its matrix elaboration in \citet[\S1]{dubrovin1977matrix}.}
\GSfield{Present status}{As of 10 November 2026, the result is established in the stated scope.}
\GSfield{Assessment}{literal statement: correct with missing hypotheses; publication match: exact; verification: read the result; development: annotated; confidence: high}
\end{GSitem}
\begin{GSitem}{1975-375.3}{A}{Isospectral Lax evolution}
\GSfield{Kind; source}{theorem; \GSpage{P000625}}
\GSfield{The notes}{All displayed flows are said to preserve the spectrum, since \(L_t=[L,M_N]\).}
\GSfield{Statement}{For the spectral pencil, the auxiliary operator must give the compatible pair \([\partial_t-M_N,\partial_x+U-EA]=0\), with signs chosen consistently. For a periodic smooth potential and periodic auxiliary coefficients, the monodromy evolves by conjugation, so its characteristic polynomial and Floquet multipliers are constant. On a noncompact domain, preservation of an operator spectrum additionally requires a specified realization and suitable decay or boundedness conditions.}
\GSfield{Reconstruction}{The notebook leaves the spectral meaning and domain unspecified; the periodic monodromy formulation makes one exact assertion that its calculation supports.}
\GSfield{Fate}{This hierarchy is treated in \citet[Ch.~III, \S2]{dubrovinmatveevnovikov1976nonlinear}.}
\GSfield{Present status}{As of 10 November 2026, the result is established in the stated scope.}
\GSfield{Assessment}{literal statement: correct with missing hypotheses; publication match: exact; verification: read the result; development: annotated; confidence: high}
\end{GSitem}
\begin{GSitem}{1975-375.4}{A}{Stationary equations and the polynomial matrix Lax equation}
\GSfield{Kind; source}{construction; \GSpage{P000625-P000626}}
\GSfield{The notes}{The stationary condition is \([L,M_N]=0\). The talk first writes \(\Lambda^{\prime}=[\Lambda,Q]\), but its lemma and proof reverse the commutator to \([Q,\Lambda]\), with \(Q=U-EA\).}
\GSfield{Statement}{Restrict the commuting auxiliary differential operator to the local solution space of \((D+Q)\psi=0\); elimination of derivatives yields a polynomial matrix \(\Lambda(E,x)\). The condition \([D+Q,\Lambda]=0\) is exactly \(\Lambda^{\prime}=[\Lambda,Q]\). Hence \(\det(zI-\Lambda(E,x))\) is independent of \(x\).}
\GSfield{Reconstruction}{\oursComputed\ The first sign in the notebook is retained and the two later signs are emended by direct expansion of the commutator.}
\GSfield{Fate}{The stationary matrix reduction is part of \citet[Ch.~III, \S2]{dubrovinmatveevnovikov1976nonlinear}. The sign correction is checked algebraically from the notebook's own definition of \(Q\).}
\GSfield{Present status}{As of 10 November 2026, the result is established in the stated scope.}
\GSfield{Assessment}{literal statement: false as written; publication match: exact; verification: read the result; development: annotated; confidence: high}
\end{GSitem}
\begin{GSitem}{1975-375.5}{A}{Algebraic finite-gap eigenfunction and spectral curve}
\GSfield{Kind; source}{definition; \GSpage{P000626}}
\GSfield{The notes}{A polynomial pencil is called finite-gap when an eigenfunction is meromorphic on a fixed finite-genus covering of the \(E\)-plane, with essential singularities allowed over infinity. The proposed curve is \(R(z,E)=\det(zI-\Lambda(E))=0\).}
\GSfield{Statement}{Take the normalization and compactification of the reduced irreducible spectral curve, with marked points above \(E=\infty\). On a regular component with generically simple eigenvalue, a common solution \((D+Q)\psi=0\), \(\Lambda\psi=z\psi\) defines the algebraic eigenline data; its normalized eigenfunction is meromorphic away from those points with prescribed exponential behavior there. This algebraic use of finite-gap is not itself a statement that a general matrix operator has a self-adjoint spectrum with finitely many real gaps.}
\GSfield{Reconstruction}{Reduction, normalization, generic simplicity and the distinction from a real spectral-gap assertion are reconstructed hypotheses.}
\GSfield{Fate}{For the matrix spectral-curve construction see \citet[Ch.~III, \S2]{dubrovinmatveevnovikov1976nonlinear}.}
\GSfield{Present status}{As of 10 November 2026, the result is established in the stated scope.}
\GSfield{Assessment}{literal statement: correct with missing hypotheses; publication match: exact; verification: read the result; development: annotated; confidence: high}
\end{GSitem}
\begin{GSitem}{1975-375.6}{A}{Bloch property of the common eigenfunction}
\GSfield{Kind; source}{claim; \GSpage{P000626}}
\GSfield{The notes}{If \(U(x+T)=U(x)\), the simultaneous eigenfunction is said to be a Bloch function because translation commutes with \(\Lambda\).}
\GSfield{Statement}{If both \(Q\) and the polynomial matrix \(\Lambda\) are \(T\)-periodic, their monodromy commutes with \(\Lambda(E,0)\). On each one-dimensional eigenspace at a simple point of the spectral curve, monodromy therefore acts by a scalar: \(\psi(x+T,P)=\mu(P)\psi(x,P)\). At a repeated eigenvalue one must choose a common eigenvector or allow generalized spectral data; an arbitrary vector in the repeated eigenspace need not be Bloch.}
\GSfield{Reconstruction}{The simple-eigenline qualification and exceptional-point distinction make explicit the linear algebra behind the notebook remark.}
\GSfield{Fate}{The Bloch/spectral-curve interpretation belongs to \citet[Ch.~III, \S2]{dubrovinmatveevnovikov1976nonlinear}.}
\GSfield{Present status}{As of 10 November 2026, the result is established in the stated scope.}
\GSfield{Assessment}{literal statement: correct with missing hypotheses; publication match: exact; verification: read the result; development: annotated; confidence: high}
\end{GSitem}
\begin{GSitem}{1975-375.7}{A}{Complete integrability of stationary matrix systems}
\GSfield{Kind; source}{theorem; \GSpage{P000625-P000626}}
\GSfield{The notes}{The stationary equations are called completely integrable Hamiltonian systems. Coefficients \(I_{ij}\) of \(R(z,E)\) are asserted to be commuting integrals.}
\GSfield{Statement}{For the notebook's matrix equation \(\Lambda^{\prime}=[\Lambda,U-EA]\), direct differentiation proves only that the coefficients of \(\det(zI-\Lambda(E))\) are constants of motion. A claim of Liouville integrability additionally requires a specified Poisson manifold or symplectic leaf, commuting integrals and a rank count. The notebook supplies neither the Poisson bracket nor that phase space; consequently its full Hamiltonian assertion is only partially identifiable here. No complete modern integrability theorem for unspecified stationary data is substituted for it.}
\GSfield{Reconstruction}{The exact consequence of the written Lax equation is separated from the stronger publication theorem. The latter has not been reproduced with all of its phase-space and independence hypotheses, so publication matching remains partial.}
\GSfield{Fate}{The matrix integrability theorem is stated in \citet[Ch.~III, \S2, p.~117]{dubrovinmatveevnovikov1976nonlinear}; further details appear in \citet[\S\S1--3]{dubrovin1977matrix}.}
\GSfield{Present status}{As of 10 November 2026, matrix Hamiltonian integrability theorems are established in the cited works; the notebook's unqualified version is not a fully specified theorem.}
\GSfield{Assessment}{literal statement: not determinable; publication match: partial; verification: read the result; development: annotated; confidence: high}
\end{GSitem}
\begin{GSitem}{1975-375.8}{A}{Recovering matrix potentials with a fixed spectral curve}
\GSfield{Kind; source}{problem; \GSpage{P000626}}
\GSfield{The notes}{The explicit problem asks for all finite-gap \(L\) with the same curve \(\Gamma\). The announced answer is a principal \(\pi\)-bundle over its Jacobian, where \(\pi\) is the diagonal invertible group modulo scalar matrices.}
\GSfield{Statement}{Fix the regular diagonal leading matrix, the polynomial spectral invariants, marked points and local spectral normalization. For a smooth generic spectral curve, the nonsingular rank-one spectral data parametrize an open part of its Jacobian; eigenvector normalization retains a diagonal gauge torsor \(\pi\cong(\mathbb C^*)^{n-1}\). The reconstruction generally produces meromorphic potentials. The assertion must be restricted to the regular spectral-data locus and does not identify every singular or globally smooth real isospectral set with an entire compact Jacobian.}
\GSfield{Reconstruction}{The open-locus and meromorphic qualifications, and the distinction between quotient and framed data, are explicit modern reconstructions. Dubrovin's published quotient notation is \(M_\Gamma/\pi=J(\Gamma)\); the open regular locus here describes the nonsingular-potential chart used in this reconstruction, not a verbatim restriction claimed for his entire algebraic parameter space.}
\GSfield{Fate}{The answer is in \citet[Ch.~III, \S2, p.~117]{dubrovinmatveevnovikov1976nonlinear} and \citet[theorem on p.~269]{dubrovin1977matrix}. The latter was received 7 December 1976 and printed in volume 11 (1977); these dates must not be replaced by its later translation copyright.}
\GSfield{Present status}{As of 10 November 2026, the result is established in the stated scope.}
\GSfield{Assessment}{literal statement: correct with missing hypotheses; publication match: partial; verification: read the result; development: annotated; confidence: high; problem \hyperref[prob:GP-1975-11]{GP-1975-11}}
\end{GSitem}

\GStalk{376}{1975-11-10}{I. M. Gelfand}{P000626-P000629}{Geometry and mathematical preparation for physicists}
The heading and informal lecture occur on P000627, continuing at the top of P000628 before Polyakov begins. An ant on a cube motivates curvature and torsion, followed by a programme for geometry and relativity on lattices, a physics reading list, and mathematics proposed for physicists. Bare subject names are recorded as a curriculum remark rather than converted into nonexistent conjectures.

\begin{GSitem}{1975-376.1}{E}{Development of loops, curvature and torsion}
\GSfield{Kind; source}{remark; \GSpage{P000627}}
\GSfield{The notes}{An ant moving on a cube and a curve developed from a plane are used to distinguish a displacement from a rotation, labelled torsion and curvature.}
\GSfield{Statement}{For an affine connection with solder form \(\theta\) and connection form \(\omega\), torsion is \(T=d\theta+\omega\wedge\theta\), while curvature is \(\Omega=d\omega+\omega\wedge\omega\). Infinitesimal affine development relates their translational and rotational effects. An arbitrary finite loop's developed displacement is not by itself a torsion tensor; the intrinsic Levi-Civita geometry of a Euclidean polyhedral surface is torsion-free, with curvature concentrated in angle deficits at vertices.}
\GSfield{Reconstruction}{The differential-form formulation and the warning about finite-loop development are explanatory reconstruction. The sketch does not specify a connection with nonzero torsion.}
\GSfield{Fate}{This is an informal exposition of connection geometry. \citet[abstract and polyhedral construction]{regge1961relativity} gives the relevant pre-talk polyhedral curvature setting, but does not verify every assertion of the ant analogy.}
\GSfield{Present status}{As of 10 November 2026, the result is established in the stated scope.}
\GSfield{Assessment}{literal statement: not determinable; publication match: thematic only; verification: publisher metadata checked; development: annotated; confidence: high}
\end{GSitem}
\begin{GSitem}{1975-376.2}{E}{Characteristic classes and basic geometry on polyhedra}
\GSfield{Kind; source}{question; \GSpage{P000627}}
\GSfield{The notes}{Gelfand asks how to build basic geometry, including characteristic classes, on higher-dimensional polyhedra.}
\GSfield{Statement}{The notes specify no category of polyhedra, tangential structure, or desired definition. For a triangulated closed smooth \(n\)-manifold there is a precise existing instance: the mod-two sum of all \(p\)-simplices of its first barycentric subdivision represents the Poincare dual of \(w_{n-p}(TM)\). This example is evidence about the programme, not an identification of a new theorem being proposed.}
\GSfield{Reconstruction}{The restricted manifold example is supplied to locate the discussion; no assertion is made for arbitrary singular polyhedra.}
\GSfield{Fate}{The relevant pre-talk instance is \citet[Theorem 1 and its mod-two corollary, pp.~511--512]{halperintoledo1972stiefel}; the same construction is discussed explicitly in Fuchs's intervention on P000635.}
\GSfield{Present status}{As of 10 November 2026, the result is established in the stated scope.}
\GSfield{Assessment}{literal statement: correct with missing hypotheses; publication match: partial; verification: read the result; development: annotated; confidence: high}
\end{GSitem}
\begin{GSitem}{1975-376.3}{E}{General relativity on a lattice}
\GSfield{Kind; source}{problem; \GSpage{P000627}}
\GSfield{The notes}{The notebook explicitly says: \ru{Задача: сделать общую относительность на решетках}.}
\GSfield{Statement}{One can assign edge lengths to a simplicial complex, use flat metrics on its top-dimensional simplices, and concentrate curvature on codimension-two hinges. In four dimensions the Regge action is proportional to the sum of hinge areas times deficit angles. A mathematically definite new problem would have to specify a discretization, approximation class, limit, and convergence criterion; none is given here.}
\GSfield{Reconstruction}{The simplicial realization is a historical comparison, not evidence that Gelfand specified that model or a convergence theorem.}
\GSfield{Fate}{\citet[abstract]{regge1961relativity} had already proposed a polyhedral formulation of general relativity in 1961. The opened publisher record confirms that prior programme; this broad seminar request is not assigned a new open-problem identifier.}
\GSfield{Present status}{As of 10 November 2026, the historical request has established realizations; the notes do not specify a particular remaining convergence problem.}
\GSfield{Assessment}{literal statement: not determinable; publication match: partial; verification: publisher metadata checked; development: annotated; confidence: high}
\end{GSitem}
\begin{GSitem}{1975-376.4}{E}{Reciprocal mathematics and physics curricula}
\GSfield{Kind; source}{remark; \GSpage{P000627-P000628}}
\GSfield{The notes}{The physics list names Feynman--Hibbs, Feynman's lectures and statistical mechanics, Landau--Lifshitz, Berestetskii--Lifshitz--Pitaevskii, and Dirac. The mathematics list includes representations, pseudodifferential and Fourier integral operators, Riemann--Roch, singularities and topological cycles.}
\GSfield{Statement}{This is a pedagogical proposal to teach established mathematical and physical subjects across the two communities. It states no common theorem, prerequisite graph, or new research question. The separate functional equation for a polynomial is indexed in the next item.}
\GSfield{Reconstruction}{No individual textbook edition is reconstructed from abbreviated names, and bare topic names are not inflated into conjectures.}
\GSfield{Fate}{The notebook is the direct evidence for the curriculum. Its sole displayed mathematical theorem is independently matched below to \citet{bernstein1972continuation}.}
\GSfield{Present status}{As of 10 November 2026, this remains a historical teaching programme, with no theorem-level status to resolve.}
\GSfield{Assessment}{literal statement: correct; publication match: thematic only; verification: publisher metadata checked; development: annotated; confidence: high}
\end{GSitem}
\begin{GSitem}{1975-376.5}{E}{The polynomial functional equation}
\GSfield{Kind; source}{theorem; \GSpage{P000628}}
\GSfield{The notes}{The example under singularities is \(L P^\lambda=a(\lambda)P^{\lambda-1}\).}
\GSfield{Statement}{For a nonzero polynomial \(P\in\mathbb C[x_1,\ldots,x_n]\), there are a nonzero polynomial \(b(s)\in\mathbb C[s]\) and a differential operator \(Q(s)\) with coefficients in \(\mathbb C[x,s]\) satisfying \(Q(s)P^{s+1}=b(s)P^s\) in the formal localization module. Taking \(s=\lambda-1\) gives the notebook's convention. Allowing polynomial dependence of the operator on the exponent is essential.}
\GSfield{Reconstruction}{The coefficient ring, formal meaning of the powers and shift of the exponent are supplied.}
\GSfield{Fate}{\citet[Theorem 1\(^{\prime}\), p.~26 of the Russian original]{bernstein1972continuation} supplies the polynomial-parameter differential functional equation; the paper appeared in 1972, before the seminar.}
\GSfield{Present status}{As of 10 November 2026, the result is established in the stated scope.}
\GSfield{Assessment}{literal statement: correct with missing hypotheses; publication match: exact; verification: read the result; development: annotated; confidence: high}
\end{GSitem}

\GStalk{377}{1975-11-10}{A. M. Polyakov}{P000627-P000632}{Quark representations, color and the Wilson-loop confinement problem}
Polyakov begins on P000628 after Gelfand's reading programme and finishes on P000630. The discussion moves from hadron representations and quark statistics to color gauge invariance, infinitely heavy probes, rectangular Wilson loops and a Gaussian Abelian comparison. Representation identities, phenomenological approximations, established constructions and the unsolved four-dimensional confinement question are separated. A reference to the penultimate issue of \emph{Theoretical and Mathematical Physics} lacks an author, title or year sufficiently specific to identify a paper.

\begin{GSitem}{1975-377.1}{E}{Quark content, spin and resonance shorthand}
\GSfield{Kind; source}{claim; \GSpage{P000628}}
\GSfield{The notes}{The notes introduce three spin-\(1/2\) quarks \(u,d,s\), label baryons spin \(1/2\) and mesons spin zero, and describe resonances as orbital excitations.}
\GSfield{Statement}{In the simplest constituent model, conventional baryons have valence content \(qqq\), while conventional mesons have \(q\bar q\). Their spins are obtained by coupling constituent spins and orbital angular momentum; baryons need not have spin \(1/2\), and mesons need not have spin zero. An orbital-excitation model describes some resonances but is not a definition of every hadronic resonance.}
\GSfield{Reconstruction}{The spin labels are restricted to examples; their literal universal reading is false even for the spin-\(3/2\) baryon discussed later on the same page.}
\GSfield{Fate}{The pre-talk review \citet[\S I and Appendix I]{levinfrankfurt1968quark} already includes spin-\(3/2\) baryons and spin-one mesons. These are old model distinctions, not a later correction discovered by this seminar.}
\GSfield{Present status}{As of 10 November 2026, the result is established in the stated scope.}
\GSfield{Assessment}{literal statement: false as written; publication match: partial; verification: read the result; development: annotated; confidence: high}
\end{GSitem}
\begin{GSitem}{1975-377.2}{E}{Flavor tensor products for mesons and baryons}
\GSfield{Kind; source}{theorem; \GSpage{P000628}}
\GSfield{The notes}{The page lists \(\bar3\otimes3=1\oplus8\), \(3^{\otimes3}=1\oplus8\oplus8\oplus10\), and \(3\otimes3=\bar3\oplus6\).}
\GSfield{Statement}{For the complex defining representation \(V\) of \(SU(3)\), \(V^*\otimes V=\mathbb C\oplus\mathfrak{sl}_3\), \(\Lambda^2V\cong V^*\), and \(V^{\otimes3}\cong\mathbb C\oplus\mathfrak{sl}_3^{\oplus2}\oplus\operatorname{Sym}^3V\). The displayed dimensions are \(9=1+8\), \(9=3+6\), and \(27=1+8+8+10\). These are exact representation identities; assigning them to flavor multiplets is a physical model.}
\GSfield{Reconstruction}{The vector-space notation distinguishes an exact tensor decomposition from empirical approximate flavor symmetry.}
\GSfield{Fate}{The quark-model context predates the talk: \citet[pp.~214--215; publisher metadata]{gellmann1964schematic}. The identities themselves follow directly from tensor symmetry and are explicitly present in the notebook.}
\GSfield{Present status}{As of 10 November 2026, the result is established in the stated scope.}
\GSfield{Assessment}{literal statement: correct; publication match: exact; verification: publisher metadata checked; development: annotated; confidence: high}
\end{GSitem}
\begin{GSitem}{1975-377.3}{E}{The additive-quark total-cross-section ratio}
\GSfield{Kind; source}{claim; \GSpage{P000628}}
\GSfield{The notes}{For high-energy baryon--baryon and baryon--meson scattering, the ratio is recorded as \(3/2\), to about ten percent; particle and antiparticle cross sections are said to coincide.}
\GSfield{Statement}{If a high-energy additive constituent model assigns the same effective scattering strength to each valence quark or antiquark, with negligible multiple-scattering and wavefunction corrections, constituent counting gives \(\sigma(BB)/\sigma(BM)\approx9/6=3/2\). Equality of particle and antiparticle total cross sections concerns an asymptotic scattering statement under additional analyticity and growth assumptions, not an exact equality at all energies. The stated ten-percent agreement is a historical estimate without a specified dataset.}
\GSfield{Reconstruction}{The factorization assumptions and separation from an asymptotic theorem are supplied; no present experimental error bound is inferred.}
\GSfield{Fate}{\citet[Introduction, p.~108, and \S VII]{levinfrankfurt1968quark} explicitly discusses this ratio under a Pomeranchuk-limit assumption before 1975.}
\GSfield{Present status}{As of 10 November 2026, the result is established in the stated scope.}
\GSfield{Assessment}{literal statement: correct with missing hypotheses; publication match: partial; verification: read the result; development: annotated; confidence: high}
\end{GSitem}
\begin{GSitem}{1975-377.4}{B}{The historical absence of pentaquarks}
\GSfield{Kind; source}{remark; \GSpage{P000628}}
\GSfield{The notes}{The notes say that a \(qqqq\bar q\) configuration has not been observed.}
\GSfield{Statement}{As a dated empirical remark this concerns the evidence available to the speaker in November 1975. It is not a representation-theoretic prohibition: \(3^{\otimes4}\otimes\bar3\) contains color singlets. A later observation of resonances consistent with such minimal valence content does not by itself determine a compact five-quark rather than molecular internal structure.}
\GSfield{Reconstruction}{The minimal-content and internal-structure distinction is supplied to prevent the historical sentence being read as a theorem.}
\GSfield{Fate}{\citet[abstract]{lhcb2015pentaquark}, first submitted 13 July 2015, reported hidden-charm structures consistent with pentaquark states. This changes the historical observational context; it does not settle their microscopic binding mechanism.}
\GSfield{Present status}{As of 10 November 2026, the 1975 non-observation statement is historical and cannot be repeated as a present exclusion.}
\GSfield{Assessment}{literal statement: correct; publication match: partial; verification: read the result; development: annotated; confidence: high}
\end{GSitem}
\begin{GSitem}{1975-377.5}{E}{Color, fermion statistics and singlet tensors}
\GSfield{Kind; source}{construction; \GSpage{P000628-P000629}}
\GSfield{The notes}{The symmetric \(uuu\) state of \(\Delta^{++}\) motivates three colors and an antisymmetric color factor; the notes also display a colorless meson sum.}
\GSfield{Statement}{For three identical spin-\(1/2\) quarks in a symmetric spatial, spin and flavor state, the antisymmetric tensor \(\varepsilon_{abc}\) in color \(\mathbb C^3\) supplies the required fermionic antisymmetry. A color-singlet meson contraction is \(\sum_a\bar q_a q_a\), while a baryon contraction is \(\varepsilon_{abc}q_aq_bq_c\). The middle term of the notebook's meson sum has two barred factors and must be emended.}
\GSfield{Reconstruction}{The fully antisymmetric color tensor is the modern formulation of the notebook's proposal. Historical paraquark and three-triplet models are precursors, not literally identical to modern QCD.}
\GSfield{Fate}{Pre-talk precedents are \citet[publisher record]{greenberg1964paraquark} and \citet[abstract]{hannambu1965triplet}; both dates and the latter's double-\(SU(3)\) model were checked.}
\GSfield{Present status}{As of 10 November 2026, the result is established in the stated scope.}
\GSfield{Assessment}{literal statement: false as written; publication match: partial; verification: publisher metadata checked; development: annotated; confidence: high}
\end{GSitem}
\begin{GSitem}{1975-377.6}{E}{Approximate flavor symmetry and exact color gauge invariance}
\GSfield{Kind; source}{claim; \GSpage{P000628-P000629}}
\GSfield{The notes}{The discussion contrasts approximately valid flavor \(SU(3)\) with exact color \(SU(3)\), suggests four flavors, and contains the isolated phrase that the Lagrangian cannot be invariant under flavor times color.}
\GSfield{Statement}{For quark fields in the fundamental color representation, a common flavor-independent kinetic term is invariant under global flavor rotations and local color gauge transformations. A nondegenerate quark-mass matrix breaks the global flavor symmetry while preserving color gauge invariance. Four flavors give an approximate global \(SU(4)\) only in a mass-degenerate idealization; adding a flavor does not change the color group. The isolated blanket denial of invariance cannot be reconciled literally with the next page without this distinction.}
\GSfield{Reconstruction}{The mass-matrix qualification is reconstructed. The numerical 20--30 percent accuracy and dating of a fourth flavor are not converted into rigorous statements.}
\GSfield{Fate}{Four-quark weak-current models already appear in \citet[abstract]{glashowiliopoulosmaiani1970weak}, published 1 October 1970. This verifies a pre-talk four-flavor precedent, not the notebook's approximate chronology or percentage.}
\GSfield{Present status}{As of 10 November 2026, the result is established in the stated scope.}
\GSfield{Assessment}{literal statement: not determinable; publication match: partial; verification: publisher metadata checked; development: annotated; confidence: high}
\end{GSitem}
\begin{GSitem}{1975-377.7}{E}{The color Yang-Mills Lagrangian and gluons}
\GSfield{Kind; source}{definition; \GSpage{P000629}}
\GSfield{The notes}{The notebook writes a Dirac kinetic term with a matrix connection and \(\operatorname{tr}F_{\mu\nu}^2/(4e^2)\), and describes eight massless gluons.}
\GSfield{Statement}{With gauge group \(SU(3)\), a connection \(A_\mu\in\mathfrak{su}(3)\), curvature \(F_{\mu\nu}=\partial_\mu A_\nu-\partial_\nu A_\mu+[A_\mu,A_\nu]\), and quarks in the defining representation, the classical gauge-invariant action contains \(\bar q(i\gamma^\mu D_\mu-m)q\) and the Yang--Mills quadratic curvature term, with signs depending on metric and Lie-algebra conventions. The eight gauge components have no gauge-invariant bare Proca mass term. This statement does not assert eight massless physical asymptotic particles.}
\GSfield{Reconstruction}{The distinction between bare gauge fields and the physical spectrum is supplied; the notebook's normalization of the coupling is retained only up to convention.}
\GSfield{Fate}{The pre-talk gauge-theory setting is explicit in \citet[\S III]{wilson1974confinement}. For the later separation of classical masslessness and a quantum mass gap see \citet[\S1]{jaffewitten2000yangmills}.}
\GSfield{Present status}{As of 10 November 2026, the result is established in the stated scope.}
\GSfield{Assessment}{literal statement: correct with missing hypotheses; publication match: partial; verification: read the result; development: annotated; confidence: high}
\end{GSitem}
\begin{GSitem}{1975-377.8}{F}{Noncommutative electrical engineering}
\GSfield{Kind; source}{remark; \GSpage{P000629}; I. M. Gelfand (intervention)}
\GSfield{The notes}{Gelfand proposes a classical version described as noncommutative electrical engineering, with a diagram.}
\GSfield{Statement}{No circuit equations, class of networks, noncommutative coefficient algebra, boundary conditions or assertion to test are specified. The strongest recoverable content is a programme to interpret classical non-Abelian gauge fields in an electrical-network language.}
\GSfield{Reconstruction}{No later noncommutative circuit theory is inferred from the metaphor. The separate Noether-type statement on the same page is indexed as 1975-377.9.}
\GSfield{Fate}{The gauge context is identifiable from \citet[\S III]{wilson1974confinement}, but no exact proposed result is identifiable from the remark; searches for the quoted phrase and Gelfand/Polyakov circuit terminology did not identify one.}
\GSfield{Present status}{As of 10 November 2026, this remains an unformalized programme remark, not a precise unsolved problem.}
\GSfield{Assessment}{literal statement: not determinable; publication match: thematic only; verification: unresolved lead; development: annotated; confidence: high}
\end{GSitem}
\begin{GSitem}{1975-377.9}{E}{Variational symmetries and conserved currents}
\GSfield{Kind; source}{theorem; \GSpage{P000629}}
\GSfield{The notes}{The line before the confinement proposal relates conserved currents to the presence of symmetry by a two-way arrow.}
\GSfield{Statement}{For a finite-order local Lagrangian with sufficiently differentiable fields and coefficients, an infinitesimal variational symmetry depending on constant parameters gives a current whose divergence is a linear combination of the Euler--Lagrange expressions and hence vanishes on solutions. Boundary terms and currents differing by identically conserved improvements must be treated consistently. A symmetry depending on arbitrary functions instead yields differential identities among the Euler--Lagrange expressions. Gauge charges further depend on boundary conditions; the bare arrow is not a universal bijection between every current and every gauge redundancy.}
\GSfield{Reconstruction}{The finite-parameter and arbitrary-function cases are separated, with the local variational and boundary hypotheses supplied. This mathematical assertion is distinct from the unformalized electrical-engineering programme in 1975-377.8.}
\GSfield{Fate}{\citet[\S1, Theorems I--II]{noether1918invariant} states the two classical results. The opened text is Tavel's 1971 English translation of the 1918 paper; its later arXiv deposit is not its first publication.}
\GSfield{Present status}{As of 10 November 2026, the result is established in the stated scope.}
\GSfield{Assessment}{literal statement: correct with missing hypotheses; publication match: partial; verification: read the result; development: annotated; confidence: high}
\end{GSitem}
\begin{GSitem}{1975-377.10}{C}{A linearly rising potential in four-dimensional pure gauge theory}
\GSfield{Kind; source}{conjecture; \GSpage{P000629-P000630}}
\GSfield{The notes}{The proposed outcome is \(E(R)\sim R\) for non-Abelian fields, contrasted with \(R^{-1}\) in the Abelian case. Colorless composites should have finite mass while an isolated color charge is excluded from the finite-energy spectrum.}
\GSfield{Statement}{A definite reconstruction is the zero-temperature infinite-volume four-dimensional pure \(SU(N)\) Yang--Mills theory, \(N\ge2\), with external fundamental charges and a renormalized rectangular Wilson loop: after removal of perimeter/self-energy terms, \(V(R)=-\lim_{T\to\infty}T^{-1}\log\langle W(R,T)\rangle\) should obey \(\lim_{R\to\infty}V(R)/R=\sigma>0\). Dynamical fundamental matter must be excluded for an asymptotic unbroken string. Noncommutativity alone, without dimension, matter content, phase and continuum hypotheses, is insufficient.}
\GSfield{Reconstruction}{The formal action explicitly contains \(d^4x\), indicating the dimension four. The notebook does not specify the continuum/regularization or matter-content conditions. The reconstructed four-dimensional pure-gauge assertion is the unresolved part; a bare infinite mass term is not the meaning of confinement.}
\GSfield{Fate}{\citet[\S IV]{wilson1974confinement} gives the pre-talk strong-coupling lattice mechanism. \citet[\S5]{jaffewitten2000yangmills} separates confinement from the existence-and-mass-gap problem; the official current description \citep{clay2026yangmills} still lists the latter as unresolved. Lattice strong coupling does not establish a continuum area law.}
\GSfield{Present status}{As of 10 November 2026, the reconstructed four-dimensional continuum area-law statement remains unproved; fixed-lattice and lower-dimensional results do not close it.}
\GSfield{Assessment}{literal statement: correct with missing hypotheses; publication match: partial; verification: read the result; development: annotated; confidence: high; problem \hyperref[prob:GP-1975-10]{GP-1975-10}}
\GSfield{Worth developing}{Worth retaining as a precise research problem, but a bounded chapter cannot supply the missing four-dimensional continuum construction and positive string tension; dimensional and screening qualifications are the useful contribution here.}
\end{GSitem}
\begin{GSitem}{1975-377.11}{E}{Infinitely heavy charges and suppression of pair creation}
\GSfield{Kind; source}{method; \GSpage{P000629}}
\GSfield{The notes}{Taking quark mass \(M\to\infty\) is used to discard quark-pair diagrams and retain paths of two external charges.}
\GSfield{Statement}{Static source worldlines arise from an infinite external-mass limit after subtracting their rest energies. Removing dynamical quark loops is a separate pure-gauge or quenched assumption, or a decoupling limit for every dynamical quark mass. Infinite mass of the two probes alone does not prevent light dynamical quarks from screening them and breaking the string.}
\GSfield{Reconstruction}{The separation of probe and sea-quark masses is a necessary reconstruction that the diagrams do not make explicit.}
\GSfield{Fate}{The loop and screening discussion is already present in \citet[\S II, especially the nearby-vacuum-loop discussion]{wilson1974confinement}.}
\GSfield{Present status}{As of 10 November 2026, the result is established in the stated scope.}
\GSfield{Assessment}{literal statement: correct with missing hypotheses; publication match: partial; verification: read the result; development: annotated; confidence: high; problem \hyperref[prob:GP-1975-10]{GP-1975-10}}
\end{GSitem}
\begin{GSitem}{1975-377.12}{E}{Wilson loops and extraction of a static energy}
\GSfield{Kind; source}{construction; \GSpage{P000629-P000630}}
\GSfield{The notes}{The field integral contains \(\operatorname{tr}\exp(i\oint A)\); its meaning is given by an ordered product. The rectangle has sides \(R,T\). One line says \(W\to ET\), whereas the next page uses \(e^{-iET}\).}
\GSfield{Statement}{The gauge-invariant loop is the trace of the path-ordered holonomy. In a Euclidean regulated formulation with a positive transfer operator, a rectangular loop has a spectral expansion \(\sum_n c_n(R)e^{-E_n(R)T}\), and nonzero overlap with the lowest state yields \(-T^{-1}\log\langle W(R,T)\rangle\to V(R)\), after source-energy subtraction. In real time the amplitude is \(e^{-iET}\), so writing it as \(e^{iW}\) gives \(W=-ET\).}
\GSfield{Reconstruction}{\oursComputed\ The phase sign is checked against \(e^{-iET}\). The Euclidean spectral interpretation, overlap hypothesis and subtraction are supplied. Path ordering is already present in the notebook's product definition and is not an omitted operation.}
\GSfield{Fate}{This construction and its physical interpretation precede the talk in \citet[\S\S II--III]{wilson1974confinement}.}
\GSfield{Present status}{As of 10 November 2026, the result is established in the stated scope.}
\GSfield{Assessment}{literal statement: false as written; publication match: partial; verification: read the result; development: annotated; confidence: high; problem \hyperref[prob:GP-1975-10]{GP-1975-10}}
\end{GSitem}
\begin{GSitem}{1975-377.13}{E}{Worldline phases and the area-law criterion}
\GSfield{Kind; source}{method; \GSpage{P000630}}
\GSfield{The notes}{The path-sum explanation uses the phase \(e^{-iET}\), treats an antiparticle as a reversed worldline, and proceeds to \(W(C)\sim\operatorname{area}(C)=RT\), hence \(E(R)\sim R\).}
\GSfield{Statement}{For a large rectangular Euclidean loop, an asymptotic law \(\log\langle W(R,T)\rangle=-\sigma RT+\text{perimeter and subleading terms}\) implies the linear static energy after the appropriate successive limits and perimeter subtraction. A worldline representation identifies the observable but does not prove its area law. The latter is an independent dynamical assertion, not a consequence of reversing one trajectory.}
\GSfield{Reconstruction}{The conditional implication and order of limits separate the valid extraction method from the unproved dynamical step.}
\GSfield{Fate}{The loop/area-law criterion was formulated in \citet[\S II]{wilson1974confinement}; its unresolved continuum use is item 1975-377.10.}
\GSfield{Present status}{As of 10 November 2026, the conditional criterion is established; proving its premise in four-dimensional continuum pure Yang--Mills theory remains the problem indexed above.}
\GSfield{Assessment}{literal statement: correct with missing hypotheses; publication match: exact; verification: read the result; development: annotated; confidence: high; problem \hyperref[prob:GP-1975-10]{GP-1975-10}}
\end{GSitem}
\begin{GSitem}{1975-377.14}{E}{Gaussian Abelian loop and the Coulomb potential}
\GSfield{Kind; source}{method; \GSpage{P000630}}
\GSfield{The notes}{The Abelian field integral becomes the exponential of a double line integral of a propagator; the rectangular-loop expression is a term proportional to \(T\) plus a term proportional to \(e^2T/R\).}
\GSfield{Statement}{For a centered Gaussian Euclidean gauge field with covariance \(D_{\mu\nu}\) and a conserved regularized loop current \(J\), \(\langle e^{ie(A,J)}\rangle=\exp[-e^2(J,DJ)/2]\). In four-dimensional Maxwell theory, removing the divergent self-energies of two opposite static charges yields \(V(R)=-e^2/(4\pi R)\) in rationalized units. Coincident-point integrals require a regulator; the sign and factor in the notebook depend on propagator and real/Euclidean-time conventions.}
\GSfield{Reconstruction}{The regulator, charge signs, units and explicit covariance identity are reconstructed. The \(R^{-1}\) law assumes three spatial dimensions and a massless noncompact Maxwell field.}
\GSfield{Fate}{The Abelian comparison is part of the pre-talk analysis in \citet[\S II]{wilson1974confinement}.}
\GSfield{Present status}{As of 10 November 2026, the result is established in the stated scope.}
\GSfield{Assessment}{literal statement: correct with missing hypotheses; publication match: partial; verification: read the result; development: annotated; confidence: high}
\end{GSitem}
\begin{GSitem}{1975-377.15}{E}{Why the non-Abelian integral is not Gaussian}
\GSfield{Kind; source}{remark; \GSpage{P000630}}
\GSfield{The notes}{The final comparison points to the commutator term in \(F\), and also to multiple extrema for compact Abelian variables.}
\GSfield{Statement}{In a non-Abelian gauge theory, \(F=dA+A\wedge A\), so the action contains cubic and quartic interactions in a local gauge potential and is not a Gaussian functional. Compact lattice \(U(1)\) has a nonlinear periodic plaquette action even though its group is Abelian. Thus Gaussian Maxwell calculations cannot decide the infrared phase of either interacting non-Abelian theory or arbitrary compact Abelian models.}
\GSfield{Reconstruction}{The distinction between a linear Abelian continuum field and a compact lattice model makes the scope of the comparison explicit.}
\GSfield{Fate}{Both compact Abelian and non-Abelian lattice formulations occur in \citet[\S\S III--IV]{wilson1974confinement}; their strong-coupling behavior is not exclusive to non-Abelian groups.}
\GSfield{Present status}{As of 10 November 2026, the result is established in the stated scope.}
\GSfield{Assessment}{literal statement: correct with missing hypotheses; publication match: partial; verification: read the result; development: annotated; confidence: high}
\end{GSitem}

\GStalk{378}{}{-}{P000630-P000632}{Year heading: 1976}
This catalogue record marks the change of calendar year. Its overlapping page links include the end of the November 1975 discussion and the dated 9 February 1976 talk, but provide no independent mathematical assertion. It receives zero occurrence items; all following item identifiers retain academic-year prefix 1975.


\GStalk{379}{1976-02-09}{I. M. Krichever}{P000630-P000634}{Baker-Akhiezer construction of Zakharov-Shabat equations}
The dated talk occupies P000631--P000632 and the first two lines of P000633. It passes from formal exponential solutions to algebraic spectral data and the uniqueness argument constructing differential operators. Gelfand asks an elementary commuting-operator problem and raises the distinction between a formal series and an analytic function. Drinfeld begins a separate algebraization immediately below on P000633.

\begin{GSitem}{1975-379.1}{A}{Zero-curvature compatibility for two auxiliary equations}
\GSfield{Kind; source}{definition; \GSpage{P000631}}
\GSfield{The notes}{The equation is \([L_1-\partial_y,L_2-\partial_t]=0\), with scalar differential operators in \(x\), monic leading terms, and constant next coefficients. Gelfand emphasizes the symmetry of the auxiliary variables.}
\GSfield{Statement}{For sufficiently differentiable operator coefficients, compatibility of \(\psi_y=L_1\psi\), \(\psi_t=L_2\psi\) is \(\partial_yL_2-\partial_tL_1=[L_1,L_2]\). The same operator identity makes sense for matrices or operators in other variables, but existence of a common solution requires suitable analytic or formal initial-value conditions. The displayed equation alone is the zero-curvature identity; it does not by itself construct a nonlinear solution.}
\GSfield{Reconstruction}{The equivalent sign convention is obtained by expanding the commutator. The original scalar differential setting is distinguished from Gelfand's broader operator analogy.}
\GSfield{Fate}{\citet[Theorem 1 and Corollary 1, pp.~291--292]{krichever1976zakharov} constructs solutions in this form. The manuscript was received 8 October 1975 and presented to the Academy on 24 October; neither date alone proves public availability before the February talk.}
\GSfield{Present status}{As of 10 November 2026, the result is established in the stated scope.}
\GSfield{Assessment}{literal statement: correct; publication match: exact; verification: read the result; development: annotated; confidence: high}
\end{GSitem}
\begin{GSitem}{1975-379.2}{A}{Formal wave-function recursion and its normalization}
\GSfield{Kind; source}{construction; \GSpage{P000631-P000632}}
\GSfield{The notes}{The ansatz is \(\psi=e^{kx+k^ny}\sum_{j\ge0}\xi_jk^{-j}\), \(\xi_0=1\). The written normalization is \(\xi_j(0,y,t)=0\) for \(j>1\).}
\GSfield{Statement}{For \(L=D^n+\sum_{i=0}^{n-2}u_iD^i\), formal coefficient comparison in \((L-\partial_y)\psi=0\) recursively determines \(\partial_x\xi_j\). Prescribing \(\xi_j(0,y,t)\) for every \(j\ge1\) fixes the integration functions. A nonzero constant \(D^{n-1}\) coefficient requires the corresponding \(k^{n-1}y\) term in the exponential phase; it is incompatible with the unmodified displayed phase.}
\GSfield{Reconstruction}{The normalization bound is emended to \(j\ge1\), and either the next coefficient must vanish or the phase must change. This is a formal calculation and asserts no convergence.}
\GSfield{Fate}{The more general polynomial phase in \citet[conditions 1--2 preceding Theorem 1]{krichever1976zakharov} makes the missing phase term explicit.}
\GSfield{Present status}{As of 10 November 2026, the result is established in the stated scope.}
\GSfield{Assessment}{literal statement: false as written; publication match: partial; verification: read the result; development: annotated; confidence: high}
\end{GSitem}
\begin{GSitem}{1975-379.3}{A}{Recovering an operator from a normalized formal wave}
\GSfield{Kind; source}{method; \GSpage{P000631-P000633}}
\GSfield{The notes}{Conversely, every normalized series is said to determine an operator up to a residual of order \(k^{-1}\); the final proof applies uniqueness to that residual.}
\GSfield{Statement}{For a normalized formal wave \(e^{kx+k^ny}(1+\sum_{j\ge1}\xi_jk^{-j})\), comparison of nonnegative powers of \(k\) uniquely gives a monic \(L_n\), with no \(D^{n-1}\) term, for which \(e^{-kx-k^ny}(L_n-\partial_y)\psi=O(k^{-1})\). If the wave also satisfies the nonspecial Baker--Akhiezer conditions below, the residual has the same allowed finite poles and zero leading normalization; uniqueness makes it identically zero. For an arbitrary formal wave only the residual statement follows.}
\GSfield{Reconstruction}{The two distinct stages, formal cancellation and global uniqueness, are separated to avoid silently assuming the analytic conclusion.}
\GSfield{Fate}{This is the proof mechanism of \citet[Theorem 1, pp.~291--292]{krichever1976zakharov}.}
\GSfield{Present status}{As of 10 November 2026, the result is established in the stated scope.}
\GSfield{Assessment}{literal statement: correct with missing hypotheses; publication match: exact; verification: read the result; development: annotated; confidence: high}
\end{GSitem}
\begin{GSitem}{1975-379.4}{E}{Commuting scalar operators of orders two and three}
\GSfield{Kind; source}{problem; \GSpage{P000631}; I. M. Gelfand (intervention)}
\GSfield{The notes}{Gelfand's school problem asks when \(L_1=-D^2+u\) and \(L_2=D^3+v_1D+v_2\) commute, with an answer in elliptic functions or elliptic curves.}
\GSfield{Statement}{On an interval with smooth complex coefficients, direct comparison of powers of \(D\) gives \(v_1=-3u/2+c_1\), \(v_2=-3u^{\prime}/4+c_2\), and \(u^{(3)}-6uu^{\prime}+4c_1u^{\prime}=0\). Integrating twice gives \((u^{\prime})^2=2u^3-4c_1u^2+2c_3u+c_4\). The solutions are elliptic functions on a nonsingular cubic, or their rational/trigonometric/hyperbolic degenerations, restricted away from poles.}
\GSfield{Reconstruction}{\oursComputed\ The coefficient comparison and integration constants were independently checked. This reconstructs a classical exercise, including degenerate cubics, and makes no priority claim.}
\GSfield{Fate}{\citet[\S3, equations (2)--(3)]{mulase1990geometric} gives this classical calculation in an equivalent normalization and traces the order-two/order-three case to Wallenberg (1903). The historical attribution is secondary within that research article; the displayed calculation was read and independently expanded here.}
\GSfield{Present status}{As of 10 November 2026, the result is established in the stated scope.}
\GSfield{Assessment}{literal statement: correct with missing hypotheses; publication match: partial; verification: read the result; development: annotated; confidence: high; problem \hyperref[prob:GP-1975-12]{GP-1975-12}}
\end{GSitem}
\begin{GSitem}{1975-379.5}{A}{Algebraic operator versus merely formal asymptotics}
\GSfield{Kind; source}{definition; \GSpage{P000631-P000632}}
\GSfield{The notes}{An operator is called algebraic when its formal wave extends meromorphically to a compact curve away from one marked point, with poles independent of \(x,y,t\). Gelfand stresses that an asymptotic expansion usually diverges.}
\GSfield{Statement}{The intended data are a smooth projective curve \(X\), a marked point \(P\), a local coordinate \(z=k^{-1}\), and a fixed pole divisor away from \(P\). The normalized wave must actually define a holomorphic function near the puncture after removal of its exponential singularity, and a meromorphic function elsewhere. A formal Laurent series does not determine such a global meromorphic function without convergence and continuation conditions.}
\GSfield{Reconstruction}{Convergence of the exponential-normalized series and fixed divisor are imposed as parts of this definition, not derived from the coefficient recursion.}
\GSfield{Fate}{The analytic data appear explicitly before \citet[Theorem 1]{krichever1976zakharov}.}
\GSfield{Present status}{As of 10 November 2026, the result is established in the stated scope.}
\GSfield{Assessment}{literal statement: correct with missing hypotheses; publication match: exact; verification: read the result; development: annotated; confidence: high}
\end{GSitem}
\begin{GSitem}{1975-379.6}{F}{An unspecified quasianalyticity analogue}
\GSfield{Kind; source}{question; \GSpage{P000632}; I. M. Gelfand (intervention)}
\GSfield{The notes}{After contrasting smooth Taylor asymptotics and analytic convergence, Gelfand asks \ru{Аналог квазианалитичности?}}
\GSfield{Statement}{No function class, growth sequence, asymptotic sector, or proposed uniqueness implication is specified. The strongest recoverable content is a request to understand when the formal wave data determine an analytic or algebraic object; there is no unique theorem or quantified conjecture to reconstruct.}
\GSfield{Reconstruction}{No Denjoy--Carleman condition, Borel summability theorem or differential-Galois criterion is substituted for this sentence.}
\GSfield{Fate}{The analytic construction in \citet{krichever1976zakharov} assumes specified spectral data and does not answer the unformulated converse. Searches for the quoted terminology with Krichever and Gelfand produced no exact identifiable statement.}
\GSfield{Present status}{As of 10 November 2026, the remark remains too underspecified to assign an open-problem status.}
\GSfield{Assessment}{literal statement: not determinable; publication match: thematic only; verification: unresolved lead; development: annotated; confidence: high}
\end{GSitem}
\begin{GSitem}{1975-379.7}{A}{Existence and uniqueness of the Baker-Akhiezer wave}
\GSfield{Kind; source}{theorem; \GSpage{P000632}}
\GSfield{The notes}{The theorem attributed in the notes to Akhiezer allows a fixed set of poles and says that \(\psi e^{-kx-k^ny}=1+o(k^{-1})\), with uniqueness up to proportionality.}
\GSfield{Statement}{Let \(X\) be a smooth compact complex curve of genus \(g\), \(P\in X\), and \(D\) an effective nonspecial divisor of degree \(g\) disjoint from \(P\). Fix a local coordinate \(k^{-1}\) and polynomial exponential phase. For times in the open set on which the normalized interpolation problem is nonsingular, there is a unique wave meromorphic on \(X\setminus\{P\}\), with poles bounded by \(D\), whose exponential-normalized germ is \(1+O(k^{-1})\). With leading coefficient fixed to one, there is no remaining scalar ambiguity.}
\GSfield{Reconstruction}{Nonspeciality, degree, open time domain and the large-O correction are supplied. The source's Akhiezer attribution is retained historically; its stated multiparameter version is matched to Krichever's publication.}
\GSfield{Fate}{\citet[conditions 1--2 and Theorem 1]{krichever1976zakharov} gives the matching construction with a general-position divisor and a \(1+\sum\xi_jk^{-j}\) expansion.}
\GSfield{Present status}{As of 10 November 2026, the result is established in the stated scope.}
\GSfield{Assessment}{literal statement: false as written; publication match: exact; verification: read the result; development: annotated; confidence: high}
\end{GSitem}
\begin{GSitem}{1975-379.8}{A}{The elliptic example requires a meromorphic correction factor}
\GSfield{Kind; source}{construction; \GSpage{P000632}}
\GSfield{The notes}{The elliptic illustration writes an exponential involving the Weierstrass zeta function.}
\GSfield{Statement}{The Weierstrass zeta function is quasi-periodic, so \(e^{x\zeta(z)}\) alone is generally not a single-valued function on the elliptic curve. A Baker--Akhiezer function needs a sigma- or theta-function ratio whose quasi-period multiplier cancels that of the exponential, together with the prescribed fixed pole and normalization. This is the genus-one instance of the preceding construction.}
\GSfield{Reconstruction}{The displayed exponential is treated as the essential-singularity factor, not a complete global formula.}
\GSfield{Fate}{The finite-pole and global-meromorphic requirements in \citet[conditions 1--2]{krichever1976zakharov} identify exactly what the isolated exponential omits.}
\GSfield{Present status}{As of 10 November 2026, the result is established in the stated scope.}
\GSfield{Assessment}{literal statement: false as written; publication match: partial; verification: read the result; development: annotated; confidence: high}
\end{GSitem}

\GStalk{380}{1976-02-09}{V. G. Drinfeld}{P000632-P000635}{Algebraization by filtered modules}
Drinfeld begins beneath the end of Krichever's proof on P000633 and continues through P000634. He replaces an embedding of the punctured-curve function ring into differential operators by a filtered module and a connection-like operator. Gelfand's reciprocity and division-ring remarks are preserved as separate, underspecified programmes. The construction directly anticipates Drinfeld's 1977 note, which acknowledges the seminar context.

\begin{GSitem}{1975-380.1}{E}{Differential operator ring over a differential algebra}
\GSfield{Kind; source}{definition; \GSpage{P000633}}
\GSfield{The notes}{\(B\) is a commutative unital complex algebra with derivation \(b\mapsto b^{\prime}\); the ring \(B[D]\) has \(Db-bD=b^{\prime}\).}
\GSfield{Statement}{For any commutative unital \(\mathbb C\)-algebra \(B\) and derivation \(\delta\), the Ore extension \(B[D;\delta]\) consists of finite sums \(\sum b_iD^i\), with multiplication determined by \(Db=bD+\delta(b)\). This definition remains meaningful when \(B\) is not a field; existence of a division ring of fractions requires extra hypotheses.}
\GSfield{Reconstruction}{The Ore-extension notation and the limitation about fractions are supplied; no analytic ellipticity is implicit in the ring definition.}
\GSfield{Fate}{The algebra is used explicitly in \citet[\S2]{drinfeld1977commutative}. The defining construction is background, distinct from the new geometric interpretation.}
\GSfield{Present status}{As of 10 November 2026, the result is established in the stated scope.}
\GSfield{Assessment}{literal statement: correct; publication match: exact; verification: read the result; development: annotated; confidence: high}
\end{GSitem}
\begin{GSitem}{1975-380.2}{F}{The claim that all commutative subrings come from smooth curves}
\GSfield{Kind; source}{claim; \GSpage{P000633}}
\GSfield{The notes}{With \(A=H^0(X\setminus\{P\},\mathcal O_X)\), the notes say that no other commutative subrings occur and attribute this to Shabat.}
\GSfield{Statement}{The smooth-curve formulation is too broad without restrictions. For example, the commutative subring \(\mathbb C[D^2,D^3]\subset\mathbb C[D]\) has a cuspidal spectrum; it is not the function ring of a smooth affine curve. Geometric classifications must allow singular spectral curves and suitable torsion-free sheaves, or restrict to subrings with a specified smooth spectral curve and regularity properties.}
\GSfield{Reconstruction}{The cusp example is a direct counterexample to the unrestricted literal reading; it does not disprove a more qualified theorem that the speaker may have intended. \ours\ The cusp obstruction is elementary: \(t\) is integral over \(\mathbb C[t^2,t^3]\), since it satisfies \(T^2-t^2=0\), but it does not lie in that ring. Thus the ring is not normal, whereas the coordinate ring of a smooth affine curve is normal.}
\GSfield{Fate}{\citet[\S2]{drinfeld1977commutative} treats a fixed smooth pointed curve, while \citet[\S2, Main Theorem]{mulase1990geometric} explicitly allows singular curves. Searches did not identify an exact Shabat theorem with the notebook's stated quantifiers.}
\GSfield{Present status}{As of 10 November 2026, the literal claim is false and the intended attribution remains unresolved; this is not evidence of an open classification problem.}
\GSfield{Assessment}{literal statement: false as written; publication match: partial; verification: unresolved lead; development: annotated; confidence: medium}
\end{GSitem}
\begin{GSitem}{1975-380.3}{F}{A noncommutative reciprocity programme}
\GSfield{Kind; source}{remark; \GSpage{P000633}; I. M. Gelfand (intervention)}
\GSfield{The notes}{Gelfand describes the goal as a noncommutative reciprocity law.}
\GSfield{Statement}{No base field, arithmetic objects, reciprocity map or functorial conditions are stated. The connection with commutative subrings is a guiding analogy, not a formulated reciprocity theorem.}
\GSfield{Reconstruction}{No later Langlands statement or Drinfeld-module theorem is read back into the phrase.}
\GSfield{Fate}{\citet[\S3]{drinfeld1977commutative} does pass to characteristic-\(p\) elliptic modules, but this thematic link does not identify or solve the unspecified reciprocity request.}
\GSfield{Present status}{As of 10 November 2026, the remark has no recoverable theorem-level formulation.}
\GSfield{Assessment}{literal statement: not determinable; publication match: thematic only; verification: unresolved lead; development: annotated; confidence: high}
\end{GSitem}
\begin{GSitem}{1975-380.4}{F}{Extensions of Gelfand-Kirillov division rings}
\GSfield{Kind; source}{question; \GSpage{P000633}; I. M. Gelfand (intervention)}
\GSfield{The notes}{Gelfand calls the real problem finding all extensions of Gelfand--Kirillov division rings, adding that polynomiality is an obstacle.}
\GSfield{Statement}{The notebook does not say whether an extension means an embedding into a fixed division algebra, a finite-dimensional division-ring extension, an extension of the center, or a change in the number of Weyl pairs. Those choices give different problems. No exact classification question can be extracted without choosing information absent from the page.}
\GSfield{Reconstruction}{The item is deliberately not identified with the fixed-rank Dixmier conjecture or with the Gelfand--Kirillov conjecture for enveloping algebras.}
\GSfield{Fate}{The historical division-ring framework is \citet[Introduction and \S1]{gelfandkirillov1966corps}. Searches combining Drinfeld, extensions and Gelfand--Kirillov division rings did not resolve the missing quantifiers.}
\GSfield{Present status}{As of 10 November 2026, this is an unresolved identification of a question, not a justified claim that a particular mathematical problem is open.}
\GSfield{Assessment}{literal statement: not determinable; publication match: thematic only; verification: unresolved lead; development: annotated; confidence: high}
\end{GSitem}
\begin{GSitem}{1975-380.5}{A}{Elliptic homomorphisms and their rank}
\GSfield{Kind; source}{definition; \GSpage{P000633}}
\GSfield{The notes}{Every nonzero image \(\varphi(a)\) is required to have an invertible leading coefficient. The notes derive \(\operatorname{ord}\varphi(a)=n\deg_P(a)\) and call \(n\) the rank.}
\GSfield{Statement}{Let \(X\) be a smooth connected projective complex curve, \(P\in X\), and \(A=H^0(X\setminus\{P\},\mathcal O_X)\). A unital \(\mathbb C\)-homomorphism \(\varphi:A\to B[D]\) is a Krichever homomorphism if every nonzero image has unit leading coefficient and some image has positive order. Then \(\operatorname{ord}\varphi(a)=n(-v_P(a))\) for nonconstant \(a\), with an integer \(n>0\). Because the pole-order semigroup is cofinite, \(n\) equals the greatest common divisor of the positive operator orders.}
\GSfield{Reconstruction}{The nontrivial positive-order condition, smooth connected curve and meaning of degree are made explicit. Elliptic here is algebraic terminology, not a real principal-symbol positivity condition.}
\GSfield{Fate}{The definition and rank lemma occur in \citet[\S2, preceding Proposition 2]{drinfeld1977commutative}, received 14 April 1976 and published in 1977.}
\GSfield{Present status}{As of 10 November 2026, the result is established in the stated scope.}
\GSfield{Assessment}{literal statement: correct with missing hypotheses; publication match: exact; verification: read the result; development: annotated; confidence: high}
\end{GSitem}
\begin{GSitem}{1975-380.6}{A}{A homomorphism encoded by a filtered module}
\GSfield{Kind; source}{construction; \GSpage{P000633-P000634}}
\GSfield{The notes}{Take \(M=B[D]\), let \(B\) act on the left and \(A\) on the right through \(\varphi\), filter by operator order, and let \(\nabla\) be left multiplication by \(D\).}
\GSfield{Statement}{The filtration has \(M_{-1}=0\), \(M=\bigcup_{i\ge0}M_i\), and each \(M_i/M_{i-1}\) is free of rank one over \(B\). The operator \(\nabla\) commutes with \(A\), satisfies \(\nabla(bm)=b\nabla m+b^{\prime}m\), and induces isomorphisms on successive graded pieces. A nonzero \(a\in A\) shifts filtration by \(n\deg_P(a)\) and gives a graded isomorphism. Conversely, these axioms and a generator of \(M_0\) identify \(M\) with \(B[D]\) and recover the right multiplication homomorphism, up to invertible scalar gauge.}
\GSfield{Reconstruction}{Exhaustiveness, the initial zero piece and local freeness are stated explicitly; the converse follows by induction on filtration degree.}
\GSfield{Fate}{\citet[Proposition 2 and its proof]{drinfeld1977commutative} gives the precise module correspondence.}
\GSfield{Present status}{As of 10 November 2026, the result is established in the stated scope.}
\GSfield{Assessment}{literal statement: correct with missing hypotheses; publication match: exact; verification: read the result; development: annotated; confidence: high}
\end{GSitem}
\begin{GSitem}{1975-380.7}{A}{From order filtrations to sheaves and a pole at infinity}
\GSfield{Kind; source}{method; \GSpage{P000634}}
\GSfield{The notes}{The filtered \(A\otimes B\)-module is reinterpreted as a sheaf on \(X\) with a filtration, and \(\nabla\) as a connection singular at infinity.}
\GSfield{Statement}{A compatible exhaustive good filtration on a finite-type \(A\otimes B\)-module allows extension across \(P\) by lattices determined by pole order. For rank one and a nonsingular spectral curve, the regular geometric case gives a line bundle and an \(\mathcal O_X\)-linear relative operator \(\nabla:L\to L(P)\) obeying the \(B\)-Leibniz rule. Arbitrary filtered modules without the freeness, growth and multiplication-isomorphism conditions need not define this geometric object.}
\GSfield{Reconstruction}{The word connection means differentiation in the coefficient parameter, not a flat holomorphic connection on all tangent directions of the curve.}
\GSfield{Fate}{The exact extension machinery is \citet[Proposition 1, Corollary 1, and Proposition 2]{drinfeld1977commutative}; the March talks unpack its rank-one case.}
\GSfield{Present status}{As of 10 November 2026, the result is established in the stated scope.}
\GSfield{Assessment}{literal statement: correct with missing hypotheses; publication match: exact; verification: read the result; development: annotated; confidence: high}
\end{GSitem}

\GStalk{381}{1976-02-23}{S. P. Novikov}{P000633-P000636}{Brodskii radial spectral inversion; Fuchs on characteristic cycles}
P000635 explicitly describes Novikov reporting work of M. Brodskii of Minsk; the matching author is M. A. Brodskii. The radial spectral problem, semiclassical action and Abel inversion match his 1976 two-page note, received on 18 June 1975. Receipt establishes submission chronology, not prior public circulation. The B classification records that the named author differs from the reporter; it does not mean that Brodskii developed an idea of Novikov's after the seminar. Fuchs then intervenes on Stiefel--Whitney cycles before Losik begins at the bottom of the page.

\begin{GSitem}{1975-381.1}{B}{Recovering a radial profile from its joint frequency array}
\GSfield{Kind; source}{problem; \GSpage{P000635}}
\GSfield{The notes}{For an elastic sphere, recover the medium parameter from all \(\lambda_{k\ell}\). The written equation contains both a \(\lambda^2 f^2\) term and a separate eigenvalue on the right, with an inconsistent sign.}
\GSfield{Statement}{The identified scalar radial model is \(y^{\prime\prime}+[\lambda_{k\ell}^2f(r)^2-\ell(\ell+1)/r^2]y=0\), \(0<r<1\). Require the regular solution at zero and a fixed real homogeneous separated condition \(\alpha y(1)+\beta y^{\prime}(1)=0\), \((\alpha,\beta)\ne(0,0)\). For the monotone reconstruction assume \(f>0\), \(f(1)=1\), and \(q(r)=rf(r)\) strictly increasing from zero to one; smoothness and \(q^{\prime}>0\) suffice for the formulas below. The data retain both radial and angular indices. This scalar equation does not recover every Lam'e parameter of an arbitrary elastic solid.}
\GSfield{Reconstruction}{The corrected equation and endpoint conditions come from the identified publication. Positivity and strict smooth monotonicity delimit a regular version; the notebook writes only monotonicity and omits boundary conditions.}
\GSfield{Fate}{\citet[equations (1)--(2) and Theorem 2]{brodskii1976spectral} is the exact publication match for the reported scalar problem, published in volume 31(4), 1976. The earliest public version is not established: a related 1975 torsional-oscillation paper is recorded as an unresolved primary-source lead in the audit. The B label records authorship relative to the reporter, with this chronology qualification.}
\GSfield{Present status}{As of 10 November 2026, the result is established in the stated scope.}
\GSfield{Assessment}{literal statement: false as written; publication match: exact; verification: read the result; development: annotated; confidence: medium; problem \hyperref[prob:GP-1975-13]{GP-1975-13}}
\end{GSitem}
\begin{GSitem}{1975-381.2}{B}{Semiclassical spectral sampling of the radial action}
\GSfield{Kind; source}{method; \GSpage{P000635}}
\GSfield{The notes}{With \(q=rf\) monotone, the notes use a turning point \(H(\tau)\), an action integral, and pairs \(\tau_{k\ell}=\sqrt{\ell(\ell+1)}/\lambda_{k\ell}\), \(m_{k\ell}=(k+1/4+C)/\sqrt{\ell(\ell+1)}\).}
\GSfield{Statement}{For \(\ell\ge1\), write \(L_\ell=\sqrt{\ell(\ell+1)}\), \(H=q^{-1}\), and \[m(\tau)=\pi^{-1}\int_{H(\tau)}^1\sqrt{\tau^{-2}f(r)^2-r^{-2}}\,dr,\qquad 0<\tau<1.\] In the smooth one-turning-point regime, the leading quantization relation samples this action as \(m(\tau_{k\ell})=(k+\delta_{k\ell})/L_\ell+o(1)\) along joint high-frequency limits in compact interior \(\tau\)-ranges. The bounded phase \(\delta_{k\ell}\) depends on the endpoint condition; it is not a universal constant. Dense limiting ratios recover the continuous action curve.}
\GSfield{Reconstruction}{The factor \(1/\pi\), bounded boundary phase, limiting interpretation and interior restriction are supplied. Individual finite-index pairs are not asserted to lie exactly on a smooth dispersion graph.}
\GSfield{Fate}{\citet[Theorem 1 and equation (4)]{brodskii1976spectral} supplies the asymptotic sampling, with the boundary correction retained explicitly. Exact overlap with the bibliographically attested 1975 precursor remains unverified; no priority date is inferred from the 1975 receipt stamp.}
\GSfield{Present status}{As of 10 November 2026, the result is established in the stated scope.}
\GSfield{Assessment}{literal statement: false as written; publication match: exact; verification: read the result; development: annotated; confidence: medium; problem \hyperref[prob:GP-1975-13]{GP-1975-13}}
\end{GSitem}
\begin{GSitem}{1975-381.3}{E}{The Abel inversion and its missing normalization}
\GSfield{Kind; source}{method; \GSpage{P000635}}
\GSfield{The notes}{The inverse \(\psi=q^{-1}\) is written as an exponential with coefficient 2 and kernel \((\tau^2-\eta^2)^{-1/2}\); the lower endpoint is absent. The preceding action integral lacks \(1/\pi\).}
\GSfield{Statement}{For positive smooth \(f\) with \(q(0)=0\), \(q(1)=1\), \(q^{\prime}>0\), let \(M(\tau)=\tau m(\tau)\) with the normalized action of the previous item. Then for \(0<\eta<1\), \[\log\psi(\eta)=2\int_\eta^1\frac{M^{\prime}(\tau)}{\sqrt{\tau^2-\eta^2}}\,d\tau.\] If the notebook's unnormalized action is retained, the multiplier must instead be \(2/\pi\). Knowing \(\psi\) recovers \(q\) by inversion and \(f(r)=q(r)/r\).}
\GSfield{Reconstruction}{\oursComputed\ The lower endpoint and normalization are emendations checked by composing the two Abel kernels and with the constant profile. The complete short derivation and a reproducible numerical check are supplied in the research log and check script.}
\GSfield{Fate}{\citet[Theorem 2, equations (5)--(6)]{brodskii1976spectral} gives this application. The Abel inversion mechanism is classical; \citet[Theorems 2.2.1 and 2.2.3, historical Remark 2.2.2]{paternainsalouhlmann2022geometric} distinguishes the older Herglotz travel-time inversion from Brodskii's spectral sampling.}
\GSfield{Present status}{As of 10 November 2026, the result is established in the stated scope.}
\GSfield{Assessment}{literal statement: false as written; publication match: partial; verification: read the result; development: annotated; confidence: high; problem \hyperref[prob:GP-1975-13]{GP-1975-13}}
\end{GSitem}
\begin{GSitem}{1975-381.4}{B}{Compatibility of an overdetermined spectral array}
\GSfield{Kind; source}{question; \GSpage{P000635}}
\GSfield{The notes}{The source calls the inverse problem overdetermined and asks when it is solvable; it mentions density and a one-to-one relation between spectral parameters.}
\GSfield{Statement}{For the fixed radial boundary problem above, a necessary condition is that its joint high-frequency samples approach a continuous monotone action/dispersion graph compatible with a positive increasing inverse profile. Such leading asymptotic compatibility is not sufficient for an arbitrary entire array to equal the exact spectrum of that profile: finitely many modified eigenvalues leave the limiting graph unchanged but may violate exact spectral constraints.}
\GSfield{Reconstruction}{The last observation separates a range theorem for a leading action transform from a range theorem for all exact frequencies. The notebook does not specify a complete topology, boundary model or list of sufficient conditions.}
\GSfield{Fate}{\citet[corollary after Theorem 2]{brodskii1976spectral} states a necessary condition, not a full necessary-and-sufficient spectral-range theorem. No claim that the whole range question was solved is made here.}
\GSfield{Present status}{As of 10 November 2026, the identified publication establishes the stated necessary compatibility and inversion results. The notes are insufficient to isolate an additional exact residual range problem, so no C classification is manufactured from the omitted endpoint data.}
\GSfield{Assessment}{literal statement: not determinable; publication match: partial; verification: read the result; development: annotated; confidence: high; problem \hyperref[prob:GP-1975-13]{GP-1975-13}}
\end{GSitem}
\begin{GSitem}{1975-381.5}{E}{Barycentric cycles representing Stiefel-Whitney classes}
\GSfield{Kind; source}{theorem; \GSpage{P000635}; D. B. Fuchs (intervention)}
\GSfield{The notes}{Fuchs observes that the sum of all edges of an arbitrary triangulation need not be a mod-two cycle, whereas it is a cycle after first barycentric subdivision; the dual classes are Stiefel--Whitney classes.}
\GSfield{Statement}{Let \(K\) be a smooth finite triangulation of a closed smooth \(n\)-manifold \(M\), and let \(K^{\prime}\) be its first barycentric subdivision. For each \(0\le p\le n\), the sum \(c_p\) of all \(p\)-simplices of \(K^{\prime}\), with coefficients in \(\mathbb F_2\), is a cycle and \([c_p]=w_{n-p}(TM)\cap[M]\). In particular, the all-edge cycle represents the homology class dual to \(w_{n-1}\), not generally to \(w_1\). For manifolds with boundary a relative formulation is needed.}
\GSfield{Reconstruction}{The dimension indexing and closed-manifold hypotheses are reconstructed. The notebook drawings support the parity motivation but do not replace the characteristic-class theorem.}
\GSfield{Fate}{The precise mod-two statement was published in \citet[Theorem 1 and its corollary, pp.~511--512]{halperintoledo1972stiefel}; their remarks also discuss the PL extension.}
\GSfield{Present status}{As of 10 November 2026, the result is established in the stated scope.}
\GSfield{Assessment}{literal statement: correct with missing hypotheses; publication match: exact; verification: read the result; development: annotated; confidence: high}
\end{GSitem}
\begin{GSitem}{1975-381.6}{E}{Parity of the simplex count and Euler characteristic}
\GSfield{Kind; source}{claim; \GSpage{P000635}; D. B. Fuchs (intervention)}
\GSfield{The notes}{Fuchs concludes that modulo two the sum of simplex counts equals the alternating sum, hence the Euler characteristic.}
\GSfield{Statement}{For a finite triangulation \(K\), \(\sum_i f_i(K)\equiv\sum_i(-1)^if_i(K)=\chi(|K|)\pmod2\). The vertices of \(K^{\prime}\) correspond to nonempty simplices of \(K\), so its all-vertex zero-cycle has augmentation \(\chi(M)\bmod2\). For a closed \(n\)-manifold this is \(\langle w_n(TM),[M]\rangle\). Finiteness is essential to the elementary numerical sum.}
\GSfield{Reconstruction}{The link between vertices of the subdivision and original simplices explains the otherwise compressed final line.}
\GSfield{Fate}{\citet[\S1, discussion of \(C_0\), p.~512]{halperintoledo1972stiefel} identifies this zero-dimensional case with the Euler characteristic.}
\GSfield{Present status}{As of 10 November 2026, the result is established in the stated scope.}
\GSfield{Assessment}{literal statement: correct; publication match: exact; verification: read the result; development: annotated; confidence: high}
\end{GSitem}

\GStalk{382}{1976-02-23}{M. V. Losik}{P000634-P000638}{The affine bundle of connections and its universal connection}
Losik begins at the bottom of P000635 and continues through P000637. Starting from frames and a principal connection form, he constructs the finite-dimensional affine bundle whose sections are connections, then the tautological connection on the pulled-back principal bundle. Gabrielov identifies the model vector bundle. The construction predates the talk in Garc'ia's 1972 paper; a mistaken base space in the notes is corrected.

\begin{GSitem}{1975-382.1}{E}{Frame bundle and vertical Maurer-Cartan identification}
\GSfield{Kind; source}{definition; \GSpage{P000635-P000636}}
\GSfield{The notes}{For a real rank-\(n\) bundle \(E\to X\), a frame \(z\) is multiplied on the right by \(A\in GL(n,\mathbb R)\). Differentiating the orbit map identifies the vertical tangent space with the Lie algebra.}
\GSfield{Statement}{For a smooth rank-\(n\) real vector bundle, its frame bundle \(P=\operatorname{Fr}(E)\) is a principal right \(GL(n,\mathbb R)\)-bundle. At \(z\in P_x\), the fundamental-vector map \(\alpha_z:\mathfrak{gl}_n\to V_zP\), \(\xi\mapsto\left.\frac{d}{dt}\right|_0 z\exp(t\xi)\), is an isomorphism; write \(\beta_z=\alpha_z^{-1}\). The action is free and transitive within each fiber, not on the whole frame bundle.}
\GSfield{Reconstruction}{The right-action convention and distinction between a fiber and the total space are supplied explicitly.}
\GSfield{Fate}{This principal-bundle framework underlies \citet[\S\S2--3]{garcia1972connections}; the source paper predates the seminar by four years.}
\GSfield{Present status}{As of 10 November 2026, the result is established in the stated scope.}
\GSfield{Assessment}{literal statement: correct; publication match: exact; verification: read the result; development: annotated; confidence: high}
\end{GSitem}
\begin{GSitem}{1975-382.2}{E}{Connection forms and equivariant horizontal complements}
\GSfield{Kind; source}{definition; \GSpage{P000636}}
\GSfield{The notes}{A connection form \(\omega\) restricts to \(\beta_z\) vertically and satisfies \(R_A^*\omega=\operatorname{Ad}(A^{-1})\omega\). Its kernel is the horizontal space.}
\GSfield{Statement}{For a smooth principal right \(G\)-bundle, a principal connection is a \(\mathfrak g\)-valued one-form \(\omega\) satisfying \(\omega(\xi^\#)=\xi\) and \(R_g^*\omega=\operatorname{Ad}_{g^{-1}}\omega\). Equivalently, \(TP=H\oplus VP\) with a smooth \(G\)-equivariant horizontal subbundle. For a frame bundle this corresponds to a linear connection on \(E\); the induced tangent-bundle connection is not automatically torsion-free.}
\GSfield{Reconstruction}{The word affine in the notes is read as a connection on the tangent/frame bundle, without adding a torsion-free condition.}
\GSfield{Fate}{Compare the connection-form and horizontal-distribution formulations used in \citet[\S3, proof of Theorem 3]{garcia1972connections}.}
\GSfield{Present status}{As of 10 November 2026, the result is established in the stated scope.}
\GSfield{Assessment}{literal statement: correct; publication match: exact; verification: read the result; development: annotated; confidence: high}
\end{GSitem}
\begin{GSitem}{1975-382.3}{E}{Affine fibers of pointwise connection data}
\GSfield{Kind; source}{construction; \GSpage{P000636-P000637}}
\GSfield{The notes}{At \(z\), the space \(C_z\) of forms restricting to \(\beta_z\) is an affine plane parallel to the space \(\Theta_z\) of forms vanishing on vertical vectors.}
\GSfield{Statement}{Define \(C_z=\{\eta\in\operatorname{Hom}(T_zP,\mathfrak g):\eta|_{V_zP}=\beta_z\}\) and \(\Theta_z=\operatorname{Hom}(T_zP/V_zP,\mathfrak g)\). Then \(C_z\) is an affine space modeled on \(\Theta_z\); subtraction of two elements is well defined, but there is no distinguished zero. The union \(C(P)\to P\) is a smooth affine bundle with its induced \(G\)-action.}
\GSfield{Reconstruction}{The affine-torsor distinction is explicit: choosing one connection identifies fibers with vector spaces, but the identification is not canonical.}
\GSfield{Fate}{The equivalent jet construction occurs in \citet[\S2, Theorems 1--2]{garcia1972connections}.}
\GSfield{Present status}{As of 10 November 2026, the result is established in the stated scope.}
\GSfield{Assessment}{literal statement: correct with missing hypotheses; publication match: exact; verification: read the result; development: annotated; confidence: high}
\end{GSitem}
\begin{GSitem}{1975-382.4}{E}{Connections as sections of an affine bundle over the base}
\GSfield{Kind; source}{theorem; \GSpage{P000637}}
\GSfield{The notes}{The quotient \(B(P)=C(P)/G\) is formed; connections correspond to its sections.}
\GSfield{Statement}{For a smooth principal \(G\)-bundle \(P\to X\), let \(\mathcal C(P)=C(P)/G\). Smooth sections \(s:X\to\mathcal C(P)\) correspond bijectively to principal connections on \(P\): the equivariant lift chooses \(\omega_z\in C_z\) at every frame. Its affine model is \(T^*X\otimes\operatorname{ad}P\), where \(\operatorname{ad}P=P\times_G\mathfrak g\). For \(P=\operatorname{Fr}(E)\), this is \(T^*X\otimes\operatorname{End}(E)\).}
\GSfield{Reconstruction}{The symbol \(\mathcal C(P)\) avoids confusing the quotient with the affine bundle over \(P\). The adjoint bundle, not a fixed global copy of \(\mathfrak g\), is necessary.}
\GSfield{Fate}{\citet[Theorem 2, p.~236]{garcia1972connections} proves the section/connection correspondence.}
\GSfield{Present status}{As of 10 November 2026, the result is established in the stated scope.}
\GSfield{Assessment}{literal statement: correct with missing hypotheses; publication match: exact; verification: read the result; development: annotated; confidence: high}
\end{GSitem}
\begin{GSitem}{1975-382.5}{E}{The tautological universal connection and its correct base}
\GSfield{Kind; source}{construction; \GSpage{P000637}}
\GSfield{The notes}{The notes call \(C(P)\) a principal bundle with a canonical connection, but state its base as \(P\).}
\GSfield{Statement}{The principal bundle is \(C(P)\to\mathcal C(P)=C(P)/G\), naturally isomorphic to the pullback of \(P\to X\) along \(\mathcal C(P)\to X\). If \(p:C(P)\to P\), define \(\Omega_{(z,\eta)}(V)=\eta(dp(V))\). This reproduces fundamental vertical fields and is \(G\)-equivariant, hence is a principal connection. Pullback by the equivariant lift of a section \(s\) recovers precisely the connection represented by \(s\). The map to \(P\) has affine fibers and is not the asserted principal \(G\)-bundle.}
\GSfield{Reconstruction}{The base-space correction and evaluation formula make the universal property precise; they are the intrinsic content of the displayed construction.}
\GSfield{Fate}{This is the canonical connection of \citet[Theorem 3, pp.~238--240]{garcia1972connections}, expressed in pointwise connection data instead of jets.}
\GSfield{Present status}{As of 10 November 2026, the result is established in the stated scope.}
\GSfield{Assessment}{literal statement: false as written; publication match: exact; verification: read the result; development: annotated; confidence: high}
\end{GSitem}
\begin{GSitem}{1975-382.6}{E}{Difference of connections as an endomorphism-valued one-form}
\GSfield{Kind; source}{remark; \GSpage{P000637}; A. M. Gabrielov (intervention)}
\GSfield{The notes}{Gabrielov identifies the affine model, first writing a fixed Lie-algebra factor and then \(T^*X\otimes\operatorname{Hom}(E,E)\).}
\GSfield{Statement}{If \(\nabla_1,\nabla_2\) are linear connections on the same smooth vector bundle \(E\), their difference is \(C^\infty(X)\)-linear in the section and is a one-form in its vector argument: \(\nabla_1-\nabla_2\in\Omega^1(X;\operatorname{End}E)\). Conversely, adding any such one-form to a connection produces another connection. A fixed global \(\mathfrak{gl}_n\) coefficient space is valid only after a global trivialization of the adjoint bundle.}
\GSfield{Reconstruction}{The final corrected bundle expression is retained; the crossed-out trivial coefficient bundle is not silently promoted to a global assertion.}
\GSfield{Fate}{The affine structure is part of \citet[\S2, Theorem 2]{garcia1972connections}; its local Christoffel-symbol version is standard background.}
\GSfield{Present status}{As of 10 November 2026, the result is established in the stated scope.}
\GSfield{Assessment}{literal statement: correct; publication match: exact; verification: read the result; development: annotated; confidence: high}
\end{GSitem}

\GStalk{383}{1976-03-01}{S. P. Novikov}{P000636-P000640}{Pointed curves, pole orders and the rank-one correspondence}
The catalogue attributes this record to Novikov, but the dated manuscript heading on P000638 explicitly names Drinfeld (Volodya). The mathematical discussion continues through P000641. This provenance discrepancy is retained rather than silently correcting the catalogue. Gelfand and Krichever contribute curve examples and pole-order remarks; S. Gelfand gives a separate divisor explanation on P000639. The main line passes from rings of functions on a pointed curve to filtered modules and acyclic line bundles.

\begin{GSitem}{1975-383.1}{E}{Pointed-curve rings and their genus-zero and elliptic examples}
\GSfield{Kind; source}{construction; \GSpage{P000638}; Gelfand (intervention; initial only)}
\GSfield{The notes}{For a smooth projective curve and marked point, the ring consists of rational functions regular off that point. The examples are a genus-zero conic and \(y^2=x(x-1)(x-2)\), with pole orders two and three for \(x,y\). The latter are identified literally with \(\wp,\wp^{\prime}\).}
\GSfield{Statement}{Let \(X\) be a smooth connected projective complex curve, \(P\in X\), and \(A=H^0(X\setminus\{P\},\mathcal O_X)\). For \((X,P)=(\mathbb P^1,\infty)\), \(A=\mathbb C[z]\). For the smooth projective completion of the displayed cubic and its point at infinity, \(A=\mathbb C[x,y]/(y^2-x(x-1)(x-2))\), with pole orders two and three. With \((\wp^{\prime})^2=4\wp^3-4\wp\), its correctly normalized uniformization is \(x=1+\wp\), \(y=\wp^{\prime}/2\). A genus-one curve has no nonconstant function with a unique simple pole, because that would define a degree-one map to \(\mathbb P^1\).}
\GSfield{Reconstruction}{\oursComputed\ The normalization follows by substituting \(x=1+w\) into the cubic. The degree-one-map argument supplies the reason behind the notebook's torus remark; no new uniformization theorem is asserted.}
\GSfield{Fate}{These are classical curve examples. The later exposition \citet[Ch.~II]{forster1981riemann} is cited for context; only its publisher record and author's contents were checked.}
\GSfield{Present status}{As of 10 November 2026, the result is established in the stated scope.}
\GSfield{Assessment}{literal statement: false as written; publication match: partial; verification: publisher metadata checked; development: annotated; confidence: high}
\end{GSitem}
\begin{GSitem}{1975-383.2}{E}{Recovering the pointed curve from its punctured function ring}
\GSfield{Kind; source}{question; \GSpage{P000638}; V. G. Drinfeld (manuscript heading)}
\GSfield{The notes}{The ring is said to determine the curve and the point uniquely, followed by the question how to describe all such rings.}
\GSfield{Statement}{The rings in question are precisely the finitely generated integrally closed complex domains of transcendence degree one whose smooth projective completion has exactly one point outside its affine spectrum. An isomorphism of such complex algebras identifies their function fields and affine curves, hence their unique smooth projective completions and the omitted points. The one-point condition is necessary: a general smooth affine curve may have several punctures.}
\GSfield{Reconstruction}{The finite-generation, normality, transcendence-degree and one-point conditions are explicit reconstruction of the broad algebraic question. No intrinsic finite list of generators and relations is being claimed.}
\GSfield{Fate}{This is the classical affine/projective correspondence for nonsingular curves; \citet[\S1]{drinfeld1977commutative} uses the fixed pointed-curve ring, without claiming the underlying correspondence as new.}
\GSfield{Present status}{As of 10 November 2026, the result is established in the stated scope.}
\GSfield{Assessment}{literal statement: correct with missing hypotheses; publication match: thematic only; verification: read the result; development: annotated; confidence: high}
\end{GSitem}
\begin{GSitem}{1975-383.3}{E}{Weierstrass gaps and the residue-class formula}
\GSfield{Kind; source}{theorem; \GSpage{P000638}; I. M. Krichever (intervention)}
\GSfield{The notes}{Krichever says that a general point has least pole order \(g+1\); a hyperelliptic branch point permits order two. If \(n\) is the least positive order and \(s+\mu_sn\) the least order in residue class \(s\), then \(\sum_{s=1}^{n-1}\mu_s=g\).}
\GSfield{Statement}{For a smooth connected projective complex curve of genus \(g\), the pole-order semigroup \(H(P)\subset\mathbb Z_{\ge0}\) has exactly \(g\) positive gaps. At a non-Weierstrass point they are \(1,\ldots,g\). If \(n\) is its least positive nongap and \(s+\mu_sn\) is the least nongap congruent to \(s\pmod n\), then the gaps in that residue class are \(s,s+n,\ldots,s+(\mu_s-1)n\); consequently the asserted sum equals \(g\). On a smooth hyperelliptic curve of genus at least one, a branch point has a degree-two function with its only pole there. The hyperelliptic polynomial must have distinct roots.}
\GSfield{Reconstruction}{The residue-class count explains the formula directly once the classical gap theorem is used. The genus and smoothness qualifications exclude degenerations hidden by the abbreviated equation.}
\GSfield{Fate}{This is classical Riemann--Roch/Weierstrass-gap theory; see the later textbook context \citet[Ch.~II, \S\S16--18]{forster1981riemann}. The full textbook was not accessible.}
\GSfield{Present status}{As of 10 November 2026, the result is established in the stated scope.}
\GSfield{Assessment}{literal statement: correct with missing hypotheses; publication match: exact; verification: publisher metadata checked; development: annotated; confidence: high}
\end{GSitem}
\begin{GSitem}{1975-383.4}{E}{Generators at an ordinary point and the missing generator}
\GSfield{Kind; source}{claim; \GSpage{P000639}; V. G. Drinfeld (manuscript heading)}
\GSfield{The notes}{The notes say there are \(g\) generators, then list orders \(g+1,g+2,\ldots,2g+1\).}
\GSfield{Statement}{At a non-Weierstrass point of a genus-\(g\) smooth projective complex curve, choose functions with unique poles of orders \(g+1,\ldots,2g+1\), one for each order. These \(g+1\) functions generate \(A\). Indeed those integers generate the nongap semigroup: an order above \(2g+1\) can repeatedly be reduced by \(g+1\). A monomial with the same pole order as a given function can be scaled to cancel its leading term; induction reduces the remainder to a constant. This asserts a generating set, not an unconditional minimal number of algebra generators.}
\GSfield{Reconstruction}{\ours\ The short leading-term induction is supplied to justify the emended count. It also includes genus zero, when the single generator has order one.}
\GSfield{Fate}{The claim is an elementary consequence of the classical gap structure recalled in \citet[Ch.~II]{forster1981riemann}; the induction here is explicit reconstruction.}
\GSfield{Present status}{As of 10 November 2026, the result is established in the stated scope.}
\GSfield{Assessment}{literal statement: false as written; publication match: partial; verification: publisher metadata checked; development: annotated; confidence: high}
\end{GSitem}
\begin{GSitem}{1975-383.5}{A}{Elliptic homomorphisms and the pole/order rank revisited}
\GSfield{Kind; source}{definition; \GSpage{P000639}; V. G. Drinfeld (manuscript heading)}
\GSfield{The notes}{The differential algebra \(B[D]\), invertible leading coefficients, and the integer relation \(\operatorname{ord}\varphi(a)=n\deg_P a\) are recalled.}
\GSfield{Statement}{For the pointed-curve ring \(A\), a commutative differential complex algebra \((B,\delta)\), and the Ore relation \(Db-bD=\delta(b)\), require \(\varphi:A\to B[D]\) to be unital, nontrivial in positive order, and to have unit leading coefficient for every nonzero image. Then its orders are a fixed positive integer multiple of pole orders, as in item 1975-380.5. The rank is that integer, not the degree of a chosen generator alone.}
\GSfield{Reconstruction}{The nonzero qualifier and positive-order requirement, omitted in the shorthand, are retained from the precise definition.}
\GSfield{Fate}{\citet[\S2, definition and rank lemma]{drinfeld1977commutative} gives the matching formulation.}
\GSfield{Present status}{As of 10 November 2026, the result is established in the stated scope.}
\GSfield{Assessment}{literal statement: correct with missing hypotheses; publication match: exact; verification: read the result; development: annotated; confidence: high}
\end{GSitem}
\begin{GSitem}{1975-383.6}{E}{The order-two/order-three commuting example recalled}
\GSfield{Kind; source}{remark; \GSpage{P000639}; I. M. Gelfand (intervention)}
\GSfield{The notes}{Gelfand interjects that commuting differential operators of orders two and three exist.}
\GSfield{Statement}{The operators \(D^2-2u\) and \(D^3-3uD-3u^{\prime}/2\) commute when \(u^{(3)}=12uu^{\prime}\); elliptic functions satisfying \((u^{\prime})^2=4u^3-g_2u-g_3\), and their degenerations, provide examples. Coefficients are taken locally away from poles. This is the earlier school exercise in another sign convention.}
\GSfield{Reconstruction}{The explicit pair makes the bare existence remark precise; the example is classical.}
\GSfield{Fate}{\citet[\S3, equations (2)--(3)]{mulase1990geometric} records the example and its early history.}
\GSfield{Present status}{As of 10 November 2026, the result is established in the stated scope.}
\GSfield{Assessment}{literal statement: correct; publication match: exact; verification: read the result; development: annotated; confidence: high; problem \hyperref[prob:GP-1975-12]{GP-1975-12}}
\end{GSitem}
\begin{GSitem}{1975-383.7}{A}{Leading coefficient and dependence on the local parameter}
\GSfield{Kind; source}{claim; \GSpage{P000639}; V. G. Drinfeld (manuscript heading)}
\GSfield{The notes}{For \(a=z^{-k}+\cdots\), the notes give \(\varphi(a)=b^kD^{nk}+\cdots\) and say that \(b\) is independent of \(k\) and of the chosen local parameter.}
\GSfield{Statement}{With \(z\) fixed and the rank hypotheses above, a unique \(b\in B^*\) gives the normalized leading coefficients \(b^k\). The pole-order semigroup is cofinite, and multiplicativity makes the leading-coefficient homomorphism extend to its group completion \(\mathbb Z\). However, if \(\widetilde z=cz+O(z^2)\), \(c\ne0\), then \(\widetilde b=c^{-1}b\). Thus the scalar coefficient is not coordinate invariant. The corresponding tangent-direction datum, with the appropriate reciprocal coordinate factor, is intrinsic.}
\GSfield{Reconstruction}{\oursComputed\ The rescaling follows by replacing the normalized function by \(c^{-k}a\); the notebook's claim of independence of the parameter is emended.}
\GSfield{Fate}{The intrinsic principal-part datum is used in \citet[Corollary 2]{drinfeld1977commutative}.}
\GSfield{Present status}{As of 10 November 2026, the result is established in the stated scope.}
\GSfield{Assessment}{literal statement: false as written; publication match: partial; verification: read the result; development: annotated; confidence: high}
\end{GSitem}
\begin{GSitem}{1975-383.8}{E}{Gauge equivalence and local removal of a scalar connection term}
\GSfield{Kind; source}{definition; \GSpage{P000639}; V. G. Drinfeld (manuscript heading)}
\GSfield{The notes}{Two homomorphisms are equivalent under conjugation by \(b\in B^*\); the notes write \(D\sim D+b^{\prime}/b\), with an arbitrary function obtainable as a logarithmic derivative among germs.}
\GSfield{Statement}{For a smooth nonvanishing complex function germ \(h\), \(h^{-1}Dh=D+h^{-1}h^{\prime}\), while \(hDh^{-1}=D-h^{-1}h^{\prime}\). The equivalence relation allows either orientation. Every smooth complex germ \(a\) is a logarithmic derivative of a unit germ, by solving \(h^{\prime}=ah\) with \(h(0)\ne0\). This local statement does not imply a periodic or global gauge exists.}
\GSfield{Reconstruction}{The conjugation orientation is made explicit. There is no sign correction to the notebook's symmetric equivalence symbol.}
\GSfield{Fate}{The equivalence relation appears in \citet[\S2]{drinfeld1977commutative}; its scalar differential calculation is elementary background.}
\GSfield{Present status}{As of 10 November 2026, the result is established in the stated scope.}
\GSfield{Assessment}{literal statement: correct; publication match: exact; verification: read the result; development: annotated; confidence: high}
\end{GSitem}
\begin{GSitem}{1975-383.9}{A}{Rank-one homomorphisms from a leading coefficient and a Jacobian point}
\GSfield{Kind; source}{theorem; \GSpage{P000639}; V. G. Drinfeld (manuscript heading)}
\GSfield{The notes}{For germs of smooth functions on the line, rank-one homomorphisms up to equivalence correspond to their leading coefficient and a point of the Jacobian.}
\GSfield{Statement}{Fix \((X,P,z)\) and let \(B\) be the algebra of smooth complex-valued germs at \(x=0\), with \(D=d/dx\). Rank-one Krichever homomorphisms \(A\to B[D]\), modulo unit scalar conjugation, correspond to pairs \((L_0,b)\), where \(b\in B^*\) and \(L_0\in\operatorname{Pic}^{g-1}(X)\) satisfies \(H^0(X,L_0)=H^1(X,L_0)=0\). The allowed line bundles form the complement of the theta divisor. Identifying this torsor with \(\operatorname{Jac}(X)=\operatorname{Pic}^0(X)\) requires a reference line bundle of degree \(g-1\). The construction is local in the coefficient parameter.}
\GSfield{Reconstruction}{The acyclic open subset, torsor identification, and germ setting are essential qualifications; an arbitrary point of the full Jacobian is not admissible.}
\GSfield{Fate}{\citet[Corollary 2]{drinfeld1977commutative} gives this precise correspondence, attributing it to Krichever.}
\GSfield{Present status}{As of 10 November 2026, the result is established in the stated scope.}
\GSfield{Assessment}{literal statement: correct with missing hypotheses; publication match: exact; verification: read the result; development: annotated; confidence: high; problem \hyperref[prob:GP-1975-14]{GP-1975-14}}
\end{GSitem}
\begin{GSitem}{1975-383.10}{A}{Recovering the homomorphism from a cyclic filtered module}
\GSfield{Kind; source}{construction; \GSpage{P000640}; V. G. Drinfeld (manuscript heading)}
\GSfield{The notes}{The module construction is repeated in rank one. The first quotient is written \(M_i/M_{i+1}\), despite the increasing filtration; subsequent displays use \(M_i/M_{i-1}\). A generator \(e\in M_0\) recovers the differential operator realization.}
\GSfield{Statement}{Use an exhaustive filtration \(0=M_{-1}\subset M_0\subset\cdots\), with each \(M_i/M_{i-1}\) a free rank-one \(B\)-module. Require the graded multiplication by \(a\ne0\) to shift by \(\deg_Pa\) isomorphically, and an \(A\)-linear operator \(\nabla\) with the \(B\)-Leibniz rule and isomorphisms between adjacent graded pieces. For a generator \(e\in M_0\), \(e,\nabla e,\ldots,\nabla^ie\) form a \(B\)-basis of \(M_i\), identifying \(M\) with \(B[D]\). The commuting \(A\)-action is right multiplication by the recovered \(\varphi(a)\); another generator changes only its scalar gauge class.}
\GSfield{Reconstruction}{The backward quotient is corrected, dimension over a possibly non-field is replaced by free rank, and the inductive meaning of the notebook's put \(D=\nabla(e)\) is specified.}
\GSfield{Fate}{\citet[Proposition 2, proof]{drinfeld1977commutative} formalizes this argument.}
\GSfield{Present status}{As of 10 November 2026, the result is established in the stated scope.}
\GSfield{Assessment}{literal statement: false as written; publication match: exact; verification: read the result; development: annotated; confidence: high}
\end{GSitem}
\begin{GSitem}{1975-383.11}{A}{Acyclic line bundles produce the pole filtration}
\GSfield{Kind; source}{construction; \GSpage{P000640-P000641}; V. G. Drinfeld (manuscript heading)}
\GSfield{The notes}{The filtration corresponds to a line bundle with vanishing zeroth and first cohomology, and \(M_i\) consists of sections with pole order at most \(i+1\) at \(P\).}
\GSfield{Statement}{For a line bundle \(L\) on the smooth projective complex curve \(X\), assume \(H^0(X,L)=H^1(X,L)=0\). Then \(\deg L=g-1\). Put \(M=H^0(X\setminus\{P\},L)\), \(M_i=H^0(X,L((i+1)P))\) for \(i\ge-1\). The successive exact sequences at \(P\) and the vanishing of \(H^1(X,L)\) give \(\dim M_i=i+1\) and \(\dim M_i/M_{i-1}=1\) for \(i\ge0\). Over a coefficient algebra, use a locally free family with the corresponding cohomology and base-change conditions, not a tensor product of a fixed fiber assumed without justification.}
\GSfield{Reconstruction}{The shift by one is retained exactly from the scan. The family qualification separates the displayed complex-vector-space argument from its relative version.}
\GSfield{Fate}{\citet[Proposition 1 and Corollary 1]{drinfeld1977commutative} supplies the relative sheaf statement.}
\GSfield{Present status}{As of 10 November 2026, the result is established in the stated scope.}
\GSfield{Assessment}{literal statement: correct with missing hypotheses; publication match: exact; verification: read the result; development: annotated; confidence: high}
\end{GSitem}
\begin{GSitem}{1975-383.12}{A}{Reconstructing the line bundle on two affine charts}
\GSfield{Kind; source}{method; \GSpage{P000641}; V. G. Drinfeld (manuscript heading)}
\GSfield{The notes}{For nonconstant \(a_0\in A\), the cover is \((X\setminus P)\cup(X\setminus\operatorname{zeros}(a_0))\). On the second chart sections are \(m/a_0^k\), with \(m\in M_{k\deg(a_0)-1}\).}
\GSfield{Statement}{Let \(M\) have the rank-one exhaustive filtration and multiplication isomorphisms of the preceding items, and write \(d=\deg_Pa_0>0\). Over \(U=X\setminus\{P\}\) use the module \(M\); over \(V=X\setminus\operatorname{zeros}(a_0)\) use the union of fractions \(a_0^{-k}M_{kd-1}\) for \(k\ge0\), inside \(M[a_0^{-1}]\). On \(U\cap V\) both localize to \(M[a_0^{-1}]\), so they glue. The filtration axioms give the line bundle extending across \(P\); the bound \(kd-1\) precisely cancels the allowed pole of order \(kd\).}
\GSfield{Reconstruction}{The common localization and union are supplied to make the notebook's chart formula unambiguous; local freeness relies on the earlier filtration hypotheses.}
\GSfield{Fate}{Compare the reconstruction proof in \citet[Proposition 1]{drinfeld1977commutative}.}
\GSfield{Present status}{As of 10 November 2026, the result is established in the stated scope.}
\GSfield{Assessment}{literal statement: correct with missing hypotheses; publication match: exact; verification: read the result; development: annotated; confidence: high}
\end{GSitem}

\GStalk{384}{1976-03-01}{S. I. Gelfand}{P000638-P000643}{Divisors and the analytic Jacobian}
The explicitly labelled S. Gelfand intervention is the short explanation of degree-zero divisor classes at the foot of P000639. The next page labels the analytic Jacobian explanation only G.; that initial does not by itself distinguish S. I. from I. M. Gelfand. Its mathematical content is indexed here as the continuation of the explanatory intervention, with the attribution uncertainty stated. Drinfeld's filtered-module and sheaf construction, including all of P000641, remains with record 383.

\begin{GSitem}{1975-384.1}{E}{Degree-zero divisor classes and line bundles}
\GSfield{Kind; source}{definition; \GSpage{P000639}}
\GSfield{The notes}{S. Gelfand explains line bundles of degree zero as divisors \(\sum n_iQ_i\), \(\sum n_i=0\), modulo divisors of meromorphic functions.}
\GSfield{Statement}{For a smooth connected projective complex curve, every line bundle admits a nonzero meromorphic section. Sending a divisor \(D\) to \(\mathcal O_X(D)\) identifies divisors modulo principal divisors with \(\operatorname{Pic}(X)\), and identifies degree-zero classes with \(\operatorname{Pic}^0(X)\). Changing the meromorphic section changes its divisor by a principal divisor, so its class is unchanged.}
\GSfield{Reconstruction}{The all-line-bundle statement is included only to explain the degree-zero restriction; it does not identify line bundles of different degree.}
\GSfield{Fate}{This is classical divisor theory; \citet[Ch.~II, \S20 and Ch.~III, \S29]{forster1981riemann} is a later reference whose contents were checked.}
\GSfield{Present status}{As of 10 November 2026, the result is established in the stated scope.}
\GSfield{Assessment}{literal statement: correct; publication match: exact; verification: publisher metadata checked; development: annotated; confidence: high}
\end{GSitem}
\begin{GSitem}{1975-384.2}{E}{The period lattice and the analytic Jacobian}
\GSfield{Kind; source}{definition; \GSpage{P000640}; Gelfand (initial only; attribution uncertain)}
\GSfield{The notes}{The analytic definition uses \(g\) holomorphic differentials and \(2g\) period vectors, yielding \(\mathbb C^g/\Lambda\). On a smooth elliptic cubic, \(dx/y\) has no singularities even at branch points. Only some complex tori arise in this manner.}
\GSfield{Statement}{For a compact connected Riemann surface of genus \(g\), \(V=H^0(X,\Omega_X^1)\) has complex dimension \(g\). Integration embeds \(H_1(X,\mathbb Z)\) as a lattice \(\Lambda\subset V^*\) of rank \(2g\); \(\operatorname{Jac}(X)=V^*/\Lambda\), canonically isomorphic to \(\operatorname{Pic}^0(X)\). On the nonsingular completion of \(y^2=x(x-a)(x-b)\), \(ab(a-b)\ne0\), \(dx/y\) extends holomorphically and without zeros. A branch coordinate such as \(y\), rather than \(x\), verifies regularity there. Jacobians carry a principal polarization; arbitrary complex tori need not be Jacobians.}
\GSfield{Reconstruction}{The dual vector space, lattice rank, smoothness condition and polarization qualification are supplied. The note does not formulate a solution of the Schottky problem.}
\GSfield{Fate}{The classical analytic construction is the subject of \citet[Ch.~II, \S\S20--21]{forster1981riemann}; verification here is bibliographic.}
\GSfield{Present status}{As of 10 November 2026, the result is established in the stated scope.}
\GSfield{Assessment}{literal statement: correct with missing hypotheses; publication match: exact; verification: publisher metadata checked; development: annotated; confidence: high}
\end{GSitem}
\begin{GSitem}{1975-384.3}{E}{Jacobi inversion converts linear motion to moving divisors}
\GSfield{Kind; source}{method; \GSpage{P000640}; Gelfand (initial only; attribution uncertain)}
\GSfield{The notes}{A straight line on the torus produces motion of \(g\) points on the curve.}
\GSfield{Statement}{Fix a base point \(Q_0\). The Abel map \(\operatorname{Sym}^gX\to\operatorname{Jac}(X)\), \(Q_1+\cdots+Q_g\mapsto[Q_1+\cdots+Q_g-gQ_0]\), is generically one-to-one: a degree-\(g\) line bundle with \(h^0=1\) has a unique effective divisor. Thus a path \(Z(t)=Z_0+tU\) avoiding the exceptional/special-divisor locus locally determines an unordered moving divisor. Choosing and smoothly following individual ordered points further requires avoiding collisions. The inversion need not be globally single-valued in ordered points.}
\GSfield{Reconstruction}{The generic-locus and collision restrictions make precise the drawing's motion statement. No differential equation for the individual points is asserted in the notes.}
\GSfield{Fate}{This is the classical Jacobi inversion problem; compare \citet[Ch.~II, \S21]{forster1981riemann}.}
\GSfield{Present status}{As of 10 November 2026, the result is established in the stated scope.}
\GSfield{Assessment}{literal statement: correct with missing hypotheses; publication match: exact; verification: publisher metadata checked; development: annotated; confidence: high}
\end{GSitem}

\GStalk{385}{1976-03-15}{V. G. Drinfeld}{P000641-P000645}{Rank-one sheaves and the tangent direction of the Jacobian flow}
The dated talk starts on P000642, recapitulates the filtered-module correspondence, and identifies its rank-one geometric data and Jacobian velocity. It ends in the upper portion of P000644, before the separate heading A. Polyakov. Gelfand's remark about an operator whose higher powers are differential is retained as an independent item; the later sheaf argument requires smooth function germs, a fixed leading coefficient and an acyclic initial line bundle.

\begin{GSitem}{1975-385.1}{A}{Rank-one homomorphisms and filtered modules revisited}
\GSfield{Kind; source}{construction; \GSpage{P000642}}
\GSfield{The notes}{The talk restates the rank-one order condition, scalar-gauge equivalence, exhaustive filtration and the two graded isomorphism conditions.}
\GSfield{Statement}{Fix the smooth pointed projective curve and differential complex algebra \(B\). The rank-one Krichever homomorphism conditions and the filtered-module axioms of 1975-383.10 are equivalent, modulo a choice of generator of the free rank-one module \(M_0\). Explicitly, \(M=B[D]\), \(a\) acts by right multiplication with \(\varphi(a)\), and \(\nabla\) is left multiplication by \(D\). The nonzero images must have unit leading coefficient; zero has no such coefficient.}
\GSfield{Reconstruction}{The equality sign is restored in the opening definition \(A=H^0(X\setminus\{P\},\mathcal O_X)\), where the notes write membership. All subsequent axioms require the whole ring.}
\GSfield{Fate}{The published comparison is \citet[Proposition 2]{drinfeld1977commutative}.}
\GSfield{Present status}{As of 10 November 2026, the result is established in the stated scope.}
\GSfield{Assessment}{literal statement: correct with missing hypotheses; publication match: exact; verification: read the result; development: annotated; confidence: high}
\end{GSitem}
\begin{GSitem}{1975-385.2}{E}{A nondifferential operator whose every higher power is differential}
\GSfield{Kind; source}{claim; \GSpage{P000642}; I. M. Gelfand (intervention)}
\GSfield{The notes}{Gelfand states that there is an operator whose second, third and all subsequent powers are differential operators, although the first is not.}
\GSfield{Statement}{Interpret operator in the ring of formal pseudodifferential operators over \(\mathbb C(x)\). Set \(P=D^2-2/x^2\), \(R=D^3-3x^{-2}D+3x^{-3}\). They commute and satisfy \(R^2=P^3\). Hence \(Q=RP^{-1}\) obeys \(Q^2=P\), \(Q^3=R\); since every integer at least two is a nonnegative combination of two and three, all \(Q^m\), \(m\ge2\), are differential. But \(Q\) is not differential: its order and leading coefficient would force \(Q=D+a\), and comparison of the \(D\)-coefficient in \(Q^2=P\) gives \(a=0\), a contradiction. The same example works with smooth germs near any nonzero point.}
\GSfield{Reconstruction}{\oursComputed\ The rational degeneration and algebraic verification supply a precise example; the source does not specify the operator category or display the example. This is a verification of classical formulas, not a new existence result.}
\GSfield{Fate}{\citet[\S3, equations (2)--(3), and \S4]{mulase1990geometric} gives the commuting family and formal fractional-power framework.}
\GSfield{Present status}{As of 10 November 2026, the result is established in the stated scope.}
\GSfield{Assessment}{literal statement: correct with missing hypotheses; publication match: partial; verification: read the result; development: annotated; confidence: high}
\end{GSitem}
\begin{GSitem}{1975-385.3}{E}{The gap count recalled in residue-class form}
\GSfield{Kind; source}{theorem; \GSpage{P000642-P000643}}
\GSfield{The notes}{The least nongap \(n\), the integers \(\mu_i\), their sum \(g\), and the existence of exactly \(g\) missing pole orders are repeated.}
\GSfield{Statement}{For a smooth connected projective complex curve of genus \(g\) and a point \(P\), the pole-order semigroup has \(g\) positive gaps. With \(n\) its least positive element and \(i+\mu_i n\) the least element in residue class \(i\), the count is \(g=\sum_{i=1}^{n-1}\mu_i\), as in 1975-383.3. This is a statement about functions having no poles away from \(P\), not about unrestricted meromorphic functions.}
\GSfield{Reconstruction}{The sole-pole condition is made explicit; the equivalence follows by counting each residue class.}
\GSfield{Fate}{This repeats the classical gap theory discussed above; \citet[Ch.~II]{forster1981riemann}.}
\GSfield{Present status}{As of 10 November 2026, the result is established in the stated scope.}
\GSfield{Assessment}{literal statement: correct; publication match: exact; verification: publisher metadata checked; development: annotated; confidence: high}
\end{GSitem}
\begin{GSitem}{1975-385.4}{A}{The relative acyclic line bundle and its simple-pole connection}
\GSfield{Kind; source}{construction; \GSpage{P000643}}
\GSfield{The notes}{The module data, apart from \(\nabla\), specify a line bundle over \(X\otimes B\) with zero cohomology. The remaining operator maps \(L\) to \(L(P)\), is \(\mathcal O_X\)-linear, and satisfies the \(B\)-Leibniz rule.}
\GSfield{Statement}{Let \(X_B=X\times\operatorname{Spec}B\). In rank one, the geometric datum is an invertible \(\mathcal O_{X_B}\)-module \(L\) with \(H^0(X_B,L)=H^1(X_B,L)=0\), a trivialization of its rank-one restriction along \(P\) up to the equivalence implicit in the filtered data, and an \(\mathcal O_X\)-linear map \(\nabla:L\to L(P)\) satisfying \(\nabla(bs)=b\nabla s+\delta(b)s\). Its induced polar map \(L/L(-P)\to L(P)/L\) must be an isomorphism. This last condition is the geometric unit-leading-coefficient condition.}
\GSfield{Reconstruction}{The polar isomorphism is stated explicitly. For the later dynamical interpretation, \(B\) specializes to smooth germs; arbitrary differential algebras do not automatically possess a value-at-zero map.}
\GSfield{Fate}{\citet[Proposition 2 and proof of Corollary 2]{drinfeld1977commutative} gives this formulation.}
\GSfield{Present status}{As of 10 November 2026, the result is established in the stated scope.}
\GSfield{Assessment}{literal statement: correct with missing hypotheses; publication match: exact; verification: read the result; development: annotated; confidence: high}
\end{GSitem}
\begin{GSitem}{1975-385.5}{A}{The reciprocal polar coefficient of the connection}
\GSfield{Kind; source}{claim; \GSpage{P000643}}
\GSfield{The notes}{After choosing a local frame \(e\), the notes write \(\nabla e/e=1/(b_0z)+\cdots\), where \(b_0\) is the leading coefficient of the homomorphism.}
\GSfield{Statement}{Fix the curve coordinate \(z\), and a local generator \(e\) of \(L\) near \(P\). If \(\nabla e=(b_0^{-1}z^{-1}+O(1))e\), \(b_0\in B^*\), then the principal-part datum is \(\beta=b_0^{-1}[z^{-1}]\). Multiplication by \(z^{-m}\) on the associated graded corresponds to \(b_0^m\nabla^m\), giving leading coefficient \(b_0^m\) for the operator attached to a normalized pole of order \(m\). A regular change of frame contributes only a regular logarithmic-derivative term.}
\GSfield{Reconstruction}{The reciprocal factor is retained from the manuscript, not changed to \(b_0\); the associated-graded argument explains it.}
\GSfield{Fate}{The datum \(\beta\) is exactly the generator used in \citet[Corollary 2]{drinfeld1977commutative}.}
\GSfield{Present status}{As of 10 November 2026, the result is established in the stated scope.}
\GSfield{Assessment}{literal statement: correct; publication match: exact; verification: read the result; development: annotated; confidence: high}
\end{GSitem}
\begin{GSitem}{1975-385.6}{A}{Initial value and leading coefficient in the rank-one inverse correspondence}
\GSfield{Kind; source}{theorem; \GSpage{P000643}}
\GSfield{The notes}{The theorem is abbreviated to a point of the Jacobian in general position. The construction then uses a natural map \(B\to\mathbb C\), value at zero, to obtain \(L_0\).}
\GSfield{Statement}{For \(B\) the algebra of smooth complex-valued germs at zero, evaluation \(B\to\mathbb C\) pulls the family back to its initial line bundle \(L_0\). Fixing a unit leading-coefficient germ \(b_0\), the gauge classes are parametrized by acyclic \(L_0\in\operatorname{Pic}^{g-1}(X)\). If \(b_0\) is not fixed, it is additional data and the Jacobian point alone is insufficient. For an arbitrary differential algebra \(B\), neither evaluation at zero nor this initial-value classification is automatic.}
\GSfield{Reconstruction}{The opening general differential-algebra notation is narrowed to the germ setting when evaluation and uniqueness of the path are invoked.}
\GSfield{Fate}{\citet[Corollary 2]{drinfeld1977commutative} states both pieces of data explicitly.}
\GSfield{Present status}{As of 10 November 2026, the result is established in the stated scope.}
\GSfield{Assessment}{literal statement: correct with missing hypotheses; publication match: exact; verification: read the result; development: annotated; confidence: high; problem \hyperref[prob:GP-1975-14]{GP-1975-14}}
\end{GSitem}
\begin{GSitem}{1975-385.7}{E}{The Picard group and the tangent space of the Jacobian}
\GSfield{Kind; source}{definition; \GSpage{P000643}}
\GSfield{The notes}{The tangent space is identified with \(H^1(X,\mathcal O_X)\), but the notes identify the whole Jacobian with \(H^1(X,\mathcal O_X^*)\).}
\GSfield{Statement}{For a smooth compact connected complex curve, \(H^1(X,\mathcal O_X^*)=\operatorname{Pic}(X)\) includes all integer degrees. Its identity component is \(\operatorname{Pic}^0(X)=\operatorname{Jac}(X)\). The analytic exponential sequence gives \(\operatorname{Pic}^0(X)\cong H^1(X,\mathcal O_X)/\operatorname{im}H^1(X,\mathbb Z)\), with the conventional \(2\pi i\) factor according to the chosen exponential. Its tangent space at the identity, and by translation at any point, is \(H^1(X,\mathcal O_X)\). A fixed-degree Picard variety is a torsor for this group.}
\GSfield{Reconstruction}{The missing degree-zero qualification is corrected; removing the star is interpreted as linearizing transition functions, not as a literal cohomology operation.}
\GSfield{Fate}{The tangent-space identification is used in \citet[Corollary 2, proof]{drinfeld1977commutative}; the surrounding Picard theory is classical.}
\GSfield{Present status}{As of 10 November 2026, the result is established in the stated scope.}
\GSfield{Assessment}{literal statement: false as written; publication match: exact; verification: read the result; development: annotated; confidence: high}
\end{GSitem}
\begin{GSitem}{1975-385.8}{A}{The principal-part boundary map gives the Jacobian velocity}
\GSfield{Kind; source}{theorem; \GSpage{P000643-P000644}}
\GSfield{The notes}{The class of \(1/z\) under the boundary map from \(\mathcal O_X(P)/\mathcal O_X\) is denoted \(\eta_P\), and \(\ell^{\prime}(x)=\eta_P/b_0(x)\). The final explanation identifies its line with the tangent to the Abel image of the curve.}
\GSfield{Statement}{Let \(\delta_P:H^0(X,\mathcal O_X(P)/\mathcal O_X)\to H^1(X,\mathcal O_X)\) be the connecting homomorphism and put \(\eta_P=\delta_P([z^{-1}])\). For the smooth-germ rank-one data above, choose the Cech and divisor sign conventions consistently. The corresponding Picard path satisfies \(\ell^{\prime}(x)=b_0(x)^{-1}\eta_P\), hence locally \(\ell(x)=\ell(0)+\eta_P\int_0^x b_0(s)^{-1}\,ds\) in a lift to the universal cover. For \(g\ge1\), the line spanned by \(\eta_P\) is the tangent image of \(X\to\operatorname{Jac}(X)\), \(Q\mapsto[Q-P]\); a chosen tangent coordinate fixes its scale and convention fixes its sign. For genus zero the Jacobian and this vector vanish.}
\GSfield{Reconstruction}{\oursComputed\ Integrating the stated velocity gives the displayed local lift. The genus-zero exception and sign convention are explicit reconstruction; no new flow theorem is claimed.}
\GSfield{Fate}{\citet[Corollary 2, proof]{drinfeld1977commutative} identifies the derivative with the boundary image of the principal part.}
\GSfield{Present status}{As of 10 November 2026, the result is established in the stated scope.}
\GSfield{Assessment}{literal statement: correct with missing hypotheses; publication match: exact; verification: read the result; development: annotated; confidence: high; problem \hyperref[prob:GP-1975-14]{GP-1975-14}}
\end{GSitem}

\GStalk{386}{1976-03-15}{A. M. Polyakov}{P000643-P000649}{Canonical quantization of Yang–Mills fields}
The mathematical text starts on P000644 after the conclusion of the preceding algebro-geometric discussion. Polyakov passes from the Yang--Mills Lagrangian to constrained Hamiltonian mechanics, separates longitudinal and transverse electric fields in Coulomb gauge, and proposes averaging the resulting charge kernel in the quantum vacuum. Several numerical factors and commutator signs are inconsistent in the manuscript; the reconstructed formulas below use one convention throughout. P000647 points ahead to the continuation on 5 April.

\begin{GSitem}{1975-386.1}{E}{Yang–Mills potentials, curvature and action}
\GSfield{Kind; source}{definition; \GSpage{P000644}}
\GSfield{The notes}{The notes take traceless anti-Hermitian \(2\times2\) matrices \(A_\mu\), write \(A_\mu\in SU(2)\) followed by ``Lie algebra'', and define \(F_{\mu\nu}=\partial_\mu A_\nu-\partial_\nu A_\mu+g[A_\mu,A_\nu]\).}
\GSfield{Statement}{A connection on a trivial \(SU(2)\)-bundle is a one-form \(A=A_\mu dx^\mu\) with \(A_\mu\in\mathfrak{su}(2)\). Its curvature is \(F=dA+gA\wedge A\). Use the positive invariant inner product \(\langle X,Y\rangle=-\operatorname{tr}(XY)\) on the Lie algebra. The Euclidean action is \(S_E=\frac14\int\langle F_{\mu\nu},F_{\mu\nu}\rangle\,d^4x\); Lorentzian contraction instead uses the spacetime metric.}
\GSfield{Fate}{These are the classical definitions; see \citet[Section 1]{jaffewitten2000yangmills}. The source's positive trace formula and anti-Hermitian convention cannot both give a nonnegative Euclidean action.}
\GSfield{Present status}{As of November 2026, the reconstructed statement is standard.}
\GSfield{Assessment}{literal statement: false as written; publication match: partial; verification: read the result; development: annotated; confidence: high}
\end{GSitem}
\begin{GSitem}{1975-386.2}{E}{Legendre transformation and Gauss constraint}
\GSfield{Kind; source}{construction; \GSpage{P000644-P000645}}
\GSfield{The notes}{The notes set \(e_i=\dot A_i-D_iA_0\), observe the absence of \(\dot A_0\), and introduce \(h_i=\frac12\varepsilon_{ijk}F_{jk}\). Both displayed Hamiltonians put a factor \(1/2\) on the \(A_0\)-term.}
\GSfield{Statement}{With \(L=\frac12\int(\langle e,e\rangle-\langle h,h\rangle)\), the canonical momenta are \(e_i\) and \(e_0=0\). Integration by parts, with boundary terms zero, gives
\[H=\frac12\int(\langle e,e\rangle+\langle h,h\rangle)\,d^3x-\int\langle A_0,D_ie_i\rangle\,d^3x.\]
Variation in \(A_0\) imposes \(D_ie_i=0\). The coefficient of this last term is one, with the conventions just fixed.}
\GSfield{Reconstruction}{\ours\ The displayed result follows directly from \(\dot A=e+DA_0\). A positive Lie-algebra inner product replaces the ambiguous bare trace. No global Hamiltonian phase space is asserted without boundary and function-space choices.}
\GSfield{Fate}{Canonical constraint reduction predates the seminar: \citet[abstract]{arnowittfickler1962quantization}; the reconstructed Hamiltonian also agrees with \citet[Section 2, equations (1)--(5)]{reinhardt2018coulomb}. The 1962 paper's metadata and abstract, rather than its full derivation, were accessible.}
\GSfield{Present status}{As of November 2026, the reconstructed statement is standard.}
\GSfield{Assessment}{literal statement: false as written; publication match: partial; verification: read the result; development: annotated; confidence: high}
\end{GSitem}
\begin{GSitem}{1975-386.3}{E}{Propagation of the Gauss constraint}
\GSfield{Kind; source}{question; \GSpage{P000645}}
\GSfield{The notes}{The speaker asks whether \(D_ie_i=0\) continues to hold and writes a Poisson-bracket condition for \(\rho=D_ie_i\).}
\GSfield{Statement}{On the canonical phase space, set \(G(\omega)=\int\langle\omega,D_ie_i\rangle\). The gauge-invariant reduced Hamiltonian \(H_0=\frac12\int(e^2+h^2)\) obeys \(\{G(\omega),H_0\}=0\). Consequently Gauss's constraint propagates in temporal gauge; with an arbitrary multiplier \(A_0\), its evolution is an adjoint rotation and preserves the zero constraint surface.}
\GSfield{Reconstruction}{\ours\ The source's bracket is the negative of the convention \(\{A_i,e_j\}=\delta_{ij}\delta\), although its Hamilton equations use that latter convention. The vanishing assertion is unaffected; the sign must be fixed when generating gauge transformations.}
\GSfield{Fate}{This is part of the pre-existing first-order canonical formulation; \citet[abstract, discussion of constraint conservation]{arnowittfickler1962quantization}. The proof here is the elementary implication of gauge invariance, not an attribution of a newly announced theorem.}
\GSfield{Present status}{As of November 2026, the reconstructed statement is standard.}
\GSfield{Assessment}{literal statement: correct with missing hypotheses; publication match: partial; verification: publisher metadata checked; development: annotated; confidence: high}
\end{GSitem}
\begin{GSitem}{1975-386.4}{E}{Gauge covariance and transverse coordinates}
\GSfield{Kind; source}{construction; \GSpage{P000645}}
\GSfield{The notes}{The transformation \(A_i^\Omega=\Omega^{-1}A_i\Omega+g^{-1}\Omega^{-1}\partial_i\Omega\) preserves \(\operatorname{tr}h^2\). The notes then write every field in terms of \(a_i\) with \(\partial_i a_i=0\) and a gauge matrix.}
\GSfield{Statement}{For smooth \(\Omega:\mathbb R^3\to SU(2)\), \(F_{ij}^\Omega=\Omega^{-1}F_{ij}\Omega\), so the magnetic energy is gauge invariant. Coulomb gauge \(\partial_i a_i=0\) is a local slice where its linearization \(M=-\partial_iD_i(a)\) is invertible after boundary conditions and residual constant transformations are fixed.}
\GSfield{Reconstruction}{A local slice, not a globally unique decomposition, is the justified reading. Distinct gauge-equivalent transverse fields may exist.}
\GSfield{Fate}{The Coulomb formulation remains useful, with its gauge-fixing determinant and Gribov restriction made explicit in \citet[Section 2]{reinhardt2018coulomb}.}
\GSfield{Present status}{As of November 2026, the reconstructed statement is standard.}
\GSfield{Assessment}{literal statement: correct with missing hypotheses; publication match: partial; verification: read the result; development: annotated; confidence: high}
\end{GSitem}
\begin{GSitem}{1975-386.5}{E}{Longitudinal electric field and the Coulomb kernel}
\GSfield{Kind; source}{construction; \GSpage{P000646}}
\GSfield{The notes}{The notes split \(e_i=P_i+\partial_i\Psi\), \(\partial_iP_i=0\), and call the longitudinal energy Coulomb energy. The commutator on the right of the Gauss equation changes order twice.}
\GSfield{Statement}{Fix \(D_i=\partial_i+g[a_i,\cdot]\) and \(D_ie_i=\rho\). Then
\[D_i\partial_i\Psi=\rho-g[a_i,P_i]=:\rho_{\rm tot}.\]
For transverse \(a\), \(M=-\partial D=-D\partial\), so \(\Psi=-M^{-1}\rho_{\rm tot}\) and
\[H_C=\tfrac12\langle\rho_{\rm tot},M^{-1}(-\Delta)M^{-1}\rho_{\rm tot}\rangle\]
when \(M\) is self-adjoint and invertible on the chosen domain. The cross term \(\int P_i\partial_i\Psi\) vanishes by integration by parts.}
\GSfield{Reconstruction}{\ours\ The sign is forced by the source's definitions, rather than chosen to imitate a later convention for the gluonic charge. The notes omit \(g\) in the first version and reverse the bracket in the later versions.}
\GSfield{Fate}{The source explicitly credits Schwinger. His 1962 radiation-gauge paper is an identified historical match at abstract level; the modern kernel is displayed in \citet[Section 2, equation (5)]{reinhardt2018coulomb}.}
\GSfield{Present status}{As of November 2026, the reconstructed statement is standard.}
\GSfield{Assessment}{literal statement: false as written; publication match: partial; verification: read the result; development: annotated; confidence: high}
\end{GSitem}
\begin{GSitem}{1975-386.6}{E}{Recovery of the Abelian Coulomb energy}
\GSfield{Kind; source}{claim; \GSpage{P000646}}
\GSfield{The notes}{For \(g=0\) the notes write \(\Delta\Psi=\rho\) and identify the energy with a double integral having kernel \(1/|x-y|\).}
\GSfield{Statement}{For a sufficiently decaying smooth charge on \(\mathbb R^3\), \(\Psi(x)=-\int\rho(y)/(4\pi|x-y|)\,dy\) solves \(\Delta\Psi=\rho\). Thus
\[\tfrac12\int|\nabla\Psi|^2\,dx=\tfrac1{8\pi}\iint\frac{\rho(x)\rho(y)}{|x-y|}\,dx\,dy.\]
Point charges require subtraction or regularization of divergent self energies.}
\GSfield{Reconstruction}{\ours\ This is an explicit reconstruction using \(-\Delta(4\pi|x|)^{-1}=\delta_0\); the notebook suppresses normalization constants. It follows directly by Green's identity.}
\GSfield{Fate}{The elementary electrostatic calculation supplies the exact normalization behind the source formula.}
\GSfield{Present status}{As of November 2026, the reconstructed statement is standard.}
\GSfield{Assessment}{literal statement: correct with missing hypotheses; publication match: none found; verification: read the result; development: annotated; confidence: high}
\end{GSitem}
\begin{GSitem}{1975-386.7}{E}{Perturbation of the non-Abelian Gauss equation}
\GSfield{Kind; source}{method; \GSpage{P000646}}
\GSfield{The notes}{An inverse-Laplacian expansion and diagrams with zero, one and two insertions of \(a_i\) are written; the first correction contains no \(g\). The accompanying historical remark names Khriplovich and Gross.}
\GSfield{Statement}{For \(P=0\), formally write \(L_a\Psi=\sum_i[a_i,\partial_i\Psi]\). Then
\[(\Delta+gL_a)^{-1}\rho=\Delta^{-1}\rho-g\Delta^{-1}L_a\Delta^{-1}\rho+O(g^2).\]
For \(P\ne0\) one must also expand the source \(\rho-g[a_i,P_i]\). Analytic convergence requires an operator topology in which \(\|g\Delta^{-1}L_a\|<1\); no such estimate is supplied.}
\GSfield{Reconstruction}{\ours\ This is the formal Neumann expansion of the preceding reconstructed equation. The list of physicists does not by itself specify a beta function, renormalization scheme, or an asymptotic-freedom theorem.}
\GSfield{Fate}{Schwinger's pre-existing Coulomb-gauge program is identifiable \citep[abstract]{schwinger1962relativistic}; the particular omitted perturbative calculation attributed to Khriplovich and Gross is not identifiable from this page.}
\GSfield{Present status}{As of November 2026, the reconstructed statement is standard.}
\GSfield{Assessment}{literal statement: false as written; publication match: thematic only; verification: publisher metadata checked; development: annotated; confidence: medium}
\end{GSitem}
\begin{GSitem}{1975-386.8}{E}{Canonical quantization in Coulomb gauge}
\GSfield{Kind; source}{construction; \GSpage{P000647}}
\GSfield{The notes}{The commutator is abbreviated as \([a(x),p(y)]=\delta(x-y)\).}
\GSfield{Statement}{For transverse components, the formal equal-time relation is
\[[a_i^a(x),P_j^b(y)]=i\hbar\delta^{ab}\left(\delta_{ij}-\frac{\partial_i\partial_j}{\Delta}\right)\delta(x-y).\]
A quantum Hamiltonian further requires regularization, ordering and the gauge-fixed inner product. The formula alone is not a construction of continuum Yang--Mills theory.}
\GSfield{Fate}{The Schrödinger representation and the Faddeev--Popov density are described in \citet[Section 2]{reinhardt2018coulomb}. The transverse projector is forced by \(\partial_i a_i=\partial_iP_i=0\).}
\GSfield{Present status}{As of November 2026, the reconstructed statement is standard.}
\GSfield{Assessment}{literal statement: false as written; publication match: partial; verification: read the result; development: annotated; confidence: high}
\end{GSitem}
\begin{GSitem}{1975-386.9}{C}{Vacuum averaging of the charge interaction}
\GSfield{Kind; source}{method; \GSpage{P000646-P000647}}
\GSfield{The notes}{The kernel \(F(x,y,a)\) replaces \(1/|x-y|\), and \(\delta E=\int\rho(x)\langle F(x,y,a)\rangle\rho(y)\,dx\,dy\) is proposed.}
\GSfield{Statement}{For specified external color sources, a regulated Coulomb-gauge vacuum expectation of \(M^{-1}(-\Delta)M^{-1}\) defines an instantaneous color-Coulomb contribution. Whether the full static-source ground-state energy grows linearly at large separation is the confinement question, not a consequence of taking this expectation.}
\GSfield{Fate}{The approach has a substantial later literature; \citet[Sections 2--3]{reinhardt2018coulomb}. The relation to confinement requires separating instantaneous energy from the fully relaxed static potential; \citet[Section 5]{jaffewitten2000yangmills} distinguishes confinement from the Yang--Mills existence and mass-gap problem.}
\GSfield{Present status}{As of November 2026, this fragment supplies a formal route to the confinement problem GP-1975-10, not its solution. The notes specify neither a regulated vacuum nor a limiting construction.}
\GSfield{Assessment}{literal statement: correct with missing hypotheses; publication match: partial; verification: read the result; development: annotated; confidence: high; problem \hyperref[prob:GP-1975-10]{GP-1975-10}}
\GSfield{Worth developing}{A precise comparison between the instantaneous Coulomb kernel and gauge-invariant static energy remains useful for understanding what this calculation can establish.}
\end{GSitem}

\GStalk{387}{1976-03-22}{I. M. Gelfand}{P000647-P000649}{A probability distribution arising in biology}
A short intervention introduces what is now called the generalized Poisson distribution and sets its normalization as an exercise. The date and text are on P000648. The formula already appeared in print before this meeting; the notebook does not identify the biological model or the unnamed second example.

\begin{GSitem}{1975-387.1}{E}{The generalized Poisson distribution}
\GSfield{Kind; source}{definition; \GSpage{P000648}}
\GSfield{The notes}{For \(0\leq\alpha<1\), \(P_n=e^{-x_0-\alpha n}x_0(x_0+\alpha n)^{n-1}/n!\); at \(\alpha=0\) the distribution is Poisson.}
\GSfield{Statement}{Let \(x_0>0\), \(0\leq\alpha<1\), and \(n\in\mathbb Z_{\ge0}\). Set \(P_0=e^{-x_0}\) and use the displayed formula for \(n\ge1\). These numbers are nonnegative. The parameter range avoids the finite-support difficulties of negative deformation parameters.}
\GSfield{Fate}{This matches \citet[Section 3, equation (3.1)]{consuljain1973poisson}, received in May 1971, revised in October 1972 and published in November 1973. The biological provenance in this intervention remains unspecified.}
\GSfield{Present status}{As of November 2026, the reconstructed statement is standard.}
\GSfield{Assessment}{literal statement: correct with missing hypotheses; publication match: exact; verification: read the result; development: annotated; confidence: high}
\end{GSitem}
\begin{GSitem}{1975-387.2}{E}{Normalization of the generalized Poisson law}
\GSfield{Kind; source}{question; \GSpage{P000648}}
\GSfield{The notes}{The explicit exercise is to verify \(\sum_{n\ge0}P_n=1\).}
\GSfield{Statement}{Let \(T(z)\) be the nonnegative solution near zero of \(T=z\exp(\alpha(T-1))\). Lagrange inversion gives
\[\sum_{n\ge0}P_nz^n=\exp(x_0(T(z)-1)).\]
For \(0\leq\alpha<1\), continuity from below and the minimal fixed point at \(z=1\) give \(T(1)=1\), hence total mass one.}
\GSfield{Reconstruction}{\ours\ The generating-function proof is supplied here. For \(\alpha=0\), \(T=z\). For \(0<\alpha<1\), the function \(h(t)=e^{\alpha(t-1)}-t\) has \(h(1)=0\) and \(h'(t)=\alpha e^{\alpha(t-1)}-1<0\) on \([0,1]\), so it has no fixed-point zero below one. Monotone iteration from zero for \(0\le z\le1\) identifies the nonnegative branch and its limit.}
\GSfield{Fate}{Normalization belongs to the already published distribution \citep[Sections 2--3]{consuljain1973poisson}. Thus this is an expository exercise, not a problem first solved after the seminar.}
\GSfield{Present status}{As of November 2026, the reconstructed statement is standard.}
\GSfield{Assessment}{literal statement: correct with missing hypotheses; publication match: exact; verification: read the result; development: annotated; confidence: high; problem \hyperref[prob:GP-1975-22]{GP-1975-22}}
\end{GSitem}

\GStalk{388}{1976-03-22}{A. V. Roiter}{P000647-P000650}{Indecomposable modules and root systems}
Roiter discusses the first Brauer--Thrall theorem, chains of indecomposable submodules and a comparison of lengths. The discussion then turns to integral quadratic forms, finite and affine Dynkin diagrams, reflections and arithmetic root criteria. Kirillov's drawing of Dynkin diagrams is an intervention within this record. A final star-shaped form raises the problem of enumerating roots by reflections.

\begin{GSitem}{1975-388.1}{E}{The first Brauer–Thrall theorem}
\GSfield{Kind; source}{theorem; \GSpage{P000648}}
\GSfield{The notes}{A finite-dimensional algebra has either finitely many indecomposable modules or indecomposables of arbitrarily large dimension. Roiter is named, followed by an Auslander generalization to Artinian rings.}
\GSfield{Statement}{For a finite-dimensional algebra \(\Lambda\) over a field, bounded dimensions of finite-dimensional indecomposable modules imply that there are only finitely many isomorphism classes of such modules. The corresponding formulation for Artin algebras uses composition length. The phrase ``Artinian rings'' needs the precise finiteness setting of the theorem being invoked.}
\GSfield{Fate}{Roiter's article was submitted on 20 March 1968 and published that year \citep{roiter1968brauerthrall}. A later proof within Artin-algebra representation theory is given by \citet[Section 6]{ringel2006foundation}; no claim about arbitrary left Artinian rings is inferred from that proof.}
\GSfield{Present status}{As of November 2026, the reconstructed statement is standard.}
\GSfield{Assessment}{literal statement: correct with missing hypotheses; publication match: exact; verification: read the result; development: annotated; confidence: high}
\end{GSitem}
\begin{GSitem}{1975-388.2}{B}{Infinite chains of indecomposable modules}
\GSfield{Kind; source}{claim; \GSpage{P000648}}
\GSfield{The notes}{The notebook asks for or asserts an infinite chain \(A_1\subset A_2\subset\cdots\) with indecomposable terms in the infinite case.}
\GSfield{Statement}{If an Artin algebra is of infinite representation type, there is an infinite sequence of finite-length indecomposables and proper monomorphisms \(M_1\hookrightarrow M_2\hookrightarrow\cdots\). The chain can be chosen to consist of Gabriel--Roiter inclusions; its lengths are unbounded.}
\GSfield{Fate}{A precise strengthened form is \citet[Section 7.3, Theorem]{ringel2006foundation}, published in 2006 after lectures in 2003--2005. The source does not specify the inclusions, and its exact contemporary publication was not identified.}
\GSfield{Present status}{As of November 2026, the reconstructed statement is standard.}
\GSfield{Assessment}{literal statement: correct with missing hypotheses; publication match: partial; verification: read the result; development: annotated; confidence: high}
\end{GSitem}
\begin{GSitem}{1975-388.3}{B}{Bounding composition length by indecomposable-chain length}
\GSfield{Kind; source}{theorem; \GSpage{P000648}}
\GSfield{The notes}{Besides composition length \(\ell(M)\), the notes define \(\bar\ell(M)\) as the largest number of terms in a chain of indecomposable submodules. Their quantifier order is \(\forall N\,\exists n:\bar\ell(M)\le n\Rightarrow\ell(M)\le N\).}
\GSfield{Statement}{The substantive uniform bound is: for each Artin algebra \(\Lambda\) and each \(n\), there exists \(N(\Lambda,n)\) such that \(\bar\ell(M)\le n\) implies \(\ell(M)\le N(\Lambda,n)\), for indecomposable finite-length \(M\). If \(p,q\) are the maximum lengths of indecomposable projectives and injectives, one may take \(N=(pq)^{n-1}\): a Gabriel--Roiter filtration starts with a simple module, and each inclusion multiplies length by at most \(pq\).}
\GSfield{Reconstruction}{\ours\ The implication printed in the notes is not this bound; it is too weak to express the intended control. The exponent is a direct deduction from the cited filtration and inclusion bound, not a formula attributed to the notebook.}
\GSfield{Fate}{Use \citet[Section 1.5 and Section 3.1, Corollary]{ringel2006foundation}. These establish the ingredients in the Artin-algebra setting; this is a later rigorous match to the intended estimate, not an exact identification of Roiter's wording.}
\GSfield{Present status}{As of November 2026, the reconstructed statement is standard.}
\GSfield{Assessment}{literal statement: not determinable; publication match: partial; verification: read the result; development: annotated; confidence: high}
\end{GSitem}
\begin{GSitem}{1975-388.4}{E}{Integral forms and finite or affine Dynkin diagrams}
\GSfield{Kind; source}{definition; \GSpage{P000648-P000649}; A. V. Roiter; A. A. Kirillov (intervention)}
\GSfield{The notes}{The form \(q(x)=\sum a_ix_i^2+\sum_{i<j}b_{ij}x_ix_j\) has \(a_i>0\), \(b_{ij}\le0\), and \(a_i\mid b_{ij}\). Positivity is associated with Dynkin diagrams; nonnegativity with extended diagrams. Kirillov draws the finite families.}
\GSfield{Statement}{Assume integer coefficients, divisibility at both endpoints, and connected coefficient graph. The matrix \(C_{ii}=2\), \(C_{ij}=b_{ij}/a_i\) is a symmetrizable generalized Cartan matrix. Positive definiteness gives finite type; positive semidefiniteness with a nontrivial radical gives affine type. Nonnegativity alone includes finite type, and disconnected forms require a componentwise statement.}
\GSfield{Reconstruction}{Connectedness and singularity distinguish the affine assertion. The cited source directly treats the simply laced representation-theoretic part; the Cartan-matrix translation of the displayed form is an explicit reconstruction.}
\GSfield{Fate}{The simply laced correspondence and its use in quiver classification predate the meeting: \citet[Section 3]{bernsteingelfandponomarev1973coxeter}. This connects thematically to GP-1972-301 and GP-1972-302, but no repeated quiver-classification question is written here. The nonsimply-laced integral-form conditions and the connected nonnegative case are addressed in \citet[Section 1 and Lemma 6]{roiter1978roots}.}
\GSfield{Present status}{As of November 2026, the reconstructed statement is standard.}
\GSfield{Assessment}{literal statement: correct with missing hypotheses; publication match: partial; verification: read the result; development: annotated; confidence: high}
\end{GSitem}
\begin{GSitem}{1975-388.5}{E}{Reflections and roots of simply laced forms}
\GSfield{Kind; source}{construction; \GSpage{P000649}}
\GSfield{The notes}{A reflection changes one coordinate and preserves \(q\); roots are obtained from coordinate vectors. In the equal-coefficient case, \(q(x)=1\) is the root criterion.}
\GSfield{Statement}{For \(a_i=1\), define \(s_i(x)=x-(2x_i+\sum_{j\ne i}b_{ij}x_j)e_i\). For a finite simply laced Dynkin form, the integral solutions of \(q(x)=1\) are precisely the real roots \(\{w(e_i):w\in\langle s_1,\ldots,s_n\rangle\}\). In affine type the same norm condition characterizes real roots; nonzero imaginary roots have \(q=0\) and must be listed separately.}
\GSfield{Fate}{The finite-type reflection argument is already in \citet[Section 3]{bernsteingelfandponomarev1973coxeter}. The word ``roots'' in the notes is orbit-defined, and should not silently be enlarged to all imaginary roots. The affine norm criterion also follows from \citet[Proposition 4]{roiter1978roots}.}
\GSfield{Present status}{As of November 2026, the reconstructed statement is standard.}
\GSfield{Assessment}{literal statement: correct with missing hypotheses; publication match: partial; verification: read the result; development: annotated; confidence: high}
\end{GSitem}
\begin{GSitem}{1975-388.6}{A}{A divisibility criterion for roots}
\GSfield{Kind; source}{claim; \GSpage{P000649}}
\GSfield{The notes}{For the more general form, the notebook states that \(x\) is a root if and only if \(q(x)\) divides every \(a_ix_i\).}
\GSfield{Statement}{Let \(q=\sum_i a_ix_i^2+\sum_{i<j}b_{ij}x_ix_j\) have positive integral \(a_i\), nonpositive \(b_{ij}\), and \(a_i,a_j\mid b_{ij}\). Assume its connected diagram is finite or affine Dynkin. A vector \(x\in\mathbb Z^n\) is a real root, that is an orbit of a coordinate vector under the integral reflections, if and only if \(q(x)>0\) and \(q(x)\mid a_ix_i\) for every \(i\). The positivity condition excludes the radical in affine type. This criterion does not hold without restrictions for arbitrary indefinite forms.}
\GSfield{Reconstruction}{Finite or connected affine type is supplied from the surrounding discussion, and \(q(x)>0\) is made explicit. A later star-shaped example with five or more arms is indefinite and must not inherit this criterion automatically.}
\GSfield{Fate}{An exact source is \citet[Proposition 4, p.205; Lemma 6, p.206]{roiter1978roots}. Roiter proves the criterion when every proper coordinate restriction is positive; the finite and connected affine cases satisfy this hypothesis. Russian publication was in 1978, English translation in 1980. The notebook of 22 March 1976 is an earlier verified witness to this claim.}
\GSfield{Present status}{As of November 2026, the stated finite/connected affine criterion is established. The initial source-identification gap was resolved by the 1978 article; indefinite extensions require its further hypotheses.}
\GSfield{Assessment}{literal statement: correct with missing hypotheses; publication match: exact; verification: read the result; development: annotated; confidence: high}
\end{GSitem}
\begin{GSitem}{1975-388.7}{E}{Roots of the star-shaped quadratic form}
\GSfield{Kind; source}{question; \GSpage{P000649}}
\GSfield{The notes}{The final example is \(q=x_0^2+\cdots+x_n^2-x_0(x_1+\cdots+x_n)\); the question is to obtain all roots by elementary reflections, illustrated by a word-changing analogy.}
\GSfield{Statement}{The orbit of the coordinate vectors under the reflections attached to this star is a precise enumerable set of real roots. Completing squares gives \(q=\sum_{i=1}^n(x_i-x_0/2)^2+(1-n/4)x_0^2\): the cases \(n\le3\) are finite, \(n=4\) is affine \(\widetilde D_4\), and \(n\ge5\) is indefinite. An exhaustive finite list is therefore not an appropriate interpretation for all \(n\).}
\GSfield{Reconstruction}{\ours\ The square completion and the distinction between finite enumeration and an orbit algorithm are supplied here. The notes do not state a bound on word length or demand an algorithm deciding all imaginary roots.}
\GSfield{Fate}{The reflection machinery is already available in \citet[Sections 2--3]{bernsteingelfandponomarev1973coxeter}. This question is understood as an exercise in that machinery; no separate unresolved classification conjecture is inferred.}
\GSfield{Present status}{As of November 2026, the reconstructed statement is standard.}
\GSfield{Assessment}{literal statement: correct with missing hypotheses; publication match: partial; verification: read the result; development: annotated; confidence: high}
\end{GSitem}

\GStalk{389}{1976-03-22}{A. A. Kirillov}{P000648-P000650}{Berezin’s analysis with anticommuting variables}
Kirillov presents superfunctions, integration and transformations, then graded Lie algebras and a Lie correspondence. Gelfand adds a nilpotent Taylor formula. A final question asks for deformations in the direction of a Poisson bracket. The unnamed 64-dimensional algebra and its particle interpretation cannot be identified from the supplied text.

\begin{GSitem}{1975-389.1}{E}{Superfunctions and split supermanifolds}
\GSfield{Kind; source}{definition; \GSpage{P000649}}
\GSfield{The notes}{Functions depend smoothly on commuting variables \(x\) and polynomially on anticommuting variables \(\xi\). A vector-bundle description is attributed to Rudakov.}
\GSfield{Statement}{A smooth real supermanifold of dimension \(p|q\) is a locally ringed space locally isomorphic to \((U,C^\infty_U\otimes\Lambda^\bullet\mathbb R^q)\). A split example is \((M,\Gamma(\Lambda E^*))\) for a rank-\(q\) vector bundle \(E\to M\). The local model does not itself select a canonical global splitting.}
\GSfield{Fate}{The sheaf definition was published by \citet[pp. 505--506]{berezinleites1975supermanifolds} in 1975, following submission on 24 March 1975. The attribution to Rudakov is preserved as the seminar's attribution, without inferring a separate publication.}
\GSfield{Present status}{As of November 2026, the reconstructed statement is standard.}
\GSfield{Assessment}{literal statement: correct with missing hypotheses; publication match: partial; verification: read the result; development: annotated; confidence: high}
\end{GSitem}
\begin{GSitem}{1975-389.2}{E}{Berezin integration}
\GSfield{Kind; source}{definition; \GSpage{P000649}}
\GSfield{The notes}{Integration is specified by \(\int d\xi=0\) and \(\int\xi,d\xi=1\), extending to products of anticommuting variables.}
\GSfield{Statement}{For an ordered set of odd generators, Berezin integration is coefficient extraction: \(\int d\xi_q\cdots d\xi_1\,f\) is the coefficient of \(\xi_1\cdots\xi_q\). Changing the order changes signs. On a supermanifold, coordinate-invariant integration requires the Berezinian transformation law, rather than an ordinary volume form.}
\GSfield{Fate}{The integration framework predates this meeting; see \citet[Section 5]{berezinkac1970lie} and \citet[pp. 507--508]{berezinleites1975supermanifolds}.}
\GSfield{Present status}{As of November 2026, the reconstructed statement is standard.}
\GSfield{Assessment}{literal statement: correct with missing hypotheses; publication match: partial; verification: read the result; development: annotated; confidence: high}
\end{GSitem}
\begin{GSitem}{1975-389.3}{E}{Coordinate changes with nilpotent terms}
\GSfield{Kind; source}{construction; \GSpage{P000649}}
\GSfield{The notes}{The examples include changes of an even coordinate by a product of odd coordinates, such as \(x\mapsto x+\xi_1\xi_2\), as well as transformations of the odd variables.}
\GSfield{Statement}{A morphism is determined locally by parity-preserving substitutions: even coordinates are sent to even superfunctions and odd coordinates to odd ones. A local isomorphism requires invertibility of the ordinary even Jacobian and of the linear odd part after nilpotents are discarded.}
\GSfield{Fate}{These are the local morphisms in the sheaf-theoretic construction of \citet[pp. 505--506]{berezinleites1975supermanifolds}. Nilpotent corrections explain why a bundle picture does not exhaust the coordinate transformations.}
\GSfield{Present status}{As of November 2026, the reconstructed statement is standard.}
\GSfield{Assessment}{literal statement: correct with missing hypotheses; publication match: partial; verification: read the result; development: annotated; confidence: high}
\end{GSitem}
\begin{GSitem}{1975-389.4}{F}{An unnamed 64-dimensional algebra}
\GSfield{Kind; source}{remark; \GSpage{P000649}}
\GSfield{The notes}{The notes mention a 64-dimensional algebra that is said to describe particles better than \(SU(3)\), without naming its generators or representation.}
\GSfield{Statement}{Dimension alone does not determine a Lie superalgebra, its parity dimensions, real form, representation or physical model. No definite mathematical identification follows.}
\GSfield{Fate}{Targeted searches on Berezin, Kirillov, 64-dimensional superalgebras and particle models did not identify a source matching all of these details.}
\GSfield{Present status}{As of November 2026, this remains an unresolved identification in the historical record.}
\GSfield{Assessment}{literal statement: not determinable; publication match: none found; verification: unresolved lead; development: annotated; confidence: low}
\end{GSitem}
\begin{GSitem}{1975-389.5}{E}{Lie superalgebras and associative envelopes}
\GSfield{Kind; source}{definition; \GSpage{P000649}}
\GSfield{The notes}{The bracket is graded skew-symmetric, satisfies graded Jacobi, and is obtained from multiplication by a signed commutator; a converse is mentioned.}
\GSfield{Statement}{Over a field of characteristic zero, a Lie superalgebra \(\mathfrak g=\mathfrak g_0\oplus\mathfrak g_1\) satisfies \([x,y]=-(-1)^{|x||y|}[y,x]\) and the graded Jacobi identity. Every associative superalgebra has this bracket. Conversely \(\mathfrak g\) embeds in its universal enveloping superalgebra, whose associated graded algebra is \(S(\mathfrak g_0)\otimes\Lambda(\mathfrak g_1)\).}
\GSfield{Reconstruction}{The precise converse is an enveloping-algebra statement, not an assertion that every Lie superalgebra already has an associative multiplication on its own vector space.}
\GSfield{Fate}{The definitions and representation framework are present in \citet[Sections 3--4]{berezinkac1970lie}, published in 1970. The enveloping formulation spells out the otherwise ambiguous converse.}
\GSfield{Present status}{As of November 2026, the reconstructed statement is standard.}
\GSfield{Assessment}{literal statement: correct with missing hypotheses; publication match: partial; verification: read the result; development: annotated; confidence: high}
\end{GSitem}
\begin{GSitem}{1975-389.6}{E}{The formal Lie correspondence in supergeometry}
\GSfield{Kind; source}{theorem; \GSpage{P000649}}
\GSfield{The notes}{A ``Lie theorem'' is attributed to Berezin and Kac ``from Kyiv''. No triangularization statement is written.}
\GSfield{Statement}{Over characteristic zero, finite-dimensional Lie superalgebras correspond to formal supergroup laws, with homomorphisms respected. This is a local/formal Lie correspondence; it does not assert that every representation of a solvable Lie superalgebra can be triangularized.}
\GSfield{Reconstruction}{The identity of Kac is G. I. Kac, coauthor of the 1970 article, rather than an automatic identification with V. G. Kac. The theorem is interpreted using that precise paper.}
\GSfield{Fate}{\citet[Section 3.5, Theorem 4]{berezinkac1970lie} proves the mutually inverse functors. This is an exact and pre-existing match to the attributed Lie correspondence.}
\GSfield{Present status}{As of November 2026, the reconstructed statement is standard.}
\GSfield{Assessment}{literal statement: correct with missing hypotheses; publication match: exact; verification: read the result; development: annotated; confidence: high}
\end{GSitem}
\begin{GSitem}{1975-389.7}{E}{Gelfand’s nilpotent Taylor formula}
\GSfield{Kind; source}{remark; \GSpage{P000649}; I. M. Gelfand}
\GSfield{The notes}{Gelfand writes \(p(t+\xi)=p(t)+p'(t)\xi\).}
\GSfield{Statement}{For a polynomial \(p\) over a characteristic-zero commutative coefficient ring and a commuting nilpotent element \(\xi\) with \(\xi^2=0\), the binomial formula gives \(p(t+\xi)=p(t)+p'(t)\xi\). The same finite Taylor substitution is meaningful for smooth superfunctions with nilpotent increments. If \(t\) is even and \(\xi\) odd, this substitution by itself is not a parity-preserving even-coordinate change.}
\GSfield{Fate}{The identity is elementary; the parity qualification separates the algebraic Taylor example from the preceding definition of supermanifold morphism.}
\GSfield{Present status}{As of November 2026, the reconstructed statement is standard.}
\GSfield{Assessment}{literal statement: correct with missing hypotheses; publication match: none found; verification: read the result; development: annotated; confidence: high}
\end{GSitem}
\begin{GSitem}{1975-389.8}{B}{Deformation in the direction of a super-Poisson bracket}
\GSfield{Kind; source}{question; \GSpage{P000650}}
\GSfield{The notes}{The notes end by asking to find deformations of the function superalgebra in the direction of its Poisson bracket.}
\GSfield{Statement}{A precise even-bracket interpretation asks for an associative parity-preserving product \(f*g=fg+\sum_{r\ge1}\hbar^rB_r(f,g)\) such that the first graded commutator is \(\hbar\{f,g\}+O(\hbar^2)\). For a split smooth supermanifold built from a symplectic base and a vector bundle with a nondegenerate symmetric fibre metric, the associated Rothstein even super-Poisson bracket has such a differential star product.}
\GSfield{Reconstruction}{Parity of the bracket, formal versus convergent deformation and nondegeneracy are absent from the notebook. The stated theorem is a definite later answer to an explicitly restricted interpretation, not a claim that all possible readings were equivalent.}
\GSfield{Fate}{\citet[Theorems 1.3--1.4 and Corollary 1.1]{bordemann2000superpoisson}, first circulated on 30 March 2000, constructs the star product and identifies its first bracket. Odd brackets require a different formulation and are not covered by this identification.}
\GSfield{Present status}{As of November 2026, the stated even symplectic interpretation is solved; the notes do not specify a unique more general problem.}
\GSfield{Assessment}{literal statement: correct with missing hypotheses; publication match: partial; verification: read the result; development: annotated; confidence: high; problem \hyperref[prob:GP-1975-23]{GP-1975-23}}
\end{GSitem}

\GStalk{390}{1976-03-22}{Ya. G. Sinai}{P000648-P000651}{Transfer matrices, the Ising model and square ice}
Sinai explains how a two-dimensional lattice partition function becomes a trace of a transfer matrix. He connects the limiting free energy to the largest eigenvalue, recalls Onsager's solution and its anticommuting-operator structure, and mentions Lieb's square-ice entropy calculation. The surviving mathematical text is on P000650.

\begin{GSitem}{1975-390.1}{E}{The row transfer matrix for the Ising model}
\GSfield{Kind; source}{construction; \GSpage{P000650}}
\GSfield{The notes}{Rows are vectors \(x,y\in\{\pm1\}^n\), and \(A(x,y)=\exp\{\beta[J_1\sum x_ix_{i-1}+J_2\sum x_iy_i]\}\), with periodic indices.}
\GSfield{Statement}{For an \(n\times R\) torus with nearest-neighbour horizontal and vertical couplings and no external field, the partition function is \(Z_{n,R}=\operatorname{tr}(A_n^R)\). Assigning all horizontal interactions to the first row produces the nonsymmetric matrix in the notes; splitting them equally gives a similar symmetric matrix. Each bond must be counted once.}
\GSfield{Fate}{The finite-width transfer-operator problem is solved in \citet[Sections II--V]{onsager1944crystal}, received 4 October 1943 and published in 1944. The identity \(Z=\operatorname{tr}A^R\) follows immediately by summing cyclic row sequences.}
\GSfield{Present status}{As of November 2026, the reconstructed statement is standard.}
\GSfield{Assessment}{literal statement: correct with missing hypotheses; publication match: exact; verification: read the result; development: annotated; confidence: high}
\end{GSitem}
\begin{GSitem}{1975-390.2}{E}{Free energy and the thermodynamic limit}
\GSfield{Kind; source}{method; \GSpage{P000650}}
\GSfield{The notes}{The probability is proportional to \(e^{-\beta H}\), and the notes identify free energy through \(\lim R^{-2}\log Z\), with singularities associated with the dominant eigenvalue.}
\GSfield{Statement}{For a finite system the Gibbs probability is \(e^{-\beta H}/Z\). At fixed width \(n\), positivity gives \(\lim_{R\to\infty}R^{-1}\log Z_{n,R}=\log\lambda_{\max}(A_n)\). A two-dimensional pressure requires a further width limit, or a joint limit with both side lengths tending to infinity. Finite positive transfer matrices have a simple dominant eigenvalue analytic in finite real couplings; a phase singularity arises only in an infinite-size limit.}
\GSfield{Reconstruction}{The source uses the word free energy for the dimensionless pressure; physical free energy per site is its negative divided by \(\beta\). The distinction between fixed width and a two-dimensional limit is essential.}
\GSfield{Fate}{\citet[abstract and Section IX]{onsager1944crystal} explicitly separates finite strips from the infinite crystal. The normalization here is fixed by the elementary Gibbs definition.}
\GSfield{Present status}{As of November 2026, the reconstructed statement is standard.}
\GSfield{Assessment}{literal statement: correct with missing hypotheses; publication match: partial; verification: read the result; development: annotated; confidence: high}
\end{GSitem}
\begin{GSitem}{1975-390.3}{E}{Onsager’s zero-field Ising solution}
\GSfield{Kind; source}{theorem; \GSpage{P000650}}
\GSfield{The notes}{The eigenvalues of the transfer matrix are attributed to Onsager, with a singularity at a distinguished temperature.}
\GSfield{Statement}{For the infinite anisotropic square-lattice Ising ferromagnet with \(J_1,J_2>0\) and zero magnetic field, the critical inverse temperature is characterized by \(\sinh(2\beta_cJ_1)\sinh(2\beta_cJ_2)=1\). The exact free energy is nonanalytic there; the internal energy is continuous and the specific heat has a logarithmic divergence.}
\GSfield{Fate}{These conclusions are in the original \citet[abstract and Section IX]{onsager1944crystal}. No statement about a nonzero magnetic field is warranted by the notes.}
\GSfield{Present status}{As of November 2026, the reconstructed statement is standard.}
\GSfield{Assessment}{literal statement: correct with missing hypotheses; publication match: exact; verification: read the result; development: annotated; confidence: high}
\end{GSitem}
\begin{GSitem}{1975-390.4}{E}{Anticommuting operators in the transfer-matrix solution}
\GSfield{Kind; source}{method; \GSpage{P000650}}
\GSfield{The notes}{Operators \(A_i\) anticommute and the transfer matrix is written schematically as \(\exp(\sum c_{ij}A_iA_j)\).}
\GSfield{Statement}{A Clifford algebra organizes the zero-field transfer calculation. For periodic spin boundary conditions, fermion-parity sectors must be treated separately, and a scalar normalization may multiply the exponential of a quadratic Clifford expression. One should not read the schematic equality as a global sector-free normal form for every transfer matrix.}
\GSfield{Fate}{The underlying algebra appears in \citet[Sections IV--VII]{onsager1944crystal}. The notebook compresses normalization and boundary-sector information, so the publication match is structural rather than a literal equality.}
\GSfield{Present status}{As of November 2026, the reconstructed statement is standard.}
\GSfield{Assessment}{literal statement: correct with missing hypotheses; publication match: partial; verification: read the result; development: annotated; confidence: high}
\end{GSitem}
\begin{GSitem}{1975-390.5}{E}{Lieb’s entropy of square ice}
\GSfield{Kind; source}{claim; \GSpage{P000650}}
\GSfield{The notes}{Lieb is credited with finding the entropy for a model of ice.}
\GSfield{Statement}{For the equal-weight square-ice model, each vertex has two incoming and two outgoing arrows. In the thermodynamic limit its number of states per vertex is \(W=(4/3)^{3/2}\), giving entropy per vertex \(k_B\log W\). This identifies the two-dimensional six-vertex model; it is not an exact entropy formula for three-dimensional physical ice.}
\GSfield{Fate}{The detailed result is \citet[abstract and final entropy formula]{lieb1967squareice}, published on 5 October 1967 after a shorter April announcement. The original full text is a source target; publisher metadata and the result stated in the abstract are verified.}
\GSfield{Present status}{As of November 2026, the reconstructed statement is standard.}
\GSfield{Assessment}{literal statement: correct with missing hypotheses; publication match: exact; verification: publisher metadata checked; development: annotated; confidence: high}
\end{GSitem}

\GStalk{391}{1976-03-29}{I. M. Krichever}{P000649-P000654}{Higher-rank commuting operators; Shnirelman on spectral phases}
The catalogue names Krichever, but the manuscript changes speakers on P000651: Krichever and Drinfeld discuss the inverse construction of commuting ordinary differential operators, then Shnirelman discusses spectral phases, closed geodesics and the wave trace through P000653. Arnold's question about nodal sets occurs within the latter discussion. These changes are retained as speaker fields, not promoted to invented catalogue records. The first part continues the preceding algebro-geometric line; the second revisits 1973-333.18--20.

\begin{GSitem}{1975-391.1}{E}{The spectral curve and joint-eigenfunction sheaf}
\GSfield{Kind; source}{construction; \GSpage{P000651}}
\GSfield{The notes}{For commuting scalar ordinary differential operators \(L_1,L_2\), the joint eigenvalues satisfy \(P(E_1,E_2)=0\). Coprime orders give a line bundle; orders four and six suggest a rank-two bundle.}
\GSfield{Statement}{A commutative ring \(\mathcal A\) of ordinary scalar differential operators of positive order has an algebraic spectral curve, with a distinguished point at infinity. In the generic smooth case, common eigenspaces define a vector bundle of rank \(r\), the greatest common divisor of the orders of all nonzero operators in \(\mathcal A\). Coprime orders force \(r=1\); the gcd of a selected pair need not equal the rank of a larger ring.}
\GSfield{Reconstruction}{The hypotheses are local smooth/analytic coefficients with nonvanishing leading coefficients, followed by a suitable normalization; singular spectral curves require torsion-free sheaves. The source's two selected orders do not by themselves fix the rank of every containing ring.}
\GSfield{Fate}{\citet[Theorem 1.1 and Section 1]{krichever1978commutative} recalls the pre-existing spectral-curve construction and clarifies the spectral data. The 1978 paper distinguishes this direct construction from the inverse problem.}
\GSfield{Present status}{As of November 2026, the reconstructed statement is standard.}
\GSfield{Assessment}{literal statement: correct with missing hypotheses; publication match: exact; verification: read the result; development: annotated; confidence: high}
\end{GSitem}
\begin{GSitem}{1975-391.2}{E}{Counting eigenvectors over a fixed spectral value}
\GSfield{Kind; source}{claim; \GSpage{P000651}}
\GSfield{The notes}{After saying that the coprime case is one-dimensional, the notes say that the number of common eigenfunctions is \(n\), with distinct eigenvalues.}
\GSfield{Statement}{If \(\operatorname{ord}L_1=n\), then \(\ker(L_1-E_1)\) is locally \(n\)-dimensional. The commuting \(L_2\) acts on this space. For generic \(E_1\) and rank one, it has \(n\) distinct eigenvalues, each with a one-dimensional joint eigenspace. Thus there are \(n\) points above a generic \(E_1\), not an \(n\)-dimensional eigenspace at one point \((E_1,E_2)\).}
\GSfield{Reconstruction}{This separates the degree of the spectral projection from the rank of the eigenfunction bundle, resolving the apparent contradiction on the page.}
\GSfield{Fate}{The solution-space construction used in the proof of \citet[Theorem 1.1]{krichever1978commutative} supplies precisely this distinction.}
\GSfield{Present status}{As of November 2026, the reconstructed statement is standard.}
\GSfield{Assessment}{literal statement: correct with missing hypotheses; publication match: partial; verification: read the result; development: annotated; confidence: high}
\end{GSitem}
\begin{GSitem}{1975-391.3}{A}{The higher-rank inverse construction}
\GSfield{Kind; source}{question; \GSpage{P000651}}
\GSfield{The notes}{The inverse problem is explicitly marked ``not done'', with orders four and six as the rank-two example: construct commuting differential operators from the bundle.}
\GSfield{Statement}{Given an appropriate smooth pointed spectral curve, rank-\(r\) bundle data, a local parameter and a compatible framing/asymptotic prescription at the marked point, reconstruct a commutative ring of scalar ordinary differential operators. A bare bundle without this extra data is not sufficient. In generic position, rank \(r\) also involves \(r-1\) functions of the independent variable after normalization.}
\GSfield{Fate}{\citet[Theorem 2.3]{krichever1978commutative}, submitted 15 February 1977 and published in 1978, completes the generic smooth-curve construction. Its introduction expressly corrects errors in the 1977 implementation. This is a subsequent answer under stated generic hypotheses, not an unrestricted classification of all singular cases.}
\GSfield{Present status}{As of November 2026, the generic inverse correspondence described here is established; the notebook’s unrestricted shorthand must retain the spectral-data qualifications.}
\GSfield{Assessment}{literal statement: correct with missing hypotheses; publication match: exact; verification: read the result; development: annotated; confidence: high; problem \hyperref[prob:GP-1975-26]{GP-1975-26}}
\end{GSitem}
\begin{GSitem}{1975-391.4}{A}{Drinfeld’s functional parameters and gauge normalization}
\GSfield{Kind; source}{claim; \GSpage{P000651}; V. G. Drinfeld}
\GSfield{The notes}{Drinfeld says \(d=(m,n)\) leaves \(d-1\) arbitrary functions; conjugation removes one coefficient, while a leading coefficient remains. Powers \(Q^2,Q^3\) already give order-four/order-six examples with two functions.}
\GSfield{Statement}{After local change of independent variable and scalar conjugation, a monic order-\(r\) auxiliary operator can be normalized to \(\partial^r+v_{r-2}(x)\partial^{r-2}+\cdots+v_0(x)\), leaving \(r-1\) functions. Spectral curve and bundle parameters remain in addition. For an arbitrary order-two \(Q=a_2\partial^2+a_1\partial+a_0\), scalar conjugation can remove the first-order term locally; if coordinate changes are not also factored out, \(a_2\) remains alongside the potential.}
\GSfield{Reconstruction}{The phrase ``specifying a pair iff specifying functions'' suppresses finite-dimensional spectral data and the equivalence relation. The example \(Q^2,Q^3\) explains why the choice of normalization matters.}
\GSfield{Fate}{The normalized \(r-1\)-function statement is made precise by \citet[Theorem 2.3]{krichever1978commutative}. It applies to the ring rank, not automatically to the gcd of any chosen pair in a larger ring.}
\GSfield{Present status}{As of November 2026, the reconstructed statement is standard.}
\GSfield{Assessment}{literal statement: correct with missing hypotheses; publication match: partial; verification: read the result; development: annotated; confidence: high; problem \hyperref[prob:GP-1975-26]{GP-1975-26}}
\end{GSitem}
\begin{GSitem}{1975-391.5}{E}{Laplacian convention and Weyl frequencies}
\GSfield{Kind; source}{definition; \GSpage{P000651}; A. I. Shnirelman}
\GSfield{The notes}{Shnirelman takes a compact analytic Riemannian manifold, writes \(\Delta u_k+\lambda_k^2u_k=0\), defines \(\Delta\) through spherical averages, and invokes \(\lambda_k\sim c k^{1/n}\).}
\GSfield{Statement}{Let \(M\) be a closed smooth Riemannian \(n\)-manifold and use the nonpositive Laplace--Beltrami operator. Then \(-\Delta u_k=\lambda_k^2u_k\), with \(\lambda_k\ge0\), counted with multiplicity. Weyl's law is
\[N(\Lambda)\sim\frac{\omega_n\operatorname{Vol}(M)}{(2\pi)^n}\Lambda^n.\]
For normalized averages over small geodesic spheres, \(\Delta f=2n\lim_{\tau\to0}(\operatorname{avg}_{S_\tau}f-f)/\tau^2\).}
\GSfield{Reconstruction}{Closedness supplies the missing boundary hypothesis. The factor \(2n\) is obtained by a normal-coordinate Taylor expansion. Analyticity is stronger than needed for the Weyl law.}
\GSfield{Fate}{The relevant positive elliptic spectral setting already appears in \citet[opening hypotheses]{duistermaatguillemin1975spectrum}; the sphere normalization is reconstructed directly. Publication metadata, rather than the original full proof, was read.}
\GSfield{Present status}{As of November 2026, the reconstructed statement is standard.}
\GSfield{Assessment}{literal statement: correct with missing hypotheses; publication match: partial; verification: publisher metadata checked; development: annotated; confidence: high}
\end{GSitem}
\begin{GSitem}{1975-391.6}{E}{The rectangular flat-torus spectrum}
\GSfield{Kind; source}{construction; \GSpage{P000651}; A. I. Shnirelman}
\GSfield{The notes}{The example uses a flat torus and frequencies \(\lambda^2=n^2/a^2+m^2/b^2\); one displayed eigenvalue equation reverses the earlier sign of \(\Delta\).}
\GSfield{Statement}{On \(\mathbb R^2/(2\pi a\mathbb Z\times2\pi b\mathbb Z)\), the functions \(e^{i(nx/a+my/b)}\) have eigenvalues \(n^2/a^2+m^2/b^2\) for \(-\Delta\). Multiplicities and the zero mode must be retained when listing the frequencies. The positive-sign equation is for \(-\Delta\), not the previously defined \(\Delta\).}
\GSfield{Fate}{This is an elementary separation-of-variables calculation. It supplies a test case for the spectral-phase question, but its highly degenerate closed geodesics do not satisfy the later isolated-orbit hypotheses.}
\GSfield{Present status}{As of November 2026, the reconstructed statement is standard.}
\GSfield{Assessment}{literal statement: correct with missing hypotheses; publication match: none found; verification: read the result; development: annotated; confidence: high}
\end{GSitem}
\begin{GSitem}{1975-391.7}{C}{Abel-renormalized distributions of spectral phases}
\GSfield{Kind; source}{question; \GSpage{P000651-P000652}; A. I. Shnirelman}
\GSfield{The notes}{The question concerns \(a\lambda_k\pmod{2\pi}\). The proposed distribution subtracts its constant Fourier mode and multiplies an exponentially smoothed spectral sum by \(\varepsilon\).}
\GSfield{Statement}{Let \((M,g)\) be a closed compact Riemannian manifold, and let \(\lambda_k\ge0\) be the eigenvalues of \(\sqrt{-\Delta_g}\), counted with multiplicity. With \(\delta_{\rm per}\) the periodic delta distribution and \(a>0\), define
\[\Phi_{a,\varepsilon}(\mu)=\varepsilon\sum_k e^{-\varepsilon\lambda_k}\left(\delta_{\rm per}(\mu-a\lambda_k)-\frac1{2\pi}\right).\]
Ask whether \(\Phi_{a,\varepsilon}\) has a limit in \(\mathcal D'(\mathbb R/2\pi\mathbb Z)\), and how its Fourier coefficients depend on closed geodesics. This is a signed renormalized distribution, not the probability distribution of the first \(N\) phases. The notebook's intended generic setting has isolated closed geodesics with pairwise incommensurable primitive unoriented lengths; nondegeneracy of every contributing iterate is an additional modern hypothesis needed for the quoted principal trace formula.}
\GSfield{Reconstruction}{The precise topology and the sign convention in the Fourier transform are supplied here. The \(m\ne0\) coefficient is \(\varepsilon(2\pi)^{-1}\operatorname{tr}e^{-(\varepsilon+ima)\sqrt{-\Delta}}\); its individual limit and distributional convergence are different questions. The exclusion of iterates and orientation reversal from the incommensurability assertion is necessary: otherwise that assumption is impossible. The source's nondegeneracy abbreviation is not deciphered with certainty.}
\GSfield{Fate}{The closed-orbit interpretation is supported by the wave-trace expansion, already developed in 1975 and explicitly recalled by \citet[Introduction]{zelditch1998wave}. No located source proves the entire global distributional assertion in exactly the notebook's generality. The chapter supplies a precise Fourier convergence criterion, the complete circle calculation and a sphere example outside the generic assumptions; it does not establish the general geometric limit.}
\GSfield{Present status}{As of November 2026, the fixed-Fourier-mode interpretation is identified; a theorem with sufficient uniform bounds for the notebook’s full Abel limit was not located in the targeted searches.}
\GSfield{Assessment}{literal statement: correct with missing hypotheses; publication match: partial; verification: read the result; development: developed in this chapter; confidence: medium; problem \hyperref[prob:GP-1975-24]{GP-1975-24}}
\GSfield{Worth developing}{Uniform control of all Fourier modes is the remaining analytical step; the bounded circle model and the degenerate sphere example clarify the necessary scope.}
\end{GSitem}
\begin{GSitem}{1975-391.8}{E}{Closed-geodesic resonances in spectral phases}
\GSfield{Kind; source}{claim; \GSpage{P000652}; A. I. Shnirelman}
\GSfield{The notes}{Under a genericity assumption, the notes say \(\Phi_a\ne0\) precisely when \(a\) is a rational multiple of a closed-geodesic length.}
\GSfield{Statement}{At a fixed nonzero Fourier index \(m\), a possible contribution occurs when \(|m|a\) is a closed-orbit length. For finitely many isolated nondegenerate orbits of that length, the leading coefficient is a sum of their principal wave invariants. Nonvanishing requires this sum not to cancel. ``Distinct incommensurable lengths'' can only refer to primitive unoriented geodesics: iterates and orientation reversal already force relations.}
\GSfield{Reconstruction}{This is a justified local interpretation of the resonance statement. It does not assert a global distributional ``if and only if'' without the convergence and noncancellation hypotheses of the preceding item.}
\GSfield{Fate}{The singularities and their orbital coefficients belong to the pre-existing wave-trace theory; \citet[Introduction]{zelditch1998wave}. The exact Shnirelman publication for this Abel formulation remains unidentified.}
\GSfield{Present status}{As of November 2026, the reconstructed statement is standard.}
\GSfield{Assessment}{literal statement: correct with missing hypotheses; publication match: partial; verification: read the result; development: annotated; confidence: medium; problem \hyperref[prob:GP-1975-24]{GP-1975-24}}
\end{GSitem}
\begin{GSitem}{1975-391.9}{E}{The symplectic return map and its spectral types}
\GSfield{Kind; source}{definition; \GSpage{P000652}; A. I. Shnirelman}
\GSfield{The notes}{A closed geodesic has a Poincaré map; its multipliers are listed as reciprocal real pairs, unit-circle pairs, and complex quadruples, with the last type associated with dimension three.}
\GSfield{Statement}{The linearized return map acts symplectically on a transverse space of dimension \(2n-2\). Its nonreal, non-unit multipliers occur in quadruples \(\lambda,\bar\lambda,\lambda^{-1},\bar\lambda^{-1}\); real and unit-circle pairs give hyperbolic and elliptic blocks. A quadruple requires \(n\ge3\), not \(n=3\) exclusively. Nondegeneracy at an iterate means \(\det(I-P_\gamma^k)\ne0\).}
\GSfield{Fate}{The block decomposition and nondegenerate setting are used in \citet[Section 1]{zelditch1998wave}. The reciprocal-pair assertion also follows directly from symplectic linear algebra.}
\GSfield{Present status}{As of November 2026, the reconstructed statement is standard.}
\GSfield{Assessment}{literal statement: correct with missing hypotheses; publication match: partial; verification: read the result; development: annotated; confidence: high}
\end{GSitem}
\begin{GSitem}{1975-391.10}{E}{Elliptic and hyperbolic orbit amplitudes}
\GSfield{Kind; source}{claim; \GSpage{P000652}; A. I. Shnirelman}
\GSfield{The notes}{The coefficients are written with denominators \(\sin(k\alpha/2)\), \(\sinh(k\alpha/2)\), and \((-1)^k/\sinh(k\alpha/2)\) for three multiplier types.}
\GSfield{Statement}{The magnitude of a nondegenerate block's contribution is proportional to \(|\det(I-P^k)|^{-1/2}\), with a separate Maslov phase. For multipliers \(e^{\pm i\alpha}\) it is \((2|\sin(k\alpha/2)|)^{-1}\); for \(e^{\pm\alpha}\), \((2|\sinh(k\alpha/2)|)^{-1}\). For \(-e^{\pm\alpha}\), even \(k\) gives the hyperbolic sine denominator, while odd \(k\) gives \(2\cosh(k\alpha/2)\).}
\GSfield{Reconstruction}{\ours\ These three magnitudes are obtained directly by multiplying \((1-\lambda^k)(1-\lambda^{-k})\). A factor \((-1)^k\) cannot replace the sine/cosine distinction in the negative-real case.}
\GSfield{Fate}{The determinant and phase organization is the metaplectic form recalled in \citet[Introduction, equation (0.1)]{zelditch1998wave}. The block calculation corrects the literal manuscript formula.}
\GSfield{Present status}{As of November 2026, the reconstructed statement is standard.}
\GSfield{Assessment}{literal statement: false as written; publication match: partial; verification: read the result; development: annotated; confidence: high}
\end{GSitem}
\begin{GSitem}{1975-391.11}{E}{Regularizing the elliptic delta comb}
\GSfield{Kind; source}{construction; \GSpage{P000652}; A. I. Shnirelman}
\GSfield{The notes}{A regularized sum of deltas at half-integer multiples of an angle is written with the damping factor and constant subtraction inside the argument of \(\delta\). The preceding line explicitly requires the angle to be poorly approximable by rationals.}
\GSfield{Statement}{One coherent reconstruction is
\[D_{\alpha,\varepsilon}=\sum_{r\ge0}e^{-\varepsilon r}\left(\delta_{\rm per}(\mu-(r+\tfrac12)\alpha)-\frac1{2\pi}\right).\]
For \(m\ne0\), its Fourier coefficient is \((2\pi)^{-1}e^{-im\alpha/2}/(1-e^{-\varepsilon-im\alpha})\). If \(|1-e^{-im\alpha}|^{-1}\) grows at most polynomially, these coefficients converge to a distribution; a Diophantine condition on \(\alpha/2\pi\) suffices. Irrationality alone does not supply that bound.}
\GSfield{Reconstruction}{\ours\ The placement of damping and subtraction is emended to make a standard Abel sum. The Fourier coefficients follow from a geometric series and show exactly where a small-divisor hypothesis enters. The notebook does supply a poor-approximation hypothesis; the polynomial Fourier bound makes its required role precise.}
\GSfield{Fate}{The stated regularization and convergence criterion are derived here directly. No exact publication match to the manuscript’s misplaced parentheses was found.}
\GSfield{Present status}{As of November 2026, the reconstructed statement is standard.}
\GSfield{Assessment}{literal statement: false as written; publication match: none found; verification: read the result; development: annotated; confidence: high}
\end{GSitem}
\begin{GSitem}{1975-391.12}{F}{Hyperbolic continuation and an elliptic-function sketch}
\GSfield{Kind; source}{remark; \GSpage{P000653}; A. I. Shnirelman}
\GSfield{The notes}{The hyperbolic case is illustrated with poles and a reference to a Weierstrass function, but period and scale normalizations are not recorded.}
\GSfield{Statement}{A hyperbolic-block sequence has exponentially decaying coefficients such as \(1/\sinh(k\alpha/2)\). It therefore defines a periodic analytic function in a nonzero complex strip. A meromorphic expression in terms of elliptic functions would additionally require its periods, residues and additive normalization.}
\GSfield{Fate}{The sketch does not identify a unique Weierstrass function or an exact identity. Searches on Shnirelman, spectral phases, hyperbolic multipliers and Weierstrass functions did not recover that identity.}
\GSfield{Present status}{As of November 2026, the exact formula remains an identification gap; exponential convergence of the reconstructed Fourier series is elementary.}
\GSfield{Assessment}{literal statement: not determinable; publication match: none found; verification: unresolved lead; development: annotated; confidence: low}
\end{GSitem}
\begin{GSitem}{1975-391.13}{E}{Holomorphic half-wave trace versus meromorphic continuation}
\GSfield{Kind; source}{claim; \GSpage{P000653}; A. I. Shnirelman}
\GSfield{The notes}{For \(f(z)=\sum e^{i\lambda_kz}\), the notes claim downward analytic continuation with simple poles at geodesic lengths and identify residues with Fourier coefficients.}
\GSfield{Statement}{For \(\operatorname{Im}z>0\), the sum converges normally on compact subsets by Weyl's bound. Its boundary value is the distributional half-wave trace. At an isolated nondegenerate nonzero closed-orbit length, its leading singularity has simple-pole type, but lower terms generally include powers times logarithms. At zero the leading singularity has order \(n\). These facts do not yield global meromorphic continuation through the real axis.}
\GSfield{Fate}{The singularity expansion in \citet[Introduction, first displayed formula]{zelditch1998wave} contains the logarithmic terms absent from the notebook. This continues the wave-trace discussion of occurrences 1973-333.18--20.}
\GSfield{Present status}{As of November 2026, the reconstructed statement is standard.}
\GSfield{Assessment}{literal statement: false as written; publication match: partial; verification: read the result; development: annotated; confidence: high}
\end{GSitem}
\begin{GSitem}{1975-391.14}{E}{Sturm’s bound for linear combinations}
\GSfield{Kind; source}{theorem; \GSpage{P000651}; A. I. Shnirelman}
\GSfield{The notes}{The discussion recalls a one-dimensional bound on the zeros of a linear combination of the first eigenfunctions and compares it with polynomials.}
\GSfield{Statement}{For a regular real Sturm--Liouville problem on a compact interval with separated self-adjoint boundary conditions, order the simple eigenvalues increasingly. A nonzero linear combination of eigenfunctions with indices \(m,\ldots,n\) has at least \(m-1\) sign changes and at most \(n-1\) interior zeros, with multiplicities treated in the precise Sturm version. In particular the span of the first \(n\) eigenfunctions is a Chebyshev space in the interior.}
\GSfield{Fate}{The original theorem was announced in 1833 and published in 1836. A modern proof is \citet[Theorems 1.4 and2.15]{berardhelffer2018sturm}: first circulated 26 June 2017, revised16 October 2018, and published in the 2020 journal issue. The distinction between circulation and issue year is retained.}
\GSfield{Present status}{As of November 2026, the reconstructed statement is standard.}
\GSfield{Assessment}{literal statement: correct with missing hypotheses; publication match: exact; verification: read the result; development: annotated; confidence: high}
\end{GSitem}
\begin{GSitem}{1975-391.15}{E}{Courant’s nodal-domain bound}
\GSfield{Kind; source}{theorem; \GSpage{P000651}; A. I. Shnirelman}
\GSfield{The notes}{The multidimensional comparison invokes the bound for an individual Laplace eigenfunction.}
\GSfield{Statement}{For a real eigenfunction associated to the \(n\)-th eigenvalue of a self-adjoint Laplacian on a connected compact manifold, or a bounded domain with standard Dirichlet or Neumann conditions, the number of nodal domains is at most \(n\), with eigenvalues counted with multiplicity. This theorem applies to a single eigenspace; it does not automatically extend to sums of different eigenvalues.}
\GSfield{Fate}{The classical theorem and indexing convention are restated in \citet[Theorem 1.1]{berardhelffer2018courant}. The distinction from an extended assertion is the issue raised in the next intervention.}
\GSfield{Present status}{As of November 2026, the reconstructed statement is standard.}
\GSfield{Assessment}{literal statement: correct with missing hypotheses; publication match: exact; verification: read the result; development: annotated; confidence: high}
\end{GSitem}
\begin{GSitem}{1975-391.16}{B}{Arnold’s Harnack-type question for eigenfunction sums}
\GSfield{Kind; source}{question; \GSpage{P000651}; V. I. Arnold}
\GSfield{The notes}{Arnold asks about transferring Harnack-type assertions to linear combinations of eigenfunctions. No precise numerical bound is recorded.}
\GSfield{Statement}{A natural definite comparison is the extended Courant property: a nonzero sum of eigenfunctions with eigenvalues at most \(\lambda_n\) should have no more nodal domains than the first index of \(\lambda_n\). This property is false in general. For example, the Neumann equilateral triangle admits a constant plus a second eigenfunction with three nodal domains, exceeding the first index two.}
\GSfield{Reconstruction}{The extended Courant property is an explicit reconstruction of one plausible Harnack-type question, not a verbatim conjecture of Arnold in these notes. Other algebraic-geometric bounds may remain meaningful.}
\GSfield{Fate}{\citet[Statement 1.2 and Section 3]{berardhelffer2018courant}, circulated in 2017 and published in 2018, gives this precise comparison and counterexample. It does not settle every unspecified Harnack-type inequality.}
\GSfield{Present status}{As of November 2026, the specified extended Courant assertion is disproved. The general wording in the seminar is broader and cannot be assigned that single fate without the reconstruction caveat.}
\GSfield{Assessment}{literal statement: correct with missing hypotheses; publication match: partial; verification: read the result; development: annotated; confidence: medium; problem \hyperref[prob:GP-1975-25]{GP-1975-25}}
\end{GSitem}
\begin{GSitem}{1975-391.17}{E}{Ground-state positivity and the polynomial analogy}
\GSfield{Kind; source}{claim; \GSpage{P000651}; A. I. Shnirelman}
\GSfield{The notes}{The first eigenfunction is described as positive everywhere; successive eigenfunctions are compared with a linear function and a conic before the nodal-set questions.}
\GSfield{Statement}{For the scalar Laplacian on a connected closed manifold the lowest eigenvalue is zero and its eigenfunctions are constant, so a ground state can be chosen strictly positive. For a connected bounded domain with Dirichlet conditions, the first eigenfunction is strictly positive in the interior after choosing its sign, and vanishes on the boundary. The comparisons of later eigenfunctions with linear and quadratic polynomials are an analogy, not a general classification of their nodal sets.}
\GSfield{Reconstruction}{Connectedness, the boundary condition and the choice of sign are supplied; ``everywhere'' must be restricted to the interior in the Dirichlet case.}
\GSfield{Fate}{This is the standard ground-state comparison underlying the nodal discussion; \citet[Introduction and Theorem 1.1]{berardhelffer2018courant}. The later-eigenfunction analogy does not add a theorem beyond the precise Sturm and Courant assertions.}
\GSfield{Present status}{As of November 2026, the reconstructed statement is standard.}
\GSfield{Assessment}{literal statement: correct with missing hypotheses; publication match: partial; verification: read the result; development: annotated; confidence: high}
\end{GSitem}

\GStalk{392}{1976-04-05}{A. M. Polyakov}{P000652-P000655}{Lattice Yang–Mills fields and their partition function}
The continuation announced on P000647 appears on P000654. Polyakov replaces continuum connections by group-valued link variables, forms plaquette products and integrates the Boltzmann weight against product Haar measure. This is an expository return to the Wilson-lattice material of 1974-345.1--8. No new explicit confinement conjecture is written on this page.

\begin{GSitem}{1975-392.1}{E}{Continuum action and its coupling convention}
\GSfield{Kind; source}{definition; \GSpage{P000654}}
\GSfield{The notes}{The action is \((4e^2)^{-1}\operatorname{tr}\int A_{\mu\nu}^2\), immediately followed by a definition of \(F_{\mu\nu}\); the commutator has no coupling constant.}
\GSfield{Statement}{After absorbing the coupling into the connection, \(F=dA+A\wedge A\) and \(S_E=-(4e^2)^{-1}\int\operatorname{tr}F_{\mu\nu}F_{\mu\nu}\,d^4x\) for anti-Hermitian matrices. Stationarity gives \(D^\mu F_{\mu\nu}=0\). If the matrices commute, the curvature reduces to the Abelian electromagnetic expression.}
\GSfield{Reconstruction}{The curvature symbol in the integrand and Euclidean sign are emended; the changed coupling convention is legitimate. The Euler--Lagrange equation assumes compactly supported variations or vanishing boundary terms.}
\GSfield{Fate}{This is the continuum starting point for \citet[Sections I--II]{wilson1974confinement}, already published in 1974.}
\GSfield{Present status}{As of November 2026, the reconstructed statement is standard.}
\GSfield{Assessment}{literal statement: false as written; publication match: partial; verification: read the result; development: annotated; confidence: high}
\end{GSitem}
\begin{GSitem}{1975-392.2}{E}{Link variables and plaquette action}
\GSfield{Kind; source}{construction; \GSpage{P000654}}
\GSfield{The notes}{Each directed link carries \(B_{x,\mu}\in SU(2)\); the four-link plaquette product enters a class function \(\chi\), called ``any character'' and then described as constant on conjugacy classes.}
\GSfield{Statement}{For link variables \(U_{x,\mu}\), set
\[U_p=U_{x,\mu}U_{x+a_\mu,\nu}U^{-1}_{x+a_\nu,\mu}U^{-1}_{x,\nu}.\]
Under \(U_{x,\mu}\mapsto\Omega_x^{-1}U_{x,\mu}\Omega_{x+a_\mu}\), the plaquette is conjugated at \(x\). Hence \(\sum_p\phi(U_p)\) is gauge invariant for every class function \(\phi\). A character is a special class function. The Wilson choice is proportional to \(\sum_p(1-\frac12\operatorname{Re}\operatorname{tr}U_p)\).}
\GSfield{Reconstruction}{An arbitrary class function gives an invariant lattice action but does not automatically have the desired continuum Yang--Mills limit. The subtraction and small-plaquette quadratic normalization distinguish the Wilson choice.}
\GSfield{Fate}{The link and plaquette construction is \citet[Section III]{wilson1974confinement}, received 12 June 1974 and published 15 October 1974. The source is therefore expository, not a new 1976 discretization.}
\GSfield{Present status}{As of November 2026, the reconstructed statement is standard.}
\GSfield{Assessment}{literal statement: correct with missing hypotheses; publication match: exact; verification: read the result; development: annotated; confidence: high}
\end{GSitem}
\begin{GSitem}{1975-392.3}{E}{Finite-lattice Haar integral and correlations}
\GSfield{Kind; source}{definition; \GSpage{P000654}}
\GSfield{The notes}{The notes write \(Z=\int\prod_{x,\mu}dB_{x,\mu}e^{-S(B)}\), count \(4n\) group integrations and \(12n\) variables, and leave a heading for correlation functions.}
\GSfield{Statement}{On a finite periodic four-dimensional lattice with \(n\) sites, there are \(4n\) independent positively oriented links. Normalized Haar measure on each \(SU(2)\) factor defines \(Z\); for a real continuous action, \(0<Z<\infty\). Expectations are \(\langle O\rangle=Z^{-1}\int Oe^{-S}\prod dU\). The unreduced configuration space has dimension \(12n\); the quotient by gauge transformations does not have that dimension.}
\GSfield{Fate}{\citet[Section III]{wilson1974confinement} provides the gauge-invariant lattice integral. Existence at finite lattice spacing and volume follows from compactness and does not establish an infinite-volume or continuum limit.}
\GSfield{Present status}{As of November 2026, the reconstructed statement is standard.}
\GSfield{Assessment}{literal statement: correct with missing hypotheses; publication match: exact; verification: read the result; development: annotated; confidence: high}
\end{GSitem}
\begin{GSitem}{1975-392.4}{E}{Wick rotation and the continuum functional integral}
\GSfield{Kind; source}{method; \GSpage{P000654}}
\GSfield{The notes}{The schematic replacement \(t\to it\), \(e^{iS}\to e^{-S}\), is followed by ``in the continuous case---a functional integral''. No measure is constructed.}
\GSfield{Statement}{Imaginary-time continuation motivates the Euclidean Boltzmann factor after the action and integration contour are transformed consistently. A rigorous reconstruction of a Lorentzian theory requires additional analytic and positivity properties. The finite product Haar integral is an actual measure; the displayed continuum functional integral remains formal here.}
\GSfield{Fate}{The Euclidean lattice strategy predates the talk \citep[Section III]{wilson1974confinement}; the continuum existence problem is stated separately by \citet[Section 4]{jaffewitten2000yangmills}.}
\GSfield{Present status}{As of November 2026, the reconstructed statement is standard.}
\GSfield{Assessment}{literal statement: correct with missing hypotheses; publication match: partial; verification: read the result; development: annotated; confidence: high}
\end{GSitem}

\GStalk{393}{1976-04-12}{D. B. Fuchs}{P000653-P000657}{Computer calculations in Hamiltonian Lie-algebra cohomology}
Fuchs describes the polynomial cochain model, its finite weight decomposition, the exotic Gelfand--Kalinin--Fuchs classes and Perchik's constant-term Euler series. The two definite questions are to compute the whole cohomology and to prove nonvanishing of the relative Euler characteristic in unbounded weights. Buchstaber's short intervention on P000656 concerns coalgebra structures arising from multiplication and averaging. Its placement is preserved within this catalogue record. The constant-term calculation is developed with a finite, explicitly checked scope in Section~\ref{dev:perchik}.

\begin{GSitem}{1975-393.1}{E}{The Hamiltonian cochain complex in polynomial symbols}
\GSfield{Kind; source}{definition; \GSpage{P000655}}
\GSfield{The notes}{The opening calls the algebra formal divergence-free vector fields, then writes alternating polynomials in pairs \((\xi_i,\eta_i)\), with differential built from \(\langle\xi_i,\eta_j\rangle-\langle\xi_j,\eta_i\rangle\).}
\GSfield{Statement}{The displayed complex is the continuous Chevalley--Eilenberg complex of
\[\mathfrak{ham}_{2n}=\mathbb R[[p_1,\ldots,p_n,q_1,\ldots,q_n]]/\mathbb R\]
with the Poisson bracket. A cochain is represented by a polynomial alternating under exchange of argument pairs and vanishing when any entire pair is zero. Writing \(z_i=(\xi_i,\eta_i)\) and \(\omega(z_i,z_j)=\langle\xi_i,\eta_j\rangle-\langle\xi_j,\eta_i\rangle\), the source convention is
\[(dP)(z_1,\ldots,z_{q+1})=\sum_{i<j}(-1)^{i+j-1}\omega(z_i,z_j)P(z_i+z_j,z_1,\ldots,\widehat z_i,\ldots,\widehat z_j,\ldots,z_{q+1}).\]
The order of the remaining arguments is unchanged.}
\GSfield{Reconstruction}{The quotient by constants supplies the omitted reduced-cochain condition. In two variables divergence-free formal fields are Hamiltonian; in higher even dimension the two Lie algebras differ. The Poisson factor, not the opening shorthand, identifies this complex.}
\GSfield{Fate}{The algebra and continuous weight convention are explicitly fixed in \citet[Section 2]{kotschickmorita2009gkf}. The displayed differential is the dual of the Poisson bracket, so its square is zero by Jacobi.}
\GSfield{Present status}{As of November 2026, the reconstructed statement is standard.}
\GSfield{Assessment}{literal statement: false as written; publication match: partial; verification: read the result; development: annotated; confidence: high}
\end{GSitem}
\begin{GSitem}{1975-393.2}{C}{Computing the full Hamiltonian cohomology}
\GSfield{Kind; source}{question; \GSpage{P000655}}
\GSfield{The notes}{The underlined task is to find the cohomology of the displayed complex, said to remain undone even for \(n=1\).}
\GSfield{Statement}{Determine the continuous cohomology \(H^*(\mathfrak{ham}_{2n};\mathbb R)\), including its graded dimensions and multiplicative structure. The first unresolved case intended here is \(\mathfrak{ham}_2\), not the Lie algebra of all formal vector fields in one variable.}
\GSfield{Fate}{Later computations and factorizations give substantial finite-weight information. \citet[p.151 and Theorem 1.2]{mikami2016metoki} still states the infinite-dimensionality question and proves the Metoki-class factorization. Stable-rank graph-homology calculations concern a different limiting object \citep[Remark 2.1]{moritasakasaisuzuki2015computations}.}
\GSfield{Present status}{As of 10 November 2026, targeted searches located no complete fixed-rank computation or resolution of the infinitude question. The latest explicitly verified open-status statement used here is from 2016; this search result is not a proof of present-day nonexistence.}
\GSfield{Assessment}{literal statement: correct with missing hypotheses; publication match: partial; verification: read the result; development: annotated; confidence: medium; problem \hyperref[prob:GP-1975-21]{GP-1975-21}}
\GSfield{Worth developing}{The finite weight complexes permit exact calculations, but a full answer needs a mechanism valid in unbounded weights.}
\end{GSitem}
\begin{GSitem}{1975-393.3}{E}{The all-vector-field comparison}
\GSfield{Kind; source}{construction; \GSpage{P000655}}
\GSfield{The notes}{The notes say that taking polynomials linear in \(\eta\) gives the cochain complex of \(W_n\).}
\GSfield{Statement}{Fibre-linear Hamiltonians \(H_v(p,q)=\sum_i p_i v_i(q)\) form a Lie subalgebra isomorphic, up to the fixed Poisson-sign convention, to the Lie algebra \(W_n\) of formal vector fields on \(\mathbb R^n\). Restricting Hamiltonian cochains to this subalgebra retains the symbol component linear in each fibre variable. This is a restriction map of complexes; those polynomial components need not form a subcomplex of the unrestricted Hamiltonian-symbol complex without projection.}
\GSfield{Reconstruction}{The restriction interpretation makes the informal sentence precise and prevents the next sentence, ``not done even for \(n=1\)'', from being misapplied to \(W_1\).}
\GSfield{Fate}{The cohomology of \(W_n\) had already been calculated by \citet[Theorem 0.3]{gelfandfuchs1970formal}, received 4 November 1969 and published in 1970. The full original English article was opened; the present distinction follows from the displayed fibre-linear embedding.}
\GSfield{Present status}{As of November 2026, the reconstructed statement is standard.}
\GSfield{Assessment}{literal statement: correct with missing hypotheses; publication match: partial; verification: read the result; development: annotated; confidence: high}
\end{GSitem}
\begin{GSitem}{1975-393.4}{E}{Finite bigraded complexes and off-diagonal contraction}
\GSfield{Kind; source}{method; \GSpage{P000655}}
\GSfield{The notes}{The total degrees in \(\xi\) and \(\eta\) are \(m_1+q\) and \(m_2+q\). Each \(K_{m_1,m_2}\) is finite-dimensional, and \(m_1\ne m_2\) is acyclic.}
\GSfield{Statement}{The differential preserves the two shifted degrees. For fixed \(n,m_1,m_2\), only finitely many exterior monomials occur: the weight-minus-one translation space is finite-dimensional, the quadratic weight-zero space is finite-dimensional, and every other factor has positive weight. The Hamiltonian \(\sum_i p_iq_i\) acts on this complex by the scalar \(m_1-m_2\), up to sign. Cartan's identity gives a contracting homotopy on a nonzero eigenspace, hence
\[H^*(\mathfrak{ham}_{2n})=\bigoplus_m H^*(K_{m,m}).\]}
\GSfield{Reconstruction}{\ours\ The finiteness argument and contraction are supplied explicitly. The formal-cochain topology is continuous cohomology; an unrestricted algebraic dual would be a different object.}
\GSfield{Fate}{The finite weight construction agrees with \citet[Section 2]{kotschickmorita2009gkf}. The source already states the reduction; the argument here explains why it is valid.}
\GSfield{Present status}{As of November 2026, the reconstructed statement is standard.}
\GSfield{Assessment}{literal statement: correct with missing hypotheses; publication match: partial; verification: read the result; development: annotated; confidence: high}
\end{GSitem}
\begin{GSitem}{1975-393.5}{E}{Low weights and the exotic Gelfand–Kalinin–Fuchs classes}
\GSfield{Kind; source}{theorem; \GSpage{P000655}}
\GSfield{The notes}{For \(n=1\), the notes give acyclicity at \(m=1,2,3\) and nonzero cohomology in degrees 7 and 10 at \(m=4\), followed by a machine-memory failure at \(m=5\).}
\GSfield{Statement}{In total Hamiltonian weight \(w=2m\), the exotic relative Gelfand--Kalinin--Fuchs class lies in \(H^7(\mathfrak{ham}_2,\mathfrak{sp}_2)_8\). The corresponding absolute complex has classes in degrees 7 and 10. The statement about exhausted computer memory is a report about a calculation, not a cohomological vanishing or nonvanishing assertion. The absolute complexes at weights \(2,4,6\) are acyclic.}
\GSfield{Fate}{The exotic class was published in Russian in 1972. Its later factorization is recalled in \citet[Theorem 1.1]{mikami2016metoki}: a degree-five, weight-ten class for fields vanishing at the origin is multiplied by the transverse symplectic class. Theorem 1.2 similarly factors the Metoki class of degree nine and weight fourteen, first reported in 1999.}
\GSfield{Present status}{As of November 2026, the reconstructed statement is standard.}
\GSfield{Assessment}{literal statement: correct with missing hypotheses; publication match: exact; verification: read the result; development: annotated; confidence: high}
\end{GSitem}
\begin{GSitem}{1975-393.6}{E}{Relative cohomology and the omitted zero-weight exception}
\GSfield{Kind; source}{construction; \GSpage{P000655}}
\GSfield{The notes}{A subcomplex \(K'_{m,m}\) is introduced with \(H(K_{m,m})=H(K'_{m,m})\otimes H(S^3)\); its Euler characteristic is initially denoted \(\chi_m^*\), then loses the star.}
\GSfield{Statement}{Here \(K'\) is relative to quadratic Hamiltonians \(\mathfrak{sp}_2\cong\mathfrak{sl}_2\). For \(m\ne0\),
\[H^q(K_{m,m})\cong H^q(K'_{m,m})\oplus H^{q-3}(K'_{m,m}).\]
At weight zero, absolute groups occur in degrees0 and7, relative groups in degrees0 and4. The nontrivial Euler series is therefore relative; the absolute Euler characteristics vanish.}
\GSfield{Fate}{The zero-weight exception is explicit in \citet[Section 1 and Theorem 2.1]{mohammadzadehfuchs2014cohomology}.}
\GSfield{Present status}{As of November 2026, the reconstructed statement is standard.}
\GSfield{Assessment}{literal statement: false as written; publication match: exact; verification: read the result; development: annotated; confidence: high}
\end{GSitem}
\begin{GSitem}{1975-393.7}{E}{Perchik’s constant-term Euler series}
\GSfield{Kind; source}{construction; \GSpage{P000655-P000656}}
\GSfield{The notes}{Perchik is credited with a generating function whose constant term in \(x\) is \(t^{-2}+2-t^8-t^{14}-t^{22}+\cdots+10t^{70}+\cdots\). The general row of factors contains an ambiguous middle exponent of \(t\).}
\GSfield{Statement}{Put \(E_m=\chi H^*(\mathfrak{ham}_2,\mathfrak{sp}_2)_{2m}\). Then
\[\sum_{m\ge-1}E_mt^{2m}=\frac12\operatorname{CT}_x\left[(1-t^{-1}x^{-1})(1-t^{-1}x)(1-x^{-2})(1-x^2)\prod_{\ell\ge1}\prod_{j=0}^{\ell+2}(1-t^\ell x^{-\ell-2+2j})\right].\]
The infinite product has a well-defined formal Laurent expansion in \(t\) with lower bound \(-2\). The linear Hamiltonians supply the two negative-weight factors; deleting them gives a different Lie algebra, \(\mathfrak{ham}_2^0\).}
\GSfield{Reconstruction}{The symbol \(\operatorname{CT}_x\) formalizes ``discard everything containing \(x\)''. Every factor in the weight-\(\ell\) row has \(t^\ell\). Section~\ref{dev:perchik} gives an independent invariant-character derivation and finite cutoff proof.}
\GSfield{Fate}{Harvard's dissertation list dates Perchik's dissertation to 1975 \citep{perchik1975dissertation}; this documents earlier work, but neither a precise public-availability date nor a theorem locator in the dissertation was verified. The expository classification E is therefore provisional, based on that documented dissertation; exact pre-seminar public circulation remains unestablished. The article appeared in 1976. \citet[Section 5]{kotschickmorita2009gkf} distinguishes the relative constructions; its Remark 9 has an inconsistent superscript for the full series, which is not followed here.}
\GSfield{Present status}{As of November 2026, the relative constant-term formula is established in the later primary literature; its exact pre-seminar circulation remains unverified.}
\GSfield{Assessment}{literal statement: correct with missing hypotheses; publication match: exact; verification: read the result; development: developed in this chapter; confidence: medium}
\end{GSitem}
\begin{GSitem}{1975-393.8}{C}{Infinitely many nonzero relative Euler characteristics}
\GSfield{Kind; source}{question; \GSpage{P000655-P000656}}
\GSfield{The notes}{Fuchs asks to prove that \(\chi_m\ne0\) for arbitrarily large \(m\). The question follows the relative-complex notation and Perchik's product.}
\GSfield{Statement}{For the coefficients \(E_m\) of the preceding fixed-rank relative Euler series, prove or disprove that \(\{m\ge1:E_m\ne0\}\) is infinite. A positive answer implies infinitely many nonzero weight components of Hamiltonian cohomology; the converse does not follow, since even and odd cohomology can cancel in an Euler characteristic.}
\GSfield{Fate}{The exact coefficient-infinitude question persists in \citet[Section 2]{mohammadzadehfuchs2014cohomology}. The weaker cohomology-infinitude question remains explicit in \citet[p.151]{mikami2016metoki}. This chapter establishes only a finite truncation and exact coefficient verification.}
\GSfield{Present status}{As of 10 November 2026, no resolution of this exact coefficient-infinitude question was located in the recorded targeted searches. The chapter’s computation through m=150 supplies finite evidence and no claim about all later weights.}
\GSfield{Assessment}{literal statement: correct with missing hypotheses; publication match: partial; verification: read the result; development: developed in this chapter; confidence: medium; problem \hyperref[prob:GP-1975-20]{GP-1975-20}}
\GSfield{Worth developing}{A closed asymptotic or structural argument for the constant-term series would turn finite evidence into an infinitude theorem.}
\end{GSitem}
\begin{GSitem}{1975-393.9}{E}{The numerical lower bound of 57 cohomology classes}
\GSfield{Kind; source}{claim; \GSpage{P000656}}
\GSfield{The notes}{The expansion through \(t^{70}\) is said to prove the existence of at least 57 nontrivial cohomology classes.}
\GSfield{Statement}{For each finite relative weight complex, \(\sum_q\dim H^q_{2m}\ge|E_m|\). The computed coefficients satisfy
\[\sum_{m=-1}^{35}|E_m|=57.\]
Thus the total relative cohomology has dimension at least 57 when these weights, including weight zero and the negative weight, are included. The positive weights alone give 54. The number does not mean 57 positive-weight primitive exotic generators.}
\GSfield{Reconstruction}{\oursComputed\ The chapter independently reproduces this finite sum and \(E_{35}=10\), then extends the verified range through \(m=150\). The inequality is the elementary comparison of the alternating and ungraded dimensions of a finite complex.}
\GSfield{Fate}{This supplies a precise interpretation and independent check of the notebook's 57. The development does not equate an Euler coefficient with a cohomological dimension or a count of algebra generators.}
\GSfield{Present status}{As of November 2026, the reconstructed statement is standard.}
\GSfield{Assessment}{literal statement: correct with missing hypotheses; publication match: partial; verification: read the result; development: developed in this chapter; confidence: high}
\end{GSitem}
\begin{GSitem}{1975-393.10}{E}{Multiplication dualized to a coproduct}
\GSfield{Kind; source}{construction; \GSpage{P000656}; V. M. Buchstaber}
\GSfield{The notes}{Buchstaber starts with \(X\times X\to X\) for a genuine group and writes a map \(C(X)\to C(X\times X)\) as a coalgebra structure.}
\GSfield{Statement}{For a compact Hausdorff topological group, multiplication gives \(\Delta f(x,y)=f(xy)\) on continuous functions, with counit evaluation at the identity. Associativity gives coassociativity. In this setting \(C(X\times X)\) is a completed tensor product; an algebraic tensor-product coalgebra is obtained instead on a suitable algebra of representative functions.}
\GSfield{Reconstruction}{The topology and tensor-product convention are supplied. Arbitrary continuous functions need not lie in the algebraic tensor product \(C(X)\otimes C(X)\).}
\GSfield{Fate}{The algebra/coalgebra viewpoint leads directly into \citet[Section 3]{buchstaber1975twovaluedI}, received 2 December 1974 and published in 1975. The notebook gives only the opening construction, not a complete theorem about all topological spaces.}
\GSfield{Present status}{As of November 2026, the reconstructed statement is standard.}
\GSfield{Assessment}{literal statement: correct with missing hypotheses; publication match: partial; verification: read the result; development: annotated; confidence: high}
\end{GSitem}
\begin{GSitem}{1975-393.11}{E}{Averaging a two-valued formal multiplication}
\GSfield{Kind; source}{construction; \GSpage{P000656}; V. M. Buchstaber}
\GSfield{The notes}{The next line replaces a value of \(f(x\cdot y)\) by its mathematical expectation. The heading says that Buchstaber is reporting Adams on Buchstaber.}
\GSfield{Statement}{For a commutative two-valued formal group law over a \(\mathbb Q\)-algebra, with branches \(F_+(x,y),F_-(x,y)\), the natural linear coproduct is
\[\Delta f=\tfrac12\bigl(f(F_+(x,y))+f(F_-(x,y))\bigr).\]
The symmetric expression is interpreted through the formal law even when its branches are not separately in the base power-series ring. Associativity of the two-valued law gives coassociativity of this averaging map, which generally is not an algebra homomorphism.}
\GSfield{Reconstruction}{This is a concrete, explicitly labelled identification of the expectation notation, supported by the contemporaneous two-valued-formal-group literature. The fragment does not itself specify the number of values or the full axioms.}
\GSfield{Fate}{The construction is read in \citet[Section 3 and the generalized translations of Section 4]{buchstaber1975twovaluedI}. The bibliographical lead Adams--Liulevicius, ``Buhstaber's work on two-valued formal groups'' (1975), matches the heading, but its full text was not accessible and no detailed theorem is attributed to it here.}
\GSfield{Present status}{As of November 2026, the reconstructed statement is standard.}
\GSfield{Assessment}{literal statement: correct with missing hypotheses; publication match: partial; verification: read the result; development: annotated; confidence: medium}
\end{GSitem}

\GStalk{394}{1976-04-12}{S. P. Novikov}{P000655-P000658}{Lorentzian geometry, Einstein’s equation and matter}
Novikov begins a geometrical introduction to general relativity after Buchstaber on P000656. The discussion covers the Lorentz metric, curvature, stress--energy, variational principles, perfect fluids and the Schwarzschild extension. Gelfand's question about quantizing geometry and Novikov's dimensional estimates are recorded separately from mathematical theorems. The next meeting continues this discussion in record396.

\begin{GSitem}{1975-394.1}{E}{The Lorentz metric and proper time}
\GSfield{Kind; source}{definition; \GSpage{P000656}}
\GSfield{The notes}{Coordinates are \((x^0,\vec x)=(ct,\vec x)\), the flat metric has diagonal entries \((1,-1,-1,-1)\), and proper time is called the length of a curve.}
\GSfield{Statement}{A spacetime is a smooth four-manifold with a Lorentz metric of signature \((+---)\). Along a future-directed timelike curve \(\gamma\),
\[\tau(\gamma)=c^{-1}\int\sqrt{g(\dot\gamma,\dot\gamma)}\,ds.\]
Null curves have zero Lorentzian length. A proper-time interpretation therefore requires a timelike curve and includes the factor \(c^{-1}\) when \(x^0=ct\).}
\GSfield{Fate}{The metric interpretation is classical; \citet[Sections 3--4]{einstein1916foundation}. The modern signature and units have been fixed explicitly.}
\GSfield{Present status}{As of November 2026, the reconstructed statement is standard.}
\GSfield{Assessment}{literal statement: correct with missing hypotheses; publication match: exact; verification: read the result; development: annotated; confidence: high}
\end{GSitem}
\begin{GSitem}{1975-394.2}{E}{Curvature contractions and Einstein’s equation}
\GSfield{Kind; source}{definition; \GSpage{P000656}}
\GSfield{The notes}{The curvature is a two-form with values in the Lie algebra of \(O(n)\); contraction gives \(R_{ij}\) and \(R\). The field equation is \(R_{ij}-\tfrac12Rg_{ij}=(8\pi K/c^4)T_{ij}\), followed by its arbitrary-dimensional form.}
\GSfield{Statement}{In a Lorentz-orthonormal frame, the curvature takes values in \(\mathfrak{o}(1,3)\), not \(\mathfrak{o}(4)\). Fix a curvature-sign convention and let \(R_{ij}=R^k{}_{ikj}\), \(R=g^{ij}R_{ij}\). With zero cosmological constant, Einstein's equation is \(G_{ij}=\kappa T_{ij}\), where \(G_{ij}=R_{ij}-\tfrac12Rg_{ij}\) and \(\kappa=8\pi G/c^4\). The tensor formula exists in any dimension, but in dimension two \(G\) vanishes identically.}
\GSfield{Reconstruction}{The orthogonal Lie algebra is corrected for the displayed signature. The notes' historical date 1916 is the year of Einstein's systematic exposition, not the first formulation of the final field equation.}
\GSfield{Fate}{The contractions and field equation appear in \citet[Sections 12, 14 and 16]{einstein1916foundation}. The exposition itself cites the 1915 work.}
\GSfield{Present status}{As of November 2026, the reconstructed statement is standard.}
\GSfield{Assessment}{literal statement: false as written; publication match: partial; verification: read the result; development: annotated; confidence: high}
\end{GSitem}
\begin{GSitem}{1975-394.3}{E}{Stress–energy from variation of the metric}
\GSfield{Kind; source}{definition; \GSpage{P000656}}
\GSfield{The notes}{The source writes \(S=\int L\sqrt g\,d^4x\) and \(\sqrt g\,T_{ij}=\delta S/\delta g_{ij}\).}
\GSfield{Statement}{For Lorentz signature the invariant density is \(\sqrt{|\det g|}\). One consistent convention is
\[\delta S_m=\frac12\int\sqrt{|g|}\,T_{ij}\,\delta g^{ij}\,d^4x
=-\frac12\int\sqrt{|g|}\,T^{ij}\,\delta g_{ij}\,d^4x.\]
Thus \(T_{ij}=2|g|^{-1/2}\delta S_m/\delta g^{ij}\) in this convention. Reversing the sign of the matter action reverses the defining sign; it does not remove the factor two or the change of index positions.}
\GSfield{Reconstruction}{\ours\ The second equality follows by differentiating \(g^{ik}g_{kj}=\delta^i_j\). The sign convention is declared rather than inferred from the abbreviated notebook Lagrangian.}
\GSfield{Fate}{This is a normalization of the standard variational definition. The notebook formula is mathematically incomplete as an indexed variational identity.}
\GSfield{Present status}{As of November 2026, the reconstructed statement is standard.}
\GSfield{Assessment}{literal statement: false as written; publication match: none found; verification: read the result; development: annotated; confidence: high}
\end{GSitem}
\begin{GSitem}{1975-394.4}{E}{The electromagnetic action and differential form}
\GSfield{Kind; source}{construction; \GSpage{P000656}}
\GSfield{The notes}{The electromagnetic Lagrangian is written \(F_{\mu\nu}F^{\mu\nu}\); the curvature two-form is written with reversed wedge indices and no half.}
\GSfield{Statement}{For a local potential one-form \(A\), set
\[F=dA=\tfrac12F_{\mu\nu}\,dx^\mu\wedge dx^\nu,\qquad F_{\mu\nu}=\partial_\mu A_\nu-\partial_\nu A_\mu.\]
With signature \((+---)\), choose \(L_m=-\tfrac14F_{\mu\nu}F^{\mu\nu}\). The preceding variational convention gives \(T_{\mu\nu}=-F_{\mu\alpha}F_\nu{}^\alpha+\tfrac14g_{\mu\nu}F_{\alpha\beta}F^{\alpha\beta}\), whose energy density in an orthonormal frame is \((|E|^2+|B|^2)/2\).}
\GSfield{Reconstruction}{\ours\ The factors and signs are reconstructed by the definition of a differential two-form and variation of the stated Lagrangian. The source correctly observes that raising indices introduces the metric.}
\GSfield{Fate}{The generally covariant electromagnetic field is treated in \citet[Section 20]{einstein1916foundation}; the particular normalization here is made explicit independently.}
\GSfield{Present status}{As of November 2026, the reconstructed statement is standard.}
\GSfield{Assessment}{literal statement: false as written; publication match: partial; verification: read the result; development: annotated; confidence: high}
\end{GSitem}
\begin{GSitem}{1975-394.5}{E}{Vacuum Einstein equations from scalar curvature}
\GSfield{Kind; source}{method; \GSpage{P000656}}
\GSfield{The notes}{In the absence of forces, stationarity of \(\int R\sqrt g\,d^4x\) is written as the field equation.}
\GSfield{Statement}{For compactly supported variations of the inverse metric in the interior, the first variation of \(\int R\sqrt{|g|}\,d^4x\) is \(\int G_{ij}\delta g^{ij}\sqrt{|g|}\,d^4x\), plus a boundary divergence. Consequently stationary metrics satisfy \(G_{ij}=0\). In dimension four this is equivalent to \(R_{ij}=0\). On a manifold with boundary, arbitrary boundary variations require a boundary prescription and possibly a boundary term.}
\GSfield{Reconstruction}{\ours\ This is the standard integration-by-parts calculation of the Einstein--Hilbert variation, with compact support specifying which divergence is discarded. ``No forces'' is interpreted here as zero matter stress and zero cosmological constant.}
\GSfield{Fate}{The source is an expository variational principle. The stated hypotheses are needed to turn its shorthand into an Euler–Lagrange assertion.}
\GSfield{Present status}{As of November 2026, the reconstructed statement is standard.}
\GSfield{Assessment}{literal statement: correct with missing hypotheses; publication match: none found; verification: read the result; development: annotated; confidence: high}
\end{GSitem}
\begin{GSitem}{1975-394.6}{E}{The meaning of gravitational energy in the field equation}
\GSfield{Kind; source}{claim; \GSpage{P000657}}
\GSfield{The notes}{The opening line says there is no energy--momentum in the gravitational field except that belonging to other forms of energy.}
\GSfield{Statement}{The matter tensor on the right-hand side excludes a separate tensor for the gravitational field. This does not assert that gravity carries no energy. Coordinate-dependent energy complexes and conserved quantities defined using suitable global or asymptotic structure are different notions. The absence of a separate local matter tensor must not be confused with vanishing gravitational energy.}
\GSfield{Fate}{Einstein's own account introduces a gravitational energy complex and discusses conservation in \citet[Sections 15--18]{einstein1916foundation}. The notebook's abbreviated sentence needs that distinction.}
\GSfield{Present status}{As of November 2026, the reconstructed statement is standard.}
\GSfield{Assessment}{literal statement: correct with missing hypotheses; publication match: partial; verification: read the result; development: annotated; confidence: high}
\end{GSitem}
\begin{GSitem}{1975-394.7}{F}{Gelfand’s question about quantizing geometry}
\GSfield{Kind; source}{question; \GSpage{P000657}; I. M. Gelfand}
\GSfield{The notes}{Gelfand asks whether the gravitational field is quantized and suggests that, if it is, there is no metric or geometrical structure of space.}
\GSfield{Statement}{The intervention asks how a quantum description of gravity changes the status of the classical metric. No choice of observables, quantization, dynamical model or precise conclusion about geometry is specified. Quantization by itself does not logically imply that every notion of metric must disappear.}
\GSfield{Fate}{No definite mathematical conjecture can be extracted without supplying a quantum-gravity framework absent from the page. The record preserves the question as an intervention rather than assigning a solved/open status to an invented theorem.}
\GSfield{Present status}{As of November 2026, this remains an identification limitation of the notebook fragment, not an assessment that all quantum-gravity programs are one unresolved mathematical problem.}
\GSfield{Assessment}{literal statement: not determinable; publication match: none found; verification: unresolved lead; development: annotated; confidence: low}
\end{GSitem}
\begin{GSitem}{1975-394.8}{E}{Planck-scale dimensional estimates}
\GSfield{Kind; source}{claim; \GSpage{P000657}}
\GSfield{The notes}{Novikov gives scales \(10^{-33}\) cm and \(10^{-43}\) seconds, contrasts nuclear interactions at \(10^{-23}\) cm, and suggests a limit on quantum mechanics.}
\GSfield{Statement}{The dimensional combinations are \(\ell_P=\sqrt{\hbar G/c^3}\) and \(t_P=\ell_P/c\), whose orders agree with the two quoted Planck estimates. The nuclear length \(10^{-23}\) cm is inconsistent with the usual nuclear scale of order \(10^{-13}\) cm; \(10^{-23}\) seconds is instead a characteristic nuclear time order. Dimensional analysis alone does not establish a mathematical breakdown of quantum mechanics or spacetime at these values.}
\GSfield{Reconstruction}{\ours\ The powers of \(\hbar,G,c\) follow by balancing mass, length and time dimensions. The scan clearly has centimetres in the nuclear comparison; the intended correction cannot be determined uniquely between its exponent and unit.}
\GSfield{Fate}{The two Planck orders are an interpretive dimensional estimate. No theorem justifying the stated limit on quantum mechanics is supplied or inferred.}
\GSfield{Present status}{As of November 2026, the reconstructed statement is standard.}
\GSfield{Assessment}{literal statement: false as written; publication match: none found; verification: read the result; development: annotated; confidence: medium}
\end{GSitem}
\begin{GSitem}{1975-394.9}{E}{The Einstein–perfect-fluid system}
\GSfield{Kind; source}{problem; \GSpage{P000657}}
\GSfield{The notes}{The material model is \(T_{ij}=(\varepsilon+p)u_i u_j-pg_{ij}\), with an equation of state \(p=p(\varepsilon)\). The task is to find \((M^4,g)\); dust and radiation are the examples.}
\GSfield{Statement}{Take a future-directed unit timelike velocity \(g(u,u)=1\), an energy density \(\varepsilon\), and an equation of state \(p=p(\varepsilon)\). Couple this perfect-fluid stress to \(G=\kappa T\) and \(\nabla_iT^{ij}=0\). The Cauchy problem additionally requires data on a spacelike hypersurface satisfying the Einstein constraints. The cases \(p=0\) and \(p=\varepsilon/3\) are dust and isotropic radiation in three spatial dimensions.}
\GSfield{Reconstruction}{The fluid fields are part of the unknowns; specifying only a manifold and a metric omits them. No particular initial or boundary data are recorded, so this is a model problem rather than a definite new global existence conjecture.}
\GSfield{Fate}{Local Cauchy theory for dust and perfect fluids was established before this seminar; \citet[Introduction, Sections 4--6 and 9--12]{fouresbruhat1958fluids}. The precise differentiability and hyperbolicity assumptions belong to the chosen local theorem.}
\GSfield{Present status}{As of November 2026, the reconstructed statement is standard.}
\GSfield{Assessment}{literal statement: correct with missing hypotheses; publication match: partial; verification: read the result; development: annotated; confidence: high}
\end{GSitem}
\begin{GSitem}{1975-394.10}{E}{Conservation of stress gives relativistic Euler equations}
\GSfield{Kind; source}{claim; \GSpage{P000657}}
\GSfield{The notes}{The notes underline \(\nabla_kT_i{}^k=0\) and identify it with the laws of hydrodynamics.}
\GSfield{Statement}{For a smooth perfect fluid, contraction with \(u\) and orthogonal projection of \(\nabla_iT^{ij}=0\) give
\[u(\varepsilon)+(\varepsilon+p)\nabla_i u^i=0,\qquad
(\varepsilon+p)\nabla_u u^j=(g^{ji}-u^ju^i)\partial_i p.\]
For dust of positive density, the second equation says that its flow lines are timelike geodesics. On a solution of Einstein's equation, stress conservation follows from the contracted Bianchi identity.}
\GSfield{Reconstruction}{\ours\ These formulas follow by differentiating the stress tensor and using \(u_i\nabla_j u^i=0\). The signs are for the source's \((+---)\) convention.}
\GSfield{Fate}{The conservation reduction and dust equations are written in \citet[Section 4; perfect-fluid analysis in Section 9]{fouresbruhat1958fluids}. The calculation is already classical in 1976.}
\GSfield{Present status}{As of November 2026, the reconstructed statement is standard.}
\GSfield{Assessment}{literal statement: correct with missing hypotheses; publication match: exact; verification: read the result; development: annotated; confidence: high}
\end{GSitem}
\begin{GSitem}{1975-394.11}{E}{Schwarzschild geometry and its maximal extension}
\GSfield{Kind; source}{construction; \GSpage{P000657}}
\GSfield{The notes}{A point mass is schematically written as a delta source; \(SO(3)\) symmetry leads to Schwarzschild, and Kruskal is credited with specifying the whole manifold in 1960.}
\GSfield{Statement}{The Schwarzschild vacuum metric away from the central curvature singularity has
\[ds^2=(1-r_s/r)c^2dt^2-(1-r_s/r)^{-1}dr^2-r^2d\Omega^2,\qquad r_s=2GM/c^2.\]
The singularity of these coordinates at \(r=r_s\) is removable in a larger Lorentzian manifold. A delta-function right-hand side for the nonlinear Einstein tensor is not defined merely by writing a tensor-valued delta; it requires an additional distributional formulation.}
\GSfield{Reconstruction}{The source conflates an exterior vacuum geometry with a point-supported source description. The historical assertion of first specifying a manifold is narrowed to the identified maximal extension, without claiming priority over every earlier extension.}
\GSfield{Fate}{\citet[abstract]{kruskal1960maximal}, received 21 December 1959 and published 1 September 1960, explicitly removes the spherical coordinate singularity and exhibits the maximal extension. Publisher metadata and abstract were verified; the original proof was not accessible in full.}
\GSfield{Present status}{As of November 2026, the reconstructed statement is standard.}
\GSfield{Assessment}{literal statement: correct with missing hypotheses; publication match: partial; verification: publisher metadata checked; development: annotated; confidence: high}
\end{GSitem}
\begin{GSitem}{1975-394.12}{E}{The areal radius has a saddle at the bifurcation sphere}
\GSfield{Kind; source}{claim; \GSpage{P000657}}
\GSfield{The notes}{The notes say that the radial coordinate has saddles at the gravitational radii.}
\GSfield{Statement}{In a Kruskal chart normalized by \(UV=(1-r/r_s)e^{r/r_s}\), the bifurcation sphere is \(U=V=0\), \(r=r_s\). Near it,
\[r-r_s=-\frac{r_s}{e}UV+O((UV)^2).\]
Thus \(r\) has a saddle in the two transverse directions, with a critical two-sphere. At a generic horizon point with only one of \(U,V\) zero, \(dr\ne0\); the whole horizon is not a critical set.}
\GSfield{Reconstruction}{\ours\ The expansion is obtained by differentiating the displayed implicit coordinate relation. It supplies a precise local meaning for the notebook's plural ``saddles''. The areal radius is not a smooth coordinate at the bifurcation sphere.}
\GSfield{Fate}{This is a local calculation in the standard maximal extension, not a new global existence assertion.}
\GSfield{Present status}{As of November 2026, the reconstructed statement is standard.}
\GSfield{Assessment}{literal statement: correct with missing hypotheses; publication match: partial; verification: publisher metadata checked; development: annotated; confidence: high}
\end{GSitem}

\GStalk{395}{1976-04-19}{I. M. Gelfand, after L. N. Vaserstein; V. G. Drinfeld}{P000656-P000658}{Serre’s problem and Suslin’s resultant method}
Gelfand, following Vaserstein's notes, asks the projective-module and unimodular-row forms of Serre's problem. Drinfeld gives the monic-column reduction and Suslin's resultant translation lemma. Suslin's and Quillen's manuscripts were received before this meeting and published in 1976; no earlier public-circulation date is inferred from receipt alone. This April notebook is itself a witness to the oral transmission of Suslin's method. The earlier three-variable problem GP-1973-33 recurs in arbitrary dimension.

\begin{GSitem}{1975-395.1}{B}{Serre’s problem over a polynomial ring}
\GSfield{Kind; source}{question; \GSpage{P000657}}
\GSfield{The notes}{The question is whether every projective \(\mathbb C[x_1,\ldots,x_n]\)-module is free. The word projective is added in a callout.}
\GSfield{Statement}{Every finitely generated projective module over \(k[x_1,\ldots,x_n]\), for a field \(k\), is free. The finite-generation hypothesis is the standard form of Serre's problem. The source asks over \(\mathbb C\), while the solution applies to every field.}
\GSfield{Fate}{The affirmative solutions are \citet[Section III, Theorems 2--3]{suslin1976projective} and \citet{quillen1976projective}. Suslin's manuscript was received 26 February and presented12 March 1976; Quillen's was received 1 March. These are submission/presentation dates, not verified public-circulation dates. Both works were published in 1976. Suslin's theorem was read; Quillen's publisher chronology was checked.}
\GSfield{Present status}{As of November 2026, the stated finite-generation problem is solved. The earlier three-variable occurrence is subsumed by this theorem.}
\GSfield{Assessment}{literal statement: correct with missing hypotheses; publication match: exact; verification: read the result; development: annotated; confidence: high; problem \hyperref[prob:GP-1973-33]{GP-1973-33}}
\end{GSitem}
\begin{GSitem}{1975-395.2}{B}{Unimodular rows and determinant-one completion}
\GSfield{Kind; source}{definition; \GSpage{P000657}}
\GSfield{The notes}{A row \((a_1,\ldots,a_k)\) is unimodular if \(\sum_i p_i a_i=1\). The question asks for a matrix of determinant one with this first row.}
\GSfield{Statement}{Let \(A=K[x_1,\ldots,x_n]\), with \(K\) a field, and let the row length be \(r\ge2\). A unimodular row gives a split surjection \(A^r\to A\), whose kernel is finitely generated projective. Its freeness gives an invertible matrix completing the row; rescaling one other row makes the determinant one. For row length one, a unit is the only entry, and determinant one requires that unit to equal one.}
\GSfield{Reconstruction}{The split kernel explains the completion. Completing all unimodular rows establishes freeness of stably free modules; equivalence with the entire projective-module problem additionally uses stable freeness of projectives over a polynomial ring over a field.}
\GSfield{Fate}{The completion follows from \citet[Section III, Theorems 2--3]{suslin1976projective}; the determinant adjustment is elementary. The exceptional row length one is omitted in the notebook.}
\GSfield{Present status}{As of November 2026, the reconstructed statement is standard.}
\GSfield{Assessment}{literal statement: correct with missing hypotheses; publication match: partial; verification: read the result; development: annotated; confidence: high; problem \hyperref[prob:GP-1973-33]{GP-1973-33}}
\end{GSitem}
\begin{GSitem}{1975-395.3}{B}{A monic unimodular column is equivalent to its constant value}
\GSfield{Kind; source}{theorem; \GSpage{P000657}; V. G. Drinfeld}
\GSfield{The notes}{Drinfeld takes a commutative coefficient ring \(B\), a unimodular column \(f\in B[x]^n\), and assumes that \(f_1\) has leading coefficient one. He states \(f(x)\sim f(0)\).}
\GSfield{Statement}{For a commutative ring with identity, a unimodular polynomial column with a monic component lies in the same \(\mathrm{GL}_n(B[x])\)-orbit as its value at zero. For \(n\ge3\), a resultant patching argument proves this; unimodular pairs are directly completable, and the length-one case is immediate. Evaluation at zero preserves unimodularity.}
\GSfield{Reconstruction}{The relation means multiplication by an invertible polynomial matrix. The coefficient ring need not be a field. A monic component is the crucial additional hypothesis, not an assertion that arbitrary projective modules over arbitrary rings are extended.}
\GSfield{Fate}{The general-ring resultant input is \citet[Theorem 3.1 and Lemma 4.1]{monceuryengui2017suslin}, tracing Suslin's 1977 Lemmas 2.1 and 2.3. The notebook is an earlier witness to this method; the modern algorithm's simplifying infinite-field hypothesis is not imposed on the general existence lemma.}
\GSfield{Present status}{As of November 2026, the reconstructed statement is standard.}
\GSfield{Assessment}{literal statement: correct with missing hypotheses; publication match: exact; verification: read the result; development: annotated; confidence: high}
\end{GSitem}
\begin{GSitem}{1975-395.4}{E}{The resultant and its leading coefficients}
\GSfield{Kind; source}{definition; \GSpage{P000657}; V. G. Drinfeld}
\GSfield{The notes}{The resultant is introduced as the product of differences of the roots of two univariate polynomials.}
\GSfield{Statement}{For \(f=a_m\prod_{i=1}^m(x-\xi_i)\) and \(g=b_n\prod_{j=1}^n(x-\eta_j)\) over a splitting field,
\[\operatorname{Res}(f,g)=a_m^n b_n^m\prod_{i,j}(\xi_i-\eta_j).\]
The Sylvester determinant defines it over any commutative coefficient ring and gives a polynomial identity \(\operatorname{Res}(f,g)=uf+vg\) for suitable polynomials \(u,v\). The unweighted root product is correct when both polynomials are monic.}
\GSfield{Reconstruction}{The leading coefficients are restored, and the determinant definition is needed because arbitrary rings need not admit root factorizations.}
\GSfield{Fate}{The resultant and its Bézout identity are the algebraic mechanism in \citet[Lemma 4.1]{monceuryengui2017suslin}. The notebook's abbreviated root formula is insufficient for the arbitrary-ring argument.}
\GSfield{Present status}{As of November 2026, the reconstructed statement is standard.}
\GSfield{Assessment}{literal statement: false as written; publication match: partial; verification: read the result; development: annotated; confidence: high}
\end{GSitem}
\begin{GSitem}{1975-395.5}{B}{Suslin’s translation by a resultant}
\GSfield{Kind; source}{theorem; \GSpage{P000658}; V. G. Drinfeld}
\GSfield{The notes}{With \(c=\operatorname{Res}(f_1,f_2)\), the main lemma says \(f(Y_1+cY_2)\sim f(Y_1)\) for any polynomial column. The proof is credited to an explicit construction by Suslin.}
\GSfield{Statement}{Let \(B\) be a commutative ring with identity and \(f\in B[x]^n\), \(n\ge2\). With \(c=\operatorname{Res}(f_1,f_2)\), there is an invertible matrix over \(B[Y_1,Y_2]\) carrying \(f(Y_1)\) to \(f(Y_1+cY_2)\). No unimodularity assumption on this column is required for the lemma itself. The two-component construction can be chosen in \(\mathrm{SL}_2\); the remaining components are adjusted using the Bézout expression for \(c\).}
\GSfield{Reconstruction}{The quantifiers and coefficient ring are made explicit. The lemma is not the assertion that an arbitrary resultant is a unit.}
\GSfield{Fate}{The two-component matrix and identity are read in \citet[Lemma 4.1, p.117]{monceuryengui2017suslin}, which identifies this as Suslin's 1977 Lemma 2.1. The seminar attributes the construction to Suslin already in April 1976.}
\GSfield{Present status}{As of November 2026, the reconstructed statement is standard.}
\GSfield{Assessment}{literal statement: correct with missing hypotheses; publication match: exact; verification: read the result; development: annotated; confidence: high}
\end{GSitem}
\begin{GSitem}{1975-395.6}{B}{Patching translations and eliminating variables}
\GSfield{Kind; source}{method; \GSpage{P000658}; V. G. Drinfeld}
\GSfield{The notes}{The notes say one needs \(c=1\), then puts \(Y_2=-Y_1\); a change of variables makes a component monic, and variables are removed one at a time.}
\GSfield{Statement}{The missing step is to obtain several resultants \(c_i\) generating the unit ideal, after allowed operations on a monic unimodular column. Choose \(d_i\) with \(\sum_i d_i c_i=1\); successive translations by \(d_i c_iY\) combine to a translation by \(Y\). Specializing \(Y=-x\) gives equivalence to the value at zero. Over a field, a triangular change of polynomial variables can make a nonzero component monic in the last variable, after rescaling its leading coefficient, and this permits induction on the number of variables.}
\GSfield{Reconstruction}{\ours\ The composition of translations explains the unit-ideal patching. It does not assume the original pair has resultant one. The monic-changing assertion is stated over a field, where the required nonzero leading coefficient is invertible.}
\GSfield{Fate}{The needed resultant ideal is \citet[Theorem 3.1]{monceuryengui2017suslin}; Section 4 supplies a constructive version. The conclusion belongs to the 1976 solution of Serre's problem, not to an unresolved question left by the missing lines.}
\GSfield{Present status}{As of November 2026, the reconstructed statement is standard.}
\GSfield{Assessment}{literal statement: correct with missing hypotheses; publication match: partial; verification: read the result; development: annotated; confidence: high}
\end{GSitem}

\GStalk{396}{1976-04-19}{S. P. Novikov}{P000657-P000658}{Relativistic fluids, coordinate freedom and Gödel’s metric}
The final record resumes the Einstein--fluid equations, gives the sound-speed inequalities, explains the apparent equation deficit by coordinate freedom, and mentions Gödel's closed timelike curves. The four-entry signature fixes the intended spacetime dimension despite the notation M-n. The chapter boundary follows immediately after this discussion.

\begin{GSitem}{1975-396.1}{E}{Perfect-fluid data and normalization of velocity}
\GSfield{Kind; source}{definition; \GSpage{P000658}}
\GSfield{The notes}{The page writes \(M^n\) with signature \((1,-1,-1,-1)\), repeats the perfect-fluid stress and explicitly imposes \(\langle u,u\rangle=1\).}
\GSfield{Statement}{The intended setting is \(M^4\) with a Lorentz metric, a unit future-directed timelike velocity, energy density and pressure satisfying an equation of state. In units \(c=1\), the unknown fields are constrained by \(G=T\), \(T_{ij}=(\varepsilon+p)u_i u_j-pg_{ij}\), and stress conservation. The source's unit normalization supplies a hypothesis left implicit in record394.}
\GSfield{Fate}{This is the pre-existing relativistic perfect-fluid system, with the convention found in \citet[Section 9]{fouresbruhat1958fluids}. The repeated spherical-symmetry aside refers back to 1975-394.11--12.}
\GSfield{Present status}{As of November 2026, the reconstructed statement is standard.}
\GSfield{Assessment}{literal statement: false as written; publication match: exact; verification: read the result; development: annotated; confidence: high}
\end{GSitem}
\begin{GSitem}{1975-396.2}{E}{Sound speed and relativistic causality}
\GSfield{Kind; source}{claim; \GSpage{P000658}}
\GSfield{The notes}{The sound speed is \(c_s^2=\partial p/\partial\varepsilon\ge0\), with \(c_s^2\le c^2=1\). Dust has \(p=0\).}
\GSfield{Statement}{For a barotropic perfect fluid in units \(c=1\), linear acoustic disturbances in the rest frame have \(c_s^2=dp/d\varepsilon\). Nonnegative squared speed is the elementary local stability requirement; \(c_s\le1\) keeps the acoustic cone inside the light cone. Strict hyperbolicity assumptions often use \(0<c_s^2<1\) and \(\varepsilon+p>0\). The zero-speed dust case requires its own formulation, so the weak inequalities alone are not a complete well-posedness theorem.}
\GSfield{Reconstruction}{\ours\ Linearizing the energy and momentum equations about constant fluid at rest gives \(\partial_t^2\delta\varepsilon-(dp/d\varepsilon)\Delta\delta\varepsilon=0\), which fixes the sound-speed formula.}
\GSfield{Fate}{The fluid characteristic cone is identified in \citet[Section 11, p.168]{fouresbruhat1958fluids}. The source states physical inequalities, not all analytical assumptions for the coupled Cauchy problem.}
\GSfield{Present status}{As of November 2026, the reconstructed statement is standard.}
\GSfield{Assessment}{literal statement: correct with missing hypotheses; publication match: exact; verification: read the result; development: annotated; confidence: high}
\end{GSitem}
\begin{GSitem}{1975-396.3}{E}{Coordinate freedom, constraints and hyperbolic reduction}
\GSfield{Kind; source}{method; \GSpage{P000658}}
\GSfield{The notes}{The notes say there are too few equations, but gauge transformations compensate; a historical aside says Hilbert rejected the equations as incorrect.}
\GSfield{Statement}{Four coordinate freedoms and four contracted Bianchi identities prevent treating the ten Einstein equations as ten independent scalar evolution equations in an arbitrary chart. Harmonic coordinates give a quasilinear hyperbolic reduced system. If the initial Einstein constraints and harmonic gauge conditions hold, stress conservation propagates the gauge constraints, and the reduced solution solves the original equations. Geometrical uniqueness is understood up to diffeomorphism.}
\GSfield{Reconstruction}{This is the precise local PDE interpretation of the equation count. It does not substantiate the anecdote about Hilbert, whose wording is retained only as source provenance.}
\GSfield{Fate}{The smooth local Cauchy problem was established by \citet[Introduction]{fouresbruhat1952cauchy}. The constraint-propagation calculation is explicit in \citet[Section 2]{fouresbruhat1958fluids}. Both substantially predate the meeting.}
\GSfield{Present status}{As of November 2026, the reconstructed statement is standard.}
\GSfield{Assessment}{literal statement: correct with missing hypotheses; publication match: partial; verification: read the result; development: annotated; confidence: high}
\end{GSitem}
\begin{GSitem}{1975-396.4}{E}{Gödel’s spacetime and closed timelike curves}
\GSfield{Kind; source}{construction; \GSpage{P000658}}
\GSfield{The notes}{Novikov mentions a Gödel metric with a closed timelike line, without displaying a metric or matter model.}
\GSfield{Statement}{Gödel's homogeneous rotating cosmological spacetime solves Einstein's equation with pressureless dust and a nonzero cosmological constant. It contains closed timelike curves. In the usual cylindrical coordinates of the original construction, circles at sufficiently large radius are timelike. The cosmological constant must be included, or absorbed into an effective stress tensor, to compare this example with the preceding equation lacking an explicit cosmological term.}
\GSfield{Fate}{The metric, dust source and cosmological term are given in \citet[Sections 2--3]{godel1949cosmological}, published 1 July 1949; Section 3 exhibits the timelike circles. The original four-page article was read.}
\GSfield{Present status}{As of November 2026, the reconstructed statement is standard.}
\GSfield{Assessment}{literal statement: correct with missing hypotheses; publication match: exact; verification: read the result; development: annotated; confidence: high}
\end{GSitem}

\section{Developments}\label{dev:1975}
\subsection{A certified finite development of the Fuchs--Perchik question}
\label{dev:perchik}

\paragraph{The question supported by the notebook.}
On \GSpage{P000655}, Fuchs asks whether nonzero Euler characteristics occur
for arbitrarily large \(m\). The product and the expansion on
\GSpage{P000655}--\GSpage{P000656} make this a definite coefficient question,
not just a request to compute a Lie algebra's cohomology. We write \(E(m)\)
for the coefficients to avoid the notebook's switching between
\(\chi_m\) and \(\chi_m^*\). Put
\[
 S_\ell=\{-\ell-2,-\ell,\ldots,\ell+2\},\qquad
 P(t,x)=\prod_{\ell=1}^{\infty}\prod_{s\in S_\ell}(1-t^\ell x^s),
\]
and define the formal Laurent series
\begin{equation}\label{eq:perchikseries}
 F(t)=\frac12\operatorname{CT}_x
 \bigl[(1-t^{-1}x^{-1})(1-t^{-1}x)
       (1-x^{-2})(1-x^2)P(t,x)\bigr]
 =\sum_{m\geq-1}E(m)t^{2m}.
\end{equation}
Here \(\operatorname{CT}_x\) extracts the coefficient of \(x^0\);
the product is expanded first in nonnegative powers of \(t\) in \(P\),
then multiplied by the displayed finite prefactor. This prescription,
including the two negative-weight factors, is essential.

\begin{question}[The notebook question, precisely normalized]
For every integer \(M\geq0\), is there an integer \(m>M\) with
\(E(m)\ne0\)? Equivalently, is \(F(t)\) not a Laurent polynomial?
\end{question}

\paragraph{The modern target and its limits.}
Let \(\mathfrak h_2\) be the topological complex Lie algebra of formal
Hamiltonian vector fields on the symplectic plane, represented by
\(\mathbb C[[p,q]]/\mathbb C\) with the Poisson bracket. Its quadratic
Hamiltonians form \(\mathfrak{sl}_2\). Use continuous cochains, trivial
coefficients, and input weight \(d-2\) for a Hamiltonian of polynomial
degree \(d\). Then the modern interpretation is
\[
 E(m)=\sum_q(-1)^q
 \dim H^q(\mathfrak h_2,\mathfrak{sl}_2)_{2m}.
\]
The target adopted here is infinite nonvanishing of these \emph{relative}
Euler characteristics at fixed symplectic dimension two. Infinite
nonvanishing of cohomology alone would be weaker: even and odd Betti
numbers might cancel. Passing to a stable symplectic rank, deleting the
constant vector fields, or retaining only positive-degree Hamiltonians
changes the question. No quantitative asymptotic, eventual nonvanishing,
or sign law is asserted.

\paragraph{What the literature establishes.}
The coefficient formula is Perchik's theorem, published in 1976
\citep{perchik1976cohomology}. Harvard's dissertation list dates his
dissertation of the same title to 1975 \citep{perchik1975dissertation}.
That establishes earlier work,
not a precise public-circulation date; the seminar records its discussion
on 12 April 1976. \citet[Theorem~2.1 and the table on Russian pp.~68--69]
{mohammadzadehfuchs2014cohomology} restate the formula, compute through
\(m=129\), and explicitly retain the infinite-nonvanishing question.
Their absolute/relative cohomology comparison excludes weight zero.
As of 10 November 2026, the targeted searches documented with this chapter
located no later resolution of this exact fixed-rank coefficient question.
This is a bounded literature finding, not a proof that no resolution exists.

\paragraph{Why a finite development is useful.}
The notebook's missing star would turn a nontrivial question into an
identically zero absolute Euler-characteristic calculation. Its general
product row also has an uncertain middle exponent. A reproducible expansion
therefore tests a mathematically decisive normalization. The contribution
here is a cutoff proof, two independent exact algorithms, an audited finite
extension of the notebook's table, and explicit bounds on what those data
prove. These are proofs and computations supplied for this companion;
no priority claim is made for the elementary lemmas or numerical values.

\begin{lemma}[Well-definedness, integrality, and parity]\ours\quad
For every integer \(N\geq0\), all coefficients of \(F\) through
\(t^N\) are determined by retaining \(P\) through degree \(N+2\)
and including only factors with \(1\leq\ell\leq N+2\).
The series \(F\) has integer coefficients, no powers below \(-2\),
and no odd powers.
\end{lemma}
\begin{proof}
Each factor in \(P\) has positive \(t\)-degree in its nonconstant term.
Thus a coefficient through degree \(L\) uses only factors with
\(\ell\leq L\); these are finitely many, each with finitely many terms.
The finite prefactor has lowest \(t\)-degree \(-2\), so \(L=N+2\)
suffices. Terms discarded beyond \(L\) cannot return to a degree at most
\(N\).

Set \(R=(1-t^{-1}x^{-1})(1-t^{-1}x)P\). Each coefficient of \(R\)
is an integer Laurent polynomial invariant under \(x\mapsto x^{-1}\).
Writing such a coefficient as \(Q(x)=\sum_k a_kx^k\), with
\(a_{-k}=a_k\), gives
\[
 \tfrac12\operatorname{CT}_x[(1-x^{-2})(1-x^2)Q]
   =\tfrac12(2a_0-a_2-a_{-2})=a_0-a_2\in\mathbb Z.
\]
Every monomial \(t^n x^k\) in \(P\) satisfies \(n\equiv k\pmod2\),
because \(s\equiv\ell\pmod2\) in each factor. The negative-weight
factors and the two Weyl factors preserve this congruence. Taking the
constant term therefore kills every odd \(t\)-degree. Finally, the only
negative \(t\)-degrees come from the two factors of degree \(-1\).
\end{proof}

\begin{remark}[Why this is the relative Euler series]\ours\quad
This interpretation can be checked independently at the level of characters.
As an \(\mathfrak{sl}_2\)-module, the weight-\(\ell\) Hamiltonians are
\(\operatorname{Sym}^{\ell+2}\mathbb C^2\), whose weights are \(S_\ell\).
The relative exterior complex removes the entire quadratic component and
retains \(\operatorname{Sym}^{1}\mathbb C^2\) in weight \(-1\) and
the components in positive weights. Its alternating character is \(R\).
For an irreducible character
\(x^d+x^{d-2}+\cdots+x^{-d}\), the difference between its coefficients
of \(x^0\) and \(x^2\) is one for \(d=0\) and zero otherwise.
Thus the operation in the proof selects the invariant multiplicity.
The chain groups at a fixed input weight are finite: at most two
negative-weight generators can occur, and the number and weights of
positive-weight generators are bounded by the total weight plus two.
Euler--Poincare therefore applies weight by weight. This derivation also
explains why an unrestricted exterior product containing the zero-weight
Cartan direction would give the factor \(1-1\), not the desired series.
\end{remark}

\begin{proposition}[Two exact coefficient algorithms]\ours\quad
Write \(P(t,x)=\sum_{n\geq0}P_n(x)t^n\), with \(P_0=1\), and put
\[
 B_r(x)=-\sum_{d\mid r}d\sum_{s\in S_d}x^{sr/d}.
\]
Then, for \(n\geq1\),
\begin{equation}\label{eq:perchikrecurrence}
 nP_n(x)=\sum_{r=1}^n B_r(x)P_{n-r}(x).
\end{equation}
Independently, multiplying each factor \(1-t^\ell x^s\), with descending
updates in \(n\), computes the same coefficients using integer additions
and subtractions only. Both methods terminate at any prescribed cutoff.
\end{proposition}
\begin{proof}
The positive product has constant term one. Its formal logarithmic
derivative in \(\mathbb Q[x,x^{-1}][[t]]\) is
\[
 \frac{t\partial_t P}{P}
 =-\sum_{\ell\geq1}\sum_{s\in S_\ell}
       \frac{\ell t^\ell x^s}{1-t^\ell x^s}
 =-\sum_{\ell\geq1}\sum_{s\in S_\ell}\sum_{j\geq1}
       \ell t^{\ell j}x^{sj}
 =\sum_{r\geq1}B_r(x)t^r.
\]
All coefficient rearrangements are finite at a fixed \(t\)-degree.
Multiplying by \(P\) and comparing the coefficient of \(t^n\) proves
\eqref{eq:perchikrecurrence}; division by \(n\) determines \(P_n\)
uniquely in characteristic zero. For the other algorithm, if \(D_{n,k}\)
is the coefficient before one factor is applied, the new coefficient is
\(D_{n,k}-D_{n-\ell,k-s}\). Descending \(n\) ensures that the second
term still belongs to the old product. Induction over factors proves the
algorithm. The cutoff lemma proves termination and prevents omitted
factors from changing any requested coefficient. Integrality of all
divisions in the recurrence follows also from its agreement with the
integer product; the implementation checks every remainder directly.
\end{proof}

\paragraph{Independent checks of the argument.}
The cutoff was re-derived from exterior degrees: the negative component
has dimension two, hence can lower a positive weight by at most two.
The invariant-multiplicity argument above independently checks the
factor \(1/2\) and both Weyl factors. The recurrence was checked by
direct factor multiplication, rather than by another implementation of
the same recurrence. In particular, agreement is checked for every
Laurent coefficient of every \(P_n\), not just for their final constant
terms. These checks found no mathematical discrepancy between the two
algorithms.

\begin{proposition}[Finite output through \(m=150\)]\oursComputed\quad
The series in \eqref{eq:perchikseries}, computed with exact integers by
both algorithms through positive-product degree \(302\), has
\(135\) nonzero coefficients among \(E(1),\ldots,E(150)\), and
\[
 \sum_{m=1}^{150}|E(m)|=9\,701\,705\,312.
\]
In particular,
\[
 E(35)=10,\quad E(86)=-22\,044,\quad
 E(129)=7\,151\,485,\quad E(150)=1\,408\,548\,898.
\]
\end{proposition}
\begin{proof}
The executable certificate is \texttt{code/perchik\_series.py}; the complete
coefficient table and verification output are in
\texttt{computations/m150/}. The two algorithms agree for all Laurent
coefficients through \(n=302\), including \(137\,561\) checked integer
divisions in the recurrence. Reflection, parity, and the support bound
\(|k|\leq3n\) are checked for every nonzero coefficient. Separate runs
at cutoffs \(82\) and \(182\) produce identical overlapping output.
The notebook's displayed terms, including \(10t^{70}\), are recovered.
The companion verification script records the prefix comparisons and
the sums in the statement.
\end{proof}

\begin{center}
\begin{tabular}{rr@{\qquad}rr@{\qquad}rr}
\toprule
\(m\)&\(E(m)\)&\(m\)&\(E(m)\)&\(m\)&\(E(m)\)\\
\midrule
130&41831534&137&-67794447&144&-30506945\\
131&79675254&138&-188663902&145&321390415\\
132&113436381&139&-311827536&146&732255863\\
133&133851464&140&-412439649&147&1137936570\\
134&131414489&141&-459062275&148&1450364576\\
135&98557132&142&-423731634&149&1571713579\\
136&31149174&143&-282894462&150&1408548898\\
\bottomrule
\end{tabular}
\end{center}
This table is generated by the supplied program; it is not a transcription
of a literature table. It extends the notebook's range, not a claim about
the largest range ever computed.

\begin{corollary}[Exactly what nonzero Euler characteristics prove]\ours\quad
For the relative continuous cohomology above,
\[
 \sum_{m=1}^{150}\sum_q
    \dim H^q(\mathfrak h_2,\mathfrak{sl}_2)_{2m}
       \geq9\,701\,705\,312.
\]
At least \(135\) distinct positive weights in this range carry nonzero
relative cohomology. Also
\(\sum_{m=-1}^{35}|E(m)|=57\), recovering the notebook's finite
lower-bound count with the relative interpretation.
\end{corollary}
\begin{proof}
For each finite weight complex, Euler--Poincare and the triangle
inequality give \(|E(m)|\leq\sum_q\dim H^q_{2m}\). Sum this inequality
over the stated ranges. Distinct weights give distinct graded summands.
No inference is made at a weight whose Euler characteristic vanishes.
\end{proof}

\paragraph{The remaining obstruction and the next bounded step.}
The equality of two exact algorithms certifies a finite prefix. It cannot
exclude a later vanishing tail, prove infinitely many sign changes,
recover the degrees of the cohomology classes, or identify representatives.
The obstacle is cancellation in a two-variable exterior product after
projection to invariants; the large unsigned chain dimensions do not
provide a lower bound on its signed constant term. The checked literature
shows continuing computation rather than abandonment; it does not
establish a sociological explanation for the absence of a general proof.
A further bounded run could compute a longer independently certified
prefix or construct and reduce the relative differential at one selected
weight with \(E(m)=0\), thereby detecting cancellations invisible to Euler
characteristics. A resolution of GP-1975-20 needs an argument excluding
eventual cancellation for unbounded \(m\); extending the table alone
does not supply that argument.

\subsection{What the Abel spectral-phase limit requires}
\label{dev:spectralphases}

\paragraph{The recorded question and a defensible target.}
Occurrence 1975-391.7 and problem GP-1975-24 concern the expression on
\GSpage{P000652}. Let \(M\) be a closed compact Riemannian manifold,
let \(\lambda_k\geq0\) be the eigenvalues of \(\sqrt{-\Delta}\), counted
with multiplicity, and fix \(a>0\). With the periodic delta normalized to
have mass one on \([0,2\pi)\), set
\[
 \Phi_{a,\varepsilon}=\varepsilon\sum_k e^{-\varepsilon\lambda_k}
 \left(\delta_{\rm per}(\mu-a\lambda_k)-\frac1{2\pi}\right),
 \qquad \varepsilon>0.
\]
The notebook proposes a limit and relates its nonzero Fourier coefficients
to resonances with lengths of closed geodesics. It also specifies a generic
length arrangement. To make that phrase coherent, distinctness and
incommensurability must concern primitive unoriented geodesics; iterates
necessarily have commensurable lengths. Isolated, nondegenerate iterates
are the modern hypotheses needed for the simple leading wave-trace
singularity. They are not silently imposed on every manifold in the first
line of the notes. Analyticity, as recorded there, can be retained but
does not itself prove the limiting assertion.

The strongest target justified here is conditional: identify every fixed
Fourier coefficient by a local trace calculation, then prove a bound
uniform in both the Fourier index and the Abel parameter. The latter
step cannot be replaced by pointwise convergence of coefficients. The
following development isolates this obstruction, solves the complete
circle model, and checks what fails when the generic-geodesic restriction
is dropped. No general theorem under the notebook's generic hypotheses
is claimed.

\paragraph{Literature and identification limit.}
The leading singularity at a nondegenerate orbit is part of the
wave-trace theory already published in 1975
\citep{duistermaatguillemin1975spectrum}.
The first expansion in \citet[Introduction]{zelditch1998wave} explicitly
has a simple-pole leading term and logarithmic lower terms. That local
expansion supports the fixed-mode mechanism; it does not, by itself,
give a bound uniform over all sampled times \(ma\).
The latter paper appeared in 1998 and was deposited on arXiv in 2000;
the deposit is not its first-publication date.
As of 10 November 2026, the searches recorded for this chapter did not
identify a theorem establishing the entire notebook limit in the stated
generality. This is an unresolved match and analytic gap in the argument
presented here, not a claim that all spectral-phase limit problems remain
open in the literature.

\begin{lemma}[A necessary and sufficient Fourier test]\ours\quad
For \(T_\varepsilon\in\mathcal D'(\mathbb R/2\pi\mathbb Z)\),
\(0<\varepsilon\leq1\), write
\(\widehat T_\varepsilon(m)=(2\pi)^{-1}
\langle T_\varepsilon,e^{-im\mu}\rangle\).
Suppose the family is bounded on each compact parameter interval
away from zero. It converges in distributions as \(\varepsilon\downarrow0\)
if and only if every coefficient converges and there are
\(C,N,\varepsilon_0>0\) such that
\[
 |\widehat T_\varepsilon(m)|\leq C(1+|m|)^N
 \quad(m\in\mathbb Z,\ 0<\varepsilon\leq\varepsilon_0).
\]
For the notebook family, \(\widehat\Phi_{a,\varepsilon}(0)=0\), and
\[
 \widehat\Phi_{a,\varepsilon}(m)
 =\frac{\varepsilon}{2\pi}
    \sum_k e^{-(\varepsilon+ima)\lambda_k},\qquad m\ne0.
\]
\end{lemma}
\begin{proof}
The last formula follows by pairing each delta with the indicated
exponential; the subtracted constant removes only the zero mode.
For necessity, distributional convergence gives pointwise boundedness
on the Frechet space of smooth periodic test functions. The uniform
boundedness principle gives equicontinuity: for some \(C,N\),
\(|\langle T_\varepsilon,f\rangle|\leq
C\max_{j\leq N}\|f^{(j)}\|_\infty\) for small \(\varepsilon\).
Testing exponentials gives the polynomial bound and their individual
convergence. For sufficiency, the limiting coefficients have the same
polynomial bound and so define a distribution. The Fourier coefficients
of a smooth test function decay faster than any power. Consequently
\[
 \langle T_\varepsilon,f\rangle
   =2\pi\sum_{m\in\mathbb Z}
       \widehat T_\varepsilon(m)\widehat f(-m)
\]
converges by dominated convergence. This proves the assertion and keeps
the normalization explicit.
\end{proof}

\begin{proposition}[The complete circle model]\ours\quad
Let \(M\) be the circle of circumference \(2\pi R\), \(R>0\).
If \(a/(2\pi R)\) is irrational, then
\(\Phi_{a,\varepsilon}\to0\) in distributions. If
\(a/(2\pi R)=p/q\) in lowest terms with \(p,q\) positive integers, then
\[
 \Phi_{a,\varepsilon}\longrightarrow
 2R\left(\frac1q\sum_{j=0}^{q-1}
           \delta_{\rm per}(\mu-2\pi j/q)-\frac1{2\pi}\right).
\]
Including or omitting the zero eigenvalue does not change the limit.
\end{proposition}
\begin{proof}
The eigenfunctions \(e^{ikx/R}\), \(k\in\mathbb Z\), give frequencies
\(|k|/R\). Thus, for \(m\ne0\), the geometric series gives
\[
 \widehat\Phi_{a,\varepsilon}(m)
 =\frac{\varepsilon}{2\pi}
 \left(1+\frac{2}{e^{(\varepsilon+ima)/R}-1}\right).
\]
Its limit is \(R/\pi\) when \(ma/R\in2\pi\mathbb Z\), and zero
otherwise. Moreover,
\[
 \left|\widehat\Phi_{a,\varepsilon}(m)\right|
 \leq\frac{\varepsilon}{2\pi}
       \left(1+\frac{2}{e^{\varepsilon/R}-1}\right)
 \leq\frac{\varepsilon+2R}{2\pi},
\]
uniformly in \(m\). The Fourier test therefore proves convergence.
In the rational case, the resonant indices are precisely the nonzero
multiples of \(q\), which are the coefficients of the displayed atomic
comb. In the irrational case no nonzero index is resonant. A zero
eigenvalue contributes only
\(\varepsilon(\delta_{\rm per}-1/(2\pi))\), which vanishes.
\end{proof}

\begin{example}[Why one cannot omit the geometric restriction]\ours\quad
On the unit round two-sphere, with \(a=2\pi\), the notebook family has
no distributional limit. More precisely,
\[
 \widehat\Phi_{2\pi,\varepsilon}(1)
     =-\frac1{\pi\varepsilon}+O(1).
\]
This example does \emph{not} refute the generic-geodesic assertion:
the round sphere has continuous families of closed geodesics of the
same length, and their return maps are degenerate.
\end{example}
\begin{proof}
The sphere frequencies are \(\lambda_\ell=\sqrt{\ell(\ell+1)}\), with
multiplicity \(2\ell+1\), \(\ell\geq0\). For completeness, the polar
coordinate identity
\(\Delta_{\mathbb R^3}=\partial_r^2+2r^{-1}\partial_r+r^{-2}\Delta_{S^2}\)
applied to a harmonic homogeneous polynomial \(r^\ell Y\) gives
\(-\Delta_{S^2}Y=\ell(\ell+1)Y\).
The decomposition of degree-\(\ell\) homogeneous polynomials into
harmonic polynomials plus \(r^2\) times degree-\(\ell-2\) polynomials
gives dimension
\(\binom{\ell+2}{2}-\binom{\ell}{2}=2\ell+1\).
Polynomial restrictions are dense on the sphere, so these harmonics
give the full spectrum.

Put \(\delta_\ell=\ell+\tfrac12-\lambda_\ell\).
For \(\ell\geq1\), rationalization gives
\[
 0<\delta_\ell=
 \frac{1/4}{\ell+1/2+\sqrt{\ell(\ell+1)}}
       \leq\frac1{4\ell}.
\]
Uniformly for \(0<\varepsilon\leq1\),
\[
 e^{-(\varepsilon+2\pi i)\lambda_\ell}
 =-e^{-\varepsilon(\ell+1/2)}
       \bigl(1+O(\ell^{-1})\bigr).
\]
After multiplying by \(2\ell+1\), summing, and then multiplying by
\(\varepsilon/(2\pi)\), the error is \(O(1)\).
The finitely many omitted low terms do not change this bound. The main
term is
\[
 -\frac{\varepsilon}{2\pi}e^{-\varepsilon/2}
       \sum_{\ell\geq0}(2\ell+1)e^{-\varepsilon\ell}
 =-\frac{\varepsilon}{2\pi}e^{-\varepsilon/2}
       \frac{1+e^{-\varepsilon}}{(1-e^{-\varepsilon})^2}
 =-\frac1{\pi\varepsilon}+O(1).
\]
A distributional limit would have a finite first Fourier coefficient,
which this estimate excludes.
\end{proof}

\paragraph{Checks, scope, and next step.}
The circle normalization was checked independently by summing over
residue classes \(k\bmod q\): each class has Abel mass tending to
\(R/q\), doubled by the two signs of nonzero \(k\). Its Fourier
transform agrees with the geometric-series calculation. The uniform
bound follows separately from the triangle inequality for the positive
Abel sum. For the sphere, rationalization gives a uniform remainder;
the closed sum of \((2\ell+1)e^{-\varepsilon\ell}\) verifies the sign
and the coefficient of the divergence. A separate numerical check,
\texttt{code/spectral\_phase\_check.py}, compares the circle formula
with direct summation and the sphere asymptotic with successively
smaller parameters; the output is in
\texttt{computations/spectral\_phase\_check.json}. Those floating-point
values are diagnostics, not the proof of convergence or divergence.

The remaining question in the generic higher-dimensional setting is
the uniform polynomial estimate in the Fourier test, together with
the exact hypotheses under which the local trace coefficients are
well defined and do not cancel. A useful further research run would
trace Shnirelman's exact publication and establish such a bound for
one specified geometric class and a specified set of admissible
\(a\). The present calculation neither proves that global estimate
nor labels the sphere counterexample as a failure under hypotheses
it does not satisfy.

\section{Problem index}\label{problems:1975}
Earlier-year occurrence identifiers are retained as text; current-year identifiers link to their entries.

\subsection{GP-1975-01: Unitary classification of essentially normal operators modulo compacts}\label{prob:GP-1975-01}
For bounded essentially normal \(T_1,T_2\) on separable infinite-dimensional \(H\), determine whether essential spectrum and complementary Fredholm-index function classify unitary equivalence modulo compact operators. The essential-normality hypothesis is reconstructed from the BDF attribution.
\par\noindent\textit{Occurrences:} \hyperref[item:1975-365.1]{1975-365.1}.
\par\noindent\textit{Status (2026-11):} As of 10 November 2026: answered affirmatively in the BDF setting before this talk; \citet[pp.~58--128]{browndouglasfillmore1973unitary}. The notebook's unrestricted formulation is not the theorem.

\subsection{GP-1975-02: Involution for elliptic Calogero--Moser systems}\label{prob:GP-1975-02}
For reduced crystallographic root systems with Weyl-orbit couplings, construct a generically independent complete set of commuting integrals for the elliptic Hamiltonian on its regular canonical phase space; retain the free center of mass in the particle \(A\)-model.
\par\noindent\textit{Occurrences:} \hyperref[item:1975-372.7]{1975-372.7}.
\par\noindent\textit{Status (2026-11):} As of 10 November 2026: affirmative in this scope; \citet[pp.~3--5]{khastgirsasaki2001liouville} is a directly checked proof, first publicly circulated 30 May 2000. This does not certify every arbitrary-\(\varepsilon\) or nonreduced variant in the shorthand table.

\subsection{GP-1975-03: Integrability of exceptional-root Calogero--Moser systems}\label{prob:GP-1975-03}
For \(G_2,F_4,E_6,E_7,E_8\), determine Liouville integrability of the rational, hyperbolic, trigonometric and untwisted elliptic root potentials with couplings constant on Weyl orbits, away from poles.
\par\noindent\textit{Occurrences:} \hyperref[item:1975-372.15]{1975-372.15}.
\par\noindent\textit{Status (2026-11):} As of 10 November 2026: affirmative in the stated scope. The group-coupling cases occur in \citet[Theorem~1 and Corollary~1]{olshanetskyperelomov1978radial}; \citet[pp.~3--5]{khastgirsasaki2001liouville} supplies the checked general involution argument. Harmonic confinement is not included in that elliptic theorem.

\subsection{GP-1975-04: Adequate inverse data for Calogero--Moser particle systems}\label{prob:GP-1975-04}
Identify data sufficient to recover a regular particle trajectory, with center-of-mass and permutation conventions fixed. The trace-only wording in the notebook is underdetermined; a concrete corrected rational problem uses the full initial matrices \(Q_0,L_0\).
\par\noindent\textit{Occurrences:} \hyperref[item:1975-372.11]{1975-372.11}.
\par\noindent\textit{Status (2026-11):} As of 10 November 2026: the rational initial-value reconstruction is explicit by diagonalizing \(Q_0+tL_0\), or its harmonic counterpart. These formulas are read in \citet[Sections~2--3]{sasaki2006explicit}. Traces alone still do not determine matrix entries or a phase point.

\subsection{GP-1975-05: Involution of the periodic Toda integrals}\label{prob:GP-1975-05}
For the periodic \(n\)-particle Toda Hamiltonian, exhibit \(n\) generically independent conserved functions that commute for the canonical Poisson bracket, keeping the free center-of-mass degree of freedom.
\par\noindent\textit{Occurrences:} \hyperref[item:1975-372.4]{1975-372.4}.
\par\noindent\textit{Status (2026-11):} As of 10 November 2026: established. \citet[pp.~1668--1669]{kawasaki1976involution}, published 1 May 1976, is a verified publisher record devoted to the claim. Its proof and the earliest-priority chronology were not independently read in this check.

\subsection{GP-1973-53: Recursive construction of KdV Lax operators}\label{prob:GP-1973-53}
Construct odd differential operators from formal fractional powers of the Schrödinger operator so that their commutators with it have order zero, and organize the resulting compatible higher KdV flows.
\par\noindent\textit{Occurrences:} 1973-326.1, 1973-327.1, \hyperref[item:1975-373.1]{1975-373.1}.
\par\noindent\textit{Status (2026-11):} As of 10 November 2026: established in the formal differential-algebra setting; \citet[Sections~8--12]{gelfanddikii1976fractional} provides the systematic formulation. Formal fractional powers are distinguished from spectral functional calculus.

\subsection{GP-1974-10: Finite-gap spectra and stationary KdV characterization}\label{prob:GP-1974-10}
Relate real smooth periodic Schrödinger potentials with finitely many open gaps to stationary equations in the KdV hierarchy; the present occurrence is the third-order, at-most-one-gap case.
\par\noindent\textit{Occurrences:} 1974-340.7, 1974-343.2, \hyperref[item:1975-373.2]{1975-373.2}.
\par\noindent\textit{Status (2026-11):} As of 10 November 2026: the periodic stationary-to-finite-gap implication is read in \citet[Section~2, Theorem~2.1]{novikov1974periodic}. The one-gap specialization requires the lower-flow constant and allowance for degenerations; the older general problem's full occurrence list belongs to the cumulative index.

\subsection{GP-1975-11: Matrix finite-gap inverse problem at fixed spectral curve}\label{prob:GP-1975-11}
Classify the regular rank-one matrix pencils \(D+U-EA\), with fixed distinct diagonal entries of \(A\), whose stationary spectral polynomial is a fixed smooth curve, and distinguish framed potentials from their diagonal-gauge quotient.
\par\noindent\textit{Occurrences:} \hyperref[item:1975-375.8]{1975-375.8}.
\par\noindent\textit{Status (2026-11):} As of November 2026, solved on the regular algebraic-geometric locus by the Jacobian construction announced here and published in 1976--1977. Related scalar precursors are 1974-336.15 and 1974-340.8; their differently quantified problems are not merged.

\subsection{GP-1975-10: Continuum confinement and the static potential}\label{prob:GP-1975-10}
For zero-temperature four-dimensional pure \(SU(N)\) Yang--Mills theory with \(N\ge2\), construct a continuum infinite-volume theory and its renormalized external-fundamental-charge observable, and prove a positive asymptotic string tension \(\lim_{R\to\infty}V(R)/R>0\). The notes motivate this statement but leave its matter-content, phase and continuum/regularization qualifications implicit; the formal action already has \(d^4x\).
\par\noindent\textit{Occurrences:} \hyperref[item:1975-377.10]{1975-377.10}, \hyperref[item:1975-377.11]{1975-377.11}, \hyperref[item:1975-377.12]{1975-377.12}, \hyperref[item:1975-377.13]{1975-377.13}, \hyperref[item:1975-386.9]{1975-386.9}.
\par\noindent\textit{Status (2026-11):} As of November 2026, the four-dimensional continuum statement is unresolved. Wilson's 1974 fixed-lattice strong-coupling construction is a prior partial realization, and the existence/mass-gap problem is logically distinct. Record 386 supplies a related Hamiltonian observable, to be linked at integration.

\subsection{GP-1975-13: Monotone radial profile from the indexed frequency matrix}\label{prob:GP-1975-13}
Recover the positive radial coefficient \(f\) of \(y^{\prime\prime}+[\lambda_{k\ell}^2f^2-\ell(\ell+1)/r^2]y=0\), with regular origin and fixed separated outer boundary condition, from its fully indexed spectrum, assuming \(rf\) is strictly increasing and normalized at the outer radius.
\par\noindent\textit{Occurrences:} \hyperref[item:1975-381.1]{1975-381.1}, \hyperref[item:1975-381.2]{1975-381.2}, \hyperref[item:1975-381.3]{1975-381.3}, \hyperref[item:1975-381.4]{1975-381.4}.
\par\noindent\textit{Status (2026-11):} As of November 2026, the reported monotone reconstruction is identified with Brodskii1976, Theorems1–2; its Abel inversion is classical. The accompanying necessary spectral compatibility condition is not advertised as a full characterization of exact arrays.

\subsection{GP-1975-14: Rank-one pointed-curve inverse correspondence}\label{prob:GP-1975-14}
For a fixed smooth pointed projective complex curve and local parameter, classify rank-one Krichever homomorphisms from its punctured-curve function ring to ordinary differential operators with smooth germ coefficients, modulo unit scalar gauge, including the unit leading-coefficient germ as data.
\par\noindent\textit{Occurrences:} \hyperref[item:1975-383.9]{1975-383.9}, \hyperref[item:1975-385.6]{1975-385.6}, \hyperref[item:1975-385.8]{1975-385.8}.
\par\noindent\textit{Status (2026-11):} As of November 2026, the precise local correspondence is established: an acyclic degree-(g-1) initial line bundle and a unit leading coefficient determine the homomorphism. Drinfeld1977, Corollary2, gives the statement and attributes it to Krichever. This is distinct from the later higher-rank inverse problem.

\subsection{GP-1975-12: Commuting scalar operators of orders two and three}\label{prob:GP-1975-12}
Find smooth local coefficients \(u,v_1,v_2\) for which \([-D^2+u,D^3+v_1D+v_2]=0\), including degenerations of the elliptic solutions.
\par\noindent\textit{Occurrences:} \hyperref[item:1975-379.4]{1975-379.4}, \hyperref[item:1975-383.6]{1975-383.6}.
\par\noindent\textit{Status (2026-11):} As of November 2026, this is a classical solved problem, already treated in the early commuting-operator literature. It is related to GP-1974-10 (finite-gap/stationary-KdV characterization), whose broader spectral quantifiers are not merged with this exact exercise.

\subsection{GP-1975-22: Normalization of the generalized Poisson law}\label{prob:GP-1975-22}
For \(x_0>0\) and \(0\le\alpha<1\), prove that \(P_0=e^{-x_0}\), \(P_n=e^{-x_0-\alpha n}x_0(x_0+\alpha n)^{n-1}/n!\) have total mass one.
\par\noindent\textit{Occurrences:} \hyperref[item:1975-387.2]{1975-387.2}.
\par\noindent\textit{Status (2026-11):} Solved before the meeting in the generalized Poisson literature (1973); the occurrence includes an elementary generating-function reconstruction. Status checked November 2026.

\subsection{GP-1975-23: Deforming a superfunction algebra along an even Poisson bracket}\label{prob:GP-1975-23}
Construct a formal associative superproduct whose first graded commutator is a given even super-Poisson bracket. The notebook leaves the geometric and parity hypotheses unspecified; the definite version treated here is the Rothstein bracket on a split supermanifold over a symplectic base.
\par\noindent\textit{Occurrences:} \hyperref[item:1975-389.8]{1975-389.8}.
\par\noindent\textit{Status (2026-11):} The specified symplectic version has a Fedosov-type construction, circulated in 2000. This is a partial identification of the notebook question, not a claim about arbitrary odd or degenerate brackets. Status November 2026.

\subsection{GP-1975-26: Reconstructing higher-rank commuting ordinary differential operators}\label{prob:GP-1975-26}
Recover a commutative ring of scalar ordinary differential operators from generic pointed spectral-curve and rank-\(r\) bundle data, with the framing and functional parameters required at infinity; orders four and six are the recorded motivating example.
\par\noindent\textit{Occurrences:} \hyperref[item:1975-391.3]{1975-391.3}, \hyperref[item:1975-391.4]{1975-391.4}.
\par\noindent\textit{Status (2026-11):} The generic smooth-curve inverse construction was completed in Krichever’s article submitted in February 1977 and published in 1978, with corrections to the preceding implementation. Status November 2026.

\subsection{GP-1975-24: Abel limits of Laplace spectral phases}\label{prob:GP-1975-24}
Let \((M,g)\) be a closed compact Riemannian manifold, let \(\lambda_k\ge0\) be the eigenvalues of \(\sqrt{-\Delta_g}\), counted with multiplicity, and let \(a>0\). Determine when \(\varepsilon\sum_ke^{-\varepsilon\lambda_k}(\delta_{\rm per}(\mu-a\lambda_k)-(2\pi)^{-1})\) converges as a periodic distribution, and express its nonzero Fourier coefficients through closed geodesics. The source’s intended genericity assumptions are incomplete. The intended geometric setting has isolated closed geodesics and pairwise incommensurable primitive unoriented lengths; nondegeneracy of contributing iterates is supplied as a modern trace-formula hypothesis, not a certain reading of an abbreviation.
\par\noindent\textit{Occurrences:} \hyperref[item:1975-391.7]{1975-391.7}, \hyperref[item:1975-391.8]{1975-391.8}.
\par\noindent\textit{Status (2026-11):} As of November 2026, fixed-mode closed-orbit contributions are explained by wave-trace theory. A source proving the full distributional limit under exactly the notebook’s hypotheses was not located; uniform growth and noncancellation are the explicit remaining qualifications, not an assertion that the whole wave-trace problem is open. The chapter develops the Fourier convergence criterion and solves the circle model; its sphere counterexample explicitly lies outside those geometric hypotheses.

\subsection{GP-1975-25: A Harnack-type bound for sums of eigenfunctions}\label{prob:GP-1975-25}
The definite reconstructed question is whether the extended Courant nodal-domain bound holds for linear combinations of Laplace eigenfunctions below an eigenvalue cutoff. Arnold’s notebook intervention asks more broadly for a Harnack-type extension and states no exact bound.
\par\noindent\textit{Occurrences:} \hyperref[item:1975-391.16]{1975-391.16}.
\par\noindent\textit{Status (2026-11):} The specified extended Courant property is false in general; an explicit two-dimensional Neumann counterexample is given in the 2018 paper. This is a partial match to the broader historical question. Status November 2026.

\subsection{GP-1975-20: Infinitely many nonzero coefficients in Perchik’s relative Euler series}\label{prob:GP-1975-20}
For the fixed-rank plane Hamiltonian Lie algebra, determine whether the relative Euler coefficients \(E_m=\chi H^*(\mathfrak{ham}_2,\mathfrak{sp}_2)_{2m}\) are nonzero for infinitely many positive integers \(m\). The series includes the linear Hamiltonians and is not the series for \(\mathfrak{ham}_2^0\).
\par\noindent\textit{Occurrences:} \hyperref[item:1975-393.8]{1975-393.8}.
\par\noindent\textit{Status (2026-11):} No resolution was located as of 10 November 2026. The bounded development derives the product rigorously and verifies all coefficients through m=150 by two independent exact methods; this does not prove infinitude.

\subsection{GP-1975-21: Full cohomology of formal Hamiltonian vector fields}\label{prob:GP-1975-21}
Compute the continuous cohomology of \(\mathfrak{ham}_{2n}\), already for \(n=1\), with its grading and products. This is distinct from the solved cohomology calculation for the full formal-vector-field algebra \(W_n\).
\par\noindent\textit{Occurrences:} \hyperref[item:1975-393.2]{1975-393.2}.
\par\noindent\textit{Status (2026-11):} Finite-weight classes and factorizations are established, but the targeted search through 10 November 2026 located no complete fixed-rank answer. Stable-rank graph-homology calculations are not a solution of this fixed-rank problem.

\subsection{GP-1973-33: Serre’s problem for polynomial rings}\label{prob:GP-1973-33}
Is every finitely generated projective module over a polynomial ring over a field free? The earlier three-variable instance recurs in arbitrary dimension, together with unimodular-row completion.
\par\noindent\textit{Occurrences:} 1973-322.1, \hyperref[item:1975-395.1]{1975-395.1}, \hyperref[item:1975-395.2]{1975-395.2}.
\par\noindent\textit{Status (2026-11):} Quillen and Suslin solved the problem in 1976. Their manuscripts were received before the April seminar, but receipt does not establish public circulation. Suslin's Section III, Theorems 2--3, and Quillen's publisher chronology were checked. Status November 2026.

\section{Corrections to the transcription and translation}\label{corr:1975}
The evidence distinguishes a transcription correction from a mathematical emendation of wording faithfully preserved in the notebook.
\subsection*{P000603 (both; faithful error): coefficient of the deformed bracket}
\noindent\textit{Transcription reads.} Two pure derivative products; no lower terms on the plane or torus\par
\noindent\textit{Critical reading.} Retain the scan reading; a Moyal reconstruction must include all binomial mixed derivatives of the same order.\par
\noindent\textit{Evidence.} Scan inspected. Expansion of the cube of the constant Poisson bidifferential operator supplies the missing mixed terms; the abbreviation for lower terms remains uncertain.\par
\subsection*{P000605 (both; faithful error): covering-group arrow}
\noindent\textit{Transcription reads.} SO(3) points to SU(2)\par
\noindent\textit{Critical reading.} The canonical double-cover homomorphism points from SU(2) to SO(3).\par
\noindent\textit{Evidence.} The diplomatic transcription explicitly records the arrow direction; the scan confirms it. The error concerns interpretation, not copying.\par
\subsection*{P000606 (both; faithful error): local gauge coupling}
\noindent\textit{Transcription reads.} The A term has no gamma contraction and mismatched parentheses\par
\noindent\textit{Critical reading.} Use \(\bar e(i\gamma^\mu(\partial_\mu+iA_\mu)-m)e\).\par
\noindent\textit{Evidence.} The missing contraction is visible in the scan; covariance and scalar-index contraction determine the mathematical repair.\par
\subsection*{P000607 (both; faithful error): Gaussian perturbation analogy}
\noindent\textit{Transcription reads.} Natural Lebesgue measure on field space; convergent-looking real cubic perturbation\par
\noindent\textit{Critical reading.} Treat the functional integral and cubic perturbation series formally or after a specified regulator/contour.\par
\noindent\textit{Evidence.} Infinite-dimensional field space has no locally finite translation-invariant Lebesgue measure; a real cubic term defeats decay on one tail.\par
\subsection*{P000607 (both; faithful error): Maxwell term}
\noindent\textit{Transcription reads.} Plus F squared in the Lorentz-signature Lagrangian\par
\noindent\textit{Critical reading.} Use \(-F_{\mu\nu}F^{\mu\nu}/(4e^2)\) for the stated Minkowski conventions; a Euclidean plus requires an explicit change of convention.\par
\noindent\textit{Evidence.} Compare the metric on P000606 and the action \(e^{iS}\) on P000607; this is a convention-sensitive emendation.\par
\subsection*{P000607 (both; ambiguity): last model exponent}
\noindent\textit{Transcription reads.} \(ix^2(1+\varepsilon x^3)\)\par
\noindent\textit{Critical reading.} Retain the quintic phase as written; do not silently turn it into the preceding cubic phase.\par
\noindent\textit{Evidence.} Scan and both transcriptions agree; the discrepancy is internal to the notebook.\par
\subsection*{P000608 (both; faithful error): description of K}
\noindent\textit{Transcription reads.} According to Maxwell\par
\noindent\textit{Critical reading.} The displayed K solves the Dirac equation; retain Maxwell as a notebook misidentification unless it was intended to anticipate D.\par
\noindent\textit{Evidence.} The operator i gamma d-m and its matrix numerator identify the Dirac kernel.\par
\subsection*{P000608 (both; faithful error): generalized Feynman parameters}
\noindent\textit{Transcription reads.} No prefactor in either parameter integral\par
\noindent\textit{Critical reading.} Insert \(\Gamma(\sum_i\lambda_i)/\prod_i\Gamma(\lambda_i)\).\par
\noindent\textit{Evidence.} The scan was inspected. Beta-integral normalization and the check with \(a=b=1\), \(\lambda=\mu=2\) give a factor of \(6\), rather than \(1\).\par
\subsection*{P000610 (both; faithful error): BDF classification}
\noindent\textit{Transcription reads.} Arbitrary elements of the Calkin algebra\par
\noindent\textit{Critical reading.} Supply essential normality of the lifts and specify unitary equivalence modulo compacts.\par
\noindent\textit{Evidence.} Scan contains no normality qualifier; this is a missing theorem hypothesis, not a translation error.\par
\subsection*{P000610 (both; faithful error): Euler exponent}
\noindent\textit{Transcription reads.} \((n^2-3n)/2\)\par
\noindent\textit{Critical reading.} n(3n-1)/2, summed over all integers n\par
\noindent\textit{Evidence.} Scan confirms the erroneous exponent. The pentagonal-number identity fixes the intended expression.\par
\subsection*{P000610 (both; faithful error): Ext(X)}
\noindent\textit{Transcription reads.} Set of embeddings C(X) into the Calkin algebra\par
\noindent\textit{Critical reading.} Stable unitary-equivalence classes of extensions, modulo the specified split-extension relation; \(K_1(X)\) denotes odd K-homology.\par
\noindent\textit{Evidence.} Scan and both transcriptions agree. The equivalence relation is absent from the notebook.\par
\subsection*{P000611 (both; faithful error): infinitesimal SU(2) transformation}
\noindent\textit{Transcription reads.} exp(+i tau alpha/2) followed by 1-i tau alpha/2\par
\noindent\textit{Critical reading.} Use the same sign in the finite transformation and its first-order expansion.\par
\noindent\textit{Evidence.} Both signs are visible in the scan; this is an internal convention mismatch.\par
\subsection*{P000613 (both; faithful error): first tau sum}
\noindent\textit{Transcription reads.} Unweighted signed partitions with at most 24 equality signs\par
\noindent\textit{Critical reading.} Use multiplicities \(0\leq r_j\leq24\) and weights \(\prod_j\binom{24}{r_j}\), or equivalently 24 labelled copies of each part.\par
\noindent\textit{Evidence.} Scan inspected. At n=2 the written unweighted sum is -1 whereas tau(2)=-24; a run of 24 copies has 23 equality signs.\par
\subsection*{P000614 (both; faithful error): Deligne root assertion}
\noindent\textit{Transcription reads.} Purely imaginary (nonreal)\par
\noindent\textit{Critical reading.} Nonreal conjugate roots of modulus \(p^{11/2}\); generally not purely imaginary.\par
\noindent\textit{Evidence.} Russian scan says both expressions. tau(2)=-24 forces real part -12, so the parenthetical is not synonymous with the preceding words.\par
\subsection*{P000614 (both; faithful error): Galois representation target}
\noindent\textit{Transcription reads.} Representation into \(\mathbb Q_\ell^2\)\par
\noindent\textit{Critical reading.} The space is \(\mathbb Q_\ell^2\) and the homomorphism target is \(\operatorname{GL}_2(\mathbb Q_\ell)\).\par
\noindent\textit{Evidence.} This is a target/type correction required for the trace assertion, not a claim about the precise handwritten arrow.\par
\subsection*{P000614 (both; faithful error): Manin coefficient polynomial}
\noindent\textit{Transcription reads.} First coefficient 1 and second coefficient +691/6\par
\noindent\textit{Critical reading.} Use \((691/18)(\Delta^8\delta^2-\Delta^2\delta^8)-(691/6)(\Delta^6\delta^4-\Delta^4\delta^6)\).\par
\noindent\textit{Evidence.} The notebook scan confirms the printed reading. Manin1973 Section1.3 supplies the coefficient theorem; exact evaluation through n=50 supports the corrected factorized form and rejects the notebook already at n=2.\par
\subsection*{P000614 (both; faithful error): Rademacher formula}
\noindent\textit{Transcription reads.} Phase \(e^{-2\pi i n/k}\) and \(\omega_{hk}^{24}=1\)\par
\noindent\textit{Critical reading.} Use exp(-2 pi i n h/k), Dedekind-sum phase exp(pi i s(h,k)), and the normalized Rademacher formula in the item.\par
\noindent\textit{Evidence.} Scan matches transcription. DLMF27.14.9--11 supplies the complete normalization; s(1,5)=1/5 disproves the asserted phase order.\par
\subsection*{P000614 (both; faithful error): partition-series inversion}
\noindent\textit{Transcription reads.} Sum from n=1 before raising to -24\par
\noindent\textit{Critical reading.} Sum from n=0 and set p(0)=1.\par
\noindent\textit{Evidence.} The constant term is required for formal series inversion; the lower limit 1 is visible in the scan.\par
\subsection*{P000614 (both; faithful error): prime on Manin summation}
\noindent\textit{Transcription reads.} No explanation of boundary weight\par
\noindent\textit{Critical reading.} At \(\delta/\Delta=1/2\), with \(\delta^{\prime}=0\), assign weight \(1/2\).\par
\noindent\textit{Evidence.} Manin1973 p374 explicitly states this rule. It is omitted in the seminar explanation; without it n=2 already fails.\par
\subsection*{P000614 (both; faithful error): prime-power tau}
\noindent\textit{Transcription reads.} \(\tau(p^a)=\theta_p^a+\overline{\theta_p}^{\,a}\)\par
\noindent\textit{Critical reading.} Use the complete symmetric sum of degree a or the Hecke recurrence.\par
\noindent\textit{Evidence.} Scan double-underlines the erroneous formula. p=2,a=2 gives -3520 instead of -1472.\par
\subsection*{P000614 (both; faithful error): quadratic-form count}
\noindent\textit{Transcription reads.} Two forms in dimension 24\par
\noindent\textit{Critical reading.} Retain the reading; the count is false even for positive definite even unimodular lattices. The two-dimensional modular-form space is a possible intended referent.\par
\noindent\textit{Evidence.} Scan explicitly reads two in all. E8 cubed, D24 plus, and Leech provide three inequivalent examples.\par
\subsection*{P000615 (both; faithful error): Eisenstein coefficients}
\noindent\textit{Transcription reads.} Coefficients \(\sigma_3(n),\sigma_5(n)\)\par
\noindent\textit{Critical reading.} For normalized \(E_4\) and \(E_6\), use \(240\sigma_3(n),-504\sigma_5(n)\), respectively, and constant term one.\par
\noindent\textit{Evidence.} The scan suppresses normalization; the ring statement requires a consistent choice of generators.\par
\subsection*{P000615 (both; faithful error): eta prefactor}
\noindent\textit{Transcription reads.} exp(2 pi i n z/24)\par
\noindent\textit{Critical reading.} exp(2 pi i z/24)\par
\noindent\textit{Evidence.} The extra n is visible outside the product in the scan. DLMF27.14.12 gives the standard prefactor.\par
\subsection*{P000615 (both; faithful error): final Kostant formula}
\noindent\textit{Transcription reads.} psi(q) equals the character sum\par
\noindent\textit{Critical reading.} psi(q) raised to dim K equals the character sum.\par
\noindent\textit{Evidence.} Compare the same notebook formula on P000610 and Kostant2003 Theorem0.1; the exponent is omitted in this repetition.\par
\subsection*{P000617 (both; ambiguity): quadratic determinant}
\noindent\textit{Transcription reads.} Square root of det A double prime\par
\noindent\textit{Critical reading.} Use the mixed Hessian, its absolute determinant for the modulus, and a separate coherent phase.\par
\noindent\textit{Evidence.} Scan inspected. The two Hessian strokes are ambiguous; the mixed-Hessian interpretation is required by the displayed generating relation.\par
\subsection*{P000617 (both; faithful error): unitarity implication}
\noindent\textit{Transcription reads.} Unitarity implies metaplectic group; absolute value c equals 1\par
\noindent\textit{Critical reading.} Unitarity gives a U(1) extension; a further phase normalization selects the double cover.\par
\noindent\textit{Evidence.} Scan confirms the implication and the following separate choice c=plus or minus1. The distinction is mathematical, not a translation error.\par
\subsection*{P000618 (both; faithful error): first Toda potential}
\noindent\textit{Transcription reads.} Difference of exponentials with n in the exponents under a sum over k\par
\noindent\textit{Critical reading.} Use the single exponential sum displayed on P000619; the difference belongs to a force equation.\par
\noindent\textit{Evidence.} Both scans inspected. Changing n to k alone would produce a telescoping sum, not the Toda potential.\par
\subsection*{P000618 (both; ambiguity): pendulum abbreviation}
\noindent\textit{Transcription reads.} Second integral is an elliptic function\par
\noindent\textit{Critical reading.} Interpret as the next integration or inversion step; there is only one independent first integral for one degree of freedom.\par
\noindent\textit{Evidence.} The abbreviated Russian permits this reading; the energy equation gives the explicit quadrature.\par
\subsection*{P000619 (both; faithful error): classical Fermi variables}
\noindent\textit{Transcription reads.} pq+qp=1 for Grassmann variables\par
\noindent\textit{Critical reading.} Classical odd coordinates anticommute to zero; 1 is the canonical quantum anticommutator or multiplication/differentiation pairing.\par
\noindent\textit{Evidence.} Scan confirms the 1. The subsequent graded Poisson construction requires separating these meanings.\par
\subsection*{P000619 (both; ambiguity): root-system rank}
\noindent\textit{Transcription reads.} n particles associated with A\_n\par
\noindent\textit{Critical reading.} n particles have relative root system A\_(n-1), with a separate center-of-mass coordinate.\par
\noindent\textit{Evidence.} The listed roots e\_i-e\_j span the sum-zero hyperplane in R\textasciicircum{}n; the notebook does not fix its rank convention.\par
\subsection*{P000620 (both; ambiguity): A and exponential grouping}
\noindent\textit{Transcription reads.} A added after the action integral and itself beginning with an exponential\par
\noindent\textit{Critical reading.} Retain the unresolved grouping; the Schur-complement pattern suggests a determinant or Berezinian measure factor.\par
\noindent\textit{Evidence.} Scan inspected; multiple braces and the inserted second trace are visibly ambiguous. No complete emendation was verified.\par
\subsection*{P000620 (both; faithful error): second-class condition}
\noindent\textit{Transcription reads.} Nonzero pairwise brackets\par
\noindent\textit{Critical reading.} Invertibility of the complete constraint-bracket matrix on the regular constraint surface.\par
\noindent\textit{Evidence.} The scan states nonzero brackets; entrywise nonvanishing does not imply invertibility.\par
\subsection*{P000621 (both; ambiguity): fractional-power normalization}
\noindent\textit{Transcription reads.} L=-D squared+u and A equals L to the three-halves with leading D cubed\par
\noindent\textit{Critical reading.} Use the formal square root of -L with leading D, or supply the equivalent complex phase for a root of L.\par
\noindent\textit{Evidence.} Scan confirms both signs. This is a formal branch and normalization issue, not spectral functional calculus.\par
\subsection*{P000621 (both; faithful error): leading coefficients of L\_m}
\noindent\textit{Transcription reads.} u\_m and u\_(m-1) are constant\par
\noindent\textit{Critical reading.} v\_m and v\_(m-1) are constant.\par
\noindent\textit{Evidence.} Scan and both transcriptions use u after defining the second operator with coefficients v; this is a notebook variable slip.\par
\subsection*{P000621-P000624 (both; provenance): retrospective dating}
\noindent\textit{Transcription reads.} Notebook date20October1975\par
\noindent\textit{Critical reading.} Retain20October; record the discrepancy with the1977 survey’s retrospective November--December1975 account.\par
\noindent\textit{Evidence.} The notebook date is explicit and applies to both headings. Krichever1977 introduction p186 gives the later recollection.\par
\subsection*{P000622 (both; faithful error): characteristic polynomial degree}
\noindent\textit{Transcription reads.} L\_m to the n+1 plus r\_n L\_m to the n\par
\noindent\textit{Critical reading.} An n-dimensional solution space has characteristic polynomial of degree n.\par
\noindent\textit{Evidence.} Scan confirms n+1. Krichever1977 Theorem2.1 uses the n by n initial-jet matrix.\par
\subsection*{P000622 (both; faithful error): local-coordinate point}
\noindent\textit{Transcription reads.} 1/k is a local parameter near x\_0\par
\noindent\textit{Critical reading.} 1/k is a local parameter near P\_0 (the marked point of the curve).\par
\noindent\textit{Evidence.} The scan says x\_0 at right but P\_0 under the expansion; x\_0 is the independent-variable basepoint.\par
\subsection*{P000623 (both; faithful error): coefficient dependence bounds}
\noindent\textit{Transcription reads.} u\_i depends on xi\_s with 0<=s<=i-1\par
\noindent\textit{Critical reading.} The displayed bounds are incorrect; for n=2, u\_0=c\_0-2c\_2 xi\_1 prime.\par
\noindent\textit{Evidence.} Scan confirms these indices. Direct matching of the normalized exponential expansion, also Krichever1977 Section1, yields the counterexample.\par
\subsection*{P000623 (both; faithful error): gap numbers}
\noindent\textit{Transcription reads.} Some pole order s+mu\_s n in each residue class\par
\noindent\textit{Critical reading.} The least available pole order in each residue class.\par
\noindent\textit{Evidence.} Without minimality the sum of mu\_s is not the genus; the scan omits the qualifier.\par
\subsection*{P000623 (both; faithful error): generic-point operator count}
\noindent\textit{Transcription reads.} g algebraically independent commuting operators\par
\noindent\textit{Critical reading.} Retain the error; a curve ring has transcendence degree one. Algebra or module generators would be different assertions.\par
\noindent\textit{Evidence.} The scan clearly says algebraically independent. The earlier polynomial-relation theorem already contradicts the wording for g>1.\par
\subsection*{P000624 (both; faithful error): Baker--Akhiezer exponential}
\noindent\textit{Transcription reads.} Exponential of the integral without x\par
\noindent\textit{Critical reading.} Use exp((x-x\_0) Omega(P)) with a regularized second-kind integral and normalization factor.\par
\noindent\textit{Evidence.} Scan confirms the missing x. The P000622 required local asymptotic cannot hold otherwise.\par
\subsection*{P000624 (both; faithful error): common centralizer}
\noindent\textit{Transcription reads.} Two operators commuting with a third commute mutually\par
\noindent\textit{Critical reading.} Require a suitable non-scalar operator; the positive-order scalar-ODE hypothesis suffices.\par
\noindent\textit{Evidence.} The scan has no qualifier; taking the third operator to be the identity gives an immediate counterexample.\par
\subsection*{P000624 (both; ambiguity): theta-series normalization}
\noindent\textit{Transcription reads.} 2 pi i multiplying the quadratic term with an unscaled linear term\par
\noindent\textit{Critical reading.} For a-periods1 use pi i m transpose B m +2 pi i m transpose z.\par
\noindent\textit{Evidence.} Scan inspected; the crossing after2 pi i leaves a damaged coefficient, and the final linear term lacks its needed factor.\par
\subsection*{P000625-P000626 (both; faithful error): stationary Lax lemma and proof}
\noindent\textit{Transcription reads.} \(\Lambda^{\prime}=[Q,\Lambda]\)\par
\noindent\textit{Critical reading.} \(\Lambda^{\prime}=[\Lambda,Q]\)\par
\noindent\textit{Evidence.} The scans show the reversed sign at the bottom of P000625 and in P000626; the earlier equation on P000625 has the consistent order. Since \([D+Q,\Lambda]=\Lambda^{\prime}+[Q,\Lambda]\), the later signs are faithful notebook errors.\par
\subsection*{P000628 (both; faithful error): spin labels}
\noindent\textit{Transcription reads.} baryons have spin 1/2; mesons spin 0\par
\noindent\textit{Critical reading.} restrict to the stated low-spin multiplets\par
\noindent\textit{Evidence.} The scan faithfully has the shorthand labels; the same page invokes the Delta++ state. Spin-3/2 baryons and spin-1 mesons already occur in Levin and Frankfurt (1968).\par
\subsection*{P000628-P000629 (both; ambiguity): flavor times color invariance}
\noindent\textit{Transcription reads.} L cannot be invariant under SU(3) flavor times SU(3) color\par
\noindent\textit{Critical reading.} distinguish approximate global flavor symmetry from exact local color gauge invariance\par
\noindent\textit{Evidence.} The negative phrase is visible in P000628; P000629 explicitly says color symmetry is exact. The source does not warrant silently deleting the negation.\par
\subsection*{P000629 (both; faithful error): middle term of meson singlet}
\noindent\textit{Transcription reads.} \(\bar q_g\bar q_g\)\par
\noindent\textit{Critical reading.} \(\bar q_g q_g\)\par
\noindent\textit{Evidence.} The scan shows two bars. A singlet in the conjugate times defining color representation has one barred and one unbarred factor.\par
\subsection*{P000629-P000630 (both; faithful error): phase of the rectangular Wilson loop}
\noindent\textit{Transcription reads.} \(e^{iW(C)},\quad W(C)\to E(R)T\)\par
\noindent\textit{Critical reading.} \(W(C)\to-E(R)T\) in the real-time convention used on P000630\par
\noindent\textit{Evidence.} Both formulas are faithful to the scans. P000630 explicitly writes exp(-iET); a Euclidean convention instead uses exp(-VT).\par
\subsection*{P000631 (both; faithful error): normalization of formal wave coefficients}
\noindent\textit{Transcription reads.} \(\xi_i(0,y,t)=0\) for \(i>1\)\par
\noindent\textit{Critical reading.} \(\xi_i(0,y,t)=0\) for \(i\ge1\), if this is to fix every integration function\par
\noindent\textit{Evidence.} The Russian transcription and scan explicitly show \(i>1\). The recurrence leaves the integration function in \(\xi_1\) undetermined under that bound.\par
\subsection*{P000631 (both; faithful error): phase and constant next coefficient}
\noindent\textit{Transcription reads.} \(u_{n-1}=\mathrm{const}\), \(e^{kx+k^ny}\)\par
\noindent\textit{Critical reading.} require \(u_{n-1}=0\), or add \(u_{n-1}k^{n-1}y\) to the phase\par
\noindent\textit{Evidence.} The coefficient of \(k^{n-1}\) in the substituted equation is \(u_{n-1}\); the scan records no corresponding term in the phase.\par
\subsection*{P000632 (both; faithful error): Baker-Akhiezer theorem hypotheses}
\noindent\textit{Transcription reads.} arbitrary marked poles; unique up to proportionality\par
\noindent\textit{Critical reading.} effective nonspecial degree-g divisor; uniqueness after leading coefficient 1 is fixed\par
\noindent\textit{Evidence.} The omitted divisor hypotheses and the normalization conflict are mathematical omissions in the notebook; Krichever1976 conditions1–2 supply the comparison.\par
\subsection*{P000632 (both; faithful error): normalized Baker-Akhiezer asymptotic}
\noindent\textit{Transcription reads.} \(1+o(1/k)\)\par
\noindent\textit{Critical reading.} \(1+O(1/k)\)\par
\noindent\textit{Evidence.} The scan and Russian text have a small \(o\), but the preceding series explicitly permits a nonzero \(\xi_1/k\) term.\par
\subsection*{P000635 (both; faithful error): action normalization and inverse integral}
\noindent\textit{Transcription reads.} action integral has no 1/pi; inverse exponential coefficient is 2 with no lower limit\par
\noindent\textit{Critical reading.} insert 1/pi in the action and eta as lower limit; alternatively retain the action and replace 2 by 2/pi\par
\noindent\textit{Evidence.} The scan omits both 1/pi and the lower endpoint. Brodskii1976 equations(3),(5), and the independently checked Abel-kernel composition determine the correction.\par
\subsection*{P000635 (both; faithful error): radial eigenvalue equation}
\noindent\textit{Transcription reads.} \(-y^{\prime\prime}+(\lambda^2f^2-\ell(\ell+1)/r^2)y=\lambda_{k\ell}y\)\par
\noindent\textit{Critical reading.} \(y^{\prime\prime}+[\lambda_{k\ell}^2f^2-\ell(\ell+1)/r^2]y=0\)\par
\noindent\textit{Evidence.} The scan matches the diplomatic text, so this is not OCR repair. Brodskii1976 title and equation(2), and the immediately following turning-point integral, identify the intended equation.\par
\subsection*{P000635 (both; faithful error): spectral interpolation equality}
\noindent\textit{Transcription reads.} the finite spectral pairs define an exact single-valued correspondence\par
\noindent\textit{Critical reading.} read as joint high-frequency asymptotic sampling with boundary-dependent phase\par
\noindent\textit{Evidence.} The scan gives a literal equality. Brodskii1976 Theorem1 is explicitly asymptotic; its corollary supplies only a necessary compatibility condition.\par
\subsection*{P000637 (both; faithful error): base of the canonical principal bundle}
\noindent\textit{Transcription reads.} C(P) is a principal bundle with base P\par
\noindent\textit{Critical reading.} C(P) is a principal G-bundle with base C(P)/G; its map to P is affine\par
\noindent\textit{Evidence.} The scan and Russian text name P. The displayed G-orbit construction and García1972 Theorems1–3 show the correct quotient base.\par
\subsection*{P000638 (both; faithful error): elliptic uniformization}
\noindent\textit{Transcription reads.} \(x=\wp(z), y=\wp^{\prime}(z)\) for \(y^2=x(x-1)(x-2)\)\par
\noindent\textit{Critical reading.} \(x=1+\wp(z;4,0), y=\wp^{\prime}(z;4,0)/2\)\par
\noindent\textit{Evidence.} The scan has the unscaled formulas. Substitution into the Weierstrass equation supplies the necessary shift and scale.\par
\subsection*{P000638 (both; provenance): speaker in dated heading}
\noindent\textit{Transcription reads.} catalogue record 383: S. P. Novikov\par
\noindent\textit{Critical reading.} retain catalogue attribution but identify manuscript heading as Drinfeld (Volodya)\par
\noindent\textit{Evidence.} The scan explicitly reads Drinfeld (Volodya), dated 1/III-76; the catalogue and manuscript are distinct witnesses.\par
\subsection*{P000639 (both; faithful error): coordinate dependence of leading coefficient}
\noindent\textit{Transcription reads.} b is independent of k and of the local parameter\par
\noindent\textit{Critical reading.} For a fixed parameter, \(b\) is independent of \(k\); if \(\widetilde z=cz+O(z^2)\), then \(\widetilde b=c^{-1}b\).\par
\noindent\textit{Evidence.} The scan explicitly adds independence of parameter. Rescaling normalized Laurent leading terms disproves that addition.\par
\subsection*{P000639 (both; faithful error): number of generators at an ordinary point}
\noindent\textit{Transcription reads.} g generators of orders g+1,...,2g+1\par
\noindent\textit{Critical reading.} g+1 generators of those orders\par
\noindent\textit{Evidence.} The list in both the scan and transcription contains g+1 orders. The pole-order induction verifies the corrected generating set.\par
\subsection*{P000640 (both; faithful error): first filtration quotient}
\noindent\textit{Transcription reads.} \(\dim_B M_i/M_{i+1}=1\)\par
\noindent\textit{Critical reading.} \(M_i/M_{i-1}\) is free of rank one over \(B\)\par
\noindent\textit{Evidence.} The scan has denominator \(M_{i+1}\), incompatible with the displayed increasing filtration; every subsequent graded map uses \(M_{i-1}\).\par
\subsection*{P000642 (both; ambiguity): definition of the function ring}
\noindent\textit{Transcription reads.} \(A\in H^0(X-P,\mathcal O_X)\)\par
\noindent\textit{Critical reading.} \(A=H^0(X-P,\mathcal O_X)\)\par
\noindent\textit{Evidence.} The scan has a membership-shaped symbol; the whole construction uses A as this algebra, not one of its elements.\par
\subsection*{P000643 (both; faithful error): Jacobian as multiplicative sheaf cohomology}
\noindent\textit{Transcription reads.} \(J=H^1(X,\mathcal O_X^*)\)\par
\noindent\textit{Critical reading.} \(\operatorname{Pic}(X)=H^1(X,\mathcal O_X^*)\); \(J=\operatorname{Pic}^0(X)\)\par
\noindent\textit{Evidence.} The scan and both transcriptions omit the degree-zero component; the notebook simultaneously speaks of bundles of a fixed degree.\par
\subsection*{P000643 (both; faithful error): Jacobian parameterization and value at zero}
\noindent\textit{Transcription reads.} a Jacobian point alone; natural B to C for the general differential algebra\par
\noindent\textit{Critical reading.} fix the unit leading coefficient as well; value at zero applies to the stated algebra of smooth germs\par
\noindent\textit{Evidence.} The previous page supplies the leading coefficient, and Corollary2 of Drinfeld1977 retains it as independent data. A general complex differential algebra need not have an augmentation.\par
\subsection*{P000644 (both; faithful error): Lie-algebra membership}
\noindent\textit{Transcription reads.} \(A_\mu\in SU(2)\) (Lie algebra)\par
\noindent\textit{Critical reading.} \(A_\mu\in\mathfrak{su}(2)\)\par
\noindent\textit{Evidence.} The written anti-Hermitian trace-zero condition is the Lie algebra condition; scan agrees with SU(2).\par
\subsection*{P000645 (both; faithful error): Legendre transformation}
\noindent\textit{Transcription reads.} \(\tfrac12 e(e+DA_0)+\tfrac12h^2\)\par
\noindent\textit{Critical reading.} \(\tfrac12(e^2+h^2)+eDA_0\)\par
\noindent\textit{Evidence.} Direct Legendre transformation; the scan clearly retains one half on \(eDA_0\).\par
\subsection*{P000645 (both; faithful error): integrated Hamiltonian}
\noindent\textit{Transcription reads.} \(\tfrac12\int(e^2+h^2-A_0De)\)\par
\noindent\textit{Critical reading.} \(\tfrac12\int(e^2+h^2)-\int A_0De\)\par
\noindent\textit{Evidence.} Integration by parts with the preceding normalization; scan checked.\par
\subsection*{P000646 (both; faithful error): first perturbative correction}
\noindent\textit{Transcription reads.} -Delta\textasciicircum{}\{-1\}[a,partial Delta\textasciicircum{}\{-1\}rho]\par
\noindent\textit{Critical reading.} -g Delta\textasciicircum{}\{-1\}[a,partial Delta\textasciicircum{}\{-1\}rho] for P=0\par
\noindent\textit{Evidence.} Formal expansion of Delta+g ad(a)partial; no g visible in scan.\par
\subsection*{P000646 (both; faithful error): three versions of the Gauss equation}
\noindent\textit{Transcription reads.} rho+[P,a]; [a,P]; rho+g[a,P]\par
\noindent\textit{Critical reading.} rho-g[a,P], with D=partial+g[a,.]\par
\noindent\textit{Evidence.} Expanding D(P+partial Psi)=rho fixes both g and commutator order; all conflicting versions occur in scan.\par
\subsection*{P000647 (both; faithful error): canonical commutator}
\noindent\textit{Transcription reads.} [a(x),p(y)]=delta(x-y)\par
\noindent\textit{Critical reading.} i hbar times color identity and transverse spatial projector\par
\noindent\textit{Evidence.} Both variables have already been restricted to transverse fields; normalization cannot be the unconstrained scalar delta.\par
\subsection*{P000648 (both; faithful error): quantifiers in the length theorem}
\noindent\textit{Transcription reads.} for every N there exists n\par
\noindent\textit{Critical reading.} for every n there exists N(Λ,n), in the intended uniform bound\par
\noindent\textit{Evidence.} The scan confirms the printed quantifiers. The replacement is a mathematical reconstruction supported by the Gabriel–Roiter filtration bound, not a transcription repair.\par
\subsection*{P000651 (both; provenance): speaker boundary after Drinfeld}
\noindent\textit{Transcription reads.} Catalogue record names only Krichever\par
\noindent\textit{Critical reading.} Preserve Shnirelman as speaker of the spectral discussion\par
\noindent\textit{Evidence.} The manuscript explicitly underlines Shnirelman; the catalogue record is retained, with item-level speaker attribution.\par
\subsection*{P000651 (both; faithful error): sphere-average Laplacian}
\noindent\textit{Transcription reads.} (average-value)/tau\textasciicircum{}2\par
\noindent\textit{Critical reading.} 2n times that limit for the standard Laplace–Beltrami operator\par
\noindent\textit{Evidence.} Normal-coordinate Taylor expansion; the missing factor is in the notebook, not introduced by translation.\par
\subsection*{P000652 (both; faithful error): elliptic delta-comb regularization}
\noindent\textit{Transcription reads.} damping and constant placed inside delta argument\par
\noindent\textit{Critical reading.} damping multiplies the delta distribution; the constant Fourier mode is subtracted outside\par
\noindent\textit{Evidence.} Scan checked. The reconstruction has the geometric-series Fourier coefficients required by the surrounding formula.\par
\subsection*{P000652 (both; faithful error): negative hyperbolic multiplier amplitude}
\noindent\textit{Transcription reads.} (-1)\textasciicircum{}k/sinh(k alpha/2)\par
\noindent\textit{Critical reading.} sinh denominator for even k; cosh denominator for odd k, with separate phase\par
\noindent\textit{Evidence.} The scan agrees with the abbreviated formula; direct determinant calculation corrects it.\par
\subsection*{P000653 (both; faithful error): analytic continuation claim}
\noindent\textit{Transcription reads.} continues analytically downwards with simple poles\par
\noindent\textit{Critical reading.} principal simple-pole-type wave-trace singularity at nondegenerate nonzero orbit lengths; no general meromorphic continuation\par
\noindent\textit{Evidence.} The note is faithful; the standard local expansion has logarithmic lower terms.\par
\subsection*{P000654 (both; faithful error): action integrand}
\noindent\textit{Transcription reads.} \(A_{\mu\nu}^2\)\par
\noindent\textit{Critical reading.} \(F_{\mu\nu}^2\), with an explicit trace-sign convention\par
\noindent\textit{Evidence.} The scan has A in the action and F in the definition below; mathematical symbol slip.\par
\subsection*{P000654 (both; ambiguity): character terminology}
\noindent\textit{Transcription reads.} any character (function constant on conjugacy classes)\par
\noindent\textit{Critical reading.} class function; a character is a special example\par
\noindent\textit{Evidence.} The manuscript itself shifts terminology. Preserve its wording and distinguish the notions.\par
\subsection*{P000655 (both; faithful error): opening algebra name}
\noindent\textit{Transcription reads.} formal divergence-free vector fields\par
\noindent\textit{Critical reading.} Hamiltonian formal vector fields on 2n variables; these coincide with divergence-free fields only when n=1\par
\noindent\textit{Evidence.} The scan reads the divergence-free abbreviation, whereas the displayed Poisson differential identifies the Hamiltonian algebra.\par
\subsection*{P000655 (both; faithful error): tensor product with H(S\textasciicircum{}3)}
\noindent\textit{Transcription reads.} all m without exception\par
\noindent\textit{Critical reading.} the absolute/relative factorization requires m nonzero\par
\noindent\textit{Evidence.} The source omits the zero-weight exception; verified in Mohammadzadeh–Fuchs2014 Section1 and Theorem2.1.\par
\subsection*{P000655-P000656 (both; faithful error): Euler characteristic symbol}
\noindent\textit{Transcription reads.} \(\chi_m\), after first defining \(\chi_m^*\) for the relative complex\par
\noindent\textit{Critical reading.} retain the relative Euler characteristic consistently\par
\noindent\textit{Evidence.} Scan confirms loss of the star; an absolute Euler interpretation makes the series false.\par
\subsection*{P000656 (both; provenance): Buchstaber intervention}
\noindent\textit{Transcription reads.} not a separate catalogue record\par
\noindent\textit{Critical reading.} retain the change of speaker within record393\par
\noindent\textit{Evidence.} The manuscript explicitly underlines Buchstaber before Novikov begins record394.\par
\subsection*{P000656 (both; faithful error): curvature Lie algebra}
\noindent\textit{Transcription reads.} O(n) in the Lorentzian example\par
\noindent\textit{Critical reading.} O(1,3) for the stated four-dimensional signature\par
\noindent\textit{Evidence.} The notebook has O(n); the orthonormal frame group is determined by the displayed metric.\par
\subsection*{P000656 (both; faithful error): electromagnetic two-form}
\noindent\textit{Transcription reads.} \(F_{\mu\nu}\,dx^\nu\wedge dx^\mu\)\par
\noindent\textit{Critical reading.} \(\tfrac12 F_{\mu\nu}\,dx^\mu\wedge dx^\nu\)\par
\noindent\textit{Evidence.} The displayed coordinate two-form needs the antisymmetric-tensor factor and consistent index order.\par
\subsection*{P000656 (both; ambiguity): general product row}
\noindent\textit{Transcription reads.} middle factor appears to have t\textasciicircum{}2\par
\noindent\textit{Critical reading.} all factors in the row have t\textasciicircum{}ell\par
\noindent\textit{Evidence.} The overwritten scan exponent is ambiguous; weight homogeneity forces ell. This is an emendation of an uncertain reading.\par
\subsection*{P000656 (both; faithful error): volume and stress variation}
\noindent\textit{Transcription reads.} \(\sqrt g\,T_{ij}=\delta S/\delta g_{ij}\)\par
\noindent\textit{Critical reading.} sqrt(abs(det g)) and a consistent factor-two, sign and index convention\par
\noindent\textit{Evidence.} Scan confirms lower indices on both T and the varied metric; differentiation of inverse metrics gives the corrected index relation.\par
\subsection*{P000657 (both; faithful error): determinant-one row completion}
\noindent\textit{Transcription reads.} arbitrary row length k\par
\noindent\textit{Critical reading.} row length at least two, or the sole entry must already equal one\par
\noindent\textit{Evidence.} A length-one matrix has determinant equal to its only entry; direct algebraic exception.\par
\subsection*{P000657 (both; ambiguity): nuclear scale}
\noindent\textit{Transcription reads.} 10\textasciicircum{}\{-23\} cm\par
\noindent\textit{Critical reading.} probably 10\textasciicircum{}\{-13\} cm, or a nuclear time of order 10\textasciicircum{}\{-23\} seconds\par
\noindent\textit{Evidence.} Scan inspected: the unit is cm. The intended emendation between exponent and unit is not uniquely recoverable.\par
\subsection*{P000657 (both; faithful error): root product for resultant}
\noindent\textit{Transcription reads.} product of root differences\par
\noindent\textit{Critical reading.} include the leading-coefficient powers unless both polynomials are monic\par
\noindent\textit{Evidence.} The source assumes only the first component monic in the preceding theorem; the missing factors matter.\par
\subsection*{P000658 (both; faithful error): dimension and signature}
\noindent\textit{Transcription reads.} M\textasciicircum{}n with four-entry signature\par
\noindent\textit{Critical reading.} M\textasciicircum{}4 in this discussion\par
\noindent\textit{Evidence.} The scan confirms n, but both the signature and velocity have exactly four components.\par
\subsection*{P000658 (both; ambiguity): reduction to resultant one}
\noindent\textit{Transcription reads.} need c=1\par
\noindent\textit{Critical reading.} patch several resultants generating the unit ideal\par
\noindent\textit{Evidence.} The written reduction omits the unit-ideal patching supplied by the general Suslin resultant lemma.\par
